\section*{集合、常用逻辑用语}
\subsection*{A Level}
\begin{enumerate}
    \item 集合 $A = \left\{  {4,{m}^{2} + {3m}}\right\}$ 中实数 $m$ 的取值范围是\blkda[3]{$\{x | x \neq 1 \text{或} -4\}$}
    \begin{jieda}
        ${m}^{2} + {3m} \neq 4$，解得$m \in \{x | x \neq 1 \text{或} -4\}$
    \end{jieda}
    \item 已知集合 $A = \{  - 1,3,0\} ,B = \left\{  {3,{m}^{2}}\right\}$ ,若 $B \subseteq  A$ ,则实数 $m$ 的值为\blkda[4]{0}.
    \begin{jieda}
        因为 $B \subseteq  A$ ,所以 ${m}^{2} =  - 1$ (舍去) 或 ${m}^{2} = 0$，所以 $m = 0$ .
    \end{jieda}
    \item 已知全集 $U = \mathbf{R}$ ,集合 $A = \{ x \mid  x < a\} ,B = \{ x \mid   - 1 < x < 3\}$ ,且 $B \subseteq  A$ ,则实数 $a$ 的取值范围是\blkda[4]{$[3, +\infty)$}.
    \begin{jieda}
        因为 $B \subseteq  A$ ,又因为 $A = \{ x \mid  x < a\} ,B = \{ x \mid   - 1 < x < 3\}$ ,所以 $a \geq  3$ .
    \end{jieda}
    \item 若集合 $A = \left\{  {2,{3a} - 1,{a}^{2} - {4a}}\right\}$ ,且 $- 4 \in  A$ ,则 $a =$ \blkda[4]{-1 或 2}.
    \begin{jieda}
        因为 $A = \left\{  {2,{3a} - 1,{a}^{2} - {4a}}\right\}$ ,且 $- 4 \in  A$ ,
        
        所以 ${3a} - 1 =  - 4$ 或 ${a}^{2} - {4a} =  - 4$ ,
        
        当 ${3a} - 1 =  - 4$ 时, $a =  - 1$ ,此时 $A = \{ 2, - 4,5\}$ ,满足题意;
        
        当 ${a}^{2} - {4a} =  - 4$ 时, $a = 2$ ,此时 $A = \{ 2,5, - 4\}$ ,满足题意,
        
        综上所述, $a =  - 1$ 或 2 .
    \end{jieda}
    \item 已知集合 $M = \{ x \mid   - 1 < x < 3\} ,N = \{ x \mid   - 2 < x < 1\}$ ,则 $M \cup  N =$ \blkda[3]{$\{ x \mid   - 2 < x < 3\}$}.
    \begin{jieda}
        由题意可知 $M \cup  N = \{ x \mid   - 2 < x < 3\}$ .
    \end{jieda}
    \item 已知集合 $A = \{ x \mid  x < a\} ,B = \{ x \mid  x \geq  1\}$ ,且 $A \cup  B = \mathbf{R}$ ,则实数 $a$ 的取值范围是\blkda[4]{$\lbrack 1, + \infty )$}.
    \begin{jieda}
        由题意知 $A \cup  B = \mathbf{R}$ ,可得 $a \geq  1$ .
    \end{jieda}
    
    \item 设 $U$ 表示全集,下列命题中假命题有\blkda[4]{0}个
    \begin{dingautolist}{172}
        \item 若 $A \cap  B = U$ ,则 $A = B = U$
        \item 若 $A \cup  B = \varnothing$ ,则 $A = B = \varnothing$
        \item 若 $A \cap  B = \varnothing$ ,则 $\overline{A} \cup  \overline{B} = U$
        \item 若 $A \cup  B = U$ ,则 $\overline{A} \cap  \overline{B} = \varnothing$
        \item 若 $A \subseteq  B$ ,则 $\overline{A} \supseteq  \overline{B}$
        \item $\varnothing  \in  \{ \varnothing \}$
        \item $\varnothing  \subseteq  \{ \varnothing \}$
    \end{dingautolist}
    \item “实数 $a$ 、 $b$ 都是正数”的否定形式为\blkda[6]{实数 $a\text{ 、 }b$ 中至少有一个是非正数}.
    
    \item $x = 2$ 是 ${x}^{2} - {3x} + 2 = 0$ 的\blkda[3]{充分不必要}条件.
    \begin{jieda}
        ${x}^{2} - {3x} + 2 = 0$ ,解得 $x = 1$ 或 2，故 $x = 2$ 是 ${x}^{2} - {3x} + 2 = 0$ 的充分不必要条件.
    \end{jieda}
    
    \item 若 $p :  - 2 < x < 2,q : \sqrt{x} < 4$ ,则 $p$ 是 $q$ 的\blkda[3]{既非充分又非必要}条件.
    \begin{jieda}
        $\because q : \sqrt{x} < 4$ ,
        
        $\therefore q : 0 \leq  x < {16}$ ,
        
        $\because  - 2 < x < 2$ 既不能推出 $0 \leq  x < {16}$ ,也不能被 $0 \leq  x < {16}$ 推出
    \end{jieda}
    \item 不等式 $m - 2 < x < m + 1$ 成立的充分非必要条件是 $0 < x < 1$ ,则 $m$ 的取值范围是\blkda[4]{$\left\lbrack  {0,2}\right\rbrack$}.
    \begin{jieda}
        由题知$(0,1)$是 $\left( {m - 2,m + 1}\right)$ 的真子集，所以 $\left\{  \begin{array}{l} m - 2 \leq  0, \\  m + 1 \geq  1, \end{array}\right.$ 且等号不同时成立，解得 $0 \leq  m \leq  2$ 
        
        所以 $m$ 的取值范围是 $\left\lbrack  {0,2}\right\rbrack$ .
    \end{jieda}
    \item 已知命题 $p : \exists x \in  \mathbf{R},{x}^{2} - {3x} + 3 \leq  0$ ,则 $\neg p$ 为\blkda[6]{$\forall x \in  \mathbf{R},{x}^{2} - {3x} + 3 > 0$}.
\end{enumerate}
\clearpage
\subsection*{B Level}
\begin{enumerate}

    \item 下列各组集合中相等的是\blkda[4]{\ding{172}\ding{174}}
    \begin{dingautolist}{172}
        \item $P = \{ x \mid  x = {2n},n \in  \mathbf{Z}\} ,Q = \{ x \mid  x = 2\left( {n - 1}\right) ,n \in  \mathbf{Z}\}$
        \item $P = \left\{  {x \mid  {x}^{2} - 1 = 0}\right\}  ,Q = \left\{  {x\left| {x = \dfrac{1 + {\left( -1\right) }^{n}}{2}}\right. ,n \in  \mathbf{Z}}\right\}$
        \item $P = \{ x \mid  x = {2m} + 1,m \in  \mathbf{Z}\} ,Q = \{ x \mid  x = {4m} \pm  1,m \in  \mathbf{Z}\}$
        \item $P = \left\{  {y \mid  y = {x}^{2} + 1}\right\}  ,Q = \left\{  {x \mid  y = {x}^{2} + 1}\right\}$
    \end{dingautolist}
    \begin{jieda}
        \begin{dingautolist}{172}
            \item 均为偶数集
            \item $P = \{1, -1\}$，$Q = \{0, 1\}$
            \item 均为奇数集
            \item 表示元素的符号不一样，$P = [1, +\infty)$，$Q = (-\infty, +\infty)$
        \end{dingautolist}
    \end{jieda}
    \item 已知集合 $A = \left\{  {1,3,{a}^{2}}\right\}$ ,集合 $B = \{ 1,2 + a\}$ ,若 $A \cup  B = A$ ,则 $a =$ \blkda[4]{2}.
    \begin{jieda}
        因为 $A \cup  B = A$ ,所以 $2 + a = 3$ 或 ${a}^{2} = 2 + a$ .

        若 $2 + a = 3$ ,则 $a = 1$ ,此时 ${a}^{2} = 1$ ,集合 ${A}$ 中的元素不满足互异性,故 $a = 1$ 舍去.
        
        若 ${a}^{2} = 2 + a$ 则 $a =  - 1$ 或 $a = 2$ .
        
        当 $a =  - 1$ 时, ${a}^{2} = 1$ ,集合 ${A}$ 中的元素不满足互异性,故 $a =  - 1$ 舍去;
        
        当 $a = 2$ 时, $A = \{ 1,3,4\} ,B = \{ 1,3\} ,A \cup  B = A$ ,故 $a = 2$ 符合题意.
    \end{jieda}
    \item 已知 $a \in \mathbf{R}$, $M = \left\{  {x \mid  {x}^{2} - {3x} - {a}^{2} + 2 = 0,x \in  \mathbf{R}}\right\}$ 的非空真子集个数是\blkda[4]{2} 
    \begin{jieda}
        $\Delta = 9 + 4a^2 - 8 = 4a^2 + 1 > 0$，故$M$中必定有两个元素，因此$M$的非空真子集个数为$2^2 - 2 = 2$.
    \end{jieda}
    
    \item 已知 $U = \{ \left( {x,y}\right)  \mid  x,y \in  \mathbf{R}\} ,M = \left\{  {\left( {x,y}\right) \left| {\dfrac{y - 2}{x - 1} = 3}\right. }\right\}  ,N = \{ \left( {x,y}\right)  \mid  y = {3x} - 1\} ,\overline{M} \cap N =$ \blkda[4]{$\{(1,2)\}$}
    \begin{jieda}
        $M$对应于直线$y = 3x-1$上除点$(1, 2)$以外的点

        $N$对应于直线$y = 3x-1$上的点，因此$\overline{M} \cap N$对应点$(1, 2)$
    \end{jieda}
    
    \item $x^2 > 1$ 是 $\dfrac{1}{x} < 1$ 的\blkda[3]{充分不必要}条件.
    \begin{jieda}
        $x^2 > 1 \Longleftrightarrow x \in (-\infty, -1) \cup (1, +\infty)$

        $\dfrac{1}{x} < 1 \Longleftrightarrow x \in (-\infty, 0) \cup (1, +\infty)$
    \end{jieda}
    
    \item 已知 $p : x > m + 3$ 或 $x < m,q :  - 4 < x < 1$ ,且 $p$ 是 $q$ 的必要而不充分条件,则实数 $m$ 的取值范围是\blkda[3]{$m \leq   - 7$ 或 $m \geq  1$}.
    \begin{jieda}
        因为 $p$ 是 $q$ 的必要而不充分条件，所以命题 $q$ 中变量的取值集合是命题 $p$ 中变量的取值集合的真子集，即 $m + 3 \leq   - 4$ 或 $m \geq  1$ ，故 $m \leq   - 7$ 或 $m \geq  1$
    \end{jieda}
    
    \item “ $- 4 < k < 0$ ” 是 “函数 $y = k{x}^{2} - {kx} - 1$ 的值恒为负值” 的\dotfill(\kh{A})
    \xx{充分非必要条件}{必要非充分条件}{充分必要条件}{即非充分也非必要条件}
    \begin{jieda}
        \begin{dingautolist}{172}
            \item $k = 0$时，函数$y = k{x}^{2} - {kx} - 1$恒为负值
            \item $k \neq 0$时，函数$y = k{x}^{2} - {kx} - 1$恒为负值的必要条件是$k < 0$，此时函数的最大值在其顶点处取得，计算对称轴为$x = \dfrac{1}{2}$，故函数的最大值为$-\dfrac{k}{4}-1$，解得$- 4 < k < 0$
        \end{dingautolist}
        故函数$y = k{x}^{2} - {kx} - 1$恒为负值的充要条件是$- 4 < k \leq 0$，故选A.
    \end{jieda}

    \item 请写出“如果两个三角形全等,那么它们的面积相等”的逆否命题:\blkda[10]{“如果两个三角形面积不相等，那么它们必然不全等”}
    \item 如果一个命题的逆命题是真命题,那么这个命题的否命题是\dotfill(\kh{A})
    \xx{真命题}{假命题}{不一定是真命题}{不一定是假命题}
    \begin{jieda}
        逆命题和否命题互为逆否命题，真假性相同
    \end{jieda}
    \item 已知集合 $A = \left\{  {x \mid  {x}^{2} - {3x} + 2 = 0}\right\}  ,B = \left\{  {x \mid  {x}^{2} - {bx} + b - 1 = 0}\right\}  ,C = \left\{  {x \mid  {x}^{2} - {cx} + }\right.$  $2 = 0\}$ 
    \begin{enumerate}
        \item 若 $A \cup  B = A$ ,求实数 $b$ 的值
        \item 若 $A \cap  C = C$ ,求实数 $c$ 的取值范围
    \end{enumerate}
    \begin{jieda}
        \begin{enumerate}
            \item $A \cup  B = A \Longleftrightarrow B \subseteq A$，$A = \{1, 2\}$，$1 \in B, b-1 \in B$

            故$b - 1 \in \{1,2\}$，因此$b \in \{2,3\}$
            \item $A \cap  C = C \Longleftrightarrow C \subseteq A$，$A = \{1, 2\}$
            
            \begin{dingautolist}{172}
                \item $C = \varnothing$，此时$\Delta = c^2 - 8 < 0$，故$c \in (-2\sqrt{2}, 2\sqrt{2})$
                \item $C \neq \varnothing$，若$1 \in C$，则$c = 3$，解得$C = \{1,2\}$，满足要求；若$2 \in C$可得到相同结果.

                \ps 可以考虑韦达定理，可知必然当$C \neq \varnothing$时，两根为$1,2$时满足要求
            \end{dingautolist}
            综上，$c \in \{3\} \cup (-2\sqrt{2}, 2\sqrt{2})$
        \end{enumerate}
    \end{jieda}

    \item 已知全集 $U = \mathbf{R},A = \left\{  {x \mid  {x}^{2} + {px} + {12} = 0}\right\}  ,B = \left\{  {x \mid  {x}^{2} - {5x} + q = 0}\right\}$ , $\overline{A} \cap  B = \{ 2\}$ ,求 $p + q$ 的值.
    \begin{jieda}
        $\overline{A} \cap  B = \{ 2\} \Longleftrightarrow 2 \in B, 2 \notin A$，
        
        $2 \in B \Longrightarrow q = 6 \Longrightarrow B = \{2, 3\}$，又因为$2 \notin \overline{A} \cap  B$，所以$3 \in A$

        $3 \in A \Longrightarrow p = -7$，故$A = \{3,4\}$

        故$p+q = -1$
    \end{jieda}
    \item 已知集合 $A = \{ x \mid  a - 1 \leq  x \leq  {2a} + 3\} ,B = \{ x \mid   - 1 \leq  x \leq  4\}$ ,全集 $U = R$ .
    \begin{enumerate}
        \item 当 $a = 1$ 时,求 $\left( {{C}_{U}A}\right)  \cap  B$
        \item 若“ $x \in  B$ ”是“ $x \in  A$ ”的必要条件,求实数 $a$ 的取值范围
    \end{enumerate}
    \begin{jieda}
        \begin{enumerate}
            \item 当 $a = 1$ 时,集合 $A = \{ x \mid  0 \leq  x \leq  5\} ,{C}_{U}A = \{ x \mid  x < 0$ 或 $x > 5\}$，$\left( {{C}_{U}A}\right)  \cap  B = \{ x \mid   - 1 \leq  x < 0\} .$
            
            \item 若“ $x \in  B$ ”是“ $x \in  A$ ”的必要条件,则 $A \subseteq  B$
            \begin{dingautolist}{172}
                \item 当 $A = \varnothing$ 时, $a - 1 > {2a} + 3,\therefore a <  - 4$
                \item $A \neq  \varnothing$ ,则 $a \geq   - 4$ 且 $a - 1 \geq   - 1,{2a} + 3 \leq  4,\therefore 0 \leq  a \leq  \dfrac{1}{2}$
            \end{dingautolist}
            综上，$a <  - 4$ 或 $0 \leq  a \leq  \dfrac{1}{2}$ .
        \end{enumerate}
    \end{jieda}
\end{enumerate}
\clearpage
\subsection*{C Level}
\begin{enumerate}
    \item 已知集合 $A = \{ x \mid  {5x} - a \leq  0\} ,B = \{ x \mid  {6x} - b > 0\} ,a,b \in  \mathbf{N}$ ,且 $A \cap  B \cap  \mathbf{N} = \{ 2,3,4\}$ ,则整数对$(a, b)$的个数为\blkda[4]{30}
    \begin{jieda}
        $A = \left(-\infty, \dfrac{a}{5}\right]$，$B = \left(\dfrac{b}{6}, +\infty\right)$，根据$A \cap  B \cap  \mathbf{N} = \{ 2,3,4\}$可知
        \begin{dingautolist}{172}
            \item $4 \in A$，$5 \notin A$，故$\dfrac{a}{5} \in [4, 5)$，即$a \in \{20, 21, 22, 23, 24\}$
            \item $2 \in B$，$1 \notin B$，故$\dfrac{b}{6} \in [1,2)$，即$b \in \{6,7,8,9,10,11\}$
        \end{dingautolist}
        故整数对$(a, b)$的个数为$30$.
    \end{jieda}
    \item 集合 $P$ 具有性质 “若 $x \in  P$ ,则 $\dfrac{1}{x} \in  P$ ”,就称集合 $P$ 是伙伴关系的集合,集合 $A = \{  - 1,0$ , $\dfrac{1}{2},\dfrac{1}{3},1,2,3,4\}$ 的所有非空子集中具有伙伴关系的集合的个数为\blkda[4]{15}.
    \begin{jieda}
        由题意得 3 和 $\dfrac{1}{3}$ 、2 和 $\dfrac{1}{2}$ 、$-1$与自身、$1$与自身 分别为 4 组满足条件 $x$ 与 $\dfrac{1}{x}$ 的伙伴,所以满足伙伴关系的非空子集有 ${2}^{4} - 1 = {15}$ (个).
    \end{jieda}
    \item 在整数集中,被 4 除所得余数为 $k$ 的所有整数组成一个 “类”,记为 $\left\lbrack  k\right\rbrack$ ,即 $\left\lbrack  k\right\rbrack   =$  $\{ {4n} + k \mid  n \in  \mathbf{Z}\} ,k = 0,1,2,3$ . 给出下列四个结论:
    \begin{dingautolist}{172}
        \item ${2021} \in  \left\lbrack  1\right\rbrack$
        \item $- 1 \in  \left\lbrack  1\right\rbrack$
        \item ${Z} = \left\lbrack  0\right\rbrack   \cup  \left\lbrack  1\right\rbrack   \cup  \left\lbrack  2\right\rbrack   \cup  \left\lbrack  3\right\rbrack$
        \item “整数 $a\text{ 、 }b$ 属于同一‘类’” 的充要条件是 “ $a - b \in  \left\lbrack  0\right\rbrack$ ”
    \end{dingautolist}
    其中,正确结论的序号是\blkda[4]{\ding{172}\ding{174}\ding{175}}.
    \begin{jieda}
        \begin{dingautolist}{172}
        \item 因为 ${2021} = 4 \times  {505} + 1$ ,则 ${2021} \in  \left\lbrack  1\right\rbrack$
        \item  $- 1 = 4 \times  \left( {-1}\right)  + 3$ ,则 $- 1 \in  \left\lbrack  3\right\rbrack$
        \item 任意整数除以 4,余数可以且只可以是 0、1、2、3四类,则 $\mathbf{Z} = \left\lbrack  0\right\rbrack   \cup  \left\lbrack  1\right\rbrack   \cup  \left\lbrack  2\right\rbrack   \cup$ [3]
        \item 
        充分性：
        
        若整数 $a\text{ 、 }b$ 属于同一 “类”,则整数 $a\text{ 、 }b$ 被 4 除的余数相同,可设 $a = 4{n}_{1} + k,b =$  $4{n}_{2} + k$ ,其中 ${n}_{1}$ 、 ${n}_{2} \in  \mathbf{Z},k \in  \{ 0,1,2,3\}$ ,则 $a - b = 4\left( {{n}_{1} - {n}_{2}}\right)$ ,故 $a - b \in  \left\lbrack  0\right\rbrack$

        必要性：
        
        若 $a - b \in  \left\lbrack  0\right\rbrack$ ,不妨令 $a = 4{n}_{1} + {k}_{1},b = 4{n}_{2} + {k}_{2}\left( {{n}_{1},{n}_{2} \in  \mathbf{Z},{k}_{1},{k}_{2} \in  \{ 0,1,2,3\} }\right)$ ,则 $a -$  $b = 4\left( {{n}_{1} - {n}_{2}}\right)  + \left( {{k}_{1} - {k}_{2}}\right)$ ,显然 ${n}_{1} - {n}_{2} \in  \mathbf{Z}$ , $\left| {{k}_{1} - {k}_{2}}\right|  \in  \{ 0,1,2,3\}$ ,于是得 $\left| {{k}_{1} - {k}_{2}}\right|  =$ 0,所以 ${k}_{1} = {k}_{2}$ ,即整数 $a,b$ 属于同一“类”
        
        所以“整数 $a,b$ 属于同一‘类’” 的充要条件是 “ $a - b \in  \left\lbrack  0\right\rbrack$ ”
        \end{dingautolist}
    \end{jieda}
    \item 已知 $A = \left\{  {x \mid  x = {a}^{2} - {b}^{2},a \in  \mathbf{Z},b \in  \mathbf{Z},x\text{ 为正整数, }x \leq  {40}}\right\}$ ,则集合 $A$ 中共含\blkda[4]{30} 个元素
    \begin{jieda}
        $a^2-b^2 = (a+b)(a-b)$，因为$(a-b)$和$(a+b)$奇偶性相同，故当$x$只能分解为一个奇数和一个偶数相乘时，不在集合$A$内，即$2$的奇数倍（$2,6,10,14,18,22,26,30,34,38$）不在集合$A$中.
    \end{jieda}
    \item 已知集合$M = \left\{x \left| \dfrac{ax- 5}{x^2-a} < 0 \right.\right\}$，若$3 \in M, 5 \notin M$，则实数$a$的取值范围是\blkda[3]{$\left[1, \dfrac{5}{3}\right) \cup (9, 25]$}
    \begin{jieda}
        \begin{enumerate}
            \item $3 \in M \Longrightarrow \dfrac{3a - 5}{9-a} < 0$
            \item $5 \notin M \Longrightarrow \dfrac{5a - 5}{25-a} \geq 0$ 或 $a = 25$(此时$\dfrac{5a - 5}{25-a}$无意义，满足条件）
        \end{enumerate}
        求解分式不等式：
        \begin{enumerate}
            \item $\dfrac{3a - 5}{9-a} < 0 \Longrightarrow (3a-5)(9-a) < 0 \Longrightarrow a \in (-\infty, \dfrac{5}{3}) \cup (9, +\infty)$
            \item $\dfrac{5a - 5}{25-a} \geq 0 \Longrightarrow (5a-5)(25-a) \geq 0, a \neq 25 \Longrightarrow a \in [1, 25)$
        \end{enumerate}
        两个不等式的解集的交集并上$\{25\}$，得到$\left[1, \dfrac{5}{3}\right) \cup (9, 25]$，即是$a$的取值范围
    \end{jieda}
    \item 已知集合$A = \{(x, y) | |x| + 2|y| \leq 4\}$，集合$B = \left\{(x, y) \left| (x-m)^2 + y^2 = \dfrac{4}{5}\right.\right\}$，若$B \subseteq A$，则实数$m$的取值范围是\blkda[4]{$[-2,2]$}
    \begin{jieda}
        首先可以判断出来集合$A, B$都是点集，因此应先找到其在坐标系内对应的图形\\
        集合$A$所对应的区域是坐标系内以$(0, 2), (4, 0), (0, -2), (-4, 0)$为顶点的菱形区域内部（包括边界），而集合$B$对应的是以$(m, 0)$为圆心，$\dfrac{2\sqrt{5}}{5}$为半径的圆\\
        由$B \subseteq A$ 可知圆不可超出菱形的边界已知$A, B$都是关于$y$轴对称，因此先考虑$m > 0$的相切的边界场景根据三角形相似可以得到$\dfrac{4-m}{2\sqrt{5}} = \dfrac{\dfrac{2\sqrt{5}}{5}}{2}$，此时$m = 2$故$m \in [-2,2]$
    \end{jieda}
\end{enumerate}
\clearpage
\subsection*{D Level}
\begin{enumerate}
    \item 如果集合 $A = \{ 1,2,3,\cdots ,8\}$ 的子集 $B$ 满足: 若 ${2k} \in  B,k$ 为正整数,则 ${2k} + 1 \in  B$ ,则称 $B$ 为 $A$ 的“好子集”, $A$ 的“好子集”有\blkda[4]{54}个
    \begin{jieda}
        \begin{dingautolist}{172}
            \item “好子集”中不能有8
            \item 对于绑定的元素，$2,3$，$4,5$，$6,7$来讲，存在三种情况，两个数都在“好子集”中，只有其中的奇数在“好子集”中，两个数都不在“好子集”中
        \end{dingautolist}
        因此，根据乘法原理可知“好子集”有$2\times3\times3\times3 = 54$个.
    \end{jieda}
    \item 十个元素组成的集合 $M = \{ {19},{93}, - 1,0,{25}, - {78}, - {94},1,{17}, - 2\}$  $M$ 的所有非空子集记为 ${M}_{i}\left( {i = 1,2,\cdots ,{1023}}\right)$ ,每一非空子集中所有元素的乘积记为 ${m}_{i}(i = 1,2,\cdots$ , 1023),则 ${m}_{1} + {m}_{2} + \cdots  + {m}_{1023} =$ \blkda[4]{$-1$}
    \begin{jieda}
        将$M_i$分成三类
        \begin{dingautolist}{172}
            \item $\{-1\}$
            \item $S_i$满足$-1 \in S_i$，且$S_i \neq \{-1\}$
            \item $T_i$满足$-1 \notin T_i$
        \end{dingautolist}
        易知$S_i$与$T_i$一一对应，且其元素乘积互为相反数，因此在计算目标值的时候，这两部分的和为$0$，故目标值为$-1$
    \end{jieda}
    \item 设集合 $S,T,S \subseteq  {\mathbf{N}}^{ * },T \subseteq  {\mathbf{N}}^{ * },S,T$ 中,至少有两个元素,且 $S,T$ 满足: \ding{172}对于任意 $x$ , $y \in  S$ ,若 $x \neq  y$ ,都有 ${xy} \in  T$ ; \ding{173}对于任意 $x,y \in  T$ ,若 $x < y$ ,则 $\dfrac{y}{x} \in  S$ . 若 $S$ 有 4 个元素,则 $S \cup  T$ 有\blkda[4]{7}个元素.
    \begin{jieda}
        设集合 $S = \left\{  {{p}_{1},{p}_{2},{p}_{3},{p}_{4}}\right\}$ ,且 ${p}_{1} < {p}_{2} < {p}_{3} < {p}_{4}$，${p}_{1},{p}_{2},{p}_{3},{p}_{4} \in  {\mathbf{N}}^{ * }$ ,则 $p_1p_2, p_1p_3,p_1p_4, p_2p_3,p_2p_4,p_3p_4 \in  T$，则$\dfrac{p_2}{p_1},\dfrac{p_3}{p_1},\dfrac{p_4}{p_1},\dfrac{p_3}{p_2},\dfrac{p_4}{p_2},\dfrac{p_4}{p_3} \in S$，且$\dfrac{p_3p_4}{p_1p_2}, \dfrac{p_2p_4}{p_1p_3} \in S$

        易知$p_4 \geq \dfrac{p_3p_4}{p_1p_2} > \dfrac{p_4}{p_1}>\dfrac{p_4}{p_2} > \dfrac{p_4}{p_3} \geq p_1$

        故必有$\dfrac{p_3p_4}{p_1p_2} = p_4$，$ \dfrac{p_4}{p_1} = p_3$，$\dfrac{p_4}{p_2} = p_2$，$\dfrac{p_4}{p_3} = p_1$，即$S = \{p_1, p_1^2, p_1^3, p_1^4\}$

        故此时可知$T = \{p_1^3, p_1^4, p_1^5, p_1^6, p_1^7\}$，$S \cup  T = \left\{  {{p}_{1},{p}_{1}^{2},{p}_{1}^{3},{p}_{1}^{4},{p}_{1}^{5},{p}_{1}^{6},{p}_{1}^{7}}\right\}$ , 即 $S \cup  T$ 中有 7 个元素.
    \end{jieda} 
\end{enumerate}
\section*{韦达定理、不等式的性质、一元二次方程根的分布}
\subsection*{A Level}
\begin{enumerate}
    \item 已知 ${x}_{1},{x}_{2}$ 是方程 ${x}^{2} - {5x} + 3 = 0$ 的两个实数根,则 $\dfrac{{x}_{1}}{{x}_{2}} + \dfrac{{x}_{2}}{{x}_{1}}$ 的值是\blkda[4]{$\dfrac{19}{3}$}.
    \begin{jieda}
        $\because {x}_{1},{x}_{2}$ 是方程 ${x}^{2} - {5x} + 3 = 0$ 的两个实数根,
        
        $\therefore {x}_{1}{x}_{2} = \dfrac{3}{1} = 3,{x}_{1} + {x}_{2} =  - \dfrac{-5}{1} = 5$ ,
        
        $\therefore \dfrac{{x}_{1}}{{x}_{2}} + \dfrac{{x}_{2}}{{x}_{1}} = \dfrac{{x}_{1}^{2} + {x}_{2}^{2}}{{x}_{1}{x}_{2}} = \dfrac{{\left( {x}_{1} + {x}_{2}\right) }^{2} - 2{x}_{1}{x}_{2}}{{x}_{1}{x}_{2}} = \dfrac{{5}^{2} - 2 \times  3}{3} = \dfrac{19}{3}$ .
        
    \end{jieda}
    
    \item 已知函数 $f\left( x\right)  = {x}^{2} - {4x} - 1$ 的两个零点分别为 ${x}_{1}$ 和 ${x}_{2}$ ,则 ${x}_{1}^{2} + {x}_{2}^{2}$ 的值为\blkda[4]{18}.
    \begin{jieda}
        因为函数 $f\left( x\right)  = {x}^{2} - {4x} - 1$ 的两个零点分别为 ${x}_{1}$ 和 ${x}_{2}$，所以 ${x}_{1}$ 和 ${x}_{2}$ 是 ${x}^{2} - {4x} - 1 = 0$ 的两个实根，所以 ${x}_{1} + {x}_{2} = 4,{x}_{1}{x}_{2} =  - 1$，所以 ${x}_{1}^{2} + {x}_{2}^{2} = {\left( {x}_{1} + {x}_{2}\right) }^{2} - 2{x}_{2}{x}_{2} = {4}^{2} - 2 \times  \left( {-1}\right)  = {18}$ .
    \end{jieda}
    \item 若方程 ${x}^{2} - {3x} - 1 = 0$ 的两实数根分别是 ${x}_{1}$ 和 ${x}_{2}$ . 则 $\dfrac{1}{{x}_{1}} + \dfrac{1}{{x}_{2}} =$ \blkda[4]{$- 3$}, $\left| {{x}_{1} - {x}_{2}}\right|  =$ \blkda[4]{$\sqrt{13}$}.
    \begin{jieda}
        由韦达定理知 ${x}_{1} + {x}_{2} = 3,{x}_{1}{x}_{2} =  - 1$ ,
        
        所以 $\dfrac{1}{{x}_{1}} + \dfrac{1}{{x}_{2}} = \dfrac{{x}_{1} + {x}_{2}}{{x}_{1}{x}_{2}} = \dfrac{3}{-1} =  - 3$
        
        ${\left( {x}_{1} - {x}_{2}\right) }^{2} = {\left( {x}_{1} + {x}_{2}\right) }^{2} - 4{x}_{1}{x}_{2} = 9 + 4 = {13}$ ,
        
        $\left| {{x}_{1} - {x}_{2}}\right|  = \sqrt{13}$
    \end{jieda}
    \item 设 ${x}_{1},{x}_{2}$ 是方程 ${x}^{2} + x - 3 = 0$ 的两个实数根,则 ${x}_{1}^{3} + {x}_{2}^{3} =$ \blkda[4]{$-10$}.
    \begin{jieda}
        因为 ${x}_{1},{x}_{2}$ 是方程 ${x}^{2} + x - 3 = 0$ 的两个实数根，所以 ${x}_{1}^{2} + {x}_{1} - 3 = 0{x}_{2}^{2} + {x}_{2} - 3 = 0$，即 ${x}_{1}^{2} = 3 - {x}_{1},{x}_{2}^{2} = 3 - {x}_{2}$ ,且 ${x}_{1} + {x}_{2} =  - 1,{x}_{1}{x}_{2} =  - 3$ ,
        
        所以 ${x}_{1}^{3} + {x}_{2}^{3} = \left( {{x}_{1} + {x}_{2}}\right) \left( {{x}_{1}^{2} - {x}_{1}{x}_{2} + {x}_{2}^{2}}\right)  = \left( {{x}_{1} + {x}_{2}}\right) \left( {3 - {x}_{1} - {x}_{1}{x}_{2} + 3 - {x}_{2}}\right) =  - 1 \times  \left( {6 + 1 + 3}\right)  =  - {10}$
    \end{jieda}
    \item 若 $a,b,c \in  \mathbf{R},a > b$ ,则下列不等式中成立的是\blkda[4]{\ding{174}}(填上正确的序号).
    \begin{dingautolist}{172}
        \item $\dfrac{1}{a} < \dfrac{1}{b}$
        \item ${a}^{2} > {b}^{2}$
        \item $\dfrac{a}{{c}^{2} + 1} > \dfrac{b}{{c}^{2} + 1}$
        \item $a\left| c\right|  > b\left| c\right|$
    \end{dingautolist}
    \begin{jieda}
        当 $a = 1,b =  - 1,c = 0$ 时,有 $a > b$ ,但 $\dfrac{1}{a} = 1 >  - 1 = \dfrac{1}{b},{a}^{2} = 1 = {b}^{2},a\left| c\right|  = 0 = b\left| c\right|$，故\ding{172}\ding{173}\ding{175}错误

        由于 $a > b,\dfrac{1}{{c}^{2} + 1} > 0$ ,故 $\dfrac{a}{{c}^{2} + 1} > \dfrac{b}{{c}^{2} + 1}$，故\ding{174}正确
    \end{jieda}
    \item 已知 $x,y \in  \mathbf{R}$ ,则 “ $x > 0,y > 0$ ” 是 “ ${xy} > 0$ ” 的\blkda[3]{充分不必要}条件, “ ${x}^{2} + {y}^{2} > 0$ ” 是 “ $x \neq  0$ ” 的\blkda[3]{必要不充分}条件. (填“充分不必要”或“必要不充分”或“充要”或“既不充分也不必要”)
    \begin{jieda}
        由 $x > 0,y > 0$ 可得 ${xy} > 0$
        
        由 ${xy} > 0$ ,得 $x > 0,y > 0$ 或 $x < 0,y < 0$ 
        
        所以 “ $x > 0,y > 0$ ” 是 “ ${xy} > 0$ ” 的充分不必要条件；
        
        由 $x \neq  0$ 可得 ${x}^{2} + {y}^{2} > 0$
        
        由 ${x}^{2} + {y}^{2} > 0$ ,得 $x \neq  0$ 或 $y \neq  0$
        
        所以 “ ${x}^{2} + {y}^{2} > 0$ ” 是 “ $x \neq  0$ ” 的必要不充分条件.
    \end{jieda}
    \item 已知实数 $a$ 、 $b$ 满足 $a > b$ ,下列不等式中正确的是\blkda[4]{\ding{174}}(只填序号).
    \begin{dingautolist}{172}
        \item $\left| a\right|  > \left| b\right|$
        \item ${a}^{2} > {b}^{2}$
        \item ${a}^{3} > {b}^{3}$
        \item $\dfrac{a}{b} > 1$
    \end{dingautolist}
    \begin{jieda}
        对于\ding{172}\ding{173}\ding{175}，取 $a = 2,b =  - 3$ ,满足 $a > b$ ,而$\left| a\right|  = 2 < 3 = \left| b\right| ,{a}^{2} = {2}^{2} < {\left( -3\right) }^{2} = {b}^{2},\dfrac{a}{b} =  - \dfrac{2}{3} < 1$，故错误；

        对于\ding{174}，${a}^{3} - {b}^{3} = \left( {a - b}\right) \left( {{a}^{2} + {ab} + {b}^{2}}\right)  = \left( {a - b}\right) \left\lbrack  {{\left( a + \dfrac{1}{2}b\right) }^{2} + \dfrac{3}{4}{b}^{2}}\right\rbrack$ ,

        由 $a > b$ ,得 $a - b > 0,{\left( a + \dfrac{1}{2}b\right) }^{2} + \dfrac{3}{4}{b}^{2} > 0$ ,因此 ${a}^{3} - {b}^{3} > 0,{a}^{3} > {b}^{3}$
    \end{jieda}
    \item 设 $a \in  \mathbf{R}$ ,“ $a > 1$ ” 是 “ $\dfrac{1}{a} < 1$ ” 的一个\blkda[3]{充分非必要}条件(充分非必要、必要非充分、充要、既非充分又非必要)
    \begin{jieda}
        $\dfrac{1}{a} < 1 \Longleftrightarrow a \in (-\infty, 0) \cup (0, 1)$
    \end{jieda}
    \item 已知 $- 1 < x < 4,2 < y < 3$ ,则 ${3x} + {2y}$ 的取值范围是\blkda[4]{$(1,18)$}.
    \begin{jieda}
        因为 $- 1 < x < 4,2 < y < 3$，所以 $- 3 < {3x} < {12},4 < {2y} < 6$，得 ${3x} + {2y} \in  \left( {-3 + 4,{12} + 6}\right)  = \left( {1,{18}}\right)$ .
    \end{jieda}
    \item 已知 $x > 3,y > 4$ ,则 ${xy}$ 的取值范围为\blkda[4]{$\left( {{12}, + \infty }\right)$}.
    \begin{jieda}
        因为 $x > 3,y > 4 > 0$，所以 ${xy} > {3y} > 3 \times  4 = {12}$，即 ${xy}$ 的取值范围为 $\left( {{12}, + \infty }\right)$ .
    \end{jieda}
    \item 若 $\alpha ,\beta$ 满足 $- \dfrac{\pi }{2} < \alpha  \leq  \beta  \leq  \dfrac{\pi }{2}$ ,则 $\alpha  - \beta$ 的取值范围是\blkda[4]{$( - \pi ,0\rbrack$}
    \begin{jieda}
        因为 $- \dfrac{\pi }{2} < \alpha  \leq  \beta  \leq  \dfrac{\pi }{2}$ ,所以 $- \dfrac{\pi }{2} < \alpha  \leq  \dfrac{\pi }{2}, - \dfrac{\pi }{2} \leq   - \beta  < \dfrac{\pi }{2}$ ,
        
        $\therefore  - \pi  < \alpha  - \beta  < \pi$，又 $\alpha  - \beta  \leq  0$，$\therefore  - \pi  < \alpha  - \beta  \leq  0$ .
    \end{jieda}
    
    \item 日常生活中,在一杯含有 $a$ 克糖的 $b$ 克糖水中 $\left( {b > a > 0}\right)$ ,再加入 $m$ 克糖 $\left( {m > 0}\right)$ (假设全部溶解), 这杯糖水变甜了. 请根据这一事实提炼出一个不等式\blkda[6]{$b > a > 0,m > 0$ 时, $\dfrac{a}{b} < \dfrac{a + m}{b + m}$}.
    \begin{jieda}
        由题意有 $b > a > 0,m > 0$ 时, $\dfrac{a}{b} < \dfrac{a + m}{b + m}$ .
        
        理由如下: $\dfrac{a}{b} - \dfrac{a + m}{b + m} = \dfrac{a\left( {b + m}\right)  - b\left( {a + m}\right) }{b\left( {b + m}\right) } = \dfrac{m\left( {a - b}\right) }{b\left( {b + m}\right) }$ ,因为 $b > a > 0,m > 0$ ,
        
        所以 $\dfrac{a}{b} - \dfrac{a + m}{b + m} = \dfrac{m\left( {a - b}\right) }{b\left( {b + m}\right) } < 0$ ,即 $\dfrac{a}{b} < \dfrac{a + m}{b + m}$ .
    \end{jieda}
\end{enumerate}

\subsection*{B Level}
\begin{enumerate}
    \item 给出四个条件:
    \begin{dingautolist}{172}
        \item $a{c}^{2} > b{c}^{2}$
        \item $\dfrac{a}{c} > \dfrac{b}{c}$
        \item $a > \left| b\right|$
        \item $a > b - 1$
    \end{dingautolist}
    其中能成为 $a > b$ 的充分条件的是\blkda[4]{\ding{172}\ding{174}}.
    \begin{jieda}
        \begin{dingautolist}{172}
            \item 若 $a{c}^{2} > b{c}^{2}$ ,则 ${c}^{2} > 0$ ,故 $a > b$ ,即 $a{c}^{2} > b{c}^{2}$ 是 $a > b$ 的充分条件
            \item 当 $c < 0$ 时,由 $\dfrac{a}{c} > \dfrac{b}{c}$ 可得 $a < b$ ,故 $\dfrac{a}{c} > \dfrac{b}{c}$ 不是 $a > b$ 的充分条件
            \item 当 $b \geq  0$ 时, $a > \left| b\right|  = b$ ,当 $b < 0$ 时, $a > \left| b\right|  =  - b > b$ ,故 $a > \left| b\right|$ 是 $a > b$ 的充分条件
            \item 当 $a = {0.5},b = 1$ 时,满足 $a > b - 1$ ,但是不满足 $a > b$ ,故 $a > b - 1$ 不是 $a > b$ 的充分条件
        \end{dingautolist}
    \end{jieda}
    \item 已知 $- 1 \leq  x - y \leq  4, - 2 \leq  x + {6y} \leq  3$ ,则 $z = x - {8y}$ 的取值范围是\blkda[2]{$- 5 \leq  z \leq  {10}$}.
    \begin{jieda}
        设 $x - {8y} = m\left( {x - y}\right)  + n\left( {x + {6y}}\right)  = \left( {m + n}\right) x + \left( {{6n} - m}\right) y$ ，则 $\left\{  \begin{array}{l} m + n = 1 \\  {6n} - m =  - 8 \end{array}\right.$ ,故 $\left\{  \begin{array}{l} n =  - 1 \\  m = 2 \end{array}\right.$ ,
        
        因为 $- 1 \leq  x - y \leq  4, - 2 \leq  x + {6y} \leq  3$ ,则 $- 2 \leq  2\left( {x - y}\right)  \leq  8, - 3 \leq   - \left( {x + {6y}}\right)  \leq  2$ ,
        
        故 $- 5 \leq  2\left( {x - y}\right)  - \left( {x + {6y}}\right)  \leq  {10}$ 即 $- 5 \leq  z \leq  {10}$
    \end{jieda}
    \item 设 $x,y$ 为实数,满足 $3 \leq  x{y}^{2} \leq  8,4 \leq  \dfrac{{x}^{2}}{y} \leq  9$ ,则 $\dfrac{x}{{y}^{3}}$ 的最小值是\blkda[4]{$\dfrac{1}{2}$}.
    \begin{jieda}
        设 $\dfrac{x}{{y}^{3}} = {\left( x{y}^{2}\right) }^{m} \cdot  {\left( \dfrac{{x}^{2}}{y}\right) }^{n}$
        
        即 $x{y}^{-3} = {x}^{m + {2n}} \cdot  {y}^{{2m} - n}$
        
        所以 $\left\{  \begin{array}{l} m + {2n} = 1 \\  {2m} - n =  - 3 \end{array}\right.$ ,解得 $\left\{  \begin{array}{l} m =  - 1 \\  n = 1 \end{array}\right.$
        
        所以 $\dfrac{x}{{y}^{3}} = {\left( x{y}^{2}\right) }^{-1} \cdot  \left( \dfrac{{x}^{2}}{y}\right)$
        
        因为 $3 \leq  x{y}^{2} \leq  8,4 \leq  \dfrac{{x}^{2}}{y} \leq  9$ ,
        
        所以 $\dfrac{1}{8} \leq  {\left( x{y}^{2}\right) }^{-1} \leq  \dfrac{1}{3}$
        
        由不等式性质可知 $\dfrac{1}{2} \leq  {\left( x{y}^{2}\right) }^{-1} \cdot  \left( \dfrac{{x}^{2}}{y}\right)  \leq  3$
        
        即 $\dfrac{1}{2} \leq  \dfrac{x}{{y}^{3}} \leq  3$ ,当且仅当 $\left\{  \begin{array}{l} \dfrac{{x}^{2}}{y} = 4 \\  {\left( x{y}^{2}\right) }^{-1} = \dfrac{1}{8} \end{array}\right.$ 时取等号,解得 $x = {2}^{\frac{7}{5}},y = {2}^{\frac{4}{5}}$ .
        
        综上可知, $\dfrac{x}{{y}^{3}}$ 的最小值为 $\dfrac{1}{2}$ .
    \end{jieda}
    
    \item 若方程 ${x}^{2} - {8x} + m = 0$ 的两根为 ${x}_{1},{x}_{2}$ ,且 $3{x}_{1} + 2{x}_{2} = {18}$ ,则 $m =$ \blkda[4]{12}.
    \begin{jieda}
        由韦达定理得, ${x}_{1} + {x}_{2} = 8,{x}_{1} \cdot  {x}_{2} = m,\Delta  = {64} - {4m} > 0$ ，即 $\left\{  \begin{array}{l} {x}_{1} + {x}_{2} = 8 \\  3{x}_{1} + 2{x}_{2} = {18} \\  {x}_{1} \cdot  {x}_{2} = m \\  \Delta  = {64} - {4m} > 0 \end{array}\right.$ ,可解得 $\left\{  \begin{array}{l} {x}_{1} = 2 \\  {x}_{2} = 6 \\  m = {12} \end{array}\right.$
    \end{jieda}
    
    \item 已知关于 $x$ 的一元二次方程 ${x}^{2} - {6x} + \left( {{4m} + 1}\right)  = 0$ 的两个实数根为 ${x}_{1},{x}_{2}$ ,且 $\left| {{x}_{1} - {x}_{2}}\right|  = 4$ , 则实数 $m$ 的值为\blkda[4]{1}.
    \begin{jieda}
        $\because {x}_{1},{x}_{2}$ 为方程 ${x}^{2} - {6x} + \left( {{4m} + 1}\right)  = 0$ 的两个实数根,

         $\therefore {x}_{1} + {x}_{2} = 6,{x}_{1}{x}_{2} = {4m} + 1,\Delta  = {36} - 4\left( {{4m} + 1}\right)  \geq  0$ ,故 $m < 2$
        
        $\because \left| {{x}_{1} - {x}_{2}}\right|  = 4$ ,则 ${\left( {x}_{1} - {x}_{2}\right) }^{2} = {16},\therefore {\left( {x}_{1} + {x}_{2}\right) }^{2} - 4{x}_{1}{x}_{2} = {16}$ ,
        
         $\therefore {36} - 4\left( {{4m} + 1}\right)  = {16}$ ,解得 $m = 1$ .
        
        $\because m < 2,\therefore m = 1$ 符合题意.
    \end{jieda}
    
    \item 若 $a \neq  b$ ,且 ${a}^{2} - {4a} + 1 = 0$ , ${b}^{2} - {4b} + 1 = 0$ ,则 $\dfrac{1}{1 + {a}^{2}} + \dfrac{1}{1 + {b}^{2}}$ 的值为\blkda[4]{1}.
    \begin{jieda}
        由题意 $a,b$ 是方程 ${x}^{2} - {4x} + 1 = 0$ 的两个解,
        
        所以 $a + b = 4,{ab} = 1$ ,则 ${a}^{2} + {b}^{2} = {\left( a + b\right) }^{2} - {2ab} = {14}$ ,
        
        $\dfrac{1}{1 + {a}^{2}} + \dfrac{1}{1 + {b}^{2}} = \dfrac{2 + {a}^{2} + {b}^{2}}{\left( {1 + {a}^{2}}\right) \left( {1 + {b}^{2}}\right) } = \dfrac{2 + {a}^{2} + {b}^{2}}{1 + {a}^{2} + {b}^{2} + {a}^{2}{b}^{2}} = \dfrac{2 + {14}}{1 + {14} + 1} = 1$
    \end{jieda}

    
    \item 方程 ${x}^{2} - {ax} + a - 1 = 0$ 有一正一负根的充要条件是\blkda[4]{$a < 1$}
    \begin{jieda}
        ${x}^{2} - {ax} + a - 1 = 0$ 有一正一负根 $\Leftrightarrow  \left\{  {\begin{array}{l} a - 1 < 0 \\  \Delta  = {\left( -a\right) }^{2} - 4\left( {a - 1}\right)  > 0 \end{array} \Longrightarrow  a < 1}\right.$
    \end{jieda}
    \item 已知关于 $x$ 的方程 ${x}^{2} - {2x} + m - 1 = 0$ 的两个实数根同号,则实数 $m$ 的取值范围为\blkda[4]{$(1,2\rbrack$}.
    \begin{jieda}
        根据题意得到 $\left\{  \begin{array}{l} \Delta  \geq  0 \\  {x}_{1}{x}_{2} > 0 \end{array}\right.$ ,即 $\left\{  \begin{array}{l} 4 - 4\left( {m - 1}\right)  \geq  0 \\  m - 1 > 0 \end{array}\right.$ ,解得 $1 < m \leq  2$ .
    \end{jieda}
    
    \item 已知一元二次方程 $m{x}^{2} - {2x} + m + 3 = 0$ 有两个正实根,则实数 $m$ 的取值范围是\blkda[3]{$0 < m \leq  \dfrac{\sqrt{13} - 3}{2}$}. 
    \begin{jieda}
        设两个正实数根分别为 ${x}_{1},{x}_{2},\left\{  {\begin{array}{l} m \neq  0 \\  \Delta  = 4 - {4m}\left( {m + 3}\right)  \geq  0 \\  {x}_{1} + {x}_{2} = \dfrac{2}{m} > 0 \\  {x}_{1}{x}_{2} = \dfrac{m + 3}{m} > 0 \end{array}, \Longrightarrow  0 < m \leq  \dfrac{\sqrt{13} - 3}{2}}\right.$ .
    \end{jieda}

    \item 若关于 $x$ 的一元二次方程 ${x}^{2} + \left( {{3a} - 1}\right) x + a + 8 = 0$ 有两个不相等的实根 ${x}_{1},{x}_{2}$ ,且${x}_{1} < 1,{x}_{2} > 1$ . 则实数 $a$ 的取值范围为\blkda[4]{$(-\infty, -2)$}.
    \begin{jieda}
        令函数 $f\left( x\right)  = {x}^{2} + \left( {{3a} - 1}\right) x + a + 8$ ,依题意, $f\left( x\right)  = 0$ 的两个不等实根 ${x}_{1},{x}_{2}$ 满足 ${x}_{1} < 1,{x}_{2} > 1$ ,
        
        而函数 $f\left( x\right)$ 图象开口向上,因此 $f\left( 1\right)  < 0$ ,则 ${1}^{2} + \left( {{3a} - 1}\right)  \times  1 + a + 8 < 0$ ,解得 $a <  - 2$，所以实数 $a$ 的取值范围为 $a <  - 2$
    \end{jieda}
    \item 已知方程 ${x}^{2} - {2ax} + {a}^{2} - 4 = 0$ 的一个实根小于 2 ,另一个实根大于 2 ,求实数 $a$ 的取值范围\blkda[4]{$(0, 4)$}.
    \begin{jieda}
        设 $f\left( x\right)  = {x}^{2} - {2ax} + {a}^{2} - 4$ ,
        
        因为方程 ${x}^{2} - {2ax} + {a}^{2} - 4 = 0$ 的一个实根小于 2,另一个实根大于 2,
        
        则满足 $f\left( 2\right)  = {a}^{2} - {4a} < 0$ ,解得 $0 < a < 4$ ,即实数 $a$ 的取值范围为$(0,4)$.
    \end{jieda}
    \item 方程 ${x}^{2} - \left( {2 - a}\right) x + 5 - a = 0$ 的两根都大于 2 ,则实数 $a$ 的取值范围是\blkda[4]{$(- 5,-4]$}.
    \begin{jieda}
        由题意,方程 ${x}^{2} - \left( {2 - a}\right) x + 5 - a = 0$ 的两根都大于 2,
        
        令 $f\left( x\right)  = {x}^{2} - \left( {2 - a}\right) x + 5 - a$ ,
        
        可得 $\left\{  {\begin{matrix} \Delta  \geq  0 \\  f\left( 2\right)  > 0 \\  \dfrac{2 - a}{2} > 2 \end{matrix},\text{ 即 }\left\{  \begin{matrix} {a}^{2} \geq  {16} \\  a + 5 > 0 \\  2 - a > 4 \end{matrix}\right. }\right.$

        解得  $- 5 < a \leq   - 4$.
    \end{jieda}
    
\end{enumerate}
\subsection*{C Level}
\begin{enumerate}
    \item 已知函数 $f\left( x\right)  = {x}^{2} + {2mx} + {2m} + 3$ 在区间$(0,2)$内有零点,求实数 $m$ 的取值范围是\blkda[4]{$\left(-\dfrac{3}{2}, -1\right]$}
    \begin{jieda}
        \begin{dingautolist}{172}
            \item 分对称轴位置讨论 \\
            $f(x)$对称轴为$x = -m$，$f(0) = 2m+3$，$f(2) = 7+6m$，$\Delta = 4m^2 - 8m - 12$
            \begin{enumerate}[1.]
                \item 若$-m \leq 0$，则$f(x)$在$(0, 2)$单调增，故“$f(x)$在$(0, 2)$上有零点”即$f(0) < 0, f(2) > 0$，此种情况无解.
                \item $-m \geq 2$，则$f(x)$在$(0, 2)$单调减，故“$f(x)$在$(0, 2)$上有零点”即$f(0) > 0, f(2) < 0$，此种情况无解.
                \item $-m \in (0, 2)$，则$f(x)$在$(0, 2)$先减后增，“$f(x)$在$(0, 2)$上有零点”即$\Delta \geq 0$，且有$f(0) > 0$或$f(2) > 0$，此种情况解得$m \in \left(-\dfrac{3}{2}, -1\right]$
            \end{enumerate}
            综上，$m \in \left(-\dfrac{3}{2}, -1\right]$
            \item 参变分离 \\
            $f(x)$至少有一个零点在区间$(0, 2)$内，即$f(x) = 0$在$(0, 2)$内有根，即方程$m = \dfrac{-3-x^2}{2x+2}$在$(0, 2)$内有根，即方程$m = \dfrac{1}{2} \cdot \dfrac{-3-(t-1)^2}{t}$在$t \in (1, 3)$内有根，即方程$m = \dfrac{1}{2} \cdot \left[-(t+\dfrac{4}{t})+2\right]$在$t \in (1, 3)$内有根.\\
            因此$m$的取值范围即$y = \dfrac{1}{2} \cdot [-(t+\dfrac{4}{t})+2]$在$t \in (1, 3)$内的值域. \\
            $y = \dfrac{1}{2} \cdot \left[-(t+\dfrac{4}{t})+2\right]$在$(1, 2)$单调增，在$(2, 3)$单调减，因此值域为$\left(-\dfrac{3}{2}, -1\right]$，故$m \in \left(-\dfrac{3}{2}, -1\right]$
        \end{dingautolist}
    \end{jieda}
    \item 已知函数 $f\left( x\right)  = {x}^{2} - {2ax} + 5$ 在区间$(1,3)$上有零点,则实数 $a$ 的取值范围是\blkda[4]{$[\sqrt{5}, 3)$ }.
    \begin{jieda}
        \begin{dingautolist}{172}
            \item 分对称轴位置讨论

            对称轴为$x = a$，$f(1) = 6-2a$，$f(3) = 14-6a$，$\Delta = 4a^2-20$

            \begin{enumerate}[1.]
                \item 若$a \leq 1$，则$f(x)$在$(1, 3)$严增，故$f\left( x\right)$ 在区间$(1,3)$上有零点即$f(1) < 0, f(3) > 0$，此种情况无解.
                \item 若$a \geq 3$，则$f(x)$在$(1, 3)$严减，故$f\left( x\right)$ 在区间$(1,3)$上有零点即$f(1) > 0, f(3) < 0$，此种情况无解.
                \item $a \in (1, 3)$，则$f(x)$在$(1, 3)$先减后增，故$f\left( x\right)$ 在区间$(1,3)$上有零点即$\Delta \geq 0, f(1) > 0$或$f(3)> 0$，解得$a \in [\sqrt{5}, 3)$.
            \end{enumerate}
            
            \item 参变分离

            函数 $f\left( x\right)  = {x}^{2} - {2ax} + 5$ 在区间(1,3)上有零点,即 ${2a} = \dfrac{{x}^{2} + 5}{x}$ 在区间$(1,3)$上有解,即 ${2a} = x + \dfrac{5}{x}$ 在区间$(1,3)$上有解
            
            令 $g\left( x\right)  = x + \dfrac{5}{x},x \in  \left( {1,3}\right)$ ,易知 $g\left( x\right)$ 在区间 $\left( {1,\sqrt{5}}\right)$ 上单调递减,在区间 $\left( {\sqrt{5},3}\right)$ 上单调递增，所以 $g{\left( x\right) }_{\min } = g\left( \sqrt{5}\right)  = 2\sqrt{5},g{\left( x\right) }_{\max } = g\left( 1\right)  = 6$，$g\left( x\right)  \in  \left\lbrack  {2\sqrt{5},6}\right)$，故$\sqrt{5} \leq  a < 3$ .
        \end{dingautolist}
    \end{jieda}
    \item 已知方程 ${x^2} + {2mx} + {2m} + 1 = 0$ 的两不等实数根均在区间$(0,1)$内,则实数 $m$ 的取值范围为\blkda[3]{$\left( {-\dfrac{1}{2},1 - \sqrt{2}}\right)$}.
    \begin{jieda}
        \begin{dingautolist}{172}
            \item 通解法

            令 $f\left( x\right)  = {x}^{2} + {2mx} + {2m} + 1$，则函数图象开口向上,对称轴为直线 $x =  - m$ .
            
            则 $\left\{  \begin{array}{l} \Delta  > 0 \\  0 <  - m < 1 \\  f\left( 0\right)  > 0 \\  f\left( 1\right)  > 0 \end{array}\right.$ ,解得 $- \dfrac{1}{2} < m < 1 - \sqrt{2}$ .
            
            故 $m$ 的取值范围为 $\left( {-\dfrac{1}{2},1 - \sqrt{2}}\right)$ .
            \item 参变分离

            ${x^2} + {2mx} + {2m} + 1 = 0 \Longleftrightarrow m = -\dfrac{-x^2-1}{2x+2}$

            换元$t = x+1, t \in (1, 2)$，则$m = \dfrac{-t^2+2t-2}{2t}$在$t \in(1, 2)$有两根，即函数$y = m$与函数$y = -\dfrac{1}{2}\left(t + \dfrac{2}{t} - 2\right)$在$t \in (1, 2)$有两交点.

            $y = -\dfrac{1}{2}\left(t + \dfrac{2}{t} - 2\right)$在$(1, \sqrt{2})$严增，在$(\sqrt{2}, 2)$严减，故$m \in \left( {-\dfrac{1}{2},1 - \sqrt{2}}\right)$满足要求.
        \end{dingautolist}
    \end{jieda}
    \item 若二次方程 ${ax^2} - \left( {4 - {a}^{2}}\right) x + 2 = 0$ 在$(0,2)$内恰有一个实根,则实数 $a$ 的取值范围为\blkda[3]{$( - 3,0)\cup (0, 1\rbrack$}.
    \begin{jieda}
        因为是二次方程，所以$a \neq 0$，设$f(x) = {ax^2} - \left( {4 - {a}^{2}}\right) x + 2$，对称轴为$x = \dfrac{4-a^2}{a}$，$f(0) = 2$，$f(2) = 2a^2+4a-6$
        \begin{enumerate}[1.]
            \item “符合函数零点存在定理”，因为$f(0) > 0$，故此时应有$f(2) < 0$，即$a \in (-3, 0) \cup (0,1)$
            \item 右端点函数值为0，此时有$a \in \{-3, 1\}$，根据韦达定理另一根为$\dfrac{1}{a}$，故当$a = 1$时，$f(x)$在$(0, 2)$存在唯一零点.
        \end{enumerate}
        综上，$(-3,0)\cup (0, 1\rbrack$
        
        \ps 重根算作两根，因此无需验证$\Delta = 0$时，顶点是否在区间内.
    \end{jieda}

    \item 定义$max M$表示数集$M$中最大的数，则$max\left\{\dfrac{1}{a}, \dfrac{1}{b}, a^2+b^2\right\}$的最小值为\blkda[4]{$\sqrt[3]{2}$}
    \begin{jieda}
        考虑特例$\dfrac{1}{a} = \dfrac{1}{b} = a^2+b^2$，此时$a = b = \sqrt[3]{\dfrac{1}{2}}$，$max\left\{\dfrac{1}{a}, \dfrac{1}{b}, a^2+b^2\right\} = \sqrt[3]{2}$ 下面反证法证明$max\left\{\dfrac{1}{a}, \dfrac{1}{b}, a^2+b^2\right\}$取得最小值为$\sqrt[3]{2}$. \\
        假设$max\left\{\dfrac{1}{a}, \dfrac{1}{b}, a^2+b^2\right\}$的最小值为$m < \sqrt[3]{2}$
        \begin{dingautolist}{172}
            \item 设$\dfrac{1}{a} = m, \dfrac{1}{b} \leq m, a^2+b^2 \leq m$ \\
            此时$a = \dfrac{1}{m}, b \geq \dfrac{1}{m}$，故$a^2+b^2 \geq \dfrac{2}{m^2} > \sqrt[3]{2} > m$，矛盾.
            \item 设$\dfrac{1}{a} \leq m, \dfrac{1}{b} = m, a^2+b^2 \leq m$\\
            此时$a \geq \dfrac{1}{m}, b = \dfrac{1}{m}$，故$a^2+b^2 \geq \dfrac{2}{m^2} > \sqrt[3]{2} > m$，矛盾.
            \item 设$\dfrac{1}{a} \leq m, \dfrac{1}{b} \leq m, a^2+b^2 = m$ \\
            此时$a \geq \dfrac{1}{m}, b \geq \dfrac{1}{m}$，故$a^2+b^2 \geq \dfrac{2}{m^2} > \sqrt[3]{2} > m$，矛盾.
        \end{dingautolist}
    \end{jieda}
    \item 定义$max M$表示数集$M$中最大的数.设$0<a<b<c<1$，已知$b\geq 2a$且$a+b\leq 1$，则$max\left\{b-a, c-b, 1-c\right\}$的最小值为\blkda[4]{$\dfrac{1}{4}$}
    \begin{jieda}
        设$A = a, B = b-a, C = c - b, D = 1 - c$，则$A+B+C+D = 1$且有$A \leq B$或$2A+B \leq 1$，目标值为$max\left\{B, C, D\right\}$

        考虑特例$B = C = D$，则此时有$A+3B = 1$，故$A = 1-3B$
        
        因此$A \leq B$且$2A+B \leq 1$即$1-3B \leq B$且$2-5B \leq 1$，即$B \geq \dfrac{1}{4}$，故有特例$A = B = C = D = \dfrac{1}{4}$
        
        下面反证法证明此特例是使得$max\left\{B, C, D\right\}$取得最小值的例子.

        假设$max\left\{B, C, D\right\}$可取到最小值$m < \dfrac{1}{4}$
        \begin{dingautolist}{172}
            \item $B = m < \dfrac{1}{4}, C \leq m, D \leq m$

            此时$m = B \geq A = 1 - B - C - D \geq 1-3m$，即$4m \geq 1$，矛盾
            \item $B \leq m, C = m < \dfrac{1}{4}, D \leq m$

            此时$m \geq B \geq A = 1 - B - C - D \geq 1-3m$，即$4m \geq 1$，矛盾
            
            \item $B \leq m, C \leq m , D = m < \dfrac{1}{4}$
            与上一情况对称
        \end{dingautolist}
        因此$max\left\{b-a, c-b, 1-c\right\}$的最小值为$\dfrac{1}{4}$
    \end{jieda}
    \item 定义$max M$表示数集$M$中最大的数.设$0<a<b<c<1$，已知$b\geq 2a$或$a+b\leq 1$，则$max\left\{b-a, c-b, 1-c\right\}$的最小值为\blkda[4]{$\dfrac{1}{5}$}
    \begin{jieda}
        设$A = a, B = b-a, C = c - b, D = 1 - c$，则$A+B+C+D = 1$且有$A \leq B$或$2A+B \leq 1$，目标值为$max\left\{B, C, D\right\}$

        
        考虑特例$B = C = D$，则此时有$A+3B = 1$，故$A = 1-3B$
        
        因此$A \leq B$或$2A+B \leq 1$即$1-3B \leq B$或$2-5B \leq 1$，即$B \geq \dfrac{1}{4}$或$B \geq \dfrac{1}{5}$，故有特例$A = \dfrac{2}{5}, B = C = D = \dfrac{1}{5}$
        
        下面反证法证明此特例是使得$max\left\{B, C, D\right\}$取得最小值的例子.

        
        假设$max\left\{B, C, D\right\}$可取到最小值$m < \dfrac{1}{5}$
        
        
        \begin{dingautolist}{172}
            \item $B = m < \dfrac{1}{5}, C \leq m, D \leq m$

            此时$m = B \geq A = 1 - B - C - D \geq 1-3m$，即$4m \geq 1$，且有$2A + B \geq 2 - 5m > 1$，故$2A+B > 1$，矛盾.
            
            \item $B \leq m, C = m < \dfrac{1}{5}, D \leq m$

            此时$m \geq B \geq A = 1 - B - C - D \geq 1-3m$，即$4m \geq 1$，且有$2A + B \geq 2 - 5m > 1$，故$2A+B > 1$，矛盾.
            
            \item $B \leq m, C \leq m , D = m < \dfrac{1}{5}$
            与上一情况对称
        \end{dingautolist}
        
        因此$max\left\{b-a, c-b, 1-c\right\}$的最小值为$\dfrac{1}{5}$
    \end{jieda}
\end{enumerate}
\section*{解不等式（高次、分式、根式）、不等式解集与方程的根、整根问题}
\subsection*{A Level}
\begin{enumerate}
    \item $\left( {x - 2}\right) \left( {x + 3}\right)  > x - 2$ 的解集为\blkda[4]{$(-\infty, -2) \cup (2, +\infty)$}.
    \begin{jieda}
        $\left( {x - 2}\right) \left( {x + 3}\right)  > x - 2 \Longleftrightarrow\left( {x - 2}\right) \left( {x + 2}\right)  > 0$，所以 $x <  - 2$ 或 $x > 2$ .
    \end{jieda}
    \item 已知 $x \in  \mathbf{R}$ ,则不等式 ${x}^{2} - {2x} - 3 < 0$ 的解集为\blkda[4]{$(-1, 3)$}.
    \begin{jieda}
        方程 ${x}^{2} - {2x} - 3 = 0$ 的解为 $x =  - 1$ 或 $x = 3$

        故不等式 ${x}^{2} - {2x} - 3 < 0$ 的解集为 $\{ x \mid   - 1 < x < 3\}$
    \end{jieda}
    \item 不等式 $x + {x}^{3} \geq  0$ 的解集是\blkda[4]{$\lbrack 0, + \infty )$}.
    \begin{jieda}
        $x + {x}^{3} \geq  0 \Longleftrightarrow x\left( {1 + {x}^{2}}\right)  \geq  0 \Longleftrightarrow x \geq 0$
    \end{jieda}
    \item 不等式 $\left( {{x}^{2} - 4}\right) {\left( x - 6\right) }^{2} \leq  0$ 的解集是\blkda[4]{$[-2, 2] \cup \{6\}$}.
    \begin{jieda}
        $\left( {{x}^{2} - 4}\right) {\left( x - 6\right) }^{2} \leq  0 \Longleftrightarrow \left( {x - 2}\right) \left( {x + 2}\right) {\left( x - 6\right)  }^{2} \leq  0$

        穿针引线可得解集为$[-2, 2] \cup \{6\}$
    \end{jieda}
    \item 不等式 $\dfrac{x - 1}{x - 3} < 0$ 的解集为\blkda[4]{$(1,3)$}.
    \begin{jieda}
        $\dfrac{x - 1}{x - 3} < 0 \Longleftrightarrow \left( {x - 1}\right) \left( {x - 3}\right)  < 0$ ,解得 $1 < x < 3$
    \end{jieda}
    
    \item 已知集合 $A = \left\{  {x\left| {\dfrac{x}{x + 2} < 0}\right. }\right\}  ,B = \{  - 1,0,1\}$ ,则 $A \cap  B =$ \blkda[4]{$\{ - 1\}$}.
    \begin{jieda}
        $\dfrac{x}{x + 2} < 0 \Longleftrightarrow  x\left( {x + 2}\right)  < 0, \Longleftrightarrow - 2 < x < 0$ ,
        
        所以 $A = \left\{  {x\left| {\dfrac{x}{x + 2} < 0}\right. }\right\}   = \{ x \mid   - 2 < x < 0\}$ ,又 $B = \{  - 1,0,1\}$，所以 $A \cap  B = \{  - 1\}$
    \end{jieda}
    
    \item 不等式 $\dfrac{\left( {{2x} - 1}\right) {\left( x + 1\right) }^{2}}{{x}^{2} - 4} \geq  0$ 的解集为\blkda[4]{$\left( {-2,\dfrac{1}{2}}\right\rbrack   \cup  \left( {2, + \infty }\right)$}.
    \begin{jieda}
        $\dfrac{\left( {{2x} - 1}\right) {\left( x + 1\right) }^{2}}{{x}^{2} - 4} \geq  0 \Longleftrightarrow \left\{\begin{array}{l}
             {x}^{2} - 4 \neq  0  \\
             \left( {{x}^{2} - 4}\right) \left( {{2x} - 1}\right) {\left( x + 1\right) }^{2} \geq  0 
        \end{array}\right.$

        穿针引线可得解集为$\left( {-2,\dfrac{1}{2}}\right\rbrack   \cup  \left( {2, + \infty }\right)$
    \end{jieda}
    \item 不等式 $\dfrac{5 - x}{{x}^{2} - {2x} - 3} \leq   - 1$ 的解集为\blkda[4]{$\left( {-1,1\rbrack \cup \lbrack 2,3}\right)$}.
    \begin{jieda}
        $\dfrac{5 - x}{{x}^{2} - {2x} - 3} \leq   - 1 \Longleftrightarrow \dfrac{{x}^{2} - {3x} + 2}{{x}^{2} - {2x} - 3} \leq  0 \Longleftrightarrow\left\{  \begin{array}{l} \left( {x + 1}\right) \left( {x - 1}\right) \left( {x - 2}\right) \left( {x - 3}\right)  \leq  0 \\  x \neq   - 1\text{ 且 }x \neq  3 \end{array}\right.$

        穿针引线可得解集为$\left( {-1,1\rbrack \cup \lbrack 2,3}\right)$ 
    \end{jieda} 
    \item 已知函数 $y = f\left( x\right) \left( {x \in  R}\right)$ 的图象如图所示,则不等式 ${xf}\left( x\right)  < 0$ 的解集为\blkda[4]{ $\left( {-1,0}\right)  \cup  \left( {1,3}\right)$}.
    
    \begin{center}
    \includegraphics[width=0.4\textwidth]{bodyxb/src/bo_d2alapjef24c73avqjjg_2_241_1430_502_220_0.jpg}
    \end{center}
    \begin{jieda}
        不等式 ${xf}\left( x\right)  < 0$ 等价为 $\left\{  \begin{array}{l} x > 0 \\  f\left( x\right)  < 0 \end{array}\right.$ 或 $\left\{  \begin{array}{l} x < 0 \\  f\left( x\right)  > 0 \end{array}\right.$ 
        
        根据图可知 $1 < x < 3$ ,或 $- 1 < x < 0$，故不等式 ${xf}\left( x\right)  < 0$ 的解集是 $\left( {-1,0}\right)  \cup  \left( {1,3}\right)$ .
    \end{jieda}
    \item 不等式 $\sqrt{x + 3} < 2$ 的解是\blkda[4]{$\lbrack  - 3,1)$}.
    \begin{jieda}
        $\sqrt{x + 3} < 2 \Longleftrightarrow \left\{\begin{array}{l}
             x \geq -3  \\
             x+3 < 4
        \end{array}\right. \Longleftrightarrow -3 \leq  x < 1$
    \end{jieda}
    \item 不等式 $\dfrac{\sqrt{x + 1}}{x - 6} \leq  0$ 的解为\blkda[4]{$\lbrack  - 1,6)$}.
    \begin{jieda}
        $\dfrac{\sqrt{x + 1}}{x - 6} \leq  0 \Longleftrightarrow \left\{  \begin{array}{l} x + 1 \geq  0 \\  x - 6 < 0 \end{array}\right. \Longleftrightarrow - 1 \leq  x < 6$
    \end{jieda}
    \item 已知集合 $A = \left\{  {x\left| {{x}^{2} - x - 2 \leq  0}\right. }\right\}  ,B = \left\{  {x\left| {\sqrt{x - 1} \leq  2}\right. }\right\}$ ,则 $A \cap  B =$ \blkda[4]{$[1, 2]$}
    \begin{jieda}
        $B = \left\{  {x \mid  \sqrt{x - 1} \leq  2}\right\}   = \{ x \mid  0 \leq  x - 1 \leq  4\}  = \{ x \mid  1 \leq  x \leq  5\}$
        
        $A = \{ x \mid   - 1 \leq  x \leq  2\}$，$A \cap  B = \{ x \mid  1 \leq  x \leq  2\}$
    \end{jieda}
\end{enumerate}
\subsection*{B Level}
\begin{enumerate}
    \item 关于 $x$ 的不等式 $\dfrac{\left( {x + 3}\right) {\left( x - 2\right) }^{999}{\left( x + 1\right) }^{2}}{{\left( x - 1\right) }^{2015}{\left( x - 4\right) }^{2016}} \leq  0$ 的解集为\blkda[4]{$(- \infty , - 3 \rbrack  \cup  \{  - 1\}  \cup  (1,2\rbrack$}
    \begin{jieda}
        原不等式等价于: $\left\{  \begin{array}{l} \left( {x + 3}\right) {\left( x - 2\right) }^{999}{\left( x + 1\right) }^{2}{\left( x - 1\right) }^{2015}{\left( x - 4\right) }^{2016} \leq  0, \\  x - 1 \neq  0, \\  x - 4 \neq  0, \end{array}\right. $
        
        结合高次不等式奇穿偶不穿的性质求解不等式组可得原不等式的解集为: $(- \infty , - 3 \rbrack  \cup  \{  - 1\}  \cup  (1,2\rbrack$
    \end{jieda} 
    
    \item 不等式 $\sqrt{{4x} - {x}^{2}} < x-1$ 的解集是\blkda[4]{$\left(\dfrac{3+\sqrt{7}}{2}, 4\right]$}.
    \begin{jieda}
        $\sqrt{{4x} - {x}^{2}} < x-1 \Longleftrightarrow \left\{\begin{array}{l}
             {4x} - {x}^{2} \geq 0  \\
             x-1 > 0 \\
             {4x} - {x}^{2} < x^2-2x+1
        \end{array}\right. \\
        \Longleftrightarrow \left\{\begin{array}{l}
             0 \leq x \leq 4  \\
             x > 1 \\
             x \in \left(-\infty, \dfrac{3-\sqrt{7}}{2}\right) \cup \left(\dfrac{3+\sqrt{7}}{2}, +\infty \right) 
        \end{array}\right. \Longleftrightarrow x \in \left(\dfrac{3+\sqrt{7}}{2}, 4\right]$
    \end{jieda}
    \item 不等式 $\left( {x - 2}\right) \sqrt{x + 3} \geq  0$ 的解集为\blkda[4]{$\lbrack 2, + \infty ) \cup  \{  - 3\}$}.
    \begin{jieda}
        \begin{dingautolist}{172}
            \item $\sqrt{x+3} > 0$且$x-2 \geq 0$，解得$x \in [2, +\infty)$
            \item $\sqrt{x+3} = 0$且$x-2 \in \mathbf{R}$，解得$x = -3$
        \end{dingautolist}
    \end{jieda}
    \item 不等式 $\sqrt{4 - {x}^{2}} > x$ 的解集是\blkda[4]{$\lbrack  - 2,\sqrt{2})$}.
    \begin{jieda}
        $\sqrt{4 - {x}^{2}} > x \Longleftrightarrow \left\{\begin{array}{l}
             4-x^2 \geq 0 \\
             x \geq 0 \\
             4 - x^2 > x^2
        \end{array}\right.$或$\left\{\begin{array}{l}
             4-x^2 \geq 0 \\
             x < 0 
        \end{array}\right.$
    \end{jieda}
    
    \item 解关于 $x$ 的不等式: 
    \begin{enumerate}
        \item $\sqrt{x - 1} + {2x} \leq  5$ 
        \item $\dfrac{{ax} - 1}{x - 2} > \dfrac{a}{2}\left( {a \in  \mathbf{R}}\right)$
    \end{enumerate}
    \begin{jieda}
        \begin{enumerate}
            \item $\sqrt{x - 1} + {2x} \leq  5 \Longleftrightarrow \sqrt{x-1} \leq 5 - 2x \Longleftrightarrow \left \{ \begin{array}{l}
                 x \geq 1  \\
                 5-2x \geq 0 \\
                 x-1 \leq (5-2x)^2
            \end{array}\right.$，解得$x \in [1, 2]$
            \item $\dfrac{{ax} - 1}{x - 2} > \dfrac{a}{2} \Longleftrightarrow \dfrac{{ax} - 1}{x - 2} - \dfrac{a}{2} > 0 \Longleftrightarrow \dfrac{ax+2a-2}{2(x-2)} > 0 \Longleftrightarrow \left\{\begin{array}{l}
                 (ax+2a-2)(x-2) > 0 \\
                 x \neq 2
            \end{array}\right. \\
            \Longleftrightarrow (ax+2a-2)(x-2) > 0$
            \begin{dingautolist}{172}
                \item $a = 0$

                $(ax+2a-2)(x-2) > 0\Longleftrightarrow 
                 (-2)(x-2) > 0 \Longleftrightarrow x \in (-\infty, 2)$
            \item $a > 0$

            $(ax+2a-2)(x-2) > 0 \Longleftrightarrow \left(x+2-\dfrac{2}{a}\right)(x-2) > 0 $

            方程$\left(x+2-\dfrac{2}{a}\right)(x-2) = 0$的两根分别为$\dfrac{2}{a}-2$和$2$

            \begin{enumerate}[1.]
                \item $a > \dfrac{1}{2}$

                $\dfrac{2}{a}-2 < 2$，故原不等式解集为$\left(-\infty, \dfrac{2}{a}-2\right)\cup (2, +\infty)$

                
                \item $a = \dfrac{1}{2}$

                $\dfrac{2}{a}-2 = 2$，故原不等式解集为$\left(-\infty, 2\right)\cup (2, +\infty)$

                
                \item $0 < a < \dfrac{1}{2}$

                $\dfrac{2}{a}-2 > 2$，故原不等式解集为$(-\infty, 2) \cup \left(\dfrac{2}{a}-2, +\infty\right)$
            \end{enumerate}

            
            \item $a < 0$

            $ax+2a-2)(x-2) > 0 \Longleftrightarrow \left(x+2-\dfrac{2}{a}\right)(x-2) < 0$

            方程$\left(x+2-\dfrac{2}{a}\right)(x-2) = 0$的两根分别为$\dfrac{2}{a}-2$和$2$

            因为$a < 0$，所以$\dfrac{2}{a}-2 < 2$恒成立，故原不等式解集为$\left(\dfrac{2}{a}-2, 2\right)$
            
            \end{dingautolist}
        \end{enumerate}
    \end{jieda}
    \item 已知集合 $A = \left\{  {x\left| {\dfrac{{ax} - 2}{x - a} \leq  0}\right. }\right\}$ ,若 $2 \notin  A$ ,则实数 $a$ 的取值范围是\blkda[4]{$(1,2\rbrack$}.
    \begin{jieda}
        集合 $A = \left\{  {x\left| {\dfrac{{ax} - 2}{x - a} \leq  0}\right. }\right\}$ ,若 $2 \notin  A$ ,故 $\dfrac{{2a} - 2}{2 - a} > 0$ 或 $2 - a = 0$ ,解得 $1 < a \leq  2$ .
    \end{jieda}
    \item 已知 $\alpha  : \dfrac{2}{x + 1} > 1,\beta  : m \leq  x \leq  2$ ,若 $\alpha$ 是 $\beta$ 的充分条件,则实数 $m$ 的取值范围是\blkda[4]{$(-\infty, -1]$}.
    \begin{jieda}
        $\dfrac{2}{x + 1} > 1 \Longleftrightarrow \dfrac{1 - x}{x + 1} > 0\Longleftrightarrow \left( {1 - x}\right) \left( {1 + x}\right)  > 0$，解得 $- 1 < x < 1$，即 $\alpha  :  - 1 < x < 1$ 
        
        若 $\alpha$ 是 $\beta$ 的充分条件,则 $(-1, 1) \subseteq \{ x \mid  m \leq  x \leq  2\}$ ，可得 $m \in (-\infty, -1]$ .
    \end{jieda}
    
    \item 已知 “ $\left( {x - m}\right) \left( {x - m - 2}\right)  < 0$ ” 是 “ $\dfrac{{2x} + 1}{x - 1} < 0$ ” 的必要非充分条件,则实数 $m$ 的取值范围是\blkda[4]{$\left\lbrack  {-1, - \dfrac{1}{2}}\right\rbrack$}.
    \begin{jieda}
        $\left( {x - m}\right) \left( {x - m - 2}\right)  < 0$ 解得 $m < x < m + 2$，$\dfrac{{2x} + 1}{x - 1} < 0$ 解得 $- \dfrac{1}{2} < x < 1$ .
        
        由题意可得得 $\left(- \dfrac{1}{2}, 1\right) \subset (m, m + 2)$ 包含 $- \dfrac{1}{2} < x < 1$，故 $m \in  \left\lbrack  {-1, - \dfrac{1}{2}}\right\rbrack$
    \end{jieda}
    \item  已知集合 $A = \left\{  {x\left| {\dfrac{7}{x + 3} \geq  1}\right. }\right\}  ,B = \left\{  {x\left| {{x}^{2} - {3mx} + 2{m}^{2} + m - 1 < 0}\right. }\right\}$ ,若 $B \subset  A$ ,则实数 $m$ 的取值范围为\blkda[4]{$\left\lbrack  {-1,\dfrac{5}{2}}\right\rbrack$}.
    \begin{jieda}
        $\dfrac{7}{x + 3} \geq  1 \Longleftrightarrow \dfrac{7}{x + 3} - 1 \geq  0 \Longleftrightarrow \dfrac{4 - x}{x + 3} \geq  0 \Longleftrightarrow \left\{  \begin{array}{l} \left( {x + 3}\right) \left( {x - 4}\right)  \leq  0 \\  x + 3 \neq  0 \end{array}\right.$，故$A = ( - 3,4\rbrack$
        
        ${x}^{2} - {3mx} + 2{m}^{2} + m - 1 < 0 \Longleftrightarrow \left\lbrack  {x - \left( {m + 1}\right) }\right\rbrack  \left\lbrack  {x - \left( {{2m} - 1}\right) }\right\rbrack  < 0$

        \begin{dingautolist}{172}
            \item 当 $m + 1 = {2m} - 1$ ,即 $m = 2$ 时, $B = \varnothing$ ,满足 $B \subset  A$ ,符合题意
            \item 当 $m + 1 > {2m} - 1$ ,即 $m < 2$ 时, $B = \left( {{2m} - 1,m + 1}\right)$，由 $B \subset  A$ ,得 $\left\{  \begin{array}{l} {2m} - 1 \geq   - 3 \\  m + 1 \leq  4 \end{array}\right.$，解得 $- 1 \leq  m < 2$
            \item 当 $m + 1 < {2m} - 1$ ,即 $m > 2$ 时, $B = \left( {m + 1,{2m} - 1}\right)$，由 $B \subset  A$ ,得 $\left\{  \begin{array}{l} m + 1 \geq   - 3 \\  {2m} - 1 \leq  4 \end{array}\right.$ ,解得 $2 < m \leq  \dfrac{5}{2}$
        \end{dingautolist}
        综上，$m \in \left\lbrack  {-1,\dfrac{5}{2}}\right\rbrack$
    \end{jieda}
    \item ${ax} + b < 0$ 的解集为 $\left( {-\infty , - 1}\right)$ ,则 $\left( {{ax} - b}\right) \left( {x + 2}\right)  < 0$ 的解集为\blkda[4]{$(-2,1)$}.
    \begin{jieda}
         ${ax} + b < 0$ 的解集为 $\left( {-\infty , - 1}\right)$，即$\left\{\begin{array}{l}
              b = a  \\
              a > 0
         \end{array}\right.$

         故$\left( {{ax} - a}\right) \left( {x + 2}\right)  < 0 \Longrightarrow  \left( {x - 1}\right) \left( {x + 2}\right)  < 0$，解得$- 2 < x < 1$ 
    \end{jieda}
    \item 不等式 $a{x}^{2} + {bx} + 2 > 0$ 的解集是 $\left\{  {x\left| { - \dfrac{1}{2} < x < \dfrac{1}{3}}\right. }\right\}$ ,则 $a + b =$ \blkda[4]{$-14$}.
    \begin{jieda}
        不等式 $a{x}^{2} + {bx} + 2 > 0$ 的解集是 $\left\{  {x\left| { - \dfrac{1}{2} < x < \dfrac{1}{3}}\right. }\right\} \Longleftrightarrow \left\{  \begin{array}{l} a < 0 \\   - \dfrac{b}{a} =  - \dfrac{1}{6}, \\  \dfrac{2}{a} =  - \dfrac{1}{6} \end{array}\right.$ 
        
        解得$\left\{  {\begin{array}{l} a =  - {12} \\  b =  - 2 \end{array}}\right.$，故$a + b =  - {14}$
    \end{jieda}
    \item 若关于 $x$ 的不等式 $a{x}^{2} + {bx} + c < 0$ 的解集是 $\left( {-\infty , - 1}\right)  \cup  \left( {3, + \infty }\right)$ ,则关于 $x$ 的不等式$a{x}^{2} + {cx} + {2b} \geq  0$ 的解集是\blkda[4]{$\left\lbrack  {-1,4}\right\rbrack$}.
    \begin{jieda}
        不等式 $a{x}^{2} + {bx} + c < 0$ 的解集是 $\left( {-\infty , - 1}\right)  \cup  \left( {3, + \infty }\right) \Longleftrightarrow \left\{  \begin{array}{l}  
        a < 0 \\
        a-b+c = 0 \\  
        9+3b+c =  0 
        \end{array}\right.$
        
        解得 $\left\{\begin{array}{l}
             a < 0 \\
             b =  - {2a}  \\
             c =  - {3a} 
        \end{array}\right.$，故$a{x}^{2} + {cx} + {2b} \geq  0 \Longleftrightarrow a{x}^{2} - {3ax} - {4a} \geq  0$
        
        因 $a < 0$ ,故得 ${x}^{2} - {3x} - 4 \leq  0$ ,解得 $x \in \left\lbrack  {-1,4}\right\rbrack$
    \end{jieda}
    \item 关于 $x$ 的不等式 ${x}^{2} - \left( {1 + {2a}}\right) x + {2a} < 0$ 的解集中恰有 2 个整数,则实数 $a$ 的取值范围是\blkda[4]{$\left\lbrack  {-1, - \dfrac{1}{2}}\right)  \cup  \left( {\dfrac{3}{2},2}\right\rbrack$}.
    \begin{jieda}
        ${x}^{2} - \left( {1 + {2a}}\right) x + {2a} < 0 \Longleftrightarrow \left( {x - 1}\right) \left( {x - {2a}}\right)  < 0$ ,

        \begin{dingautolist}{172}
            \item 当 ${2a} = 1$ 时,不等式化为 ${\left( x - 1\right) }^{2} < 0$ ,此时无解
            \item 当 ${2a} > 1$ 时,解得 $1 < x < {2a}$ ,要使解集中恰有两个整数,则 $3 < {2a} \leq  4$ ,得 $\dfrac{3}{2} < a \leq  2$
            \item 当 ${2a} < 1$ 时,解得 ${2a} < x < 1$ ,要使解集中恰有两个整数,则 $- 2 \leq  {2a} <  - 1$ ,得 $- 1 \leq  a <  - \dfrac{1}{2}$
        \end{dingautolist}
        综上，$a \in \left\lbrack  {-1, - \dfrac{1}{2}}\right)  \cup  \left( {\dfrac{3}{2},2}\right\rbrack$ .
    \end{jieda}
    \item 已知关于 $x$ 的不等式组 $\left\{  \begin{array}{l} {x}^{2} - {2x} - 8 > 0, \\  2{x}^{2} + \left( {{2k} + 7}\right) x + {7k} < 0 \end{array}\right.$ 仅有一个整数解,则实数 $k$ 的取值范围是\blkda[4]{$ [-5, 3) \cup (4,5]$}.
    \begin{jieda}
        不等式 ${x}^{2} - {2x} - 8 > 0$ 的解集为 $\{ x \mid  x <  - 2$ 或 $x > 4\}$ .
        
        不等式 $2{x}^{2} + \left( {{2k} + 7}\right) x + {7k} < 0$ 化为 $\left( {x + k}\right) \left( {x + \dfrac{7}{2}}\right)  < 0$ .

        \begin{dingautolist}{172}
            \item 当 $k = \dfrac{7}{2}$ 时,不等式组的解集为 $\varnothing$ ,不符合题意.
            
            \item 当 $k > \dfrac{7}{2}$ 时,不等式 $2{x}^{2} + \left( {{2k} + 7}\right) x + {7k} < 0$ 的解集为 $\left\{x\left |  { - k < x <  - \dfrac{7}{2}}\right. \right\}$，由解集中的整数为$-4$,得 $- 5 \leq   - k <  - 4$ ,解得 $4 < k \leq  5$ .

            \item 当 $k < \dfrac{7}{2}$ 时,不等式 $2{x}^{2} + \left( {{2k} + 7}\right) x + {7k} < 0$ 的解集为 $\left\{  {x\left| { - \dfrac{7}{2} < x <  - k}\right. }\right\}$，由解集中的整数为$-3$,得 $- 3 <  - k \leq  5$ ,解得 $- 5 \leq  k < 3$ .
        \end{dingautolist}
        综上，$k \in [-5, 3) \cup (4,5]$.
    \end{jieda}
\end{enumerate}
\subsection*{C Level}
\begin{enumerate}
    \item 已知关于$x$的不等式组$\left \{ \begin{array}{l}
         \dfrac{x-a}{x-1+a} < 0\\
         x+3a>1 
    \end{array}\right.$的整数解恰有两个，则实数$a$的取值范围是\blkda[4]{$(1, 2]$}
    \begin{jieda}
        首先可以确定第一个不等式的解集为$(a, 1-a)$或$(1-a, a)$，因为$a$和$1-a$关于$\dfrac{1}{2}$对称，因此$\dfrac{1}{2}$必然在该不等式的解集中。而第二个不等式的解集为$(1-3a, +\infty)$
        \begin{dingautolist}{172}
            \item 若$a \in [0, 1]$，则$(a, 1-a)$或$(1-a, a)$没有整数解，故整个不等式组的解集中也没有整数解.
            \item 若$a < 0$，则$a < 0 < 1 < 1-a < 1-3a$，则整个不等式组解集为空集.
            \item 若$a > 1$，则$1-3a < 1-a < 0 < 1 < a$整个不等式组的解集为$(1-a, a)$，此时解集中必有两个整数解$0, 1$，故当$a \in (1, 2]$时，满足要求.
        \end{dingautolist}
        综上，$a \in (1, 2]$ \\
        \ps 本题画数轴可以更方便分析，本题判断出$a$和$1-a$关于$\dfrac{1}{2}$对称后可以猜测到两个整根是0和1.
    \end{jieda}
    \item 已知不等式组$ \left\{ \begin{array}{l}
         x^2-x+a-a^2 < 0  \\
         x + 2a > 1
    \end{array} \right.$的整数解恰好有两个，$a$的取值范围是\blkda[4]{$(1, 2]$}
    \begin{jieda}
        设函数$f(x) = x^2 - x$，则第一个不等式即$f(x) < f(a)$，根据函数的性质可知不等式即$\left|x-\dfrac{1}{2}\right| < \left|a - \dfrac{1}{2}\right|$，故$\dfrac{1}{2}-\left|a - \dfrac{1}{2}\right| < x < \dfrac{1}{2} + \left|a - \dfrac{1}{2}\right|$，故第一个不等式解集为$(1 - a, a)$或$(a, 1-a)$，第二个不等式解集为$(1-2a, +\infty)$
        \begin{dingautolist}{172}
            \item $a \in [0, 1]$，此时第一个不等式中不含有两个整数解，不符合要求.
            \item $a < 0$，第一个不等式解集为$(a, 1-a)$，而$1-a < 1-2a$，故不等式组解集为空集，不符合要求.
            \item $a > 1$，第一个不等式解集为$(1 - a, a)$，故不等式解集为$(1 - a, a)$，当$a \in (1, 2]$满足要求.
        \end{dingautolist}
    \end{jieda}
    \item 若关于$x$的不等式$k|x| > |x - 2|$恰有4个整数解，则实数$k$的取值范围是\blkda[4]{$\left(\dfrac{3}{5}, \dfrac{2}{3}\right]$}
        \begin{jieda}
            易知$x = 0$不是不等式的解，故$k|x| > |x - 2|$进行参变分离，同解变形$k > \left|\dfrac{x-2}{x}\right|$，即$k > \left|1 - \dfrac{2}{x}\right|$ \\
            由函数$f(x) = \left|1 - \dfrac{2}{x}\right|$的图像可知整根的集合为$\{2, 3, 4, 5\}$，故$f(5) < k \leq f(6)$，即$k \in \left(\dfrac{3}{5}, \dfrac{2}{3}\right]$
            \ps \textbf{解决整根问题的第一步往往是确认整根的集合}
        \end{jieda}
    
    \item 已知关于 $x$ 的不等式组: $\left\{  \begin{array}{l} 4{x}^{2} - {4ax} =  - {3x} - a + 1, \\  x > \dfrac{a - 1}{4} \end{array}\right.$ 有且只有一个实数解,则实数 $a$ 的取值范围是\blkda[4]{$\left( {-\infty ,1\rbrack \cup \lbrack 2, + \infty }\right)  \cup  \left\{  \dfrac{5}{4}\right\}$}
    \begin{jieda}
        由 $4{x}^{2} - {4ax} =  - {3x} - a + 1$ ,得 $4{x}^{2} - \left( {{4a} - 3}\right) x + \left( {a - 1}\right)  = 0$ ,进而得 $\lbrack x - (a -$ 1)] $\left( {{4x} - 1}\right)  = 0$ ,解得 $x = a - 1$ 或 $x = \dfrac{1}{4}$ .
        
        因为不等式组有且只有一个实数解，所以 
        
        $\left\{  \begin{array}{l} a - 1 > \dfrac{a - 1}{4}, \\  \dfrac{1}{4} \leq  \dfrac{a - 1}{4} \end{array}\right.$ 或 $\left\{  \begin{array}{l} a - 1 \leq  \dfrac{a - 1}{4}, \\  \dfrac{1}{4} > \dfrac{a - 1}{4} \end{array}\right.$ 或 $\left\{  \begin{array}{l} a - 1 = \dfrac{1}{4} \\  \dfrac{1}{4} > \dfrac{a - 1}{4}\end{array}\right.$
        
        解得 $a \in ( - \infty ,1\rbrack  \cup   \lbrack 2, + \infty ) \cup  \left\{  \dfrac{5}{4}\right\}$ .
    \end{jieda}
\end{enumerate}
\subsection*{D Level}
\begin{enumerate}
    \item 设关于 $x$ 的不等式 $\dfrac{{x}^{2} + \left( {2{a}^{2} + 2}\right) x - {a}^{2} + {4a} - 7}{{x}^{2} + \left( {{a}^{2} + {4a} - 5}\right) x - {a}^{2} + {4a} - 7} < 0$ 的解集是一些区间的并集,且这些区间的长度和 (规定: 区间$(a, b)$的长度为 $b - a$ ) 不小于 12,则实数 $a$ 的取值范围为 \dotfill(\kh{A})
    \xx{$a \leq   - 1$ 或 $a \geq  5$}{$a <  - 1$ 或 $a \geq  5$}{$a <  - 2$ 或 $a \geq  3$}{$a \leq   - 2$ 或 $a \geq  3$}
    \begin{jieda}
        因为$- {a}^{2} + {4a} - 7 < 0$，而分子分母的二次项系数均为1，因此分子分母对应的方程均有互异实根，设${x}^{2} + \left( {2{a}^{2} + 2}\right) x - {a}^{2} + {4a} - 7 = 0$的两根分别为$x_1, x_2$，${x}^{2} + \left( {{a}^{2} + {4a} - 5}\right) x - {a}^{2} + {4a} - 7 = 0$的两根分别为$x_3, x_4$，
        
        根据上述分析可知$x_1 < 0 < x_2, x_3 < 0 < x_4$，目标式为$|x_2-x_4| + |x_1-x_3|$

        $x_1 x_2 = x_3x_4 = - {a}^{2} + {4a} - 7 < 0$，因为乘积相同，故必有$x_1 < x_3, x_2 < x_4$或$x_1 > x_3, x_2 > x_4$，因此$x_2-x_4,x_1-x_3$正负性相同，故目标式可以化为$|(x_1+x_2)-(x_3+x_4)|$
        
        而根据$\left\{\begin{array}{l}
             x_1 + x_2 = \left( 2{a}^{2} + 2\right)  \\
             x_3 + x_4 = \left( {{a}^{2} + {4a} - 5}\right)
        \end{array}\right.$可知$(x_1+x_2) - (x_3+x_4) = a^2-4a+7 > 0$，故目标式为$(x_1+x_2)-(x_3+x_4) = a^2-4a+7 \geq 12$，解得$a \in (-\infty, -1] \cup [5, +\infty)$

    \end{jieda}
\end{enumerate}
\section*{基本不等式、 对勾函数}
\subsection*{A Level}
\begin{enumerate}
    \item 不等式 $a + \dfrac{1}{a} \geq  2$ 成立的充要条件是\blkda[4]{$a > 0$}.
    \begin{jieda}
        均值不等式 $a + \dfrac{1}{a} \geq  2$ 成立的条件是 $a > 0$ .
    \end{jieda}
    \item “ $a = 1$ ” 是 “ ${x}^{2} + \dfrac{a}{{x}^{2}} \geq  2$ ” 的\blkda[4]{充分不必要}条件. (填“充分非必要”、“必要非充分”、“充分必要”、“既不充分也不必要”)
    \begin{jieda}
        当 $a = 1$ 时, ${x}^{2} + \dfrac{a}{{x}^{2}} = {x}^{2} + \dfrac{1}{{x}^{2}} \geq  2\sqrt{{x}^{2} \cdot  \dfrac{1}{{x}^{2}}} = 2$，当且仅当 ${x}^{2} = \dfrac{1}{{x}^{2}} \Longrightarrow  x =  \pm  1$ 时,取等号,所以充分性成立
        
        由 ${x}^{2} + \dfrac{a}{{x}^{2}} \geq  2\sqrt{{x}^{2} \cdot  \dfrac{a}{{x}^{2}}} = 2\sqrt{a} \geq  2$ ，所以 $a \geq  1$ ,故必要性不成立,、
        
        所以 “ $a = 1$ ” 是 “ ${x}^{2} + \dfrac{a}{{x}^{2}} \geq  2$ ” 的充分不必要条件
    \end{jieda}
    \item 已知 $a,b \in  \mathbf{R}$ ,且 ${ab} > 0$ ,则下列不等式中,恒成立的序号是\blkda[4]{\ding{175}\ding{176}\ding{177}}.
    \begin{dingautolist}{172}
        \item ${a}^{2} + {b}^{2} > {2ab}$
        \item $a + b \geq  2\sqrt{ab}$
        \item $\dfrac{1}{a} + \dfrac{1}{b} > \dfrac{2}{\sqrt{ab}}$
        \item $\dfrac{b}{a} + \dfrac{a}{b} \geq  2$
        \item $\dfrac{{a}^{2} + {b}^{2}}{2} \geq  {ab}$
        \item $\dfrac{{\left( a + b\right) }^{2}}{4} \leq  \dfrac{{a}^{2} + {b}^{2}}{2}$
        \item $\dfrac{a + b}{2} \geq$  $\dfrac{ab}{a + b}$
    \end{dingautolist}
    \begin{jieda}
        易知因为 ${a}^{2} + {b}^{2} \geq  {2ab}$ 对于 $a,b \in  \mathbf{R}$ 恒成立,当且仅当 $a = b$ 时,等号成立,所以 \ding{172} 错误，\ding{176}\ding{177}正确;

        对于\ding{173},\ding{174},显然 $a < 0,b < 0$ 时,不等式均不成立,即\ding{173}和\ding{174}错误；

        对于\ding{175},因为 ${ab} > 0$ ,所以 $\dfrac{b}{a} > 0,\dfrac{a}{b} > 0$ ,由基本不等式可得 $\dfrac{b}{a} + \dfrac{a}{b} \geq  2\sqrt{\dfrac{b}{a} \cdot  \dfrac{a}{b}} = 2$ ,当且仅当 $a = b$ 成立即\ding{175}正确;

        对于\ding{178}，当$a=b=-1$时不等式不成立，故错误
    \end{jieda}
    \item  若 $0 < x < 1,0 < y < 1$ ,则 ${x}^{2} + {y}^{2}\text{ 、 }x + y\text{ 、 }{2xy}\text{ 、 }2\sqrt{xy}$ 中最大的一个是\blkda[4]{$x + y$}.
    \begin{jieda}
        $0 < x < 1,0 < y < 1$ ,由基本不等式得 ${x}^{2} + {y}^{2} \geq  {2xy};x + y \geq  2\sqrt{xy}$ ;
        
        又因为 $0 < x < 1,0 < y < 1$ ,
        
        所以 ${x}^{2} + {y}^{2} - \left( {x + y}\right)  = {x}^{2} - x + {y}^{2} - y = x\left( {x - 1}\right)  + y\left( {y - 1}\right)  < 0$ ,
        
        故 ${x}^{2} + {y}^{2} < x + y$ ,
        
        所以最大的一个是 $x + y$
    \end{jieda}
    \item 若 $x > 1$ ,则 $x + \dfrac{1}{x - 1}$ 的最小值是\blkda[4]{3}.
    \begin{jieda}
        $x + \dfrac{1}{x - 1} = x - 1 + \dfrac{1}{x - 1} + 1 \geq  2\sqrt{\left( {x - 1}\right)  \times  \dfrac{1}{x - 1}} + 1 = 3$ ,
        
        当且仅当 $x - 1 = \dfrac{1}{x - 1}$ 即 $x = 2$ 时取等号,
        
        $\therefore x = 2$ 时 $x + \dfrac{1}{x - 1}$ 取得最小值 3 .
    \end{jieda}
    \item 已知 $x < 0$ ,则函数 ${y} = {x} + \dfrac{1}{{x}}$ 的最大值是\blkda[4]{-2}.
    \begin{jieda}
        因为 $x < 0$ ,所以 $- x > 0, - \dfrac{1}{x} > 0$ ,
        
        故 $x + \dfrac{1}{x} =  - \left\lbrack  {\left( {-x}\right)  + \left( {-\dfrac{1}{x}}\right) }\right\rbrack   \leq   - 2\sqrt{\left( {-x}\right) \left( {-\dfrac{1}{x}}\right) } =  - 2$ ,
        
        当且仅当 $- x =  - \dfrac{1}{x}$ ,即 $x =  - 1$ 时,等号成立,
        
        所以 $x < 0$ 时, ${y} = {x} + \dfrac{1}{{x}}$ 的最大值是 -2 .
    \end{jieda}
    \item 已知 $a > 0$ ,则 $\dfrac{{a}^{2} + 4}{a}$ 的最小值为\blkda[4]{4}.
    \begin{jieda}
        $\dfrac{{a}^{2} + 4}{a} = a + \dfrac{4}{a} \geq  2\sqrt{a \cdot  \dfrac{4}{a}} = 4,a = 2$ 时取等号
    \end{jieda}
    \item 设 $x$ 、 $y$ 满足 $x + y = {10}$ ,且 $x$ 、 $y$ 都是正数,则 ${xy}$ 的最大值为\blkda[4]{25}.
    \begin{jieda}
        由于 $x$ 、 $y$ 都是正数,故 ${xy} \leq  {\left( \dfrac{x + y}{2}\right) }^{2} = {25}$ ,当且仅当 $x = y = 5$ 时等号成立，故 ${xy}$ 的最大值为 25
    \end{jieda}
    \item 已知正实数 $x,y$ 满足: $x + \dfrac{1}{y} = 1$ ,则 $\dfrac{x}{y}$ 的最大值为\blkda[4]{$\dfrac{1}{4}$}.
    \begin{jieda}
        $\dfrac{x}{y} = x \times  \dfrac{1}{y} \leq  \dfrac{1}{4}{\left( x + \dfrac{1}{y}\right) }^{2} = \dfrac{1}{4}$，当且仅当 $x = \dfrac{1}{y} = \dfrac{1}{2}$，即 $x = \dfrac{1}{2},y = 2$ 时取得等号
        
        故 $\dfrac{x}{y}$ 的最大值为 $\dfrac{1}{4}$
    \end{jieda}
    \item 已知正实数 $a,b$ 满足 ${ab} = 8$ ,则 $a + {2b}$ 的最小值为\blkda[4]{8}.
    \begin{jieda}
        因为 $a,b > 0,{ab} = 8$，故$b = \dfrac{8}{a}$ ，所以 $a + {2b} = a + 2 \cdot  \dfrac{8}{a} \geq  2\sqrt{a \cdot  \dfrac{16}{a}} = 8$ ,
        
        当且仅当 $a = \dfrac{16}{a}$ ,即 $a = 4,b = 2$ 时取等号，所以 $a + {2b}$ 的最小值为 8
    \end{jieda}
    
    \item 若 $a > 0,b > 0$ ,且 $a + b = 4$ ,则下列不等式恒成立的是\blkda[4]{\ding{175}}(填序号).
    \begin{dingautolist}{172}
        \item $\dfrac{1}{ab} \leq  \dfrac{1}{4}$
        \item $\dfrac{1}{a} + \dfrac{1}{b} \leq  1$
        \item $\sqrt{ab} \geq  2$
        \item ${a}^{2} + {b}^{2} \geq  8$
    \end{dingautolist}
    \begin{jieda}
        因为 $4 = a + b \geq  2\sqrt{ab}$ (当且仅当 $a = b$ 时,等号成立)，即 $\sqrt{ab} \leq  2,{ab} \leq  4,\dfrac{1}{ab} \geq  \dfrac{1}{4}$ ,故\ding{172}\ding{174}不成立;
        
        $\dfrac{1}{a} + \dfrac{1}{b} = \dfrac{a + b}{ab} = \dfrac{4}{ab} \geq  1$ ,故\ding{173}不成立;
        
        ${a}^{2} + {b}^{2} = {\left( a + b\right) }^{2} - {2ab} = {16} - {2ab} \geq  8$ ,故\ding{175}成立.
    \end{jieda}
    
    \item 已知全集 $U = R$ ,实数 $a,b$ 满足 $a > b > 0$ ,集合 $M = \left\{  {x \mid  \sqrt{ab} < x < a}\right\}  ,N = \{ x \mid  b < x < \dfrac{a + b}{2}\}$ ,则 $\overline{M} \cap  N =$ \blkda[4]{$\left\{  {x\left| {b < x \leq  \sqrt{ab}}\right. }\right\}$}
    \begin{jieda}
        由于 $\overline{M} = \left\{  {x \mid  x \leq  \sqrt{ab}\text{ 或 }x \geq  a}\right\}$ ,又 $a > b > 0,b < \sqrt{ab} < \dfrac{a + b}{2} < a$ 则
        
        $\overline{M} \cap  N = \left\{  {x \mid  b < x \leq  \sqrt{ab}}\right\}$        
    \end{jieda}
\end{enumerate}
\subsection*{B Level}
\begin{enumerate}
    \item 函数 $y = \dfrac{x + 1}{\sqrt{x} - 1}\left( {x > 1}\right)$ 的最小值为\blkda[4]{$2\sqrt{2} + 2$}
    \begin{jieda}
        $y = \dfrac{x + 1}{\sqrt{x} - 1} = \dfrac{x - 1}{\sqrt{x} - 1} + \dfrac{2}{\sqrt{x} - 1} = \sqrt{x} - 1 + \dfrac{2}{\sqrt{x} - 1} + 2 \geq  2\sqrt{\left( {\sqrt{x} - 1}\right)  \cdot  \dfrac{2}{\sqrt{x} - 1}} + 2 = 2\sqrt{2} + 2$，当且仅当 $\sqrt{x} - 1 = \dfrac{2}{\sqrt{x} - 1}$ ,即 $x = 3 + 2\sqrt{2}$ 时,等号成立.
    \end{jieda}
    \item 下列各式中, 最小值为4的是\dotfill(\kh{B})
    \xx
    {$x+2\sqrt{x}+5$}
    {$\dfrac{x^2+5}{\sqrt{x^2+1}}$}
    {$x^2+3+\dfrac{4}{x^2+3}$}
    {$x+\dfrac{4}{x}$}
    \begin{jieda}
        注意取等条件是否可以取到！
    \end{jieda}
    \item 函数 $y = \dfrac{2{x}^{2} - {3x} + 5}{x - 1}$ 在 $x \in  \left\lbrack  {\dfrac{3}{2},3}\right\rbrack$ 上的最大值和最小值的乘积为\blkda[4]{${40}\sqrt{2} + {10}$}
    \begin{jieda}
        令 $t = x - 1,x = t + 1,\because x \in  \left\lbrack  {\dfrac{3}{2},3}\right\rbrack  ,\therefore t \in  \left\lbrack  {\dfrac{1}{2},2}\right\rbrack$ ,
        
        $\therefore y = \dfrac{2{\left( t + 1\right) }^{2} - 3\left( {t + 1}\right)  + 5}{t} = \dfrac{2{t}^{2} + t + 4}{t} = 2\left( {t + \dfrac{2}{t}}\right)  + 1$ ,
        
        令 $g\left( t\right)  = 2\left( {t + \dfrac{2}{t}}\right)  + 1,t \in  \left\lbrack  {\dfrac{1}{2},2}\right\rbrack$ ,
        
        由对勾函数的性质可知,函数 $g\left( t\right)$ 在 $\left\lbrack  {\dfrac{1}{2},\sqrt{2}}\right\rbrack$ 上为减函数,在 $\left\lbrack  {\sqrt{2},2}\right\rbrack$ 上为增函数,
        
        $g\left( \dfrac{1}{2}\right)  = {10},g\left( \sqrt{2}\right)  = 4\sqrt{2} + 1,g\left( 2\right)  = 7$，故$g{\left( t\right) }_{\max } = {10},g{\left( t\right) }_{\min } = 4\sqrt{2} + 1$
        
        因此函数 $y = \dfrac{2{x}^{2} - {3x} + 5}{x - 1}$ 在 $x \in  \left\lbrack  {\dfrac{3}{2},3}\right\rbrack$ 上的最大值和最小值的乘积为 ${40}\sqrt{2} + {10}$ .
    \end{jieda}
    \item 设 $a,b > 0,a + \dfrac{1}{b} = 1$ ,则 $b + \dfrac{1}{a}$ 的最小值为\blkda[4]{4}.
    \begin{jieda}
        易知 $b + \dfrac{1}{a} = \left( {b + \dfrac{1}{a}}\right) \left( {a + \dfrac{1}{b}}\right)  = {ab} + \dfrac{1}{ab} + 2 \geq  2\sqrt{{ab} \cdot  \dfrac{1}{ab}} + 2 = 4$ ,
        
        当且仅当 ${ab} = 1$ ,即 $a = \dfrac{1}{2},b = 2$ 时取得最小值.
    \end{jieda}
    \item 已知实数 $a,b$ 满足 $- 1 < a < 1 < b$ ,且 $a + b = 2$ ,则 $\dfrac{1}{a + 1} + \dfrac{3a}{b - 1}$ 的最小值为\blkda[4]{$\sqrt{3} - 1$}.
    \begin{jieda}
    
        由 $a + b = 2$ ,所以 $\dfrac{1}{a + 1} + \dfrac{3a}{b - 1} = \dfrac{1}{a + 1} + \dfrac{3\left( {2 - b}\right) }{b - 1} = \dfrac{1}{a + 1} + \dfrac{3}{b - 1} - 3$ .

        $\dfrac{1}{a + 1} + \dfrac{3}{b - 1} - 3 = \dfrac{1}{2}\left\lbrack  {\dfrac{1}{a + 1} + \dfrac{3}{b - 1}}\right\rbrack  \left\lbrack  {\left( {a + 1}\right)  + \left( {b - 1}\right) }\right\rbrack   - 3 = \dfrac{1}{2}\left( {1 + 3 + \dfrac{b - 1}{a + 1} + \dfrac{3\left( {a + 1}\right) }{b - 1}}\right)  - 3 \geq  \dfrac{1}{2}\left( {4 + 2\sqrt{\dfrac{b - 1}{a + 1} \cdot  \dfrac{3\left( {a + 1}\right) }{b - 1}}}\right)  - 3 = \sqrt{3} - 1$
        
        当且仅当 $\dfrac{b - 1}{a + 1} = \dfrac{3\left( {a + 1}\right) }{b - 1}$ ,即 $a =  - 2 + \sqrt{3},b = 4 - \sqrt{3}$ 时,等号成立.
    \end{jieda}
    \item 已知 $a > 0,b > 0$ ,且 $\dfrac{1}{a + 1} + \dfrac{2}{b + 1} = 1$ ,则 $a + b$ 的最小值为\blkda[4]{$2\sqrt{2} + 1$}.
    \begin{jieda}
        由 $a > 0,b > 0,\dfrac{1}{a + 1} + \dfrac{2}{b + 1} = 1$，得 $a + b = \left( {a + 1}\right)  + \left( {b + 1}\right)  - 2 = \left( {\dfrac{1}{a + 1} + \dfrac{2}{b + 1}}\right) \left\lbrack  {\left( {a + 1}\right)  + \left( {b + 1}\right) }\right\rbrack   - 2 = \dfrac{b + 1}{a + 1} + \dfrac{2\left( {a + 1}\right) }{b + 1} + 1 \geq  2\sqrt{\dfrac{b + 1}{a + 1} \cdot  \dfrac{2\left( {a + 1}\right) }{b + 1}} + 1 = 2\sqrt{2} + 1$
        
        当且仅当 $\dfrac{b + 1}{a + 1} = \dfrac{2\left( {a + 1}\right) }{b + 1}$ ,即 $a = \sqrt{2},b = \sqrt{2} + 1$ 时, $a + b$ 取得最小值 $2\sqrt{2} + 1$ .
    \end{jieda}
    \item 已知 $a > 0,b > 0$ ,且 ${9a} + b = {ab}$ ,则 $a + {4b}$ 的最小值为\blkda[4]{49}.
    \begin{jieda}
        因为 $a > 0,b > 0$ 且 ${9a} + b = {ab}$ ,所以 $\dfrac{1}{a} + \dfrac{9}{b} = 1$ ,
        
        所以 $a + {4b} = \left( {a + {4b}}\right) \left( {\dfrac{1}{a} + \dfrac{9}{b}}\right)  = 1 + {36} + \dfrac{4b}{a} + \dfrac{9a}{b} \geq  {37} + 2\sqrt{\dfrac{4b}{a} \times  \dfrac{9a}{b}} = {37} + {12} = {49}$ ,
        
        当且仅当 $\dfrac{4b}{a} = \dfrac{9a}{b}$ 即 $a = 7,b = \dfrac{21}{2}$ 时取等号，所以 $a + {4b}$ 最小值为 49 .
    \end{jieda}

    \item 已知正数 $x$ 、 $y$ 满足 $x + {2y} = 1$ ,则 $\dfrac{2}{{3x} + {4y}} + \dfrac{1}{x + {3y}}$ 的最小值为\blkda[4]{$\dfrac{8}{5}$}.
    \begin{jieda}
        设$k(x + {2y}) = a(3x+4y)+b(x+3y)$，即$\left\{\begin{array}{l}
             3a+b = k  \\
             4a+3b = 2k
        \end{array}\right.$，$\left\{\begin{array}{l}
             a = 1 \\
             b = 2 \\
             k = 5
        \end{array}\right.$是一组解，故$\dfrac{2}{{3x} + {4y}} + \dfrac{1}{x + {3y}} = \dfrac{2}{{3x} + {4y}} + \dfrac{2}{2x + {6y}} = \dfrac{1}{5}\left(\dfrac{2}{{3x} + {4y}} + \dfrac{2}{2x + {6y}}\right) \left((3x+4y)+2(x+3y)\right) = \dfrac{2}{5}\left(1+1+\dfrac{{3x} + {4y}}{2x + {6y}}+\dfrac{2x + {6y}}{{3x} + {4y}}\right) \geq \dfrac{2}{5}(2+2) = \dfrac{8}{5}$ .
    \end{jieda}

    \item 若 $a$ 、 $b$ 为正实数,且 $a + {4b} + 3 = {ab}$，则 ${ab}$ 的最小值为\blkda[4]{${11} + 4\sqrt{7}$}
    \begin{jieda}
        $a + {4b} + 3 = {ab} \Longleftrightarrow ab = \dfrac{4b+3}{b-1} \cdot b$

        令$t = b-1, t > 0$，则$ab = \dfrac{(4t+7)(t+1)}{t} = 4t + \dfrac{7}{t} + 11 \geq 11+4\sqrt{7}$，当且仅当$t = \dfrac{\sqrt{7}}{2}$时取等号
    \end{jieda}
    \item 已知正实数 $x,y$ 满足 ${xy} - x - y = 3$ ,则 ${2x} + y$ 的最小值为\blkda[4]{$4\sqrt{2} + 3$}.
    \begin{jieda}
        正实数 $x,y$ 满足 ${xy} - x - y = 3$ ,故 $\left( {x - 1}\right) y = 3 + x$ ,所以 $y = \dfrac{3 + x}{x - 1}$ ,
        
        则 $\dfrac{3 + x}{x - 1} > 0$ ,又 $x > 0$ ,解得 $x > 1$ ,
        
        故 ${2x} + y = {2x} + \dfrac{3 + x}{x - 1} = {2x} + 1 + \dfrac{4}{x - 1} = 2\left( {x - 1}\right)  + \dfrac{4}{x - 1} + 3 \geq  2\sqrt{2\left( {x - 1}\right)  \cdot  \dfrac{4}{x - 1}} + 3 = 4\sqrt{2} + 3,$
        
        当且仅当 $2\left( {x - 1}\right)  = \dfrac{4}{x - 1}$ ,即 $x = 1 + \sqrt{2}$ 时,等号成立，故 ${2x} + y$ 的最小值为 $4\sqrt{2} + 3$
    \end{jieda} 
    \item 已知 $a > 0,b > 0$ ,且 ${ab} = a - b + 3$ ,则 $a + b$ 的最小值为\blkda[4]{$2\sqrt{2}$}.
    \begin{jieda}
        因为 ${ab} = a - b + 3$ ,解得: $b = \dfrac{a + 3}{a + 1} = 1 + \dfrac{2}{a + 1}$，则 $a + b = a + 1 + \dfrac{2}{a + 1} \geq  2\sqrt{2}$
        
        当且仅当 $a = \sqrt{2} - 1,b = \sqrt{2} + 1$ 时，等号成立
    \end{jieda}
    \item 设 $a > 0,b > 0$ ,且 ${5ab} + {b}^{2} = 1$ ,则 $a + b$ 的最小值为\blkda[4]{$\dfrac{4}{5}$}.
    \begin{jieda}
        因为 ${5ab} + {b}^{2} = 1$ ,所以 $a = \dfrac{1 - {b}^{2}}{5b} = \dfrac{1}{5b} - \dfrac{b}{5}$ ,

        所以 $a + b = \dfrac{1}{5b} - \dfrac{b}{5} + b = \dfrac{1}{5b} + \dfrac{4b}{5} \geq  2\sqrt{\dfrac{1}{5b} \cdot  \dfrac{4b}{5}} = \dfrac{4}{5}$ ,当且仅当 $a = \dfrac{3}{10},b = \dfrac{1}{2}$ 时,等号成立，所以 $a + b$ 的最小值为 $\dfrac{4}{5}$ .
    \end{jieda} 
\end{enumerate}
\subsection*{C Level}
\begin{enumerate}
    \item 若正实数 $x\text{ 、 }y$ 满足 ${2x} + y = 2$ ,则 $\dfrac{4{x}^{2}}{y + 1} + \dfrac{{y}^{2}}{{2x} + 2}$ 的最小值是\blkda[4]{$\dfrac{4}{5}$}.
    \begin{jieda}
        根据题意,若 ${2x} + y = 2$ ,则 $\dfrac{4{x}^{2}}{y + 1} + \dfrac{{y}^{2}}{{2x} + 2} = \dfrac{{\left( 2 - y\right) }^{2}}{y + 1} + \dfrac{{\left( 2 - 2x\right) }^{2}}{2\left( {x + 1}\right) } =$  $\dfrac{{\left\lbrack  \left( y + 1\right)  - 3\right\rbrack  }^{2}}{y + 1} + 2\dfrac{{\left\lbrack  \left( x + 1\right)  - 2\right\rbrack  }^{2}}{x + 1} = \left( {y + 1}\right)  + \dfrac{9}{y + 1} + 2\left( {x + 1}\right)  + \dfrac{16}{2\left( {x + 1}\right) } - {14} = \dfrac{9}{y + 1} +$  $\dfrac{16}{2\left( {x + 1}\right) } - 9$
        
         $\dfrac{4{x}^{2}}{y + 1} + \dfrac{{y}^{2}}{{2x} + 2} =$  $\left\lbrack  {\dfrac{9}{y + 1} + \dfrac{16}{2\left( {x + 1}\right) }}\right\rbrack  \dfrac{\left( {{2x} + 2}\right)  + \left( {y + 1}\right) }{5} - 9 \\
         = \dfrac{1}{5}\left\lbrack  {{16} + 9 + \dfrac{{18}\left( {x + 1}\right) }{y + 1} + \dfrac{8\left( {y + 1}\right) }{x + 1}}\right\rbrack   - 9 \geq \dfrac{1}{5}\left\lbrack  {{25} + 2\sqrt{\dfrac{{18}\left( {x + 1}\right) }{y + 1} \times  \dfrac{8\left( {y + 1}\right) }{x + 1}}}\right\rbrack   - 9 = \dfrac{4}{5}$ ,当且仅当 $y + 1 = \dfrac{3}{2}\left( {x + 1}\right)  = \dfrac{15}{7}$ 时等号成立, 即 $\dfrac{4{x}^{2}}{y + 1} + \dfrac{{y}^{2}}{{2x} + 2}$ 的最小值是 $\dfrac{4}{5}$
    \end{jieda}
    \item 已知关于 $x$ 的不等式 $- {x}^{2} + {6ax} - 3{a}^{2} \geq  0\left( {a > 0}\right)$ 的解集为 $\left\lbrack  {{x}_{1},{x}_{2}}\right\rbrack$ ,则 ${x}_{1} + {x}_{2} + \dfrac{3a}{{x}_{1}{x}_{2}}$ 的最小值是\blkda{$2\sqrt{6}$}
    \begin{jieda}
         不等式 $- {x}^{2} + {6ax} - 3{a}^{2} \geq  0\left( {a > 0}\right)$ 的解集为 $\left\lbrack  {{x}_{1},{x}_{2}}\right\rbrack$ ,所以 ${x}_{1}\text{、}{x}_{2}$ 是方程 ${x}^{2} - {6ax} + 3{a}^{2} = 0$ 的实数根，所以 $\left\{  \begin{array}{l} {x}_{1} + {x}_{2} = {6a} \\  {x}_{1}{x}_{2} = 3{a}^{2}\end{array}\right.$
         
         ${x}_{1} + {x}_{2} + \dfrac{3a}{{x}_{1}{x}_{2}} = {6a} + \dfrac{1}{a} \geq  2\sqrt{{6a} \cdot  \dfrac{1}{a}} = 2\sqrt{6}\text{,当且仅当}{6a} =  $
        $\dfrac{1}{a}$ ,即 $a = \dfrac{\sqrt{6}}{6}$ 时取等号，所以 ${x}_{1} + {x}_{2} + \dfrac{3a}{{x}_{1}{x}_{2}}$ 的最小值为 $2\sqrt{6}$
    \end{jieda}
    \item 若 $x \leq   - 1$ 时,不等式 $x + \dfrac{9}{x} - 2 < {k}^{2} - {9k}$ 恒成立,则 $k$ 的取值范围是\blkda[4]{$\left( {-\infty ,1}\right)  \cup  \left( {8, + \infty }\right)$}
    \begin{jieda}
         由于 ${k}^{2} - {9k} > {\left( x + \dfrac{9}{x} - 2\right) }_{\max } =  - 8$ ,所以 $k \in  \left( {-\infty ,1}\right)  \cup  \left( {8, + \infty }\right)$
    \end{jieda}
    \item 当 $0 < m < \dfrac{1}{2}$ 时,若 $\dfrac{1}{m} + \dfrac{2}{1 - {2m}} \geq  {k}^{2} - {2k}$ 恒成立,则实数 $k$ 的取值范围为\blkda[4]{$[-2,4]$}.
    \begin{jieda}
        由于 $0 < m < \dfrac{1}{2}$ ,所以 $\dfrac{1}{2} \cdot  {2m} \cdot  \left( {1 - {2m}}\right)  \leq  \dfrac{1}{2} \cdot  {\left\lbrack  \dfrac{{2m} + \left( {1 - {2m}}\right) }{2}\right\rbrack  }^{2} = \dfrac{1}{8}$ ,当且仅当 ${2m} = 1 - {2m}$ ,即 $m = \dfrac{1}{4}$ 时取等号,所以 $\dfrac{1}{m} + \dfrac{2}{1 - {2m}} = \dfrac{1}{m\left( {1 - {2m}}\right) } \geq  8$ ,因为 $\dfrac{1}{m} +$  $\dfrac{2}{1 - {2m}} \geq  {k}^{2} - {2k}$ 恒成立,所以 ${k}^{2} - {2k} - 8 \leq  0$ ,所以 $- 2 \leq  k \leq  4$
    \end{jieda}
    \item 若 $\forall x > a$ ,关于 $x$ 的不等式 ${2x} + \dfrac{2}{x - a} \geq  5$ 恒成立,则实数 $a$ 的取值范围是\blkda[4]{$\left\lbrack  {\dfrac{1}{2}, + \infty }\right)$}.
    \begin{jieda}
        若关于 $x$ 的不等式 ${2x} + \dfrac{2}{x - a} \geq  5$ 恒成立，则 ${\left( 2x + \dfrac{2}{x - a}\right) }_{\min } \geq  5$
        
        因为 $x > a$ ,故 ${2x} + \dfrac{2}{x - a} = 2\left( {x - a}\right)  + \dfrac{2}{x - a} + {2a} \geq  2\sqrt{2\left( {x - a}\right)  \times  \dfrac{2}{x - a}} + {2a} = 4 + {2a}$ ，当且仅当 $x = a + 1$ 时取等,故得 $4 + {2a} \geq  5$ ,解得 $a \geq  \dfrac{1}{2}$ .
    \end{jieda}
    \item 已知 $x > 0,y > 0$ 且 $\dfrac{1}{2}x < y < {2x},x + y = 3$ ,若 $\dfrac{1}{{2x} - y} + \dfrac{2}{{2y} - x} \geq  a$ 恒成立,则实数 $a$ 的范围是 \blkda[4]{$\left( {-\infty ,\dfrac{3 + 2\sqrt{2}}{3}}\right\rbrack$}
    \begin{jieda}
        因为 $x > 0,y > 0$ 且 $x + y = 3$ ,若 $\dfrac{1}{{2x} - y} + \dfrac{2}{{2y} - x} \geq  a$ 恒成立，则 $a \leq  {\left( \dfrac{1}{{2x} - y} + \dfrac{2}{{2y} - x}\right) }_{\min }$
        
        $\dfrac{1}{{2x} - y} + \dfrac{2}{{2y} - x} = \dfrac{1}{3}\left( {\dfrac{1}{{2x} - y} + \dfrac{2}{{2y} - x}}\right) \left( {{2x} - y + {2y} - x}\right) = \dfrac{1}{3}\left\lbrack  {3 + \dfrac{{2y} - x}{{2x} - y} + \dfrac{2\left( {{2x} - y}\right) }{{2y} - x}}\right\rbrack   \geq  \dfrac{1}{3}\left\lbrack  {3 + 2\sqrt{\dfrac{{2y} - x}{{2x} - y} \cdot  \dfrac{2\left( {{2x} - y}\right) }{{2y} - x}}}\right\rbrack   = \dfrac{3 + 2\sqrt{2}}{3}$，当且仅当 $\dfrac{{2y} - x}{{2x} - y} = \dfrac{2\left( {{2x} - y}\right) }{{2y} - x}$ ,即 $x = \sqrt{2},y = 3 - \sqrt{2}$ 时,等号成立
        
        $\therefore a \leq  \dfrac{3 + 2\sqrt{2}}{3}$ ,即实数 $a$ 的取值范围是 $\left( {-\infty ,\dfrac{3 + 2\sqrt{2}}{3}}\right\rbrack$ .
    \end{jieda}

\end{enumerate}

\section*{含绝对值不等式、不等式恒成立与有解问题}
\subsection*{A Level}
\begin{enumerate}
    \item 不等式 $\left| {x + 1}\right|  > 3$ 的解集是\blkda[4]{$\left( {-\infty , - 4}\right)  \cup  \left( {2, + \infty }\right)$}. 
    \begin{jieda}
        不等式 $\left| {x + 1}\right|  > 3$ 等价于 ${\left( x + 1\right) }^{2} > 9$ ,即 ${x}^{2} + {2x} - 8 > 0$ ,解得 $x <  - 4$ 或 $x > 2$ , 所以不等式 $\left| {x + 1}\right|  > 3$ 的解集是 $\left( {-\infty , - 4}\right)  \cup  \left( {2, + \infty }\right)$ .
    \end{jieda}
    
    \item 若不等式 $\left| {x - a}\right|  < 2\left( {a \in  \mathbf{R}}\right)$ 的解集为$(-1, t)$,则实数 $t$ 等于\blkda[4]{3}. 
    \begin{jieda}
        不等式 $\left| {x - a}\right|  < 2$ ,化为 $a - 2 < x < a + 2$ ,因此不等式 $\left| {x - a}\right|  < 2$ 的解集为 $\left( {a - 2,a + 2}\right)$ , 依题意, $\left( {-1,t}\right)  = \left( {a - 2,a + 2}\right)$ ,于是 $\left\{  \begin{array}{l} a - 2 =  - 1 \\  t = a + 2 \end{array}\right.$ ,解得 $a = 1,t = 3$ , 所以实数 $t$ 等于 3
    \end{jieda}
    \item 不等式 $\left| {1 - {4x}}\right|  > {2x}$ 的解集为\blkda[4]{$\left(-\infty, \dfrac{1}{6}\right) \cup \left(\dfrac{1}{2}, +\infty\right)$} 
    \begin{jieda}
        $\left| {1 - {4x}}\right|  > {2x} \Longleftrightarrow 1 - {4x} > {2x}$ 或 $1 - {4x} <  - {2x}$ , 解得 $x < \dfrac{1}{6}$ 或 $x > \dfrac{1}{2}$
    \end{jieda}
    \item 不等式 $\left| {{x}^{2} - 4}\right|  < x + 2$ 的解集为\blkda[4]{$(1, 3)$}. 
    \begin{jieda}
        $\left| {{x}^{2} - 4}\right|  < x + 2 \Longleftrightarrow -(x+2) < x^2-4 < (x+2) \Longleftrightarrow \left\{\begin{array}{l}
             (x+2)(x-3) < 0  \\
             (x+2)(x-1) > 0
        \end{array}\right.$，解得$ x \in (1,3)$
    \end{jieda}
    \item 不等式 $\left| {x + 2}\right|  + \left| {x - 2}\right|  \leq  4$ 的解集是\blkda[4]{$\left\lbrack  {-2,2}\right\rbrack$}.
    \begin{jieda}
        \begin{dingautolist}{172}
            \item 当 $x <  - 2$ 时, $- x - 2 + 2 - x \leq  4$ ,解得 $x \geq   - 2$ ,此时解集为空集
            \item 当 $- 2 \leq  x \leq  2$ 时, $x + 2 + 2 - x \leq  4$ ,即 $4 \leq  4$ ,符合要求,此时解集为 $\left\lbrack  {-2,2}\right\rbrack$
            \item 当 $x > 2$ 时, $x + 2 + x - 2 \leq  4$ ,解得 $x \leq  2$ ,此时解集为空集
        \end{dingautolist}
        综上，解集为 $\left\lbrack  {-2,2}\right\rbrack$
    \end{jieda} 
    \item 不等式 $\left| {{2x} + 1}\right|  - \left| {5 - x}\right|  > 0$ 的解集为\blkda[4]{$\left( {-\infty , - 6}\right)  \cup  \left( {\dfrac{4}{3}, + \infty }\right)$}
    \begin{jieda}
        由不等式 $\left| {{2x} + 1}\right|  - \left| {5 - x}\right|  > 0$，可得 ${\left( 2x + 1\right) }^{2} > {\left( 5 - x\right) }^{2}$，即 $\left( {x + 6}\right) \left( {{3x} - 4}\right)  > 0$
        
        解得 $x \in  \left( {-\infty , - 6}\right)  \cup  \left( {\dfrac{4}{3}, + \infty }\right)$ .
    \end{jieda}
    
    \item 已知 $x \in  \mathbf{R}$ ,则 “ ${x}^{3} > 8$ ” 是 “ $\left| x\right|  > 2$ ” 的\dotfill{(\kh{A})}
    \xx{充分不必要条件}{必要不充分条件}{充要条件}{既不充分也不必要条件}
    \begin{jieda}
        ${x}^{3} > 8 \Longleftrightarrow  x > 2$，$\left| x\right|  > 2 \Longleftrightarrow  x > 2$ 或 $x <  - 2$

        故“ ${x}^{3} > 8$ ” 是 “ $\left| x\right|  > 2$ ” 的充分不必要条件
    \end{jieda} 
    \item 已知 $x \in  \mathbf{R}$ ,则 “ $\left| {x + 1}\right|  + \left| {x - 1}\right|  \leq  2$ ” 是 “ $\dfrac{1}{x} > 1$ ” 的\dotfill{(\kh{B})}
    \xx{充分不必要条件}{必要不充分条件}{充要条件}{既不充分也不必要条件}
    \begin{jieda}
        $\left| {x + 1}\right|  + \left| {x - 1}\right|  \leq  2 \Longleftrightarrow x \in \left\lbrack  {-1,1}\right\rbrack$

        $\dfrac{1}{x} > 1 \Longleftrightarrow x \in (0, 1)$

        $(0, 1) \subset [-1, 1]$，故“ $\left| {x + 1}\right|  + \left| {x - 1}\right|  \leq  2$ ” 是 “ $\dfrac{1}{x} > 1$ ” 的必要不充分条件
    \end{jieda}
    
    \item 如果 $\left| {x + 1}\right|  + \left| {x + 9}\right|  > a$ 对任意实数 $x$ 总成立,那么 $a$ 的取值范围是\blkda[4]{$\left( {-\infty ,8}\right)$}.
    \begin{jieda}
        $\left| {x + 1}\right|  + \left| {x + 9}\right|  \geq  \left| {x + 1 - \left( {x + 9}\right) }\right|  = 8$ ,当且仅当 $- 9 \leq  x \leq   - 1$ 时等号成立，故 $a < 8$
    \end{jieda}
    \item 若存在实数 $x$ 使得不等式 $\left| {x + 1}\right|  + \left| {x - a}\right|  \leq  2$ 成立,则实数 $a$ 的取值范围是\blkda[4]{$\left\lbrack  {-3,1}\right\rbrack$}.
    \begin{jieda}
        因为 $\left| {x + 1}\right|  + \left| {x - a}\right|  \geq  \left| {\left( {x + 1}\right)  - \left( {x - a}\right) }\right|  = \left| {a + 1}\right|$ ,当且仅当 $\left( {x + 1}\right) \left( {x - a}\right)  \leq  0$ 时,等号成立，由题意可得 $\left| {a + 1}\right|  \leq  2$ ,解得 $- 3 \leq  a \leq  1$
    \end{jieda}
    \item 若关于 $x$ 的一元二次不等式 $2{x}^{2} - {kx} + \dfrac{3}{8} > 0$ 对于一切实数 $x$ 都成立,则实数 $k$ 的取值范围为\blkda[4]{$\left( {-\sqrt{3},\sqrt{3}}\right)$}.
    \begin{jieda}
        由题意 $\Delta  = {k}^{2} - 4 \times  2 \times  \dfrac{3}{8} < 0, - \sqrt{3} < k < \sqrt{3}$ .
    \end{jieda}
    \item 已知命题“ $\exists x \in  \mathbf{R}$ ,使 $4{x}^{2} + \left( {a - 2}\right) x + \dfrac{1}{4} \leq  0$ ”是真命题,则实数 $a$ 的取值范围是\blkda[4]{$(-\infty, 0] \cup [4, +\infty)$}. 
    \begin{jieda}
        依题意: “ $\exists x \in  \mathbf{R}$ ,使 $4{x}^{2} + \left( {a - 2}\right) x + \dfrac{1}{4} \leq  0$ ”是真命题,
        
        所以 $\Delta  = {\left( a - 2\right) }^{2} - 4 \times  4 \times  \dfrac{1}{4} = {a}^{2} - {4a} \geq  0$，解得 $a \leq  0$ 或 $a \geq  4$ .
    \end{jieda}
\end{enumerate}
\subsection*{B Level}
\begin{enumerate}
    
    \item  $x$ 的不等式 ${x}^{2} - 2\left| x\right|  - 3 < 0$ 解集是\blkda[4]{$(-3,3)$}.
    \begin{jieda}
        \begin{dingautolist}{172}
            \item 当 $x \geq  0$ 时,不等式化为 ${x}^{2} - {2x} - 3 < 0$ ,解得 $- 1 < x < 3$ ,即 $0 \leq  x < 3$
            \item 当 $x < 0$ 时,不等式化为 ${x}^{2} + {2x} - 3 < 0$ ,解得 $- 3 < x < 1$ ,即 $- 3 < x < 0$
        \end{dingautolist}
        综上，不等式的解集为$(-3,3)$
    \end{jieda}
    
    \item 命题“任意 $x \in  \left\lbrack  {-1,2}\right\rbrack$ , ${x}^{2} - {2x} - a \leq  0$ ”为真命题,则实数 $a$ 的取值范围是\blkda[4]{$\lbrack 3, + \infty )$}.
    \begin{jieda}
        任意 $x \in  \left\lbrack  {-1,2}\right\rbrack$ , ${x}^{2} - {2x} - a \leq  0$ 恒成立 $\Longleftrightarrow  {x}^{2} - {2x} \leq  a$ 恒成立,故只需
        
        ${\left( {x}^{2} - 2x\right) }_{\max } \leq  a$ ,记 $f\left( x\right)  = {x}^{2} - {2x} = {\left( x - 1\right) }^{2} - 1$ , $x \in  \left\lbrack  {-1,2}\right\rbrack$ ,易知 $f{\left( x\right) }_{\max } = f\left( {-1}\right)  = 3$ ,所以 $3 \leq  a$ .
    \end{jieda}
    \item 若存在 $x \in  \left\lbrack  {1,3}\right\rbrack$ ,使不等式 ${x}^{2} - {2ax} + a + 2 \leq  0$ 成立,则 $a$ 的取值范围为\blkda[4]{$\lbrack 2, + \infty )$}.
    \begin{jieda}
        由 ${x}^{2} - {2ax} + a + 2 \leq  0 \Longrightarrow  {x}^{2} + 2 \leq  a\left( {{2x} - 1}\right)$

        因为 $x \in  \left\lbrack  {1,3}\right\rbrack$ ,所以 ${2x} - 1 \in  \left\lbrack  {1,5}\right\rbrack$ ,令 $t = {2x} - 1 \in  \left\lbrack  {1,5}\right\rbrack  ,x = \dfrac{t + 1}{2}$ ,
        
        由 ${x}^{2} + 2 \leq  a\left( {{2x} - 1}\right)  \Longrightarrow  a \geq  \dfrac{{x}^{2} + 2}{{2x} - 1} = \dfrac{\dfrac{1}{4}\left( {{t}^{2} + {2t} + 9}\right) }{t} = \dfrac{1}{4}\left( {t + \dfrac{9}{t} + 2}\right)$ ,
        
        构造函数 $g\left( t\right)  = \dfrac{1}{4}\left( {t + \dfrac{9}{t} + 2}\right)  \geq  \dfrac{1}{4}\left( {2\sqrt{t \cdot  \dfrac{9}{t}} + 2}\right)  = 2$ ,
        
        即 $g{\left( t\right) }_{\min } = 2$ ,当且仅当 $t = 3 \in  \left\lbrack  {1,5}\right\rbrack$ 时取等号，所以 $a \geq  g{\left( t\right) }_{\min } = 2$
    \end{jieda} 
    \item 若命题“ $\forall x > 0$ ,使得 ${x^2} + {2ax} + {2a} + 3 \geq  0$ ”为假命题,则实数 $a$ 的取值范围\blkda[4]{$(-\infty, -1)$}.
    \begin{jieda}
         “ $\exists x \in  \left( {0, + \infty }\right)$ ,使得 $x^{2} + {2a}x + {2a} + 3 < 0$ ” 为真命题，即 ${2a} <  - \dfrac{x^{2} + 3}{x + 1}$ 在 $x \in  \left( {0, + \infty }\right)$ 内有解，即 ${2a} < {\left( -\dfrac{x^{2} + 3}{x + 1}\right) }_{\max }$
         
        因为 $- \dfrac{x^{2} + 3}{x + 1} =  - \dfrac{{\left( x + 1\right) }^{2} - 2\left( {x + 1}\right)  + 4}{x + 1} =  - \left( {x + 1 + \dfrac{4}{x + 1} - 2}\right) \leq - 2$（当且仅当 $x + 1 = \dfrac{4}{x + 1}$ ,即 $x = 1$ 时等号成立），
        所以 ${\left( -\dfrac{x^{2} + 3}{x + 1}\right) }_{\max } =  - 2$ ,所以 ${2a} <  - 2$ ,解得 $a <  - 1$
    \end{jieda}
    \item 命题： $\exists x \in  \left\lbrack  {1,4}\right\rbrack  ,{x^2} - \left( {{a^2} - {4a} - 1}\right) x + 4 < 0$ 的否定为真命题,则实数 $a$ 的最大值为\blkda[4]{5}.
    \begin{jieda}
        $\forall x \in  \left\lbrack  {1,4}\right\rbrack  ,{x^2} - \left( {{a^2} - {4a} - 1}\right) x + 4 \geq  0$为真命题.
                
        分离参数化简得: ${a^2} - {4a} - 1 \leq  \dfrac{{x}^{2} + 4}{x}\left( {x \in  \left\lbrack  {1,4}\right\rbrack  }\right)$ 恒成立.
        
        对 $\forall x \in  \left\lbrack  {1,4}\right\rbrack  ,\dfrac{{x^2} + 4}{x} = x + \dfrac{4}{x} \geq  2\sqrt{4} = 4$ ,当且仅当 $x = 2$ 时取得最小值 4 ,
        
        即 ${a^2} - {4a} - 1 \leq  4\therefore a \in  \left\lbrack  {-1,5}\right\rbrack  ,\therefore a$ 的最大值为 5
    \end{jieda}
    \item  已知 $f\left( x\right)  = {x}^{2} - \left( {a + 1}\right) x + a$ . 
    \begin{enumerate}
        \item 若 $f\left( x\right)  >  - \dfrac{1}{4}$ 恒成立,求实数 $a$ 的取值范围
        \item 求不等式 $f\left( x\right)  > 0$ 的解集
    \end{enumerate}
    \begin{jieda}
        \begin{enumerate}
            \item $f\left( x\right)  >  - \dfrac{1}{4}$ 恒成立即${x}^{2} - \left( {a + 1}\right) x + a + \dfrac{1}{4} > 0$ 对 $\forall x \in  \mathbf{R}$ 恒成立,
            
            故 $\Delta  = {\left( -a - 1\right) }^{2} - 4\left( {a + \dfrac{1}{4}}\right)  < 0$ ,化简得 $\Delta  = a\left( {a - 2}\right)  < 0$ ,解得 $0 < a < 2$
            \item $f\left( x\right)  = {x}^{2} - \left( {a + 1}\right) x + a > 0$ ,即 $\left( {x - a}\right) \left( {x - 1}\right)  > 0$
            \begin{dingautolist}{172}
                \item 当 $a > 1$ 时,不等式的解为 $\{ x \mid  x < 1$ 或 $x > a\}$
                \item 当 $a < 1$ 时,不等式的解为 $\{ x \mid  x < a$ 或 $x > 1\}$
                \item 当 $a = 1$ 时,不等式的解为 $\{ x \mid  x \neq  1\}$
            \end{dingautolist}
        \end{enumerate}
    \end{jieda}
    
    \item 设 $y = m{x}^{2} + \left( {1 - m}\right) x + m - 2$ .
    \begin{enumerate}
        \item 若不等式 $y \geq   - 2$ 对一切实数 $x$ 恒成立,求实数 $m$ 的取值范围
        \item 解关于 $x$ 的不等式 $m{x}^{2} + \left( {1 - m}\right) x + m - 2 < m - 1\left( {m \in  \mathbf{R}}\right)$
    \end{enumerate}
    \begin{jieda}
        \begin{enumerate}
            \item 由题设 $m{x}^{2} + \left( {1 - m}\right) x + m - 2 \geq   - 2$ ,即 $m{x}^{2} + \left( {1 - m}\right) x + m \geq  0$ 对一切实数 $x$ 恒成立

            \begin{dingautolist}{172}
                \item 当 $m = 0$ 时, $m{x}^{2} + \left( {1 - m}\right) x + m = x \geq  0$ 不恒成立
                \item 当 $m \neq  0$ 时,只需 $\left\{  \begin{array}{l} m > 0 \\  \Delta  = {\left( 1 - m\right) }^{2} - 4{m}^{2} \leq  0 \end{array}\right.$ ,可得 $m \geq  \dfrac{1}{3}$
            \end{dingautolist}
            综上, $m \geq  \dfrac{1}{3}$ .
            \item 
            原不等式即$m{x}^{2} + \left( {1 - m}\right) x - 1 < 0$
            \begin{dingautolist}{172}
                \item 当 $m = 0$ 时, $m{x}^{2} + \left( {1 - m}\right) x + m - 2 < m - 1$ ,即 $x - 2 <  - 1$ ,可得 $x < 1$
                \item $m > 0$ ,则 $\left( {x + \dfrac{1}{m}}\right) \left( {x - 1}\right)  < 0$ ,可得 $- \dfrac{1}{m} < x < 1$
                \item $m < 0$ ,则 $\left( {x + \dfrac{1}{m}}\right) \left( {x - 1}\right)  > 0$
                \begin{enumerate}[1.]
                    \item $- 1 < m < 0$ 时,可得 $x >  - \dfrac{1}{m}$ 或 $x < 1$
                    \item $m =  - 1$ 时,可得 $x \neq  1$ ,解集为 $\left( {-\infty ,1}\right)  \cup  \left( {1, + \infty }\right)$
                    \item $m <  - 1$ 时,可得 $x > 1$ 或 $x <  - \dfrac{1}{m}$
                \end{enumerate}
            \end{dingautolist}
        \end{enumerate}
    \end{jieda}
    \item 已知函数 $f\left( x\right)  = \left| x\right|  + \left| {x + a}\right|$ 
    \begin{enumerate}
        \item 若不等式 $f\left( x\right)  \geq  {2a} - 1$ 对任意的 $x \in  \mathbf{R}$ 恒成立,求实数 $a$ 的取值范围
        \item 若不等式 $f\left( x\right)  \leq  {2a} - 1$ 的解集为 $\left\lbrack  {b,b + 3}\right\rbrack$ ,求实数 $a$ 、 $b$ 的值
    \end{enumerate}
    \begin{jieda}
        \begin{enumerate}
            \item 对 $\forall x \in  \mathbf{R},f\left( x\right)  = \left| x\right|  + \left| {x + a}\right|  \geq  \left| {x - \left( {x + a}\right) }\right|  = \left| a\right|$ ,
            当且仅当 $x\left( {x + a}\right)  \leq  0$ 时取等号
            
            故原条件等价 $\left| a\right|  \geq  {2a} - 1$ ,即 $a \geq$
            ${2a} - 1$ 或 $a \leq   - \left( {{2a} - 1}\right)$ ,所以 $a \leq  1$ ,故实数 $a$ 的取值范围是 $( - \infty ,1\rbrack$ 
            \item 由 ${2a} - 1 \geq  \left| x\right|  + \left| {x + a}\right|  \geq  0$ ,可知 ${2a} - 1 \geq  0$ ,所以 $a \geq  \dfrac{1}{2}$ ,
            故 $- a < 0$ ,故 $f\left( x\right)  = \left\{  \begin{array}{ll}  - {2x} - a, & x <  - a, \\  a, &  - a \leq  x \leq  0, \\  {2x} + a, & x > 0 \end{array}\right.$ 的图像如图所
            示
            \begin{center}
                \includegraphics[width=0.5\linewidth]{bodyxb//src/yikeyilian3.png}
            \end{center}
            由图可知 $\left\{  \begin{array}{l}  - {2b} - a = {2a} - 1, \\  2\left( {b + 3}\right)  + a = {2a} - 1, \end{array}\right.$ 解得 $\left\{  \begin{array}{l} a = 2, \\  b =  - \dfrac{5}{2}。 \end{array}\right.$
        \end{enumerate}
    \end{jieda}
    \item 已知函数 $f\left( x\right)  = \left| {x + 1}\right|  - \left| {x - 2}\right|$ .
    \begin{enumerate}
        \item 求不等式 $f\left( x\right)  \geq  1$ 的解集
        \item 若不等式 $f\left( x\right)  \geq  {x}^{2} - x + m$ 的解集非空,求实数 $m$ 的取值范围
    \end{enumerate}
    \begin{jieda}
        \begin{enumerate}
            \item $f\left( x\right)  = \left| {x + 1}\right|  - \left| {x - 2}\right|  = \left\{  \begin{array}{ll}  
            - 3, & x \leq  - 1 \\  
            {2x} - 1,& - 1 <  x <  2 \\  
            3, & x \geq 2 \end{array}\right.$

                
            \begin{dingautolist}{172}
                \item $x \leq  - 1$，无解
                \item $- 1 <  x <  2$，$f\left( x\right)  \geq  1 \Longleftrightarrow x \geq  1$
                \item $x \geq 2$，$f\left( x\right)  \geq  1$恒成立
            \end{dingautolist}
            
            综上，$f\left( x\right)  \geq  1$ 的解集为 $\{ x \mid  x \geq  1\}$
            \item 原式等价于存在 $x \in  \mathbf{R}$ 使得 $f\left( x\right)  - {x}^{2} + x \geq  m$ 成立，即 $m \leq  \left\lbrack  {f\left( x\right)  - {x}^{2} + x}\right\rbrack_{max}$ ,设 $g\left( x\right)  = f\left( x\right)  - {x}^{2} + x$ .
            
            $g\left( x\right)  = \left\{  \begin{array}{ll}  
            - {x}^{2} + x - 3, & x \leq   - 1 \\ 
            - {x}^{2} + {3x} - 1, & - 1 < x < 2 \\
            - {x}^{2} + x + 3, &x \geq  2 \end{array}\right.$ 

            \begin{dingautolist}{172}
                \item 当 $x \leq   - 1$ 时, $g\left( x\right)  =  - {x}^{2} + x - 3$ ,其开口向下,对称轴方程为 $x = \dfrac{1}{2} >  - 1$，故$g\left( x\right)  \leq  g\left( {-1}\right)  =  - 1 - 1 - 3 =  - 5$
                \item 当 $- 1 < x < 2$ 时, $g\left( x\right)  =  - {x}^{2} + {3x} - 1$ ,其开口向下,对称轴方程为 $x = \dfrac{3}{2} \in  \left( {-1,2}\right)$，故$g\left( x\right)  \leq  g\left( \dfrac{3}{2}\right)  =  - \dfrac{9}{4} + \dfrac{9}{2} - 1 = \dfrac{5}{4}$
                \item 当 $x \geq  2$ 时, $g\left( x\right)  =  - {x}^{2} + x + 3$ ,其开口向下,对称轴方程为 $x = \dfrac{1}{2} < 2$，故$g\left( x\right)  \leq  g\left( 2\right)  =  - 4 + 2 + 3 = 1$
            \end{dingautolist}
            
            综上, $m \leq \dfrac{5}{4}$
        \end{enumerate}
    \end{jieda}
    \item 已知函数 $f\left( x\right)  = \left| {x - {a}^{2}}\right|  + \left| {x - {2a} + 1}\right|$
    \begin{enumerate}
        \item 当 $a = 2$ 时,求不等式 $f\left( x\right)  \geq  4$ 的解集
        \item 若 $f\left( x\right)  \geq  4$恒成立 ,求 $a$ 的取值范围
    \end{enumerate}
    \begin{jieda}
        \begin{enumerate}
        \item 当 $a = 2$ 时, $f\left( x\right)  = \left| {x - 4}\right|  + \left| {x - 3}\right| = \left\{\begin{array}{ll}
             7 - {2x}, & x \leq  3  \\
             4, & 3 < x < 4 \\
             {2x} - 7, & x \geq  4
        \end{array}\right.$ .
        \begin{dingautolist}{172}
            \item $x \leq  3$，$f\left( x\right)  \geq  4 \Longleftrightarrow x \leq  \dfrac{3}{2}$
            \item $3 < x < 4$，无解
            \item $x \geq  4$，$f\left( x\right)  \geq  4 \Longleftrightarrow x \geq  \dfrac{11}{2}$
        \end{dingautolist}
        
        综上所述: $f\left( x\right)  \geq  4$ 的解集为 $\left\{  x \left |  x \leq  \dfrac{3}{2}\right. \right.$ 或 $\left. {x \geq  \dfrac{11}{2}}\right\}$ .
        \item $f\left( x\right)  = \left| {x - {a}^{2}}\right|  + \left| {x - {2a} + 1}\right|  \geq  \left| {\left( {x - {a}^{2}}\right)  - \left( {x - {2a} + 1}\right) }\right|  = \left| {-{a}^{2} + {2a} - 1}\right|  = {\left( a - 1\right) }^{2}$ (当且仅当 ${2a} - 1 \leq  x \leq  {a}^{2}$ 时取等号),
        
        $\therefore {\left( a - 1\right) }^{2} \geq  4$ ,解得: $a \leq   - 1$ 或 $a \geq  3$
    \end{enumerate}
    \end{jieda}
    
    \item 已知 $f\left( x\right)  = \left| {x + 1}\right|  - \left| {{ax} - 1}\right|$ .
    \begin{enumerate}
        \item 当 $a = 1$ 时,求不等式 $f\left( x\right)  > 1$ 的解集
        \item 若 $x \in  \left( {0,1}\right)$ 时不等式 $f\left( x\right)  > x$ 成立,求 $a$ 的取值范围
    \end{enumerate}
    \begin{jieda}
        \begin{enumerate}
            \item 当 $a = 1$ 时, $f\left( x\right)  = \left| {x + 1}\right|  - \left| {x - 1}\right|$ ,即 $f\left( x\right)  = \left\{  \begin{array}{ll}  
            - 2, & x \leq   - 1 \\  
            {2x}, & - 1 < x < 1 \\
            2, & x \geq  1 
            \end{array}\right.$ 
                
            \begin{dingautolist}{172}
                \item $x \leq  -1$，无解
                \item $-1 < x < 1$，$f\left( x\right)  > 1 \Longleftrightarrow x > \dfrac{1}{2}$
                \item $x \geq  1$，恒成立
            \end{dingautolist}
            故不等式 $f\left( x\right)  > 1$ 的解集为 $\left\{  {x \left |  x > \dfrac{1}{2}\right.}\right\}$ .
            \item  $x \in  \left( {0,1}\right)$ 时, $\left| {x + 1}\right|  - \left| {{ax} - 1}\right|  > x \Longleftrightarrow \left| {{ax} - 1}\right|  < 1 \Longleftrightarrow -1 < ax-1 < 1 \Longleftrightarrow 0 < a  < \dfrac{2}{x}$

            当$0 < a \leq  2$时，该式恒成立.
            
        \end{enumerate}
    \end{jieda}
    \item 设函数 $f\left( x\right)  = 5 - \left| {x + a}\right|  - \left| {x - 2}\right|$ .
    \begin{enumerate}
        \item 当 $a = 1$ 时,求不等式 $f\left( x\right)  \geq  0$ 的解集
        \item 若 $f\left( x\right)  \leq  1$ 恒成立,求 $a$ 的取值范围
    \end{enumerate}
    \begin{jieda}
        \begin{enumerate}
            \item 当 $a = 1$ 时，$f\left( x\right)  = \left\{  \begin{array}{ll} 
            {2x} + 4, & x \leq   - 1 \\
            2, & - 1 < x \leq  2 \\
            - {2x} + 6, &x > 2. \end{array}\right.$
            
            可得 $f\left( x\right)  \geq  0$ 的解集为 $\{ x \mid   - 2 \leq  x \leq  3\}$ .
            \item $f\left( x\right)  \leq  1$ 等价于 $\left| {x + a}\right|  + \left| {x - 2}\right|  \geq  4$ .
            
            而 $\left| {x + a}\right|  + \left| {x - 2}\right|  \geq  \left| {a + 2}\right|$ ,且当 $x = 2$ 时等号成立. 故 $f\left( x\right)  \leq  1$ 等价于 $\left| {a + 2}\right|  \geq  4$ .
            
            由 $\left| {a + 2}\right|  \geq  4$ 可得 $a \leq   - 6$ 或 $a \geq  2$ ,所以 $a$ 的取值范围是 $\left( {-\infty , - 6}\right\rbrack   \cup  \lbrack 2, + \infty )$ .
        \end{enumerate}
    \end{jieda}
\end{enumerate}
\subsection*{C Level}
\begin{enumerate}
    \item 设 $m \in  \mathbf{R}$ ,若对于任意实数 $a$ ,都存在 $x \in  \left\lbrack  {-2,2}\right\rbrack$ 满足 $\left| {{x}^{2} - 1}\right|  + \left| {x - a}\right|  > m$ ,则 $m$ 的取值范围是\blkda{$\left( {-\infty ,5}\right)$}
    \begin{jieda}
    记 $f\left( x\right)  = \left| {{x}^{2} - 1}\right|  + \left| {x - a}\right| ,x \in  \left\lbrack  {-2,2}\right\rbrack$ 

    因为任意实数 $a$ ,都存在 $x \in  \left\lbrack  {-2,2}\right\rbrack$ 满足 $\left| {{x}^{2} - 1}\right|  + \left| {x - a}\right|  > m$ ,所以 ${\left\lbrack  f{\left( x\right) }_{\max }\right\rbrack  }_{\min } > m$

    易得 $f{\left( x\right) }_{\max } = \max \{ f\left( {-2}\right) ,f\left( 2\right) \}  = \max \{ 3 + \left| {2 - a}\right| ,3 + \left| {2 + a}\right| \}$

    根据$y_1 = 3 + \left| {2 - a}\right|$与$y_2 = 3 + \left| {2 + a}\right|$的函数图像可知$f{\left( x\right) }_{\max } = \left\{  \begin{array}{ll}  a + 5, & a \geq  0 \\  
    - a + 5, & a < 0 \end{array}\right.$ 
    
    所以当 $a = 0$ 时, ${\left\lbrack  f{\left( x\right) }_{\max }\right\rbrack  }_{\min } = 5$，所以 $m$ 的取值范围是 $\left( {-\infty ,5}\right)$ 。
    \end{jieda}
    \item 设 $a \in  \mathbf{R}$ ,若 $x > 0$ 时,均有 $\left\lbrack  {\left( {a - 1}\right) x - 1}\right\rbrack  \left( {{x}^{2} - {ax} - 1}\right)  \geq  0$ ,则实数 $a$ 的取值集合为\blkda{$\left\{  \dfrac{3}{2}\right\}$}
    \begin{jieda}
    \begin{dingautolist}{172}
        \item 当 $a = 1$ 时
        
        不等式左边为 $- {x}^{2} + x + 1 = - {\left( x - \dfrac{1}{2}\right) }^{2} + \dfrac{5}{4} \leq  \dfrac{5}{4}$ ,不满足题意
        \item 当 $a \neq  1$ 时

        构造函数 ${y}_{1} = \left( {a - 1}\right) x - 1$，${y}_{2} = {x}^{2} - {ax} - 1$ ,它们都过定点 $P\left( {0, - 1}\right)$ (如图)
        \begin{center}
            \includegraphics[width=0.5\linewidth]{bodyxb//src/yikeyilian2-9.png}
        \end{center}
        考查函数 ${y}_{1} = \left( {a - 1}\right) x - 1$ : 令 $y = 0$ ,得 $M\left( {\dfrac{1}{a - 1},0}\right)$，当$x$充分大时$y_2 > 0$，则应有$y_1 \geq 0$，所以必有$a > 1$; 
        
        考查函数 ${y}_{2} = {x}^{2} - {ax} - 1$ ,因为 $x > 0$ 时均有 $\left\lbrack  {\left( {a - 1}\right) x - 1}\right\rbrack  \left( {{x}^{2} - {ax} - 1}\right)  \geq  0$，当且仅当 ${y}_{2} = {x}^{2} - {ax} - 1$ 过点 $M\left( {\dfrac{1}{a - 1},0}\right)$时满足要求，代入得 ${\left( \dfrac{1}{a - 1}\right) }^{2} - \dfrac{a}{a - 1} - 1 = 0$
        
        解得$a = \dfrac{3}{2}$ ,或 $a = 0$ (舍去),故答案为 $\left\{  \dfrac{3}{2}\right\}$
    \end{dingautolist}
    \end{jieda}
    \item 若对任意 $x \in  \left\lbrack  {1,2}\right\rbrack$ ,均有 $\left| {{x}^{2} - a}\right|  + \left| {x + a}\right|  = \left| {{x}^{2} + x}\right|$ ,则实数 $a$ 的取值范围为\blkda[4]{$\left\lbrack  {-1,1}\right\rbrack$}.
    \begin{jieda}
        因为在绝对值三角不等式 $\left| {a + b}\right|  \leq  \left| a\right|  + \left| b\right|$ 中,当 $a,b$ 同号时有 $\left| {a + b}\right|  = \left| a\right|  + \left| b\right|$ , 又因为 $\left| {{x}^{2} + x}\right|  = \left| {\left( {{x}^{2} - a}\right)  + \left( {x + a}\right) }\right|  = \left| {{x}^{2} - a}\right|  + \left| {x + a}\right|$，所以 $\left( {{x}^{2} - a}\right) \left( {x + a}\right)  \geq  0$ 在 $x \in  \left\lbrack  {1,2}\right\rbrack$ 恒成立，
        \begin{enumerate}
            \item 方法1 

            构造函数$y_1 = x+a$，$y_2 = x^2-a$
            \twotu{\includegraphics[width=0.6\linewidth]{bodyxb//src/fenceng1.png}}{\includegraphics[width=0.6\linewidth]{bodyxb//src/fenceng2.png}}
            \begin{dingautolist}{172}
                \item $a=0$时，满足要求
                \item $a>0$时，当且仅当$\sqrt{a} \leq 1$，满足要求
                \item $a<0$时，当且仅当$-a \leq 1$，满足要求
            \end{dingautolist}
            综上，$a \in [-1, 1]$
            \item 方法2 
            
            $\left( {{x}^{2} - a}\right) \left( {x + a}\right)  \geq  0$ 在 $x \in  \left\lbrack  {1,2}\right\rbrack$ 恒成立，即$[1, 2]$是不等式的解集的子集

            \begin{dingautolist}{172}
                \item $a \in (-\infty, 0]$
                
                此时对应方程的一个实根是$-a$，故不等式的解集为$(-a, +\infty)$，当$a \in [-1, 0]$时满足$[1, 2]$是不等式的解集的子集
                \item $a \in (0, 1]$
    
                设此时方程对应的三个根分别是$-a, -\sqrt{a}, \sqrt{a}$满足$-\sqrt{a} < -a  < \sqrt{a}$，易知$\sqrt{a} \leq 1$，故必然满足$[1, 2]$是不等式的解集的子集
                
                \item $a \in (1, +\infty)$
                
                此时对应方程的三个实根分别是$-a, -\sqrt{a}, \sqrt{a}$满足$-a < -\sqrt{a} < \sqrt{a}$，故不等式的解集为$(-a, -\sqrt{a}) \cup (\sqrt{a}, +\infty)$，因为$-\sqrt{a} < 1 < \sqrt{a}$，所以不满足$[1, 2]$是不等式的解集的子集
            \end{dingautolist}
            
            综上，$a \in \left\lbrack  {-1,1}\right\rbrack$ .
        \end{enumerate}
        
    \end{jieda}
\end{enumerate}
\section*{函数的定义}
\subsection*{A Level}
\begin{enumerate}
    \item 已知 $f\left( {\dfrac{1}{2}x - 1}\right)  = {2x} + 3$ ,若 $f\left( t\right)  = 4$ ,则 $t =$ \blkda[4]{$- \dfrac{3}{4}$}.
    \begin{jieda}
        令 $\dfrac{1}{2}x - 1 = m$ ,则 $x = {2m} + 2$，所以 $f\left( m\right)  = 2\left( {{2m} + 2}\right)  + 3 = {4m} + 7$，因为 $f\left( t\right)  = 4$，所以 ${4t} + 7 = 4$ ,解得 $t =  - \dfrac{3}{4}$
    \end{jieda}
    \item 函数 $f\left( x\right)  = \dfrac{\sqrt{4 - {x}^{2}}}{x - 1}$ 的定义域为\blkda[4]{$\left\lbrack  {-2,1}\right)  \cup  (1,2\rbrack$}.
    \begin{jieda}
        要使函数有意义，须满足 $\left\{  \begin{array}{l} 4 - {x}^{2} \geq  0 \\  x - 1 \neq  0 \end{array}\right.$，解得 $- 2 \leq  x \leq  2$ 且 $x \neq  1$，故函数 $f\left( x\right)$ 的定义域为 $\left\lbrack  {-2, 1) \cup  (1,2}\right\rbrack$ .
    \end{jieda}
    \item 若函数 $y = \sqrt{m{x}^{2} + {4mx} + 3}$ 中 $x$ 的取值范围为 $\mathbf{R}$ ,则 $m$ 的取值范围是\blkda[4]{$\left\lbrack  {0,\dfrac{3}{4}}\right\rbrack$}.
    \begin{jieda}
        由已知 $m{x}^{2} + {4mx} + 3 \geq  0$ 恒成立
        \begin{dingautolist}{172}
            \item 当 $m = 0$ 时符合题意
            \item 当 $m \neq  0$ 时, $\Delta  = {16}{m}^{2} - {12m} \leq  0$，解得$0 < m \leq  \dfrac{3}{4}$
        \end{dingautolist}
        综上所述 $m \in  \left\lbrack  {0,\dfrac{3}{4}}\right\rbrack$ 
    \end{jieda}
    \item 函数 $f\left( x\right)  = \sqrt{-{x}^{2} + {ax} - 1}$ 在 $\left\lbrack  {\dfrac{1}{2},3}\right\rbrack$ 上有意义,则实数 $a$ 的取值范围为\blkda[4]{$\left\lbrack  {\dfrac{10}{3}, + \infty }\right)$}.
    \begin{jieda}
        由题意可知 $- {x}^{2} + {ax} - 1 \geq  0$ 在 $\left\lbrack  {\dfrac{1}{2},3}\right\rbrack$ 上恒成立，抛物线开口向下，则只需要端点处满足即可，因此 $\left\{  \begin{array}{l}  - \dfrac{1}{4} + \dfrac{1}{2}a - 1 \geq  0 \\   - 9 + {3a} - 1 \geq  0 \end{array}\right.  \Longrightarrow  a \geq  \dfrac{10}{3}$ ,
        
        所以满足题意的实数 $a$ 的取值范围为 $\left\lbrack  {\dfrac{10}{3}, + \infty }\right)$ .
    \end{jieda}
    \item 已知函数 $y = f\left( x\right)$ 的定义域为 $\left\lbrack  {1,4}\right\rbrack$ ,则 $f\left( {x + 2}\right)$ 的定义域为\blkda[4]{$\left\lbrack  {-1,2}\right\rbrack$}.
    \begin{jieda}
        根据题意 $f\left( x\right)$ 定义域为 $\left\lbrack  {1,4}\right\rbrack$ .
        
        则 $f\left( {x + 2}\right)$ 的定义域可得 $1 \leq  x + 2 \leq  4$ ,即 $- 1 \leq  x \leq  2$ .
        
        故 $f\left( {x + 2}\right)$ 的定义域为 $\left\lbrack  {-1,2}\right\rbrack$
    \end{jieda}
    \item 函数 $y = f\left( x\right)$ 的定义域 $\left\lbrack  {-2,4}\right\rbrack$ ,则函数 $g\left( x\right)  = f\left( x\right)  + f\left( {-x}\right)$ 的定义域为\blkda[4]{$\left\lbrack  {-2,2}\right\rbrack$}.
    \begin{jieda}
        $f(-x)$的定义域为$[-4, 2]$，函数相加定义域取交集，因此函数 $g\left( x\right)  = f\left( x\right)  + f\left( {-x}\right)$ 的定义域为 $\left\lbrack  {-2,2}\right\rbrack$ .
    \end{jieda}
    \item 已知 $M = \left\{  {y\left| {y = {x}^{2} - {4x} + 5}\right. }\right\}  ,N = \left\{  {x\left| {y = \sqrt{x + 1}}\right. }\right\}$ ,则 $M \cup  N =$ \blkda[4]{$\lbrack  - 1, + \infty )$}.
    \begin{jieda}
        $y = {x}^{2} - {4x} + 5 = {\left( x - 2\right) }^{2} + 1 \geq  1$ ,所以 $M = \lbrack 1, + \infty )$ .
        
        由 $x + 1 \geq  0$ 解得 $x \geq   - 1$ ,所以 $N = \lbrack  - 1, + \infty )$ ,
        
        所以 $M \cup  N = \lbrack  - 1, + \infty )$
    \end{jieda}
    \item 函数 $y = \dfrac{1 - x}{2 + x}\left( {-1 \leq  x \leq  1}\right)$ 的值域为\blkda[4]{$\left\lbrack  {0,2}\right\rbrack$}.
    \begin{jieda}
        由 $y = \dfrac{1 - x}{2 + x} = \dfrac{3}{2 + x} - 1$，又 $- 1 \leq  x \leq  1$ ,则 $1 \leq  2 + x \leq  3$ ,则 $1 \leq  \dfrac{3}{2 + x} \leq  3$ ,所以 $0 \leq  \dfrac{3}{2 + x} - 1 \leq  2$ ,
        
        故函数 $y = \dfrac{1 - x}{2 + x}\left( {-1 \leq  x \leq  1}\right)$ 的值域为 $\left\lbrack  {0,2}\right\rbrack$ .
    \end{jieda}
    \item 函数 $y = x + \dfrac{1}{x}$ 的值域为\blkda[4]{$\left( {-\infty , - 2\rbrack \cup \lbrack 2, + \infty }\right)$}.
    \item 函数 $f\left( x\right)  = \dfrac{1}{{x}^{2} - x + 2}$ 的值域为\blkda[4]{$\left( {0,\dfrac{4}{7}}\right\rbrack$}.
    \begin{jieda}
        因为二次函数 $y = {x}^{2} - x + 2 = {\left( x - \dfrac{1}{2}\right) }^{2} + \dfrac{7}{4}$ 的值域为 $\left\lbrack  {\dfrac{7}{4}, + \infty }\right)$ ,
        
        所以 $f\left( x\right)  = \dfrac{1}{{x}^{2} - x + 2}$ 的定义域是 $\mathbf{R}$ ,值域为 $\left( {0,\dfrac{4}{7}}\right\rbrack$ .
    \end{jieda}
    \item 函数 $y = x - \dfrac{1}{x}$ 在 $\left\lbrack  {1,2}\right\rbrack$ 上的值域为\blkda[4]{$\left\lbrack  {0,\dfrac{3}{2}}\right\rbrack$}.
    \begin{jieda}
        $\because y = x - \dfrac{1}{x}$ 在 $\left\lbrack  {1,2}\right\rbrack$ 上为增函数，则 $y = x - \dfrac{1}{x}$ 在 $\left\lbrack  {1,2}\right\rbrack$ 上的最小值为 $y = 1 - 1 = 0$ ,最大值为 $y = 2 - \dfrac{1}{2} = \dfrac{3}{2}$，即 $y \in  \left\lbrack  {0,\dfrac{3}{2}}\right\rbrack$ .
    \end{jieda}
    \item 若函数 $y = f\left( x\right)$ 的值域是 $\left\lbrack  {a,b}\right\rbrack$ ,则函数 $y = f\left( {x + 1}\right)  + 1$ 的值域为\blkda[4]{$\left\lbrack  {a + 1,b + 1}\right\rbrack$}.
    \begin{jieda}
        将函数 $y = f\left( x\right)$ 向左平移一个单位,得函数 $y = f\left( {x + 1}\right)$ ,其值域为 $\left\lbrack  {a,b}\right\rbrack$ ,
        
        继续将函数 $y = f\left( {x + 1}\right)$ 向上平移一个单位,得函数 $y = f\left( {x + 1}\right)  + 1$ ,其值域也变为 $\left\lbrack  {a + 1,b + 1}\right\rbrack$ , 故答案为 $\left\lbrack  {a + 1,b + 1}\right\rbrack$
    \end{jieda}
    \item 已知函数 $f\left( {x + 2}\right)  = 2{x}^{2} + {4x} + 3$ ,则 $f\left( x\right)  =$ \blkda[4]{$2x^2-4x+3$}
    \item 已知函数 $f\left( {\sqrt{x} + 1}\right)  = x - 4$，则 $f\left( x\right)  =$ \blkda[4]{${x}^{2} - {2x} - 3\left( {x \geq  1}\right)$}.
    \begin{jieda}
        令 $\sqrt{x} + 1 = t\left( {t \geq  1}\right)$ ,则 $x = {\left( t - 1\right) }^{2}\left( {t \geq  1}\right)$ ,
        
        于是有 $f\left( t\right)  = {\left( t - 1\right) }^{2} - 4 = {t}^{2} - {2t} - 3\left( {t \geq  1}\right)$ ,所以 $f\left( x\right)  = {x}^{2} - {2x} - 3\left( {x \geq  1}\right)$ .
    \end{jieda}
    \end{enumerate}
\subsection*{B Level}
\begin{enumerate}
    \item 已知 $f\left( x\right)$ 是一次函数，且满足 $f\left( {f\left( x\right) }\right)  = {9x} + 4$ ，则 $f\left( x\right) = $\blkda[4]{${3x} + 1$ 或 $ - {3x} - 2$}.
    \begin{jieda}
        设 $f\left( x\right)  = {kx} + b\left( {k \neq  0}\right)$ ,则 $f\left( {f\left( x\right) }\right)  = k\left( {{kx} + b}\right)  + b = {k}^{2}x + {kb} + b$ ,
        
        故 ${k}^{2}x + {kb} + b = {9x} + 4$ ,所以 $\left\{  \begin{array}{l} {k}^{2} = 9 \\  {kb} + b = 4 \end{array}\right.$，解得 $\left\{  \begin{array}{l} k = 3 \\  b = 1 \end{array}\right.$ 或 $\left\{  \begin{array}{l} k =  - 3 \\  b =  - 2 \end{array}\right.$
        
        故 $f\left( x\right)  = {3x} + 1$ 或 $f\left( x\right)  =  - {3x} - 2$ .
    \end{jieda}
    \item 若函数 $f\left( x\right)$ 满足 $f\left( x\right)  - {2f}\left( \dfrac{1}{x}\right)  = x + 2$ ,则 $f\left( x\right)  =$ \blkda[4]{$- \dfrac{1}{3}\left( {x + \dfrac{2}{x}}\right)  - 2$}.
    \begin{jieda}
        由 $f\left( x\right)  - {2f}\left( \dfrac{1}{x}\right)  = x + 2$ ,可得 $f\left( \dfrac{1}{x}\right)  - {2f}\left( x\right)  = \dfrac{1}{x} + 2$ ,
        
        联立两式消去 $f\left( \dfrac{1}{x}\right)$ ,可得 $f\left( x\right)  =  - \dfrac{1}{3}\left( {x + \dfrac{2}{x}}\right)  - 2$ .
    \end{jieda}
    \item 已知 $f\left( x\right)$ 满足 ${3f}\left( x\right)  + {2f}\left( {1 - x}\right)  = {4x}$ ,则 $f\left( x\right)$ 解析式为\blkda[4]{$f\left( x\right)  = {4x} - \dfrac{8}{5}$}.
    \begin{jieda}
        由 ${3f}\left( x\right)  + {2f}\left( {1 - x}\right)  = {4x}$ 可得 ${3f}\left( {1 - x}\right)  + {2f}\left( x\right)  = 4\left( {1 - x}\right)$ 
        
        联立两式，消去$f(1-x)$可得: $f\left( x\right)  = {4x} - \dfrac{8}{5}$
    \end{jieda}
    \item 已知函数 $f\left( x\right)  = \sqrt{x - 1}$ 的值域为 $\lbrack 1, + \infty )$ ,则其定义域是\blkda[4]{$\lbrack 2, + \infty )$}.
    \item 一系列函数若它们的解析式相同、值域相同,但其定义域不同,则称这一系列函数为 “同族函数”,若解析式为 $y = {x}^{2}$ 、值域为 $\{ 1,2\}$ 的“同族函数”共有\blkda[4]{9}个.
    \begin{jieda}
        与解析式为 $y = {x}^{2}$ ,值域是 $\{ 1,2\}$ 的“同族函数”的定义域可以为:
        
        $\{ 1,\sqrt{2}\} ,\{ 1, - \sqrt{2}\} ,\{  - 1,\sqrt{2}\} ,\{  - 1, - \sqrt{2}\} ,\{ 1, - 1,\sqrt{2}\} ,\{ 1, - 1, - \sqrt{2}\} ,$
        
        $\{ 1, - \sqrt{2},\sqrt{2}\} ,\{  - 1, - \sqrt{2},\sqrt{2}\} ,\{  - 1,1,\sqrt{2}, - \sqrt{2}\}$ ,共 9 个.
    \end{jieda}
    \item 一系列函数若它们的对应关系相同、值域相同,则称这一系列的函数为 “同族函数”. 对应关系为 $y = {x}^{2} - 1$ ,值域为\{0,3,8\}的“同族函数”有\blkda[4]{27}个.
    \begin{jieda}
        因为对应关系为 $y = {x}^{2} - 1$ ,值域为 $\{ 0,3,8\}$ ,所以由 ${x}^{2} - 1 = 0 \Longrightarrow  x =  \pm  1,{x}^{2} - 1 =$  $3 \Longrightarrow  x =  \pm  2,{x}^{2} - 1 = 8 \Longrightarrow  x =  \pm  3$ 知,要构成满足条件的函数,定义域应该由 $\pm  1, \pm  2, \pm  3$ 这三组数所构成,且 $\pm  1, \pm  2, \pm  3$ 这三组数中,每组的两个数至少要取一个,所以满足题意的定义域可以有 $3 \times  3 \times  3 = {27}$ (种),即所求的 “同族函数” 有 27 个.
    \end{jieda}
    \item 已知函数 $y = f\left( {x + 1}\right)$ 的定义域为 $\left\lbrack  {1,3}\right\rbrack$ ,则 $f\left( {x}^{2}\right)$ 的定义域为\blkda[4]{$\left\lbrack  {-2, - \sqrt{2}}\right\rbrack   \cup  \left\lbrack  {\sqrt{2},2}\right\rbrack$}.
    \begin{jieda}
        在 $y = f\left( {x + 1}\right)$ 中, $x \in  \left\lbrack  {1,3}\right\rbrack  ,\therefore x + 1 \in  \left\lbrack  {2,4}\right\rbrack  ,\therefore f\left( x\right)$ 的定义域是 $\left\lbrack  {2,4}\right\rbrack$
        
        从而在 $f\left( {x}^{2}\right)$ 中 $2 \leq  {x}^{2} \leq  4, - 2 \leq  x \leq   - \sqrt{2}$ 或 $\sqrt{2} \leq  x \leq  2$ ,
        
        $\therefore f\left( {x}^{2}\right)$ 的定义域是 $\left\lbrack  {-2, - \sqrt{2}}\right\rbrack   \cup  \left\lbrack  {\sqrt{2},2}\right\rbrack$ .
    \end{jieda}
    \item 若 $y = f\left( x\right)$ 的定义域为 $\left\lbrack  {-1,1}\right\rbrack$ ,则函数 $y = f\left( {3x}\right)  + f\left( \dfrac{x}{3}\right)$ 的定义域为\blkda[4]{$\left\lbrack  {-\dfrac{1}{3},\dfrac{1}{3}}\right\rbrack$}.
    \begin{jieda}
        因为 $y = f\left( x\right)$ 的定义域为 $\left\lbrack  {-1,1}\right\rbrack$ ,则 $\left\{  \begin{array}{l}  - 1 \leq  {3x} \leq  1 \\   - 1 \leq  \dfrac{x}{3} \leq  1 \end{array}\right.$ ,解得 $- \dfrac{1}{3} \leq  x \leq  \dfrac{1}{3}$ ,即函数
        
        $y = f\left( {3x}\right)  + f\left( \dfrac{x}{3}\right)$ 的定义域为 $\left\lbrack  {-\dfrac{1}{3},\dfrac{1}{3}}\right\rbrack$
    \end{jieda}
    \item 已知函数 $y = f\left( x\right)$ 的定义域为 $\left\lbrack  {-2,5}\right\rbrack$ ,则函数 $y = \dfrac{f\left( {{2x} - 1}\right) }{x - 1}$ 的定义域为\blkda[4]{$\left\lbrack  {-\dfrac{1}{2},1}\right)  \cup  (1,3\rbrack$}.
    \begin{jieda}
        函数 $y = f\left( x\right)$ 的定义域为 $\left\lbrack  {-2,5}\right\rbrack$，则 $\left\{  \begin{array}{l}  - 2 \leq  {2x} - 1 \leq  5 \\  x - 1 \neq  0 \end{array}\right.$ ,则 $- \dfrac{1}{2} \leq  x < 1$ 或 $1 < x \leq  3$
        
        则函数 $y = \dfrac{f\left( {{2x} - 1}\right) }{x - 1}$ 的定义域为 $\left\lbrack  {-\dfrac{1}{2},1}\right)  \cup  (1,3\rbrack$ .
    \end{jieda}
    \item 已知函数 $y = f\left( x\right)$ 的定义域为 $\left\lbrack  {a,b}\right\rbrack  ,\{ \left( {x,y}\right)  \mid  y = f\left( x\right) ,a \leq  x \leq  b\}  \cap  \{ \left( {x,y}\right)  \mid  x =$ 0\} 只有一个子集, 则\dotfill{(\kh{A})}
    \xx{${ab} > 0$}{${ab} \geq  0$}{${ab} < 0$}{${ab} \leq  0$}
    \begin{jieda}
        因为 $\{ \left( {x,y}\right)  \mid  y = f\left( x\right) ,a \leq  x \leq  b\}  \cap  \{ \left( {x,y}\right)  \mid  x = 0\}$ 只有一个子集，即该交集为空集, 说明 $\{ \left( {x,y}\right)  \mid  y = f\left( x\right) ,a \leq  x \leq  b\}$ 中没有横坐标为 0 的点,所以由函数的定义可知 $0 \notin  \left\lbrack  {a,b}\right\rbrack$ ,所以 $a\text{ 、 }b$ 同号.
    \end{jieda}
    \item 函数 $y = \sqrt{a{x}^{2} + {ax} + 1}$ 的定义域为 $\mathbf{R}$ ,则 $a \in$ \blkda[4]{$\{ a \mid  0 \leq  a \leq  4\}$}.
    \begin{jieda}
        因为任意 $x \in  \mathbf{R}$ ,根式 $\sqrt{a{x}^{2} + {ax} + 1}$ 恒有意义,所以 $a{x}^{2} + {ax} + 1 \geq  0$ 的解集为 $\mathbf{R}$，即不等式 $a{x}^{2} + {ax} + 1 \geq  0$ 在 $\mathbf{R}$ 上恒成立.

        \begin{dingautolist}{172}
            \item 当 $a = 0$ 时, $1 \geq  0$ 恒成立,满足题意
            \item 当 $a \neq  0$ 时, $\left\{  \begin{array}{l} a > 0 \\  \Delta  = {a}^{2} - {4a} \leq  0 \end{array}\right.$ ,解得 $0 < a \leq  4$
        \end{dingautolist}
        综上, $a \in  \{ a \mid  0 \leq  a \leq  4\}$
    \end{jieda}
    \item 已知函数 $f\left( x\right)  = \dfrac{x + 1}{a{x}^{2} - {2ax} + 1}$ 的定义域为 $\mathbf{R}$ ,则实数 $a$ 的取值范围是\blkda[4]{$0 \leq  a < 1$}.
    \begin{jieda}
        函数 $f\left( x\right)  = \dfrac{x + 1}{a{x}^{2} - {2ax} + 1}$ 的定义域为 $\mathbf{R}$，得 $\forall x \in  \mathbf{R},a{x}^{2} - {2ax} + 1 \neq  0$ 恒成立
        \begin{dingautolist}{172}
            \item 当 $a = 0$ 时, $1 \neq  0$ 恒成立
            \item 当 $a \neq  0$ 时, $\Delta  = 4{a}^{2} - {4a} < 0$ ,得 $0 < a < 1$
        \end{dingautolist}
        
        综上,实数 $a$ 的取值范围是 $0 \leq  a < 1$ .
    \end{jieda}
    \item 函数的图象是两条线段 (如图),它的定义域为 $\lbrack  - 1,0) \cup  (0,1\rbrack$ ,则不等式 $f\left( x\right)  - f\left( {-x}\right)  >  - 1$ 的解集为\blkda[4]{$\lbrack  - 1,0) \cup  \left( {\dfrac{1}{2},1}\right\rbrack$}.
    
    \begin{center}
    \includegraphics[width=0.2\textwidth]{bodyxb/src/bo_d2cah3v7aajc738prqf0_5_241_203_248_250_0.jpg}
    \end{center}
    \begin{jieda}
        $f(x) = \left\{\begin{array}{ll}
             x+1, & -1 \leq x < 0 \\
             x-1, & 0 < x \leq 1
        \end{array}\right.$
        读图即可写解集.
    \end{jieda}
    \item 写出一个满足: $f\left( {x + y}\right)  = f\left( x\right)  + f\left( y\right)  + {2xy}$ 的函数解析式为\blkda[4]{$f\left( x\right)  = {x}^{2}$}.
    \begin{jieda}
        $f\left( {x + y}\right)  = f\left( x\right)  + f\left( y\right)  + {2xy}$ 中,令 $x = y = 0$ ,解得 $f\left( 0\right)  = 0$ ,
        
        令 $y =  - x$ 得 $f\left( {x - x}\right)  = f\left( x\right)  + f\left( {-x}\right)  - 2{x}^{2}$ ,故 $f\left( x\right)  + f\left( {-x}\right)  = 2{x}^{2}$ ,
        
        不妨设 $f\left( x\right)  = {x}^{2}$ ,满足要求.
    \end{jieda}
\end{enumerate}
\subsection*{C Level}
\begin{enumerate}
    \item 已知函数 $f\left( x\right)  = \dfrac{x}{{ax} + b}$ ( $a,b$ 为常数,且 $a \neq  0$ ) 满足 $f\left( 2\right)  = 1$ ,方程 $f\left( x\right)  = x$ 有唯一解,则函数 $f\left( x\right)$ 的解析式为\blkda[8]{$f\left( x\right)  = \dfrac{2x}{x + 2}$ 或 $f\left( x\right)  = 1\left( {x \neq  0}\right)$}.
    \begin{jieda}
        由$f(2) = 1$，$\dfrac{2}{2a+b} = 1$，故$2a+b = 2$；$f(x) = x$即$\dfrac{x}{ax+b} = x$，即$\left(\dfrac{1}{ax+b}-1\right)x = 0$，容易判断其有两个解$x_1  = 0, x_2 = \dfrac{1-b}{a}$.
        \begin{dingautolist}{172}
            \item 两根重合，$0 = \dfrac{1-b}{a}$可得$b = 1$. 进一步推得$a = \dfrac{1}{2}$，故$f\left( x\right)  = \dfrac{2x}{x + 2}$
            \item $x_1 = 0$不在定义域内，即$b = 0$，可得$a = 1$，故$f(x) = 1(x \neq 0)$
        \end{dingautolist}
    \end{jieda}

    \item 对于函数 $f\left( x\right)  = \sqrt{a{x}^{2} + {bx}}$ ,其中 $b > 0$ ,若 $f\left( x\right)$ 的定义域与值域相同,则非零实数 $a$ 的值为\blkda[4]{$-4$}.
    \begin{jieda}
        \begin{dingautolist}{172}
            \item 若 $a > 0$，$f\left( x\right)$ 的定义域为: $D = \left( {-\infty , - \dfrac{b}{a}}\right\rbrack   \cup  \lbrack 0, + \infty )$，$f\left( x\right)$ 的值域 $A \subseteq  \lbrack 0, + \infty )$ ,故 $D \neq  A$ ,不合要求.
            \item 若 $a < 0$ ,对于正数 $b,f\left( x\right)$ 的定义域为 ${D} = \left\lbrack  {0, - \dfrac{b}{{a}}}\right\rbrack$ .
            
            由于此时 ${\left\lbrack  f\left( x\right) \right\rbrack  }_{\max } = f\left( {-\dfrac{b}{2a}}\right)  = \dfrac{b}{2\sqrt{-a}}$ ,故函数的值域 $A = \left\lbrack  {0,\dfrac{b}{2\sqrt{-a}}}\right\rbrack$ .
            
            由题意,有 $- \dfrac{b}{a} = \dfrac{b}{2\sqrt{-a}}$ ,由于 $b > 0$ ,所以 $a =  - 4$ .
        \end{dingautolist}
    \end{jieda}
    \item 已知函数 $y = {x}^{2} + {2x}$ 在闭区间 $\left\lbrack  {a,b}\right\rbrack$ 上的值域为 $\left\lbrack  {-1,3}\right\rbrack$ ,则 $a \cdot b$ 的最大值为\blkda[4]{3}.
    \begin{jieda}
    画出函数 $f\left( x\right)  = {x}^{2} + {2x}$ 的图像可知,要使其在闭区间 $\left\lbrack  {a,b}\right\rbrack$ 上的值域为 $\left\lbrack  {-1,3}\right\rbrack$
    \begin{center}
    \includegraphics[width=0.2\textwidth]{bodyxb/src/bo_d2cah3v7aajc738prqf0_7_342_938_350_343_0.jpg}
    \end{center}
    由于有且仅有 $f\left( {-1}\right)  =  - 1$ ,所以 $- 1 \in  \left\lbrack  {a,b}\right\rbrack   \Longrightarrow  a \leq   - 1 \leq  b$
    
    而 $f\left( {-3}\right)  = f\left( 1\right)  = 3$ ,所以有 $\left\lbrack  {a,b}\right\rbrack   \subseteq  \left\lbrack  {-3,1}\right\rbrack  ,a =  - 3$ 或 $b = 1$ ,
    
    因为$a < 0$，所以$a \cdot b$ 的最大值为正值时必有$b < 0$，故对于$a, b$两者均是越小，其绝对值越大，其乘积越大，因此$a \cdot b$ 的最大值为 $\left( {-3}\right)  \times  \left( {-1}\right)  = 3$
    \end{jieda}    
    \item 已知函数 $f\left( {-{x}^{2} + {4x} - 1}\right)$ 的定义域为 $\left\lbrack  {0,m}\right\rbrack$ ,则可求得函数 $f\left( {{2x} - 1}\right)$ 的定义域为 $\left\lbrack  {0,2}\right\rbrack$ ,求实数 ${m}$ 的取值范围.
    \begin{jieda}
        函数 $f\left( {{2x} - 1}\right)$ 的定义域为 $\left\lbrack  {0,2}\right\rbrack$，故$f(x)$的定义域为$[-1, 3]$
        
        令 $t =  - {x}^{2} + {4x} - 1$ , 则 $- 1 \leq  t \leq  3$ ,由题意知,当 $x \in  \left\lbrack  {0,m}\right\rbrack$ 时, $t \in  \left\lbrack  {-1,3}\right\rbrack$ ,作出函数 $t =  - {x}^{2} + {4x} - 1$ 的图象,
        
        \begin{center}
        \includegraphics[width=0.2\textwidth]{bodyxb/src/bo_d2cah3v7aajc738prqf0_6_242_849_341_342_0.jpg}
        \end{center}
        
        如图所示,由图可得,当 $x = 0$ 或 $x = 4$ 时, $t =  - 1$ ,当 $x = 2$ 时, $t = 3,\therefore 2 \leq  m \leq  4$ ,时 $t \in  \left\lbrack  {-1,3}\right\rbrack  ,\therefore$ 实数 $m$ 的取值范围是 $2 \leq  m \leq  4$
    \end{jieda}
\end{enumerate}
\subsection*{D Level}
\begin{enumerate}
    \item 函数 $f\left( x\right)  = \sqrt{1 - {x}^{2}}, - \dfrac{1}{2} \leq  x \leq  \dfrac{1}{2}$ 的图象绕着原点顺时针旋转 $\theta^{\circ} \left( \theta \in [0, 180]\right)$ ,若得到的图象仍是函数图象,则 $\theta$ 可取值的集合为\blkda[4]{$\left\lbrack  {0, 60}\right\rbrack   \cup  \left\lbrack  120, 180 \right\rbrack$}.
    \begin{jieda}
        参考图如下：
        \fourtu{\includegraphics[width=0.9\textwidth]{bodyxb/src/bo_d29ci7n7aajc738p8c80_14_157_1324_312_276_0.jpg}}{\includegraphics[width=0.9\textwidth]{bodyxb/src/bo_d29ci7n7aajc738p8c80_14_494_1324_315_276_0.jpg}}{\includegraphics[width=0.9\textwidth]{bodyxb/src/bo_d29ci7n7aajc738p8c80_14_836_1322_311_275_0.jpg}}{\includegraphics[width=0.9\textwidth]{bodyxb/src/bo_d29ci7n7aajc738p8c80_14_1174_1321_312_275_0.jpg}}
    
        $f\left( x\right)$ 的图象为如图 1 所示的一段弧,弧所在的圆的方程为: ${x}^{2} + {y}^{2} = 1$ ,其中 $A\left( {-\dfrac{1}{2},\dfrac{\sqrt{3}}{2}}\right) ,B\left( {\dfrac{1}{2},\dfrac{\sqrt{3}}{2}}\right)$ . 
        
        在图象绕原点旋转的过程中,当 $B$ 从图 1 的位置旋转到 $(1,0)$ 前(含),如图 2 所示,根据函数的定义,任意竖线与图象最多只有 1 个交点,在这个旋转过程所得的图形均为函数的图象,故 $0 \leq  \theta  \leq  60$ 满足要求. 
        
        在图象绕原点旋转的过程中,当 $B$ 从图 2 的 $(1,0)$ 位置旋转到 $x$ 轴下方,而 $A$ 在 $x$ 轴上,如图 3 所示,根据函数的定义,在这个旋转过程所得的图形不是函数的图象,故 $60 < \theta  < 120$ 不符合. 
        
        当 $A$ 旋转到 $x$ 轴下方时,如图 4 所示,根据函数的定义,在这个旋转过程所得的图形又是函数的图象了,故 $60 \leq  \theta  \leq  180$ 符合.
    \end{jieda}
\end{enumerate}

\section*{函数的值域、函数的图像}
\subsection*{A Level}
\begin{enumerate}
    \item 函数 $y = \left| {x + 1}\right|  - \left| {x - 2}\right|$ 的值域是\blkda[4]{$\left\lbrack  {-3,3}\right\rbrack$}.
    \begin{jieda}
        斜坡函数，根据三角不等式可以直接写出值域.
        
        写成分段函数亦可：$y = \left| {x + 1}\right|  - \left| {x - 2}\right|  = \left\{  \begin{array}{l}  - 3,x <  - 1 \\  {2x} - 1, - 1 \leq  x \leq  2, \\  3,x > 2 \end{array}\right.$
    \end{jieda}
    \item 函数 $y = \dfrac{1 - x}{2 + x}\left( {-1 \leq  x \leq  1}\right)$ 的值域为\blkda[4]{$\left\lbrack  {0,2}\right\rbrack$}.
    \begin{jieda}
        由 $y = \dfrac{1 - x}{2 + x} = \dfrac{3}{2 + x} - 1$ ,
        
        又 $- 1 \leq  x \leq  1$ ,则 $1 \leq  2 + x \leq  3$ ,则 $1 \leq  \dfrac{3}{2 + x} \leq  3$ ,所以 $0 \leq  \dfrac{3}{2 + x} - 1 \leq  2$ ,
        
        故函数 $y = \dfrac{1 - x}{2 + x}\left( {-1 \leq  x \leq  1}\right)$ 的值域为 $\left\lbrack  {0,2}\right\rbrack$ .
    \end{jieda}
    \item 函数 $y = \dfrac{1}{{x}^{2} + {2x} + 2}$ 的值域是\blkda[4]{$(0,1\rbrack$}.
    \begin{jieda}
        由 $f\left( x\right)  = {x}^{2} + {2x} + 2 = {\left( x + 1\right) }^{2} + 1$ ,可得 $f\left( x\right)$ 的最小值为 1
        
        故$y = \dfrac{1}{{x}^{2} + {2x} + 2}$ 的值域为 $(0,1\rbrack$ .
    \end{jieda} 
    \item 函数 $y = \dfrac{{x}^{2} - {4x} + 3}{{x}^{2} - 1}$ 的值域是\blkda[4]{$\{ y \mid  y \neq   - 1$ 且 $y \neq  1\}$}.
    \begin{jieda}
        函数 $y = \dfrac{{x}^{2} - {4x} + 3}{{x}^{2} - 1}$ 中, ${x}^{2} - 1 \neq  0$ ,则 $x \neq   - 1$ 且 $x \neq  1$ ,
        
        于是 $y = \dfrac{\left( {x - 1}\right) \left( {x - 3}\right) }{\left( {x - 1}\right) \left( {x + 1}\right) } = \dfrac{x - 3}{x + 1} = 1 - \dfrac{4}{x + 1}$ ,由 $\dfrac{4}{x + 1} \neq  0$ ,得 $y \neq  1$ ; 由 $x \neq  1$ ,得 $y \neq   - 1$ ,
        
        所以原函数的值域为 $\{ y \mid  y \neq   - 1$ 且 $y \neq  1\}$ .
    \end{jieda}
    \item 求函数 $y = \dfrac{{2x} + 1}{{x}^{2} - {2x} + 2}$ 的值域\blkda[4]{$\left\lbrack  {\dfrac{3 - \sqrt{13}}{2},\dfrac{3 + \sqrt{13}}{2}}\right\rbrack$}.
    \begin{jieda}
        设$t = 2x+1$，则$x = \dfrac{t-1}{2}$

        $y = \dfrac{t}{1 + \left(\dfrac{t-1}{2} - 1\right)^2} = \dfrac{4t}{4 + (t-3)^2} = \dfrac{4t}{t^2-6t+13}$

        \begin{dingautolist}{172}
            \item 当$t = 0$时，$y = 0$
            \item 当$t \neq 0$时，$y = \dfrac{4}{t+\dfrac{13}{t}-6}$

            $t + \dfrac{13}{t} \in (-\infty, -2\sqrt{13}] \cup [2\sqrt{13}, +\infty)$，故$t+\dfrac{13}{t}-6 \in (-\infty, -2\sqrt{13}-6] \cup [2\sqrt{13}-6, +\infty)$，故$y \in \left[\dfrac{3 - \sqrt{13}}{2}, 0\right)\cup\left(0, \dfrac{3 + \sqrt{13}}{2}\right]$
        \end{dingautolist}
        综上，$y \in \left\lbrack  {\dfrac{3 - \sqrt{13}}{2},\dfrac{3 + \sqrt{13}}{2}}\right\rbrack$
    \end{jieda}
    
    \item 函数 $y = \sqrt{-{x}^{2} + {2x} + 2}$ 的值域为\blkda[4]{$\left\lbrack  {0,\sqrt{3}}\right\rbrack$}.
    \begin{jieda}
        令 $\mu \left( x\right)  =  - {x}^{2} + {2x} + 2$ ,则 $0 \leq  \mu \left( x\right)  =  - {x}^{2} + {2x} + 2 =  - {\left( x - 1\right) }^{2} + 3 \leq  3$ ,

        所以 $0 \leq  y \leq  \sqrt{3}$ .
    \end{jieda}
    
    \item 函数 $g\left( x\right)  = x - \sqrt{1 - {2x}}$ 的值域为\blkda[4]{$\left( {-\infty ,\dfrac{1}{2}}\right\rbrack$}
    \begin{jieda}
        函数 $g\left( x\right)$ 的定义域为 $\left( {-\infty ,\dfrac{1}{2}}\right\rbrack$
        
        令 $t = \sqrt{1 - {2x}},x = \dfrac{1 - {t}^{2}}{2}\left( {t \geq  0}\right)$ ,得 $y =  - \dfrac{1}{2}{t}^{2} - t + \dfrac{1}{2}$ ,故 $y \in  \left( {-\infty ,\dfrac{1}{2}}\right\rbrack$ ,
        所以函数 $g\left( x\right)  = x - \sqrt{1 - {2x}}$ 的值域为 $\left( {-\infty ,\dfrac{1}{2}}\right\rbrack$ .
    \end{jieda}
    \item 函数 $y = \sqrt{x - 3} + \sqrt{5 - x}$ 的值域为\blkda[4]{$\left\lbrack  {\sqrt{2},2}\right\rbrack$}
    \begin{jieda}
        由题意得 $\left\{  \begin{array}{l} x - 3 \geq  0 \\  5 - x \geq  0 \end{array}\right.$ ,解得 $3 \leq  x \leq  5$ 
        
        $\because {y}^{2} = 2 + 2\sqrt{\left( {x - 3}\right) \left( {5 - x}\right) } = 2 + 2\sqrt{-{\left( x - 4\right) }^{2} + 1}$，$\therefore 2 \leq  {y}^{2} \leq  4$ 
        
        由 $y$ 的非负性知原函数的值域为 $\left\lbrack  {\sqrt{2},2}\right\rbrack$ .
    \end{jieda}
    \item 函数 $f\left( x\right)  = \sqrt{x - 1} - \sqrt{x - 3}$ 的值域为\blkda[4]{$\left( {0,\sqrt{2}}\right\rbrack$}.
    \begin{jieda}
        $f\left( x\right)  = \sqrt{x - 1} - \sqrt{x - 3} = \dfrac{2}{\sqrt{x - 1} + \sqrt{x - 3}}$ ,显然该函数在其定义域 $\lbrack 3, + \infty )$ 上单调递减, 故 $f\left( x\right)$ 的值域为 $\left( {0,\sqrt{2}}\right\rbrack$ .
    \end{jieda}
    
    \item 函数 $y = \left\{  \begin{array}{l} x + 3,x \leq  1, \\   - x + 5,x > 1 \end{array}\right.$ 的最大值是\blkda[4]{4}.
    \item 函数$y=\dfrac{x-2}{x-1}$的图像是\dotfill(\kh{B})
    \fourtu
    {\includegraphics[totalheight=4.5\baselineskip]{./bodyxb/src/bx01ch03hanshudetuxiangt8a.png}}
    {\includegraphics[totalheight=4.5\baselineskip]{./bodyxb/src/bx01ch03hanshudetuxiangt8b.png}}
    {\includegraphics[totalheight=4.5\baselineskip]{./bodyxb/src/bx01ch03hanshudetuxiangt8c.png}}
    {\includegraphics[totalheight=4.5\baselineskip]{./bodyxb/src/bx01ch03hanshudetuxiangt8d.png}}
    \begin{jieda}
        可以先确定对称中心是$(1, 1)$，然后再代入$x = 0$判断.
    \end{jieda}
    \item 函数 $f\left( x\right)  = \dfrac{x - 1}{x - 2}$ 的对称中心是\blkda[4]{$(2,1)$}.
    \item 已知 $f\left( x\right)$ 的图象恒过点$(1, - 1)$,则函数 $f\left( {x - 3}\right)$ 的图象恒过点\blkda[4]{$(4, -1)$}.
    \begin{jieda}
        因为 $f\left( x\right)$ 的图象恒过点$(1, -1)$,所以当 $x - 3 = 1$ 时, $f\left( {x - 3}\right)  =  - 1$ ,

        即函数 $f\left( {x - 3}\right)$ 的图象恒过点$(4, - 1)$.
    \end{jieda}
    \item 为了得到函数 $y = \dfrac{1}{x - 1}$ 的图象,可以把函数 $y = \dfrac{1}{x}$ 的图象右移\blkda[4]{1}个单位.
    \begin{jieda}
        为了得到函数 $y = \dfrac{1}{x - 1}$ 的图象,根据平移规律,可以把函数 $y = \dfrac{1}{x}$ 的图象右移 1 个单位
    \end{jieda}
    \item 函数 $y = x + \left| x\right|  + 1$ 向右平移 1 个单位的函数解析式为 $y =$ \blkda[4]{$x + \left| {x - 1}\right|$}.
    \begin{jieda}
        函数 $y = x + \left| x\right|  + 1$ 向右平移 1 个单位,

        可得 $y = x + \left| x\right|  + 1 = x - 1 + \left| {x - 1}\right|  + 1 = x + \left| {x - 1}\right|$ .
    \end{jieda}
    \item 把函数 $y = {\left( x - 2\right) }^{2} + 2$ 的图象向左平移 1 个单位,再向上平移 1 个单位,所得图象对应的函数的解析式是\blkda[4]{$y = {\left( x - 1\right) }^{2} + 3$}.
    \begin{jieda}
        把函数 $y = {\left( x - 2\right) }^{2} + 2$ 的图象向左平移 1 个单位,得 $y = {\left\lbrack  \left( x + 1\right)  - 2\right\rbrack  }^{2} + 2 = {\left( x - 1\right) }^{2} +$ 2 的图象,

        再向上平移 1 个单位,即得到 $y = {\left( x - 1\right) }^{2} + 2 + 1 = {\left( x - 1\right) }^{2} + 3$ 的图象
    \end{jieda}
\end{enumerate}


\subsection*{B Level}
\begin{enumerate}
    \item 将函数 $f\left( x\right)  = \dfrac{x - 1}{{2x} - {x}^{2}}$ 的图象向左平移 1 个单位长度,得到函数 $g\left( x\right)$ 的图象,则函数 $g\left( x\right)$ 的图象大致是\dotfill{(\kh{B})}
    \xx{\includegraphics[width=0.39\textwidth]{bodyxb/src/bo_d2cs50f7aajc738pvmbg_0_335_610_315_256_0.jpg}}{\includegraphics[width=0.39\textwidth]{bodyxb/src/bo_d2cs50f7aajc738pvmbg_0_867_607_319_259_0.jpg}}{\includegraphics[width=0.39\textwidth]{bodyxb/src/bo_d2cs50f7aajc738pvmbg_0_579_870_318_257_0.jpg}}{\includegraphics[width=0.39\textwidth]{bodyxb/src/bo_d2cs50f7aajc738pvmbg_0_580_1131_317_257_0.jpg}}
    \begin{jieda}
        考虑$x \Longrightarrow \infty$的符号和极限可排除$A,C$，考虑$x \Longrightarrow 1^-$的符号可选B.
    \end{jieda}
    \item 若函数 $y = f\left( x\right)$ 的值域是 $\left\lbrack  {\dfrac{1}{2},3}\right\rbrack$ ,则函数 $F\left( x\right)  = f\left( {{2x} + 1}\right)  + \dfrac{1}{f\left( {{2x} + 1}\right) }$ 的值域是\blkda[4]{$\left\lbrack  {2,\dfrac{10}{3}}\right\rbrack$}.
    \begin{jieda}
        因函数 $y = f\left( x\right)$ 的值域是 $\left\lbrack  {\dfrac{1}{2},3}\right\rbrack$ ,从而得函数 $t = f\left( {{2x} + 1}\right)$ 值域为 $\left\lbrack  {\dfrac{1}{2},3}\right\rbrack$ ,
        
        函数 $F\left( x\right)$ 变为 $y = t + \dfrac{1}{t},t \in  \left\lbrack  {\dfrac{1}{2},3}\right\rbrack$ ,由对勾函数的性质知值域是 $\left\lbrack  {2,\dfrac{10}{3}}\right\rbrack$ .
    \end{jieda}
    \item 若 $\left| {{2x} + 3}\right|  + \left| {{2x} - 1}\right|  \geq  a$ 恒成立,则 $a$ 的取值范围为\blkda[4]{$( - \infty ,4\rbrack$}
    \begin{jieda}
        因为 $a \leq  \left| {{2x} + 3}\right|  + \left| {{2x} - 1}\right|$ 恒成立，所以 $a \leq  {\left( \left| 2x + 3\right|  + \left| 2x - 1\right| \right) }_{\min }$ .
        因为 $\left| {{2x} + 3}\right|  + \left| {{2x} - 1}\right|  \geq  \left| {{2x} + 3 - \left( {{2x} - 1}\right) }\right|  = 4$ ,当且仅当 $- \dfrac{3}{2} \leq  x \leq  \dfrac{1}{2}$ 时等号成立,
        所以 ${\left( \left| 2x + 3\right|  + \left| 2x - 1\right| \right) }_{\min } = 4$ .
        所以 $a \leq  4$ .
    \end{jieda}
    \item 若不等式 $\left| {x - 4}\right|  - \left| {x - 3}\right|  \geq  a$ 对一切 $x \in  \mathbf{R}$ 恒成立,则实数 $a$ 的取值范围是\blkda[4]{$( - \infty , - 1\rbrack$}.
    \begin{jieda}
        令 $f\left( x\right)  = \left| {x - 4}\right|  - \left| {x - 3}\right|$ ,则 g那就三角不等式可知 $f{\left( x\right) }_{\min } =  - 1$
        
        因为不等式 $\left| {x - 4}\right|  - \left| {x - 3}\right|  \geq  a$ 对一切 $x \in  \mathbf{R}$ 恒成立，所以 $a \leq   - 1$ ,即实数 $a$ 的取值范围是 $( - \infty , - 1\rbrack$ .
    \end{jieda}
    \item 若 $x + \left| {x - {2c}}\right|  > 1$ 的解集为 $R$ ,则实数 $c$ 的范围为\blkda[4]{$\left( {\dfrac{1}{2}, + \infty }\right)$}.
    \begin{jieda}
        令函数 $f\left( x\right)  = x + \left| {x - {2c}}\right|  = \left\{  \begin{array}{l} {2x} - {2c}\left( {x \geq  {2c}}\right) \\  {2c}\left( {x < {2c}}\right)  \end{array}\right.$ , $\therefore f\left( x\right)$ 的最小值为 ${2c}$ . 由题意得 ${2c} > 1,\therefore c > \dfrac{1}{2}$  $\therefore$ 实数 $c$ 的范围为 $\left( {\dfrac{1}{2}, + \infty }\right)$
    \end{jieda}
    \item 函数 $y = \left| {x - 1}\right|  + \left| {x - 2}\right|  + \left| {x - 3}\right|$ 的最小值为\blkda[4]{2}. 
    \begin{jieda}
        由题意，$y = \left| {x - 1}\right|  + \left| {x - 2}\right|  + \left| {x - 3}\right|  = \left\{  \begin{array}{l} 6 - {3x},x \leq  1 \\  4 - x,1 < x \leq  2 \\  x,2 < x \leq  3 \\  {3x} - 6,x > 3 \end{array}\right.$
        
        图象如下图所示:
        
        \begin{center}
        \includegraphics[width=0.25\textwidth]{bodyxb/src/bo_d2cs50f7aajc738pvmbg_3_246_1458_446_552_0.jpg}
        \end{center}
        
        数形结合可知,当 $x = 2$ 时,函数取得最小值 $y = 2$
    \end{jieda}
    \item 已知函数 $f\left( x\right)  = \left\{  \begin{array}{l} {x}^{2},x \geq  1 \\  x + \dfrac{4}{x} - 5,0 < x < 1 \end{array}\right.$ ,则 $f\left( x\right)$ 的值域为\blkda[4]{$\left( {0, + \infty }\right)$}.
    \begin{jieda}
        因为 $f\left( x\right)  = \left\{  \begin{array}{l} {x}^{2},x \geq  1 \\  x + \dfrac{4}{x} - 5,0 < x < 1 \end{array}\right.$ ,
        
        当 $x \geq  1$ 时, $f\left( x\right)  = {x}^{2} \geq  1$ ,
        
        当 $0 < x < 1$ 时,函数 $f\left( x\right)  = x + \dfrac{4}{x} - 5$ 单调递减,故 $f\left( x\right)  > f\left( 1\right)  = 0$ ,
        
        综上,函数 $f\left( x\right)$ 的值域为 $\left( {0, + \infty }\right)$
    \end{jieda}
    \item 已知函数 $f\left( x\right)  \in  \left\lbrack  {1,5}\right\rbrack$ ,则函数 $g\left( x\right)  = f\left( x\right)  + \dfrac{1}{f\left( x\right) }$ 的值域为\blkda[4]{$\left\lbrack  {2,\dfrac{26}{5}}\right\rbrack$}.
    \item 已知函数$y=f(\abs{x})$的图像如左下图所示, 则函数$y=f(x)$的图像不可能是\dotfill(\kh{B})
    \begin{minipage}{\linewidth}
    \centering
    \includegraphics[height=5.0\baselineskip]{./bodyxb/src/bx01ch03hanshudetuxiangt9start.png}\qquad
    
\includegraphics[height=5.0\baselineskip]{./bodyxb/src/bx01ch03hanshudetuxiangt9a.png}\qquad
    
\includegraphics[height=5.0\baselineskip]{./bodyxb/src/bx01ch03hanshudetuxiangt9b.png}\qquad
    
\includegraphics[height=5.0\baselineskip]{./bodyxb/src/bx01ch03hanshudetuxiangt9c.png}\qquad
    
\includegraphics[height=5.0\baselineskip]{./bodyxb/src/bx01ch03hanshudetuxiangt9d.png}
    \end{minipage}
    \begin{jieda}
        $f(|x|)$是将$f(x)$在$y$轴右侧图像翻到左侧
    \end{jieda}
    
    \item 已知函数$f(x)=ax^3+bx^2+cx+d$的图像如下图, 则\dotfill(\kh{A})
    \xx{$b\in(-\infty,0)$}{$b\in(0,1)$}{$b\in(1,2)$}{$b\in(2,+\infty)$}
    \begin{minipage}{\linewidth}
    \centering
    \includegraphics[height=5\baselineskip]{./bodyxb/src/bx01ch03hanshudetuxiangt11.png}
    \end{minipage}
    \begin{jieda}
        首先可以判断出$a > 0$，然后根据三个零点列方程可以得到$b = -3a$或者将$f(x)$写成因式相乘的形式：$f(x) = ax(x-1)(x-2)$，也可以得到$b = -3a$
    \end{jieda}
    
\end{enumerate}


\subsection*{C Level}
\begin{enumerate}
    \item 已知函数 $y = \dfrac{{ax} + b}{{x}^{2} + 1}$ 的值域为 $\left\lbrack  {-1,4}\right\rbrack$ ,则常数 $a + b =$ \blkda[4]{$7$ 或 $-1$}.
    \begin{jieda}
        因为 $y = \dfrac{{ax} + b}{{x}^{2} + 1}$ ,所以 ${x}^{2}y - {ax} + y - b = 0$ ,
        
        $\Delta  = {a}^{2} - {4y}\left( {y - b}\right)  \geq  0$ ,即 $4{y}^{2} - {4by} - {a}^{2} \leq  0$ ,
        
        因为函数 $y = \dfrac{{ax} + b}{{x}^{2} + 1}$ 的值域为 $\left\lbrack  {-1,4}\right\rbrack$ ,
        
        所以 ${y}_{1} =  - 1,{y}_{2} = 4$ 是方程 $4{y}^{2} - {4by} - {a}^{2} = 0$ 的两个根,
        
        所以 $- 1 + 4 = b, - 1 \times  4 =  - \dfrac{{a}^{2}}{4}$ ,
        
        解得 $a = 4,b = 3$ 或 $a =  - 4,b = 3$ ,所以 $a + b = 7$ 或 $-1$ .
    \end{jieda}
    \item 函数 $y = \sqrt{{x}^{2} - {2x} + 5} + \sqrt{{x}^{2} - {4x} + {13}}$ 的值域为\blkda[4]{$\left\lbrack  {\sqrt{26}, + \infty }\right)$}
    \begin{jieda}
        由题设 $y = \sqrt{{\left( x - 1\right) }^{2} + {\left( 0 - 2\right) }^{2}} + \sqrt{{\left( x - 2\right) }^{2} + {\left( 0 + 3\right) }^{2}}$ ,
        所以所求值域化为求 $x$ 轴上点 $C$ 到 $A\left( {1,2}\right)$ 与 $B\left( {2, - 3}\right)$ 距离和的范围,
        由图知： $\left| {CA}\right|  + \left| {CB}\right|  \geq  \left| {AB}\right|  = \sqrt{26}$ ,
        所以函数值域为 $\left\lbrack  {\sqrt{26}, + \infty }\right)$ .
    \end{jieda}
    \item 若关于 $x$ 的不等式 ${x}^{2} + \left| {{2x} + m}\right|  \leq  4$ 在 $x \in  \lbrack 0, + \infty )$ 有解,则实数 $m$ 的取值范围是 \blkda[4]{$\left\lbrack  {-5,4}\right\rbrack$}
    \begin{jieda}
        原不等式可变形为 $\left| {{2x} + m}\right|  \leq  4 - {x}^{2}$ ,作出函数$f\left( x\right)  = \left| {{2x} + m}\right|$ 与 $g\left( x\right)  = 4 - {x}^{2}$ 的图像,如图所示。
        \begin{center}
            \includegraphics[width=0.5\linewidth]{bodyxb//src/yikeyilian2-12.png}
        \end{center}
        由题意,在$x \geq  0$ 时,至少有一点满足 $f\left( x\right)  \leq  g\left( x\right)$ ,当 $y =  - {2x} - m$ 与$g\left( x\right)  = 4 - {x}^{2}$ 相切时, $- {2x} - m = 4 - {x}^{2}$ ,即 ${x}^{2} - {2x} - 4 - m = 0$
        
        由 $\Delta  = 4 + 4\left( {4 + m}\right)  = 0$ ,得 $m =  - 5$ 。当 $y = {2x} + m$ 过点 $M(0$ ,4) 时, $m = 4$ ,所以根据图象可知 $- 5 \leq  m \leq  4$
    \end{jieda}
    
    \item 已知函数$y=f(x)$与$y=g(x)$在$[-2,2]$的图像如图所示: \\
    \begin{minipage}{\linewidth}
    \centering
    
\includegraphics[totalheight=8\baselineskip]{./bodyxb/src/bx01ch03hanshudeyunsuan02t1401.png}
    \qquad\qquad
    
\includegraphics[totalheight=8\baselineskip]{./bodyxb/src/bx01ch03hanshudeyunsuan02t1402.png}
    \end{minipage}
    
    判断下列命题的真假, 并说明理由: 
    \begin{enumerate}
    \item 方程$f(g(x))=0$有且仅有6个根; 
    \item 方程$g(f(x))=0$有且仅有3个根; 
    \item 方程$f(f(x))=0$有且仅有5个根; 
    \item 方程$g(g(x))=0$有且仅有4个根. 
    \end{enumerate}
    \begin{jieda}
    $f(x)=0$的根可取近似值为$0,\pm1.5$; $g(x)=0$的根可取近似值为$-1.8,0.4$. 
    \begin{enumerate}
    \item $f(g(x))=0\iff g(x)=0,\pm1.5$, 结合$g(x)$的图像可得解的个数为$3+3=6$个. 
    \item $g(f(x))=0\iff f(x)=-1.8,0.4$, 结合$f(x)$的图像可得解的个数为$1+3=4$个. 
    \item $f(f(x))=0\iff f(x)=0,\pm1.5$, 结合$f(x)$的图像可得解的个数为$3+2=5$个. 
    \item $g(g(x))=0\iff g(x)=-1.8,0.4$, 结合$g(x)$的图像可得解的个数为$2+2=4$个. 
    \end{enumerate}
    综上, 只有(2)是假命题, 其他都是真命题. 
    \end{jieda}
    
    \item 任取$x_1,x_2\in[a,b], x_1\neq x_2$, 若$f\left(\dfrac{x_1+x_2}{2}\right)>\dfrac{f(x_1)+f(x_2)}{2}$, 则称$f(x)$是$[a,b]$上的凸函数. 
    下列图像中, 可以是凸函数图像的是\dotfill(\kh{C})
    \begin{center}
        \includegraphics[width=0.9\linewidth]{bodyxb//src/ec705b924905107021a3cb8a69dcf8ac.jpg}
    \end{center}
\end{enumerate}
\clearpage
\subsection*{D Level}
\begin{enumerate}
    \item 设 $f\left( x\right)  = \left| {x + \dfrac{1}{x}}\right|  + \left| {1 - \dfrac{1}{x}}\right|$ ,若存在 $a \in  \mathbf{R}$ 使得关于 $x$ 的方程 ${\left( f\left( x\right) \right) }^{2} + {af}\left( x\right)  + b = 0$ 恰有六个解,则 $b$ 的取值范围是\blkda[4]{$\left( {4\sqrt{2} + 2, + \infty }\right)$}.
    \begin{jieda}
        $f\left( x\right)  = \left\{\begin{array}{ll}
            - x - \dfrac{2}{x} + 1, & x < 0 \\
            x + \dfrac{2}{x} - 1, & 0 < x < 1 \\
            x + 1,  & x \geq 1
        \end{array}\right. $，则 $f\left( x\right)$ 图像如图所示:
        
        \begin{center}
        \includegraphics[width=0.35\textwidth]{bodyxb/src/bo_d2bm22f7aajc738poiq0_3_244_205_1094_752_0.jpg}
        \end{center}
        由图可知：
        \begin{dingautolist}{172}
            \item 当 $x < 0$ 时, $f\left( x\right)  =  - x - \dfrac{2}{x} + 1 \geq  2\sqrt{2} + 1$
            \item 当 $x > 0$ 时, $f{\left( x\right) }_{\min } = 2$ 
        \end{dingautolist}
        
        令 $t = f\left( x\right)$ ,则 ${t}^{2} + {at} + b = 0$，故关于 $x$ 的方程 ${\left( f\left( x\right) \right) }^{2} + {af}\left( x\right)  + b = 0$ 恰有六个解应满足关于 $t$ 的方程 ${t}^{2} + {at} + b = 0$ 有两个解 ${t}_{1}\text{ 、 }{t}_{2}$（${t}_{1} < {t}_{2}$），且${t}_{1} \in  \left( {2,2\sqrt{2} + 1}\right) ,{t}_{2} \in  \left( {2\sqrt{2} + 1, + \infty }\right)$ ,
        
        令 $g\left( t\right)  = {t}^{2} + {at} + b$，则应有 $\left\{  \begin{array}{l} g\left( 2\right)  > 0 \\  g\left( {2\sqrt{2} + 1}\right) < 0 \end{array}\right.$，即$\left\{  \begin{array}{l} 4+2a+b  > 0 \\  
        9+4\sqrt{2}+(2\sqrt{2}+1)a+b < 0 \end{array}\right.$
        
        要存在 $a$ 满足条件，不等式变形有$\left\{\begin{array}{l}
             a > \dfrac{-b - 4}{2}  \\
             a < \dfrac{-b - 9 - 4\sqrt{2}}{2\sqrt{2} + 1} 
        \end{array}\right.$，则应满足 $\dfrac{-b - 4}{2} < \dfrac{-b - 9 - 4\sqrt{2}}{2\sqrt{2} + 1}$，解得 $b > 4\sqrt{2} + 2$ .
    \end{jieda}
    \item 设 $a \in  \mathbf{R}$ ,若不等式 $\left| {{x}^{2} + \dfrac{1}{x}}\right|  + \left| {{x}^{2} - \dfrac{1}{x}}\right|  + {ax} \geq  {4x} - 8$ 恒成立,则实数 $a$ 的取值范围是\dotfill{(\kh{D})}
    \xx{$\left\lbrack  {-2,{12}}\right\rbrack$}{$\left\lbrack  {-2,{10}}\right\rbrack$}{$\left\lbrack  {-4,4}\right\rbrack$}{$\left\lbrack  {-4,{12}}\right\rbrack$}
    \begin{jieda}
        $\left| {{x}^{2} + \dfrac{1}{x}}\right|  + \left| {{x}^{2} - \dfrac{1}{x}}\right|  + {ax} \geq  {4x} - 8$ 恒成立,
        即为 $\left| {{x^2} + \dfrac{1}{x}}\right|  + \left| {{x^2} - \dfrac{1}{x}}\right|  + 8 \geq  \left( {4 - a}\right) x$ 恒成立
        \begin{dingautolist}{172}
            \item 当 $x > 0$ 时,可得 $4 - a \leq  \left| {x + \dfrac{1}{{x}^{2}}}\right|  + \left| {x - \dfrac{1}{{x}^{2}}}\right|  + \dfrac{8}{x}$ 的最小值
            
            由 $\left| {x + \dfrac{1}{{x}^{2}}}\right|  + \left| {x - \dfrac{1}{{x}^{2}}}\right|  + \dfrac{8}{x} \geq  \left| {x + \dfrac{1}{{x}^{2}} + x - \dfrac{1}{{x}^{2}}}\right|  + \dfrac{8}{x} = {2x} + \dfrac{8}{x} \geq  2\sqrt{{2x} \cdot  \dfrac{8}{x}} = 8$
            
            故当且仅当 $x = 2$ 原式取得最小值 8 ,即有 $4 - a \leq  8$ ,则 $a \geq   - 4$
            \item 当 $x < 0$ 时,可得 $4 - a \geq   - \left\lbrack  {\left| {x + \dfrac{1}{{x}^{2}}}\right|  + \left| {x - \dfrac{1}{{x}^{2}}}\right|  - \dfrac{8}{x}}\right\rbrack$ 的最大值
            
            由 $\left| {x + \dfrac{1}{{x}^{2}}}\right|  + \left| {x - \dfrac{1}{{x}^{2}}}\right|   - \dfrac{8}{x} \geq \left| {x + \dfrac{1}{{x}^{2}} + x - \dfrac{1}{{x}^{2}}}\right|  - \dfrac{8}{x} =   - {2x} - \dfrac{8}{x} \geq  2\sqrt{{2x} \cdot  \dfrac{8}{x}} = 8$
            
            故当且仅当 $x =  - 2$ 原式取得最大值$-8$,即有 $4 - a \geq   - 8$ ,则 $a \leq  {12}$
        \end{dingautolist}

        综上可得 $- 4 \leq  a \leq  {12}$ .
    \end{jieda}
\end{enumerate}
\section*{单调性、奇偶性}
\subsection*{A Level}
\begin{enumerate}
    \item 函数 $y = f\left( x\right)$ 在区间 $\left\lbrack  {-2,2}\right\rbrack$ 上的图象如图所示,则此函数的增区间是\blkda[4]{$\left\lbrack  {-2,1}\right\rbrack$}.
    
    \begin{center}
    \includegraphics[width=0.3\textwidth]{bodyxb/src/bo_d2dc2ev7aajc738q2ikg_0_242_557_356_293_0.jpg}
    \end{center}
    \item 函数 $f\left( x\right)  = \left| x\right|$ 的单调递减区间为\blkda[4]{$( - \infty ,0\rbrack$}
    
    \item 函数 $f\left( x\right)  = \sqrt{{x}^{2} + {4x} - {12}}$ 的单调减区间为\blkda[4]{$\left( {-\infty , - 6}\right]$}.

    \begin{jieda}
        函数 $f\left( x\right)  = \sqrt{{x}^{2} + {4x} - {12}}$ 是由函数 $g\left( u\right)  = \sqrt{u}$ 和 $u\left( x\right)  = {x}^{2} + {4x} - {12}$ 组成的复合函数, $\because {x}^{2} + {4x} - {12} \geq  0$ ,解得 $x \leq   - 6$ 或 $x \geq  2$ ,
        
        $\therefore$ 函数 $y = f\left( x\right)$ 的定义域是 $\{ x \mid  x \leq   - 6$ 或 $x \geq  2\}$ ,
        
        因为函数 $u\left( x\right)  = {x}^{2} + {4x} - {12}$ 在 $( - \infty , - 6\rbrack$ 单调递减,在 $\lbrack 2, + \infty )$ 单调递增,
        
        而 $g\left( u\right)  = \sqrt{u}$ 在 $\lbrack 0, + \infty )$ 上单调递增,
        
        由复合函数单调性的“同增异减”,可得函数 $f\left( x\right)$ 的单调减区间 $( - \infty , - 6\rbrack$ .
    \end{jieda}
    \item 函数 $y = {x}^{2} - {2mx} + 3$ 在区间 $\left\lbrack  {1,3}\right\rbrack$ 上具有单调性,则 $m$ 的取值范围为\blkda[4]{$m \leq  1$ 或 $m \geq  3$}.
    \begin{jieda}
        二次函数 $y = {x}^{2} - {2mx} + 3$ 的对称轴为 $x = m$ ,因函数 $y = {x}^{2} - {2mx} + 3$ 在区间 $\left\lbrack  {1,3}\right\rbrack$ 上具有单调性，所以 $m \leq  1$ 或 $m \geq  3$
    \end{jieda}
    \item 已知 $f\left( x\right)$ 是定义在 $\mathbf{R}$ 上的严格增函数,且 $f\left( {{x}^{2} - 2}\right)  < f\left( {-x}\right)$ ,则 $x$ 的取值范围是\blkda[4]{$(-2,1)$}.
    \begin{jieda}
        因 $f\left( x\right)$ 是定义在 $\mathbf{R}$ 上的严格增函数,故由 $f\left( {{x}^{2} - 2}\right)  < f\left( {-x}\right)$ 可得
        
        ${x}^{2} - 2 <  - x$ ,即 ${x}^{2} + x - 2 < 0$ ,解得 $- 2 < x < 1$ .
    \end{jieda}
    \item 若 $f\left( x\right)$ 是 $\mathbf{R}$ 上单调递减的函数,且 $f\left\lbrack  {f\left( x\right) }\right\rbrack   = {4x} - 1$ ,则 $f\left( x\right)  =$ \blkda[4]{$- {2x} + 1$}.
    \begin{jieda}
        因为 $f\left( x\right)$ 是 $\mathbf{R}$ 上单调递减的函数,所以可设 $f\left( x\right)  = {kx} + b\left( {k < 0}\right)$ ,

        所以 $f\left\lbrack  {f\left( x\right) }\right\rbrack   = k \cdot  f\left( x\right)  + b = k\left( {{kx} + b}\right)  + b = {k}^{2}x + {kb} + b$ ,
        
        又因为 $f\left\lbrack  {f\left( x\right) }\right\rbrack   = {4x} - 1$ ,所以 ${k}^{2}x + {kb} + b = {4x} - 1$ 恒成立, 所以 $\left\{  \begin{array}{l} {k}^{2} = 4 \\  {kb} + b =  - 1 \end{array}\right.$ ,因为 $k < 0$ ,所以 $k =  - 2,b = 1$ . 所以 $f\left( x\right)  =  - {2x} + 1$ .
    \end{jieda}
    \item 函数 $y = \left\{  \begin{array}{l} x\left( {1 + x}\right) ,x \geq  0 \\  x\left( {1 - x}\right) ,x < 0 \end{array}\right.$ 的奇偶性为\blkda[4]{奇函数}.
    \begin{jieda}
        作出函数图像如下图所示:
        
        \begin{center}
        \includegraphics[width=0.2\textwidth]{bodyxb/src/bo_d2dc2ev7aajc738q2ikg_2_245_788_336_333_0.jpg}
        \end{center}
    \end{jieda}
    \item 已知定义在 $\left\lbrack  {a - 2,{a}^{2}}\right\rbrack$ 上的函数 $y = f\left( x\right)$ 是偶函数,则实数 $a$ 的值为\blkda[4]{$-2$ 或 $1$}.
    \begin{jieda}
        由题设,函数 $y = f\left( x\right)$ 的定义域关于原点对称，即 ${a}^{2} + a - 2 = 0$，解得 $a =  - 2$ 或 $a = 1$
    \end{jieda}
    \item 若函数 $y = f\left( x\right)$ 是偶函数,且在 $\left( {-\infty ,0}\right)$ 上是严格增函数,则 $f\left( \pi \right) \text{ 、 }f\left( {-3}\right) \text{ 、 }f\left( {-\sqrt{3}}\right)$ 的大小关系是\blkda[4]{$f\left( \pi \right)  < f\left( {-3}\right)  < f\left( {-\sqrt{3}}\right)$}.
    \begin{jieda}
        因为函数 $y = f\left( x\right)$ 是偶函数,且在 $\left( {-\infty ,0}\right)$ 上是严格增函数,
        
        所以 $f\left( \pi \right)  = f\left( {-\pi }\right)  < f\left( {-3}\right)  < f\left( {-\sqrt{3}}\right)$ .
    \end{jieda}
    
    \item 已知函数 $f\left( x\right)  = \left\{  \begin{array}{ll}  - {x}^{2} + {2x}, & x > 0, \\  0, & x = 0, \\  {x}^{2} + {mx}, & x < 0 \end{array}\right.$ 是奇函数,则 $m =$ \blkda{2}
    \begin{jieda}
        $f(1) = f(-1)$即可判断
    \end{jieda}
    
    \item 若奇函数在 $\left( {-\infty ,0}\right)$ 上是增函数,且 $f\left( {-1}\right)  = 0$ ,则使得 $f\left( x\right)  < 0$ 的取值范围是\blkda[4]{$\left( {-\infty , - 1}\right)  \cup  \left( {0,1}\right)$}.
    \begin{jieda}
        因为 $f\left( x\right)$ 在 $\left( {-\infty ,0}\right)$ 上是增函数, $f\left( {-1}\right)  = 0$ ,
        
        所以 $x <  - 1$ 时, $f\left( x\right)  < 0$ ,
        
        又因为 $f\left( x\right)$ 是奇函数,所以 $f\left( x\right)$ 在 $\left( {0, + \infty }\right)$ 上也是增函数, $f\left( 1\right)  =  - f\left( {-1}\right)  = 0$ ,
        
        所以 $0 < x < 1$ 时, $f\left( x\right)  < 0$ ,综上, $f\left( x\right)  < 0$ 的解集为 $\left( {-\infty ,0}\right)  \cup  \left( {0,1}\right)$ .
    \end{jieda}
    \item 已知函数 $f\left( x\right)$ 是定义在 $\mathbf{R}$ 上的奇函数,当 $x \in  \left( {-\infty ,0}\right)$ 时, $f\left( x\right)  = 2{x}^{3} + {x}^{2}$ ,则$f\left( 2\right)  =$ \blkda[4]{12}.
    \begin{jieda}
        函数 $f\left( x\right)$ 是定义在 $R$ 上的奇函数, $f\left( {-x}\right)  =  - f\left( x\right)$ ,则 $f\left( x\right)  =  - f\left( {-x}\right)$，$f\left( 2\right)  =  - f\left( {-2}\right)  =  - \left\lbrack  {2 \times  {\left( -2\right) }^{3} + {\left( -2\right) }^{2}}\right\rbrack  = {12}.$
    \end{jieda}
\end{enumerate}
\clearpage
\subsection*{B Level}
\begin{enumerate}
    \item 函数 $f\left( x\right)  = \dfrac{1}{\sqrt{-{x}^{2} + {4x} + 5}}$ 的单调增区间是\blkda[4]{$\lbrack 2,5)$}.
    \begin{jieda}
        由题意可知 $- {x}^{2} + {4x} + 5 > 0$ ,解得 $- 1 < x < 5$ ,即函数定义域为(-1,5),
        
        易知函数由 $y = \dfrac{1}{\sqrt{t}},t\left( x\right)  =  - {x}^{2} + {4x} + 5$ 复合而成,
        
        且 $y = \dfrac{1}{\sqrt{t}}$ 在 $\left( {0, + \infty }\right)$ 单调递减, $t\left( x\right)  =  - {x}^{2} + {4x} + 5$ 在(-1,2)单调递增,在 $\lbrack 2,5)$ 上单调递减; 利用复合函数单调性可得 $f\left( x\right)  = \dfrac{1}{\sqrt{-{x}^{2} + {4x} + 5}}$ 的单调增区间是 $\lbrack 2,5)$
    \end{jieda}
    
    \item 已知函数 $f\left( x\right)  = \sqrt{{x}^{2} + {4ax} + 7}$ 在 $\left\lbrack  {1,2}\right\rbrack$ 上是减函数,则 $a$ 的取值范围是\blkda[4]{$\left\lbrack  {-\dfrac{11}{8}, - 1}\right\rbrack$}.
    \begin{jieda}
        因为函数 $y = \sqrt{x}$ 在 $\lbrack 0, + \infty )$ 上单调递增,所以函数 $y = {x}^{2} + {4ax} + 7$ 在 $\left\lbrack  {1,2}\right\rbrack$ 上单调递减.
        
        由 $\left\{  \begin{matrix}  - \dfrac{4a}{2} \geq  2 \\  4 + {8a} + 7 \geq  0 \end{matrix}\right.$ 解得 $- \dfrac{11}{8} \leq  a \leq   - 1$ .
    \end{jieda}
    \item 若函数 $f\left( x\right)  = \left\{  \begin{array}{l} \left( {2 + {3a}}\right) x - 1,x > 1 \\  {4ax} - {x}^{2},x \leq  1 \end{array}\right.$ 在 $\mathbf{R}$ 上单调递增,则实数 $a$ 的取值范围是\blkda[4]{$\left\lbrack  {\dfrac{1}{2},2}\right\rbrack$}
    \begin{jieda}
        根据题意得 $\left\{  \begin{array}{l} 2 + {3a} \geq 0 \\  {2a} \geq  1 \\  2 + {3a} - 1 \geq  {4a} - 1 \end{array}\right.$ ,解得 $\dfrac{1}{2} \leq  a \leq  2$，所以实数 $a$ 的取值范围是 $\left\lbrack  {\dfrac{1}{2},2}\right\rbrack$ .
    \end{jieda}
    \item 若函数$y=\dfrac{ax}{2a^2-5a+2}$在$\bR$上为严格减函数, 则$a$的取值范围是\blkda[3]{$(-\infty, 0) \cup \left(\dfrac{1}{2}, 2\right)$}. 
        \begin{jieda}
        $\begin{cases}
        a < 0\\
        2a^2-5a+2>0
        \end{cases}$
        或$\begin{cases}
        a > 0 \\
        2a^2-5a+2 < 0
        \end{cases}$, 
        解得$a \in (-\infty, 0) \cup \left(\dfrac{1}{2}, 2\right)$. 
        \end{jieda}
    \item 已知函数 $f\left( x\right)$ 对一切实数 $x$ 都满足 $f\left( x\right)  + f\left( {-x}\right)  = 0$ ,且当 $x < 0$ 时, $f\left( x\right)  = 2{x}^{2} - x + 1$，则 $f\left( x\right)  =$ \blkda[4]{$\left\{  \begin{array}{l}  - 2{x}^{2} - x - 1,x > 0 \\  0,x = 0 \\  2{x}^{2} - x + 1,x < 0 \end{array}\right.$}.
    \begin{jieda}
        函数 $f\left( x\right)$ 对一切实数 $x$ 都满足 $f\left( x\right)  + f\left( {-x}\right)  = 0$ ,
        
        所以 $f\left( 0\right)  = 0$ ,
        
        设 $x > 0$ ,则 $- x < 0,f\left( {-x}\right)  = 2{x}^{2} + x + 1$ ,
        
        又因为 $f\left( x\right)  + f\left( {-x}\right)  = 0$ ,即 $f\left( x\right)  =  - f\left( {-x}\right)$ ,
        
        所以 $f\left( x\right)  =  - 2{x}^{2} - x - 1$
        
        所以 $f\left( x\right)  = \left\{  \begin{array}{l}  - 2{x}^{2} - x - 1,x > 0 \\  0,x = 0 \\  2{x}^{2} - x + 1,x < 0 \end{array}\right.$ .        
    \end{jieda}
    
    \item 已知函数 $f\left( x\right)  = \dfrac{{ax} + b}{{x}^{2} + 1}$ 是奇函数,且 $f\left( 2\right)  = \dfrac{2}{5}$ ,则 $f\left( \dfrac{1}{2}\right)  =$ \blkda[4]{$\dfrac{2}{5}$}.
    \begin{jieda}
        $f\left( x\right)  = \dfrac{{ax} + b}{{x}^{2} + 1}$ 的定义域为 $\mathbf{R}$ ,而 $f\left( x\right)$ 为奇函数,
        
        故 $f\left( 0\right)  = 0 = b$ ,而 $f\left( 2\right)  = \dfrac{2}{5}$ ,故 $\dfrac{2a}{4 + 1} = \dfrac{2}{5}$ ,故 $a = 1$ ,
        
        所以 $f\left( x\right)  = \dfrac{x}{{x}^{2} + 1}$ ,此时 $f\left( {-x}\right)  = \dfrac{-x}{{x}^{2} + 1} =  - f\left( x\right)$ ,故 $f\left( x\right)$ 为奇函数,
        
        故 $f\left( \dfrac{1}{2}\right)  = \dfrac{\dfrac{1}{2}}{\dfrac{5}{4}} = \dfrac{2}{5}$
    \end{jieda}
    
    \item 已知函数 $f\left( x\right)  = a{x}^{7} + b{x}^{3} + {x}^{2} + {cx} - {2023}$ ,且 $f\left( {10}\right)  = 6$ ,则 $f\left( {-{10}}\right)  =$ \blkda[4]{-3852}.
    \begin{jieda}
        由 $f\left( x\right)  = a{x}^{7} + b{x}^{3} + {x}^{2} + {cx} - {2023}$ ,得 $a{x}^{7} + b{x}^{3} + {cx} = f\left( x\right)  - {x}^{2} + {2023}$ ,
        
        构建函数 $h\left( x\right)  = a{x}^{7} + b{x}^{3} + {cx} = f\left( x\right)  - {x}^{2} + {2023}$ ,定义域为 $\mathbf{R}$ ,
        
        则 $h\left( {-x}\right)  = a{\left( -x\right) }^{7} + b{\left( -x\right) }^{3} + c\left( {-x}\right)  =  - \left( {a{x}^{7} + b{x}^{3} + {cx}}\right)  =  - h\left( x\right)$ ,即 $h\left( x\right)$ 是奇函数,
        
        于是 $h\left( {-{10}}\right)  =  - h\left( {10}\right)$ ,所以 $f\left( {-{10}}\right)  - {\left( -{10}\right) }^{2} + {2023} =  - \left\lbrack  {f\left( {10}\right)  - {10}^{2} + {2023}}\right\rbrack$ ,
        
        可得 $f\left( {-{10}}\right)  =  - f\left( {10}\right)  - {3846}$ ,
        
        又 $f\left( {10}\right)  = 6$ ,因此 $f\left( {-{10}}\right)  =  - {3852}$
    \end{jieda}
    
    \item 函数 $f\left( x\right)  = \dfrac{{\left( x + 1\right) }^{2} + {x}^{3}}{{x}^{2} + 1}$ 在 $\left\lbrack  {-{2020},{2020}}\right\rbrack$ 上的最大值和最小值分别为 $M,m$ ,则 $M + m =$ \blkda[4]{2}.
    \begin{jieda}
        因为 $f\left( x\right)  = \dfrac{{\left( x + 1\right) }^{2} + {x}^{3}}{{x}^{2} + 1} = \dfrac{\left( {{x}^{3} + {2x}}\right)  + \left( {{x}^{2} + 1}\right) }{{x}^{2} + 1} = \dfrac{{x}^{3} + {2x}}{{x}^{2} + 1} + 1$ ,
        
        令 $g\left( x\right)  = \dfrac{{x}^{3} + {2x}}{{x}^{2} + 1},x \in  \left\lbrack  {-{2020},{2020}}\right\rbrack$ ,
        
        定义域关于原点对称,且 $g\left( {-x}\right)  = \dfrac{-{x}^{3} - {2x}}{{x}^{2} + 1} =  - \dfrac{{x}^{3} + {2x}}{{x}^{2} + 1} =  - g\left( x\right)$ ,
        
        所以 $g\left( x\right)  = \dfrac{{x}^{3} + {2x}}{{x}^{2} + 1}$ 为奇函数,函数图象关于(0,0)对称,
        
        所以 $f\left( x\right)$ 关于(0,1)对称,所以 $M + m = 2$ .
    \end{jieda}
    
    \item 设 $a$ 为常数, $y = f\left( x\right)$ 是定义在 $\mathbf{R}$ 上的奇函数,当 $x < 0$ 时, $f\left( x\right)  = {9x} + \dfrac{{a}^{2}}{x} + 7$ ,若 $f\left( x\right)  \geq$  $a + 1$ 对一切 $x \geq  0$ 成立,则实数 $a$ 的取值范围是\blkda[4]{$\left(-\infty, - \dfrac{8}{7}\right]$}.
    \begin{jieda}
        $f\left( x\right)$ 是奇函数,当 $x < 0$ 时, $f\left( x\right)  \leq  7 - 6\left| a\right|$ ,易得当 $x > 0$ 时, $f\left( x\right)  \geq  6\left| a\right|  -$ 7,且有 $f\left( 0\right)  = 0$ ,所以 $\left\{  \begin{array}{l} 6\left| a\right|  - 7 \geq  a + 1, \\  0 \geq  a + 1, \end{array}\right.$ 解得 $a \leq   - \dfrac{8}{7}$ 
    \end{jieda}
    \item 已知函数 $y = f\left( x\right)$ 是定义在 $\mathbf{R}$ 上的奇函数,对任意两个不相等的正数 ${x}_{1}\text{ 、 }{x}_{2}$ 都有 $\dfrac{{x}_{2}f\left( {x}_{1}\right)  - {x}_{1}f\left( {x}_{2}\right) }{{x}_{1} - {x}_{2}} < 0$ ,则函数 $y = \left\{  \begin{array}{ll} \dfrac{f\left( x\right) }{x}, & x \neq  0, \\  0, & x = 0 \end{array}\right.$ \dotfill{(\kh{A})}
    \xx{是偶函数,且在 $\left( {0, + \infty }\right)$ 上严格递减}{是偶函数,且在 $\left( {0, + \infty }\right)$ 上严格递增}{是奇函数, 且是严格减函数}{是奇函数, 且是严格增函数}
    \begin{jieda}
        $\dfrac{{x}_{2}f\left( {x}_{1}\right)  - {x}_{1}f\left( {x}_{2}\right) }{{x}_{1} - {x}_{2}} < 0 \Longleftrightarrow \dfrac{\dfrac{f(x_1)}{x_1} - \dfrac{f(x_2)}{x_2}}{x_1-x_2} < 0$ ,即 $g\left( x\right)  = \dfrac{f\left( x\right) }{x}$ 在 $\left( {0, + \infty }\right)$ 上严格递减. $g\left( {-x}\right)  = \dfrac{f\left( {-x}\right) }{-x} = \dfrac{-f\left( x\right) }{-x} = \dfrac{f\left( x\right) }{x} = g\left( x\right)$ ,故 $g\left( x\right)$ 为偶函数.
    \end{jieda}
    
    \item 已知函数 $f\left( x\right)$ 是定义在 $\mathbf{R}$ 上的奇函数,且在区间 $\lbrack 0, + \infty )$ 上是减函数若 $f\left( {{2a} + 1}\right)  + f\left( 1\right)  < 0$ ,则实数 $a$ 的取值范围是\blkda[4]{$\left( {-1, + \infty }\right)$}
    \begin{jieda}
        根据题意可知$f\left( {{2a} + 1}\right)  + f\left( 1\right)  < 0 \Longleftrightarrow 2a+1+1 > 0$，故$a \in \left( {-1, + \infty }\right)$
    \end{jieda}
    
    \item 设 $f\left( x\right)$ 是定义域为 $\mathbf{R}$ 的奇函数, $g\left( x\right)$ 是定义域为 $\mathbf{R}$ 的偶函数,若函数 $f\left( x\right)  + g\left( x\right)$ 的值域为 $\lbrack 1,3)$ ,则函数 $f\left( x\right)  - g\left( x\right)$ 的值域为\blkda{$(-3, -1]$}.
    \begin{jieda}
        对$\forall x_0 \in \mathbf{R}$，$\left\{\begin{array}{l}
             f(x_0)+g(x_0) \in [1,3)  \\
             f(-x_0)+g(-x_0) \in [1,3)
        \end{array}\right.$，根据函数性质可知$f(-x_0)+g(-x_0) = -f(x_0)+g(x_0) = -(f(x_0)-g(x_0))$，故$f\left( x\right)  - g\left( x\right)$ 的值域为$(-3, -1]$.
    \end{jieda}
    
    \item 给定函数 $f\left( x\right)$ 和 $g\left( x\right)$ ,令 $h\left( x\right)  = max \{ f\left( x\right) ,g\left( x\right) \}$ ,对以下三个论断:
    \begin{dingautolist}{172}
        \item 若 $f\left( x\right)$ 和 $g\left( x\right)$ 都是奇函数,则 $h\left( x\right)$ 也是奇函数
        \item 若 $f\left( x\right)$ 和 $g\left( x\right)$ 都是非奇非偶函数,则 $h\left( x\right)$ 也是非奇非偶函数
        \item $f\left( x\right)$ 和 $g\left( x\right)$ 之一与 $h\left( x\right)$ 有相同的奇偶性
    \end{dingautolist}
    其中, 正确论断的个数为\blkda[4]{0}个.
    \begin{jieda}
        \begin{dingautolist}{172}
            \item 错误. 反例:当 $f\left( x\right)  = x,g\left( x\right)  = 0$ 时, $h\left( x\right)$ 不是奇函数.
            
            \item 错误. 反例:当 $f\left( x\right)  = \left\{  {\begin{array}{ll} x, & x \leq  0, \\  0, & x > 0, \end{array}g\left( x\right)  = \left\{  \begin{array}{ll} 0, & x \leq  0, \\   - x, & x > 0 \end{array}\right. }\right.$ 时, $h\left( x\right)  = 0$ 是奇函数.
            
            \item 错误. 反例:当 $f\left( x\right)  = x,g\left( x\right)  =  - x$ 时, $f\left( x\right)$ 与 $g\left( x\right)$ 均为奇函数, $h\left( x\right)  = \left| x\right|$ 为偶函数.
        \end{dingautolist}
    \end{jieda}
    
    \item $f\left( x\right)$ 是定义在 $\left( {0, + \infty }\right)$ 上的增函数,且对定义域内任意的 $x\text{ 、 }y$ 都有 $f\left( \dfrac{x}{y}\right)  =$  $f\left( x\right)  - f\left( y\right)$ 成立
    \begin{enumerate}
        \item 求 $f\left( 1\right)$ 的值
        \item 若 $f\left( 6\right)  = 1$ ,解不等式 $f\left( {x + 5}\right)  - f\left( \dfrac{1}{x}\right)  < 2$ 
    \end{enumerate}
    \begin{jieda}
        \begin{enumerate}
            \item 令 $x = y = 1$ ,则 $f\left( 1\right)  = 0$ 
            \item 
            根据定义可知$f\left( {x + 5}\right)  - f\left( \dfrac{1}{x}\right)  = f\left\lbrack  {\left( {x + 5}\right) x}\right\rbrack$，故应寻找函数值为2的自变量

            
            令 $x = {36},y = 6$ ,由 $f\left( 6\right)  = 1$ 得 $f\left( {36}\right)  = 2$
            
            而 $f\left\lbrack  {\left( {x + 5}\right) x}\right\rbrack  < 2 = f\left( {36}\right)$ ,且 $f\left( x\right)$ 是定义在 $\left( {0, + \infty }\right)$ 上的增函数,则 $\left( {x + 5}\right) x < 36$,所以 $- 9 < x < 4$ ,又因为 $x > 0,x + 5 > 0$ ,所以 $0 < x < 4$ 。
        \end{enumerate}        
    \end{jieda}
    
    \item  已知函数 $f\left( x\right)  = \dfrac{{ax} - b}{1 + {x}^{2}}$ 是定义在 $\left\lbrack  {-1,1}\right\rbrack$ 上的奇函数,且 $f\left( 1\right)  =  - 1$ .
    \begin{enumerate}
        \item 求函数 $f\left( x\right)$ 的解析式
        \item 判断 $f\left( x\right)$ 在 $\left\lbrack  {-1,1}\right\rbrack$ 上的单调性,并用单调性定义证明
        \item 解不等式 $f\left( {t - 1}\right)  + f\left( {t}^{2}\right)  > f\left( 0\right)$
    \end{enumerate}
    \begin{jieda}
        \begin{enumerate}
            \item 函数 $f\left( x\right)  = \dfrac{{ax} - b}{1 + {x}^{2}}$ 是定义在 $\left\lbrack  {-1,1}\right\rbrack$ 上的奇函数,
            
            $f\left( {-x}\right)  =  - f\left( x\right) ;\dfrac{-{ax} - b}{1 + {x}^{2}} =  - \dfrac{{ax} - b}{1 + {x}^{2}}$ ,解得 $b = 0$ ,
            
            $\therefore f\left( x\right)  = \dfrac{ax}{1 + {x}^{2}}$ ,而 $f\left( 1\right)  =  - 1$ ,解得 $a =  - 2$ ,
            
            $\therefore f\left( x\right)  = \dfrac{-{2x}}{1 + {x}^{2}},x \in  \left\lbrack  {-1,1}\right\rbrack$ .
            \item 函数 $f\left( x\right)  = \dfrac{-{2x}}{1 + {x}^{2}}$ 在 $\left\lbrack  {-1,1}\right\rbrack$ 上为减函数；
            
            证明如下: 任意 ${x}_{1},{x}_{2} \in  \left\lbrack  {-1,1}\right\rbrack$ 且 ${x}_{1} < {x}_{2}$ ,则
            
            $f\left( {x}_{1}\right)  - f\left( {x}_{2}\right)  = \dfrac{-2{x}_{1}}{1 + {x}_{1}^{2}} - \dfrac{-2{x}_{2}}{1 + {x}_{2}^{2}} = \dfrac{-2\left( {{x}_{1} - {x}_{2}}\right) \left( {1 - {x}_{1}{x}_{2}}\right) }{\left( {1 + {x}_{1}^{2}}\right) \left( {1 + {x}_{2}^{2}}\right) }$
            
            因为 ${x}_{1} < {x}_{2}$ ,所以 ${x}_{1} - {x}_{2} < 0$ ,又因为 ${x}_{1},{x}_{2} \in  \left\lbrack  {-1,1}\right\rbrack$ ,
            
            所以 $1 - {x}_{1}{x}_{2} > 0$ ,所以 $f\left( {x}_{1}\right)  - f\left( {x}_{2}\right)  > 0$ ,
            
            即 $f\left( {x}_{1}\right)  > f\left( {x}_{2}\right)$ ,所以函数 $f\left( {x}_{1}\right)  > f\left( {x}_{2}\right)$ 在 $\left\lbrack  {-1,1}\right\rbrack$ 上为减函数.

            \item 由题意, $f\left( {t - 1}\right)  + f\left( {t}^{2}\right)  > f\left( 0\right)$ ,又 $f\left( 0\right)  = 0$ ,所以 $f\left( {t - 1}\right)  + f\left( {t}^{2}\right)  > 0$ ,
            
            即解不等式 $f\left( {t}^{2}\right)  >  - f\left( {t - 1}\right)$ ,所以 $f\left( {t}^{2}\right)  > f\left( {1 - t}\right)$ ,
            
            所以 $\left\{  \begin{array}{l}  - 1 \leq  {t}^{2} \leq  1 \\   - 1 \leq  1 - t \leq  1,\text{ 解得 }0 \leq  t < \dfrac{\sqrt{5} - 1}{2}, \\  {t}^{2} < 1 - t \end{array}\right.$
            
            所以该不等式的解集为 $\left\lbrack  {0,\dfrac{\sqrt{5} - 1}{2}}\right)$ .
        \end{enumerate}
    \end{jieda}
\end{enumerate}
\clearpage
\subsection*{C Level}
\begin{enumerate}
    \item 若函数 $y = f\left( x\right)$ 在区间 $I$ 上是增函数,且函数 $y = \dfrac{f\left( x\right) }{x}$ 在区间 $I$ 上是减函数,则称函数 $f\left( x\right)$ 是区间 $I$ 上的 “ $H$ 函数”. 对于命题:
    \begin{dingautolist}{172}
        \item 函数 $f\left( x\right)  =  - x + 2\sqrt{x}$ 是(0,1)上的 “ $H$ 函数”
        \item 函数 $g\left( x\right)  = \dfrac{2x}{1 - {x}^{2}}$ 是(0,1)上的 “ $H$ 函数”
    \end{dingautolist}
    其中真命题有\blkda[4]{\ding{172}}
    \begin{jieda}
        \begin{dingautolist}{172}
            \item 设 $t = \sqrt{x},x \in  \left( {0,1}\right)$ 单调递增, $y =  - {t}^{2} + {2t},t \in  \left( {0,1}\right)$ 单调递增,故 $f\left( x\right)$ 在(0,1)上单调递增, $\dfrac{f\left( x\right) }{x} = \dfrac{2}{\sqrt{x}} - 1,x \in  \left( {0,1}\right)$ 单调递减,故函数 $f\left( x\right)$ 是$(0,1)$上的 “ $H$ 函数”.
            \item $g\left( x\right)  = \dfrac{2}{\dfrac{1}{x} - x},x \in  \left( {0,1}\right)$ 单调递增, $\dfrac{g\left( x\right) }{x} = \dfrac{2}{1 - {x}^{2}},x \in  \left( {0,1}\right)$ 单调递增,故函数$g\left( x\right)$ 不是$(0,1)$上的 “ $H$ 函数”.
        \end{dingautolist}
    \end{jieda}

    \item 定义在 $\mathbf{R}$ 上的函数 $f\left( x\right)$ ,若存在 $a \in  \mathbf{R}$ 且 $a \neq  0$ ,使得 $f\left( {x + a}\right)  > f\left( x\right)  + f\left( a\right)$ 恒成立, 则称 $f\left( x\right)$ 具有“ $P$ 性质”. 已知 ${f}_{1}\left( x\right)$ 是 $\mathbf{R}$ 上的增函数,且 ${f}_{1}\left( x\right)  \leq  0$ 恒成立； ${f}_{2}\left( x\right)$ 是 $\mathbf{R}$ 上的减函数,且存在 ${x}_{0} < 0$ ,使得 ${f}_{2}\left( {x}_{0}\right)  = 0$ ,则\dotfill{(\kh{A})}
    \xx{${f}_{1}\left( x\right)$ 和 ${f}_{2}\left( x\right)$ 都具有 “ $P$ 性质”}{${f}_{1}\left( x\right)$ 不具有 “ $P$ 性质”, ${f}_{2}\left( x\right)$ 具有 “ $P$ 性质”}{${f}_{1}\left( x\right)$ 具有 “ $P$ 性质”, ${f}_{2}\left( x\right)$ 不具有 “ $P$ 性质”}{${f}_{1}\left( x\right)$ 和 ${f}_{2}\left( x\right)$ 都不具有 “ $P$ 性质”}
    
    \begin{jieda}
        在 ${f}_{1}\left( x\right)$ 中,任意 $a > 0$ ,由单调递增可得 ${f}_{1}\left( {x + a}\right)  > {f}_{1}\left( x\right)$ ,又 ${f}_{1}\left( a\right)  \leq  0$ ,所以存在 $a > 0$ ,使得 ${f}_{1}\left( {x + a}\right)  > {f}_{1}\left( x\right)  + {f}_{1}\left( a\right)$ ,故 ${f}_{1}\left( x\right)$ 具有 “ $P$ 性质”; 
        
        在 ${f}_{2}\left( x\right)$ 中,设 $a = {x}_{0} < 0$ ,由单调递减可得 ${f}_{2}\left( {x + a}\right)  > {f}_{2}\left( x\right)$ ,又 ${f}_{2}\left( a\right)  = 0$ ,所以存在 $a < 0$ ,使得 ${f}_{2}\left( {x + a}\right)  > {f}_{2}\left( x\right)  + {f}_{2}\left( a\right)$ ,故 ${f}_{2}\left( x\right)$ 具有“ $P$ 性质”.
    \end{jieda}
    \item 已知函数 $f\left( x\right)  = \mid {x}^{2} + \left( {{3m} + 5}\right) \left| x\right|  + 1 \mid$ 的定义域为 $\mathbf{R}$ ,且该函数在定义域上有且仅有 8 个单调区间,则实数 $m$ 的取值范围是\dotfill{(\kh{C})}
    \xx{$\left( {-\infty , - \dfrac{5}{3}}\right)$}{$\left( {-\infty , - \dfrac{5}{3}}\right)  \cup  \left( {-1, + \infty }\right)$}{$\left( {-\infty , - \dfrac{7}{3}}\right)$}{$\left( {-\infty , - \dfrac{7}{3}}\right)  \cup  \left( {-1, + \infty }\right)$}
    \begin{jieda}
        由题易得 $f\left( x\right)  = f\left( {-x}\right)$ ,则有 $f\left( x\right)$ 为偶函数,所以 $f\left( x\right)$ 在 $\left( {0, + \infty }\right)$ 上有四个单调区间,画图可得 $y = {x}^{2} + \left( {{3m} + 5}\right) x + 1$ 必须满足有两个正的零点
        
        故应有对称轴大于 0且有两个不同的实根，即 $\left\{  \begin{array}{l}  - \dfrac{{3m} + 5}{2} > 0, \\  \Delta  = {\left( 3m + 5\right) }^{2} - 4 > 0, \end{array}\right.$ 解得 $m <  - \dfrac{7}{3}$ .
    \end{jieda}
\end{enumerate}


\section*{对称性、周期性、函数性质综合}
\subsection*{A Level}
\begin{enumerate}
    \item 函数 $f\left( x\right)  = \left| {{2x} - 1}\right|  + 2$ 的对称轴方程为\blkda[4]{$x = \dfrac{1}{2}$}.
    
    \item 函数 $f\left( x\right)  = \dfrac{1}{x} - x$ 的图象关于\blkda[4]{原点}对称.
    
    \item 函数 $f\left( x\right)  = {x}^{3} - 3{x}^{2} + {5x} - 1$ 图象的对称中心为\blkda[4]{$(1,2)$}
    \begin{jieda}
        \ps 三次函数$f(x) = ax^2+bx^2+cx+d$对称中心坐标为$\left(-\dfrac{b}{3a}, f\left(-\dfrac{b}{3a}\right)\right)$
    
        由题意设对称中心的坐标为(a, b),则有 ${2b} = f\left( {a + x}\right)  + f\left( {a - x}\right)$ 对任意 $x \in  R$ 均成立, 代入函数解析式得，${2b} = {\left( a + x\right) }^{3} - 3{\left( a + x\right) }^{2} + 5\left( {a + x}\right)  - 1 + {\left( a - x\right) }^{3} - 3{\left( a - x\right) }^{2} + 5\left( {a - x}\right)  - 1$
        
        整理得到${2b} = {\left( a + x\right) }^{3} - 3{\left( a + x\right) }^{2} + 5\left( {a + x}\right)  - 1 + {\left( a - x\right) }^{3} - 3{\left( a - x\right) }^{2} + 5\left( {a - x}\right)  - 1,$
        
        整理得到 ${2b} = \left( {{6a} - 6}\right) {x}^{2} + 2{a}^{3} - 6{a}^{2} + {10a} - 2 = 0$ 对任意 $x \in  R$ 均成立,
        
        所以 $\left\{  \begin{array}{l} {6a} - 6 = 0 \\  2{a}^{3} - 6{a}^{2} + {10a} - 2 = {2b} \end{array}\right.$ ,所以 $a = 1,b = 2$，即对称中心$(1,2)$.
    \end{jieda}
    \item 已知函数 $f\left( x\right)  = {x}^{2} + {mx}$ ,若 $\forall x \in  \mathbf{R},f\left( {1 - x}\right)  = f\left( {1 + x}\right)$ ,则 $m =$ \blkda[4]{$-2$}.
    \begin{jieda}
        因 $\forall x \in  \mathbf{R},f\left( {1 - x}\right)  = f\left( {1 + x}\right)$ ,函数 $f\left( x\right)$ 的对称轴为直线 $x = 1$ ,

        而由 $f\left( x\right)  = {x}^{2} + {mx} = {\left( x + \dfrac{m}{2}\right) }^{2} - \dfrac{{m}^{2}}{4}$ 可知其对称轴为直线 $x =  - \dfrac{m}{2}$ ,故 $- \dfrac{m}{2} = 1$ ,解得 $m =  - 2$ .
    \end{jieda} 
    \item 函数 $y = \dfrac{x - 1}{x - b}$ 的图像关于点(3, c)中心对称,则 $b + c =$ \blkda[4]{4}.
    \begin{jieda}
        函数 $y = \dfrac{x - 1}{x - b} = \dfrac{\left( {x - b}\right)  + b - 1}{x - b} = 1 + \dfrac{b - 1}{x - b}$ 的图像关于点(3, c)中心对称，所以 $\left\{  \begin{array}{l} 3 - b = 0 \\  c = 1 \end{array}\right.$ ,解得 $b = 3,c = 1$ ,
        
        所以 $b + c = 4$ .
    \end{jieda}
    \item 奇函数 $y = f\left( x\right)$ 的图像关于直线 $x = 2$ 对称, $f\left( 3\right)  = 3$ ,则 $f\left( {-1}\right)  + f\left( 0\right)  =$ \blkda[4]{$-3$}.
    \begin{jieda}
        因为函数是奇函数,所以 $f\left( 0\right)  = 0$
        
        因为函数关于直线 $x = 2$ 对称, $f\left( {4 - x}\right)  = f\left( x\right)$ ,则 $f\left( 1\right)  = f\left( 3\right)$ ,
        
        $f\left( {-1}\right)  =  - f\left( 1\right)  =  - f\left( 3\right)  =  - 3$ ,所以 $f\left( {-1}\right)  + f\left( 0\right)  =  - 3$
    \end{jieda}
    
    \item 已知函数 $f\left( x\right)$ 满足 $f\left( {x + 3}\right)  = f\left( x\right)$ . 当 $x \in  \left\lbrack  {0,2}\right\rbrack$ 时, $f\left( x\right)  = {x}^{2} + 4$ ,则 $f\left( {2024}\right)  =$ \blkda[4]{8}.
    \begin{jieda}
        因为 $f\left( {x + 3}\right)  = f\left( x\right)$ ,所以 $f\left( x\right)$ 是以 3 为周期的周期函数,
        
        所以 $f\left( {2024}\right)  = f\left( {{674} \times  3 + 2}\right)  = f\left( 2\right)  = {2}^{2} + 4 = 8$
    \end{jieda} 
    \item 设 $f\left( x\right)$ 是以 2 为最小正周期的周期函数,且当 $x \in  \left\lbrack  {0,2}\right\rbrack$ 时, $f\left( x\right)  = {\left( x - 1\right) }^{2}$ ,则$f\left( 5\right)  =$ \blkda[4]{0}, $f\left( \dfrac{9}{2}\right)  =$ \blkda[4]{$\dfrac{1}{4}$}.
    \begin{jieda}
        $f\left( 5\right)  = f\left( 1\right)  = {\left( 1 - 1\right) }^{2} = 0,f\left( \dfrac{9}{2}\right)  = f\left( {4 + \dfrac{1}{2}}\right)  = f\left( \dfrac{1}{2}\right)  = {\left( \dfrac{1}{2} - 1\right) }^{2} = \dfrac{1}{4}$
    \end{jieda} 
    \item 已知定义在 $\mathbf{R}$ 上的函数 $f\left( x\right)$ 满足 $f\left( x\right)  = f\left( \left| x\right| \right) ,f\left( {x - 1}\right)  = f\left( {x + 1}\right)$ ,且当 $x \in  \left\lbrack  {0,1}\right\rbrack$ 时, $f\left( x\right)  = {x}^{2}$ ,则 $f\left( \dfrac{3}{2}\right)  =$ \blkda{$\dfrac{1}{4}$}
    \begin{jieda}
        由题意可知,函数 $f\left( x\right)$ 满足 $f\left( x\right)  = f\left( \left| x\right| \right)$ ,故函数 $f\left( x\right)$ 是偶函数,又 $f\left( {x - 1}\right)  = f(x + 1)$ ,即周期为 2
        
        故 $f\left( \dfrac{3}{2}\right)  = f\left( {-\dfrac{1}{2}}\right)  = f\left( \dfrac{1}{2}\right)  = {\left( \dfrac{1}{2}\right) }^{2} = \dfrac{1}{4}$
    \end{jieda}
    \item 设 $f\left( x\right)$ 是定义域为 $\mathbf{R}$ 的奇函数,且 $f\left( {2 + x}\right)  = f\left( {-x}\right)$ . 若 $f\left( {-3}\right)  = 3$ ,则 $f\left( 5\right)  =$ \blkda[4]{3}.
    \begin{jieda}
        因为 $f\left( x\right)$ 是定义域为 $\mathbf{R}$ 的奇函数，则 $f\left( {2 + x}\right)  = f\left( {-x}\right)  =  - f\left( x\right)$
        
        所以 $f\left( {4 + x}\right)  =  - f\left( {2 + x}\right)  = f\left( x\right)$ ,
        
        所以 $f\left( x\right)$ 是周期为 4 的函数,则 $f\left( 5\right)  = f\left( {-3}\right)  = 3$
    \end{jieda}
    \item 已知偶函数 $f\left( x\right)$ 满足 $f\left( {x + 4}\right)  = f\left( x\right)  + {2f}\left( 2\right)$ ,则 $f\left( {2022}\right)  =$ \blkda[4]{0}.
    \begin{jieda}
        令 $x =  - 2$ ,则 $f\left( 2\right)  = f\left( {-2}\right)  + {2f}\left( 2\right)$ ,又 $f\left( {-2}\right)  = f\left( 2\right)$，所以 $f\left( 2\right)  = 0$ ,于是 $f\left( {x + 4}\right)  = f\left( x\right)  + {2f}\left( 2\right)$ 化为: $f\left( {x + 4}\right)  = f\left( x\right)$，所以 $f\left( x\right)$ 的周期 $T = 4$
        
         $f\left( {2022}\right)  = f\left( {{505} \times  4 + 2}\right)  = f\left( 2\right)  = 0$
    \end{jieda}

    \item 函数 $y = f\left( x\right)$ 的周期为 2,且当 $x \in  \lbrack  - 1,1)$ 时, $f\left( x\right)  = x$ ,则 $y = f\left( x\right)$，$x \in  \lbrack {2k} - 1,{2k} + 1)\left( {k \in  \mathbf{Z}}\right)$ 的解析式为\blkda[4]{$f\left( x\right)  = x - {2k}$}.
    \begin{jieda}
        因为函数 $y = f\left( x\right)$ 的周期为 2,当 $x \in  \lbrack  - 1,1)$ 时, $f\left( x\right)  = x$,且 $k \in  \mathbf{Z}$ ,当 ${2k} - 1 \leq  x < {2k} + 1$ 时,则 $- 1 \leq  x - {2k} < 1$，故当 $x \in  \lbrack {2k} - 1,{2k} + 1)\left( {k \in  \mathbf{Z}}\right)$ 时, $f\left( x\right)  = f\left( {x - {2k}}\right)  = x - {2k}$ .
    \end{jieda}
    
    
\end{enumerate}
\subsection*{B Level}
\begin{enumerate}
    \item 设函数 $f\left( x\right)  = \dfrac{{x}^{3} + {\left( 2x + 1\right) }^{2} + 3}{{x}^{2} + 1}$ 在区间 $\left\lbrack  {-2,2}\right\rbrack$ 上的最大值为 $M$ ,最小值为 $N$ ,则 $M + N$ 的值为\blkda[4]{8}.
    \begin{jieda}
        由 $f\left( x\right)  = \dfrac{{x}^{3} + {\left( 2x + 1\right) }^{2} + 3}{{x}^{2} + 1} = \dfrac{{x}^{3} + 4{x}^{2} + {4x} + 4}{{x}^{2} + 1} = \dfrac{{x}^{3} + {4x}}{{x}^{2} + 1} + 4$ ,
        
        设 $g\left( x\right)  = \dfrac{{x}^{3} + {4x}}{{x}^{2} + 1},x \in  \left\lbrack  {-2,2}\right\rbrack$ ,
        
        则 $g\left( {-x}\right)  = \dfrac{-{x}^{3} - {4x}}{{\left( -x\right) }^{2} + 1} =  - \dfrac{{x}^{3} + {4x}}{{x}^{2} + 1} =  - g\left( x\right)$ ,
        
        所以函数 $g\left( x\right)$ 在 $\left\lbrack  {-2,2}\right\rbrack$ 上为奇函数,
        
        所以 $g{\left( x\right) }_{\max } + g{\left( x\right) }_{\min } = 0$
        
        由题意,得 $\left\{  \begin{array}{l} M = g{\left( x\right) }_{\max } + 4 \\  N = g{\left( x\right) }_{\min } + 4 \end{array}\right.$ ,
        
        所以 $M + N = g{\left( x\right) }_{\max } + g{\left( x\right) }_{\min } + 8 = 8$ .
    \end{jieda}
    \item 设定义在 $\mathbf{R}$ 上的函数 $f\left( x\right)$ 在 $( - \infty ,2\rbrack$ 单调递减,且 $f\left( {x + 2}\right)$ 为偶函数,若 $m > 0,n > 0$ , ${2m} \neq  n$ ,且有 $f\left( {2m}\right)  = f\left( n\right)$ ,则 $\dfrac{1}{m} + \dfrac{2m}{n}$ 的最小值为\blkda[4]{$\dfrac{1}{2} + \sqrt{2}$}. 
    \begin{jieda}
        $f\left( {x + 2}\right)$ 为偶函数,则 $f\left( {-x + 2}\right)  = f\left( {x + 2}\right)$ ,则 $f\left( x\right)$ 的对称轴为 $x = 2$ ,
        
        函数 $f\left( x\right)$ 在 $( - \infty ,2\rbrack$ 单调递减,则函数 $f\left( x\right)$ 在 $\lbrack 2, + \infty )$ 单调递增,
        
        若 $m > 0,n > 0,{2m} \neq  n$ ,且有 $f\left( {2m}\right)  = f\left( n\right)$ ,
        
        则 $\dfrac{{2m} + n}{2} = 2$ ,即 ${2m} + n = 4,{2m} = 4 - n$ ,
        
        $\therefore \dfrac{1}{m} + \dfrac{2m}{n} = \dfrac{1}{m} + \dfrac{4 - n}{n} = \dfrac{1}{m} + \dfrac{4}{n} - 1 = \dfrac{1}{4}\left( {{2m} + n}\right) \left( {\dfrac{1}{m} + \dfrac{4}{n}}\right)  - 1$
        
        $= \dfrac{1}{4}\left( {{2m} + n}\right) \left( {\dfrac{1}{m} + \dfrac{4}{n}}\right)  - 1 = \dfrac{1}{4}\left( {6 + \dfrac{n}{m} + \dfrac{8m}{n}}\right)  - 1$
        
        $\geq  \dfrac{1}{4}\left( {6 + 2\sqrt{\dfrac{n}{m} \cdot  \dfrac{8m}{n}}}\right)  - 1 = \dfrac{1}{2} + \sqrt{2},$
        
        当且仅当 $\dfrac{n}{m} = \dfrac{8m}{n}$ 且 ${2m} + n = 4,m > 0,n > 0$ ,即 $m = 2\left( {\sqrt{2} - 1}\right) ,n = 4\left( {2 - \sqrt{2}}\right)$ 时,等号成立,
        
        故 $\dfrac{1}{m} + \dfrac{2m}{n}$ 的最小值为 $\dfrac{1}{2} + \sqrt{2}$ .
    \end{jieda}
    \item 已知奇函数 $f\left( x\right)$ 的定义域为 $\mathbf{R},f\left( {x + 3}\right)  =  - f\left( {-x}\right)$ ,且 $f\left( 2\right)  = 0$ ,则 $f\left( x\right)$ 在 $\left\lbrack  {0,6}\right\rbrack$ 上的零点个数的最小值为\blkda[4]{9}.
    \begin{jieda}
        由 $f\left( {x + 3}\right)  =  - f\left( {-x}\right)$ ,可得 $f\left( x\right)$ 的图象关于点 $\left( {\dfrac{3}{2},0}\right)$ 对称,
        
        又 $f\left( x\right)$ 是奇函数,所以 $f\left( {x + 3}\right)  =  - f\left( {-x}\right)  = f\left( x\right)$ ,
        
        则 $f\left( x\right)$ 的周期为 3,所以 $f\left( 0\right)  = f\left( 3\right)  = f\left( 6\right)  = 0$ ,
        
        $f\left( 5\right)  = f\left( 2\right)  = 0,f\left( 4\right)  = f\left( 1\right)  = f\left( {-2}\right)  =  - f\left( 2\right)  = 0$ ,
        
        而 $f\left( {1.5}\right)  = f\left( {-{1.5}}\right)  =  - f\left( {1.5}\right)$ ,则 $f\left( {1.5}\right)  = f\left( {4.5}\right)  = 0$ .
        
        故 $f\left( x\right)$ 在 $\left\lbrack  {0,6}\right\rbrack$ 上的零点个数的最小值为 9 .
    \end{jieda}
    \item 已知函数 $f\left( x\right)$ 的定义域为 $\mathbf{R}$ ,若 $f\left( {x + 1}\right)  - 2$ 为奇函数,且 $f\left( {1 - x}\right)  = f\left( {3 + x}\right)$ ,则 $f\left( {2023}\right)  =$ \blkda[4]{2}.
    \begin{jieda}
        因为 $f\left( {x + 1}\right)  - 2$ 为奇函数,则 $f\left( {-x + 1}\right)  - 2 =  - \left\lbrack  {f\left( {x + 1}\right)  - 2}\right\rbrack$ ,
        
        所以, $f\left( {1 + x}\right)  + f\left( {1 - x}\right)  = 4$ ,
        
        在等式 $f\left( {1 + x}\right)  + f\left( {1 - x}\right)  = 4$ 中,令 $x = 0$ ,可得 ${2f}\left( 1\right)  = 4$ ,解得 $f\left( 1\right)  = 2$ ,
        
        又因为 $f\left( {1 - x}\right)  = f\left( {3 + x}\right)$ ,则 $f\left( {1 + x}\right)  + f\left( {3 + x}\right)  = 4$ 
        
        所以, $f\left( {x + 3}\right)  + f\left( {x + 5}\right)  = 4$
        
        因此 $f\left( {x + 5}\right)  = f\left( {x + 1}\right)$ ,即 $f\left( {x + 4}\right)  = f\left( x\right)$ ,
        
        所以,函数 $f\left( x\right)$ 为周期函数,且该函数的周期为 4 ,
        
        所以, $f\left( {2023}\right)  = f\left( {4 \times  {505} + 3}\right)  = f\left( 3\right)  = 4 - f\left( 1\right)  = 2$ .
    \end{jieda}
    \item 已知函数 $y = f\left( x\right)$ 是定义在 $\mathbf{R}$ 上的奇函数,且满足 $f\left( {2 + x}\right)  =  - f\left( {2 - x}\right) ,f\left( 1\right)  = 1$ ,则 $f\left( 1\right)  + f\left( 2\right)  + \cdots  + f\left( {2023}\right)  =$ \blkda[4]{0}.
    \begin{jieda}
        因为函数 $y = f\left( x\right)$ 是定义在 $\mathbf{R}$ 上的奇函数,则 $f\left( {-x}\right)  =  - f\left( x\right)$ ,
        
        因为 $f\left( {2 + x}\right)  =  - f\left( {2 - x}\right)  = f\left( {x - 2}\right)$ ,即 $f\left( {x + 4}\right)  = f\left( x\right)$ ,
        
        所以,函数 $f\left( x\right)$ 为周期函数,且周期为 4,则 $f\left( 4\right)  = f\left( 0\right)  = 0$ ,
        
        在等式 $f\left( {2 + x}\right)  =  - f\left( {2 - x}\right)$ 中,令 $x = 0$ ,可得 $f\left( 2\right)  =  - f\left( 2\right)$ ,所以, $f\left( 2\right)  = 0$ ,
        
        因为 $f\left( 3\right)  = f\left( {-1}\right)  =  - f\left( 1\right)  =  - 1$ ,则 $f\left( 1\right)  + f\left( 2\right)  + f\left( 3\right)  + f\left( 4\right)  = 1 + 0 - 1 + 0 = 0$ ,
        
        因为 ${2023} = 4 \times  {505} + 3$ ,
        
        所以, $f\left( 1\right)  + f\left( 2\right)  + \cdots  + f\left( {2023}\right)  = {505} \times  \left\lbrack  {f\left( 1\right)  + f\left( 2\right)  + f\left( 3\right)  + f\left( 4\right) }\right\rbrack   + f\left( 1\right)  + f\left( 2\right)  + f\left( 3\right)$
        
        $= {505} \times  0 + 1 + 0 - 1 = 0$ .
    \end{jieda}
    \item 已知函数 $y = f\left( x\right)$ 对于 $x \in  R$ 恒有 $f\left( {2 - x}\right)  + f\left( x\right)  = 2$ ,若 $f\left( x\right)$ 与函数 $g\left( x\right)  = \dfrac{x}{x - 1}$ 的图像的点交为 $\left( {{x}_{1},{y}_{1}}\right) ,\left( {{x}_{2},{y}_{2}}\right) \ldots \left( {{x}_{n},{y}_{n}}\right)$ ,则 $\left( {{x}_{1} + {y}_{1}}\right)  + \left( {{x}_{2} + {y}_{2}}\right)  + \ldots  + \left( {{x}_{n} + {y}_{n}}\right)  =$ \blkda[4]{${2n}$}
    \begin{jieda}
        因为函数 $y = f\left( x\right)$ 对于 $x \in  R$ 恒有 $f\left( {2 - x}\right)  + f\left( x\right)  = 2$，所以函数 $y = f\left( x\right)$ 的图像关于点(1,1)对称;
        
        $g\left( x\right)  = \dfrac{x}{x - 1} = 1 + \dfrac{1}{x - 1}$ 的图像关于点(1,1)对称,
        
        所以当 $\left( {{x}_{1},{y}_{1}}\right)$ 为 $y = f\left( x\right)$ 和 $y = g\left( x\right)$ 的图像的交点时,点 $\left( {2 - {x}_{1},2 - {y}_{1}}\right)$ 也是 $y = f\left( x\right)$ 和 $y = f\left( x\right)$ 的图像的交点.
        
        所以 $\left( {{x}_{1} + {y}_{1}}\right)  + \left( {{x}_{2} + {y}_{2}}\right)  + \ldots  + \left( {{x}_{n} + {y}_{n}}\right)$
        
        $= \left( {{x}_{1} + {x}_{2} + \cdots  + {x}_{n}}\right)  + \left( {{y}_{1} + {y}_{2} + \cdots {y}_{n}}\right)$
        
        $= \left( {{x}_{1} + {x}_{2} + \cdots  + {x}_{n}}\right)  + \left( {{y}_{1} + {y}_{2} + \cdots {y}_{n}}\right)$
        
        $= 2 \times  \dfrac{n}{2} + 2 \times  \dfrac{n}{2} = {2n}$
    \end{jieda}
    \item 已知定义在 $\mathbf{R}$ 上的函数 $f\left( x\right)$ 满足 $f\left( {-x}\right)  = f\left( x\right) ,f\left( {x + 2}\right)  = f\left( {2 - x}\right)$ ,当 $x \in  \left\lbrack  {0,1}\right\rbrack$ 时，$f\left( x\right)  = {x}^{3} - {3x}$ ,则 $f\left( {2023}\right)  =$ \blkda[4]{$-2$}.
    \begin{jieda}
        由已知可知 $f\left( x\right)$ 是偶函数，且关于$x = 2$对称，因此可以推出周期为4
        
        故 $f\left( {2023}\right)  = f\left( {-1}\right)  = f\left( 1\right)  = {\left( 1\right) }^{3} - 3 \times  1 =  - 2$ .
    \end{jieda}
    \item 已知函数 $f\left( x\right)$ 的定义域为 $\mathbf{R}$ ,且 $f\left( {x + y}\right)  + f\left( {x - y}\right)  = f\left( x\right) f\left( y\right) ,f\left( 1\right)  = 1$ ,则$\mathop{\sum }\limits_{{k = 1}}^{{2024}}f\left( k\right)  =$ \blkda[4]{0}.
    \begin{jieda}
        依题意, $f\left( 1\right)  = 1,f\left( {x + y}\right)  + f\left( {x - y}\right)  = f\left( x\right) f\left( y\right)$ ,
        
         令 $y = 1$ 得 $f\left( {x + 1}\right)  + f\left( {x - 1}\right)  = f\left( x\right) f\left( 1\right)  = f\left( x\right)$ ,
        
        所以 $f\left( {x + 1}\right)  = f\left( x\right)  - f\left( {x - 1}\right)$ ,则 $f\left( {x + 2}\right)  = f\left( {x + 1}\right)  - f\left( x\right)$ ,
        
        $f\left( {x + 3}\right)  = f\left( {x + 2}\right)  - f\left( {x + 1}\right)  =  - f\left( x\right) ,$
        
        
        
        所以 $f\left( {x + 6}\right)  = f\left( {x + 3 + 3}\right)  =  - f\left( {x + 3}\right)  =  - \left( {-f\left( x\right) }\right)  = f\left( x\right)$ ,
        
        所以 $f\left( x\right)$ 是周期为 6 的周期函数.
        
        令 $x = 1,y = 0$ ,则 $f\left( 1\right)  + f\left( 1\right)  = f\left( 1\right) f\left( 0\right) ,f\left( 0\right)  = 2$ ,
        
        $f\left( 2\right)  = f\left( 1\right)  - f\left( 0\right)  =  - 1,f\left( 3\right)  = f\left( 2\right)  - f\left( 1\right)  =  - 2,$
        
        $f\left( 4\right)  = f\left( 3\right)  - f\left( 2\right)  =  - 1,f\left( 5\right)  = f\left( 4\right)  - f\left( 3\right)  = 1$ ,
        
        $f\left( 6\right)  = f\left( 5\right)  - f\left( 4\right)  = 2$ ,所以 $f\left( 1\right)  + f\left( 2\right)  + f\left( 3\right)  + f\left( 4\right)  + f\left( 5\right)  + f\left( 6\right)  = 0$ ,
        
        因为 ${2024} = {337} \times  6 + 2$ ,所以 $\mathop{\sum }\limits_{{k = 1}}^{{2024}}f\left( k\right)  = f\left( 1\right)  + f\left( 2\right)  = 0$ .
    \end{jieda}
    \item 设 $f\left( x\right)$ 为定义在整数集上的函数, $f\left( 1\right)  = 1,f\left( 2\right)  = 0,f\left( {-1}\right)  < 0$ ,对任意的整数 $x,y$ 均有 $f\left( {x + y}\right)  = f\left( x\right) f\left( {1 - y}\right)  + f\left( {1 - x}\right) f\left( y\right)$ . 则 $f\left( {55}\right)  =$ \blkda[4]{$-1$}.
    \begin{jieda}
        令 $x = y = 1$ ,则 $f\left( 2\right)  = f\left( 1\right) f\left( 0\right)  + f\left( 0\right) f\left( 1\right)  = {2f}\left( 0\right)  = 0,\therefore f\left( 0\right)  = 0$ ;
        
        令 $x = 2,y =  - 1$ ,则 $f\left( 1\right)  = {f}^{2}\left( 2\right)  + {f}^{2}\left( {-1}\right)  = {f}^{2}\left( {-1}\right)  = 1$ ,又 $f\left( {-1}\right)  < 0,\therefore f\left( {-1}\right)  =  - 1$ ;
        
        令 $y = 1$ ,则 $f\left( {x + 1}\right)  = f\left( x\right) f\left( 0\right)  + f\left( {1 - x}\right) f\left( 1\right)  = f\left( {1 - x}\right) ,\therefore f\left( x\right)$ 关于直线 $x = 1$ 对称;
        
        令 $y =  - x$ ,则 $f\left( 0\right)  = f\left( x\right) f\left( {1 + x}\right)  + f\left( {1 - x}\right) f\left( {-x}\right)  = \left\lbrack  {f\left( x\right)  + f\left( {-x}\right) }\right\rbrack  f\left( {1 + x}\right)  = 0$ ,
        
        $\because f\left( {1 + x}\right)  = 0$ 不恒成立, $\therefore f\left( x\right)  + f\left( {-x}\right)  = 0$ 恒成立, $\therefore f\left( x\right)$ 为奇函数,
        
        $\because f\left( {x + 2}\right)  = f\left( {-x}\right)  =  - f\left( x\right) ,\therefore f\left( {x + 4}\right)  =  - f\left( {x + 2}\right)  = f\left( x\right)$ ,
        
        $\therefore f\left( x\right)$ 是周期为 4 的周期函数, $\therefore f\left( {55}\right)  = f\left( {4 \times  {14} - 1}\right)  = f\left( {-1}\right)  =  - 1$ .
    \end{jieda}
    
    \item 已知 $f\left( x\right)$ 是定义域为 $\mathbf{R}$ 的偶函数, $f\left( {5.5}\right)  = 2,g\left( x\right)  = \left( {x - 1}\right) f\left( x\right)$ 若 $g\left( {x + 1}\right)$ 是偶函数,则 $g\left( {-{0.5}}\right)  =$\blkda[4]{3}
    \begin{jieda}
        由 $g\left( {x + 1}\right)$ 为偶函数,则 $g\left( x\right)$ 关于 $x = 1$ 对称,即 $g\left( x\right)  = g\left( {2 - x}\right)$ ,即 $\left( {x - 1}\right) f\left( x\right)  = (1 -x)f\left( {2 - x}\right)$ ,即 $f\left( x\right)  + f\left( {2 - x}\right)  = 0$ ,所以, $f\left( x\right)$ 关于(1,0)对称,又 $f\left( x\right)$ 是定义域为 $\mathbf{R}$ 的偶函数,则$f\left( x\right)  =  - f\left( {2 - x}\right)  =  - f\left( {x - 2}\right)$ ,从而 $f\left( {x - 4}\right)  = f\left\lbrack  {\left( {x - 2}\right)  - 2}\right\rbrack   =  - f\left( {x - 2}\right)  =  - \left\lbrack  {-f\left( x\right) }\right\rbrack   = f\left( x\right)$，即 $f\left( {x - 4}\right)  = f\left( x\right)$ ,所以 $f\left( x\right)$ 周期为 4 ,所以 $f\left( {5.5}\right)  = f\left( {1.5}\right)  = f\left( {-{2.5}}\right)  = f\left( {2.5}\right)  = 2$ ,所以 $g\left( {-{0.5}}\right)  = g\left( {2.5}\right)  = {1.5f}\left( {2.5}\right)  = 3$
    \end{jieda}
    
    \item 设函数 $f\left( x\right)$ 的定义域为 $\mathbf{R},f\left( {x + 1}\right)$ 为奇函数, $f\left( {x + 2}\right)$ 为偶函数,当 $x \in  \left\lbrack  {1,2}\right\rbrack$ 时, $f\left( x\right)  = {kx} + b$ 若 $f\left( 0\right)  + f\left( 3\right)  = 8$ ,则 $f\left( \dfrac{2021}{2}\right)  =$ \blkda{$-4$}
    \begin{jieda}
    $f(x)$关于$(1, 0)$中心对称，关于$x = 2$轴对称，故其周期为$4$.

    因为$f(x)$关于$(1, 0)$中心对称，定义域为$\mathbf{R}$，所以$f(1) = 0$，故$k+b = 0$

    因为$f(x)$关于$x = 2$轴对称，故$f(3) = f(1) = 0$，故$f(0) = 8$
    
    因为$f(x)$关于$(1, 0)$中心对称，$f(2) = -f(0) = -8$，故$2k+b = -8$

    综上，可得$\left\{\begin{array}{l}
         k = -8 \\
         b = 8
    \end{array}\right.$
    
    因为$f(x)$周期为4，所以$f\left(\dfrac{2021}{2}\right) = f\left(\dfrac{5}{2}\right) = f\left(\dfrac{3}{2}\right) = -4$
    \end{jieda}
    \item 设函数 $f\left( x\right)$ 的定义域为 $\mathbf{R},f\left( {x + 1}\right)$ 为奇函数, $f\left( {x + 2}\right)$ 为偶函数,当 $x \in  \left\lbrack  {1,2}\right\rbrack$ 时, $f\left( x\right)  = a{x}^{2} + b$ . 若 $f\left( 0\right)  + f\left( 3\right)  = 6$ ,则 $f\left( \dfrac{2025}{2}\right)  =$\blkda[4]{$\dfrac{5}{2}$}
    \begin{jieda}
        $f(x)$关于$(1, 0)$中心对称，关于$x = 2$轴对称，故其周期为$4$.

        因为$f(x)$关于$(1, 0)$中心对称，定义域为$\mathbf{R}$，所以$f(1) = 0$，故$a+b = 0$

        因为$f(x)$关于$x = 2$轴对称，故$f(3) = f(1) = 0$，故$f(0) = 6$
        
        因为$f(x)$关于$(1, 0)$中心对称，$f(2) = -f(0) = -6$，故$4a+b = -6$

        综上，可得$\left\{\begin{array}{l}
             a = -2 \\
             b = 2
        \end{array}\right.$
        
        因为$f(x)$周期为4，所以$f\left(\dfrac{2025}{2}\right) = f\left(\dfrac{1}{2}\right) = -f\left(\dfrac{3}{2}\right) = \dfrac{5}{2}$
    \end{jieda}
\end{enumerate}
\subsection*{C Level}
\begin{enumerate}
    \item 已知定义在$\mathbf{R}$上的偶函数$f(x)$满足$f(x) = -f(\pi - x)$，当$x \in \left[0, \dfrac{\pi}{2}\right)$时$f(x) = \sqrt{x}$，则函数$y = f(x)$在$\left[-\dfrac{3\pi}{2}, \dfrac{5\pi}{2}\right]$上的零点个数为\blkda{9}
    \begin{jieda}
        $f(x) = -f(\pi - x)$即$f(x)$关于$\left(\dfrac{\pi}{2}, 0\right)$中心对称，又因为$f(x)$是偶函数，所以$2\pi$是$f(x)$的一个周期. 因为定义域为$\mathbf{R}$，所以$f\left(\dfrac{k\pi}{2}\right) = 0, k \in \mathbf{Z}$；结合图像可知有9个零点.
    \end{jieda}
    \item 已知函数 $f\left( x\right)$ 是 $\mathbf{R}$ 上的偶函数,对于任意 $x \in  \mathbf{R}$ 都有 $f\left( {x + 6}\right)  = f\left( x\right)  + f\left( 3\right)$ 成立,当 ${x}_{1}$ 、 ${x}_{2} \in  \left\lbrack  {0,3}\right\rbrack$ ,且 ${x}_{1} \neq  {x}_{2}$ 时,都有 $\dfrac{f\left( {x}_{1}\right)  - f\left( {x}_{2}\right) }{{x}_{1} - {x}_{2}} > 0$ . 给出以下三个命题:
    \begin{dingautolist}{172}
        \item 直线 $x =  - 6$ 是函数 $f\left( x\right)$ 图象的一条对称轴
        \item 函数 $f\left( x\right)$ 在区间 $\left\lbrack  {-9, - 6}\right\rbrack$ 上为增函数
        \item 函数 $f\left( x\right)$ 在区间 $\left\lbrack  {-9,9}\right\rbrack$ 上恰有五个零点.    
    \end{dingautolist}
    以上命题中, 正确的序号是\blkda[4]{\ding{172}}.
    \begin{jieda}
        令 $x =  - 3$ ,有 $f\left( {-3 + 6}\right)  = f\left( {-3}\right)  + f\left( 3\right)$，即$f(-3) = 0$
        
        又 $f\left( x\right)$ 是 $\mathbf{R}$ 上的偶函数,所以 $f\left( 3\right)  = 0$，则有 $f\left( {x + 6}\right)  = f\left( x\right)$ ,所以 $f\left( x\right)$ 的周期为 6 .
        \begin{center}
        \includegraphics[width=0.6\textwidth]{bodyxb/src/bo_d29ci7n7aajc738p8c80_29_760_1405_737_185_0.jpg}
        \end{center}
        \begin{dingautolist}{172}
            \item 函数为偶函数,则函数关于 $y$ 轴对称,又由函数的周期为 6,则直线 $x =  - 6$ 是函数 $f\left( x\right)$ 图象的一条对称轴
            \item 当 ${x}_{1},{x}_{2} \in  \left\lbrack  {0,3}\right\rbrack$ ,且 ${x}_{1} \neq  {x}_{2}$ 时,都有 $\dfrac{f\left( {x}_{1}\right)  - f\left( {x}_{2}\right) }{{x}_{1} - {x}_{2}} > 0$ ,则函数 $y = f\left( x\right)$ 在 [0,3]上为增函数,根据函数的对称性和周期性等性质,所以函数 $y = f\left( x\right)$ 在 $\left\lbrack  {-9, - 6}\right\rbrack$ 上为减函数
            \item 在 $f\left( {x + 6}\right)  = f\left( x\right)  + f\left( 3\right)$ 式中,令 $x = 0$ ,又因为 $f\left( 3\right)  = 0$ ,函数 $y = f\left( x\right)$ 在 $\left\lbrack  {0,3}\right\rbrack$ 上为严格增函数,故 $f\left( 0\right)  \neq  0$ ,故函数 $y = f\left( x\right)$ 在 $\left\lbrack  {-9,9}\right\rbrack$ 上零点个数一定是偶数个
        \end{dingautolist}
    \end{jieda}
    \item 定义域为实数集的偶函数 $f\left( x\right)$ 满足 $f\left( {x + 1}\right)  = f\left( {x - 1}\right) ,x \in  \mathbf{R}$ 恒成立,若当 $x \in  [2 , 3]$ 时, $f\left( x\right)  = x$ ,给出如下四个结论:
    \begin{dingautolist}{172}
        \item 函数 $f\left( x\right)$ 的图象关于直线 $x =  - 4$ 对称
        \item 对任意实数 $a$ ,关于 $x$ 的方程 $f\left( x\right)  - \left| {x - a}\right|  = 0$ 一定有解
        \item 若存在实数 $a$ ,使得关于 $x$ 的方程 $f\left( x\right)  - \left| {x - a}\right|  = 0$ 有一个根为 2,则此方程所有根之和为$-20$
        \item 若关于 $x$ 的不等式 $f\left( x\right)  - \left| {x - a}\right|  < 0$ 在区间 $\lbrack 0, + \infty )$ 上恒成立,则 $a$ 有最大值
    \end{dingautolist}
    其中,所有正确结论的编号是\blkda{\ding{172}\ding{173}}.
    \begin{jieda}
        \begin{dingautolist}{172}
            \item 因为函数 $f\left( x\right)$ 满足 $f\left( {x + 1}\right)  = f\left( {x - 1}\right)$ , 所以函数 $f\left( x\right)$ 为周期为 2 的函数,又因为函数为偶函数, 图象关于 $y$ 轴对称,所以 $f\left( x\right)$ 也关于直线 $x =  - 4$ 对称
            \item 方程 $f\left( x\right)  - \left| {x - a}\right|  = 0$ 一定有解,即函数 $f\left( x\right)$ 的图象与函数 $y = \left| {x - a}\right|$ 的图象一定有交点. 在同一个直角坐标系中作出两个函数 $y = \left| x\right|$ 与 $y = f\left( x\right)$ 的图象如图所示. 由图象可得,将 $y =  \mid  x \mid$ 左右平移后一定会与函数 $y = f\left( x\right)$ 相交
            \begin{center}
            \includegraphics[width=0.3\textwidth]{bodyxb/src/bo_d29ci7n7aajc738p8c80_30_1077_1575_433_334_0.jpg}
            \end{center}
            \item 根据图象可知若 $x = 2$ 为 $y = f\left( x\right)$ 与 $y = \left| {x - a}\right|$ 的一个交点，则$a = 0$或$4$
            
            当 $a = 0$ 时, $y = \left| x\right|$ 与 $y = f\left( x\right)$ 的图象都关于 $y$ 轴对称,所有交点的横坐标之和为 0

            
            当 $a = 4$ 时, $y = \left| x\right|$ 与 $y = f\left( x\right)$ 的图象都关于 $x = 4$ 轴对称，每一组关于$x = 4$对称的两根之和为$8$，而有两个函数有无数个交点.
            \item 若关于 $x$ 的不等式 $f\left( x\right)  - \left| {x - a}\right|  < 0$ 在区间 $\lbrack 0, + \infty )$ 上恒成立,即 $\left| {x - a}\right|  >$  $f\left( x\right)$ 恒成立,通过图象的平移可知, $a <  - 2$ ,即实数 $a$ 没有最大值
        \end{dingautolist}
    \end{jieda}
\end{enumerate}
\subsection*{D Level}
\begin{enumerate}
    \item 函数 $f\left( x\right)  = \dfrac{x}{x + 1} + \dfrac{x + 1}{x + 2} + \dfrac{x + 2}{x + 3} + \cdots  + \dfrac{x + {2021}}{x + {2022}}$ 的图像的对称中心为\blkda[4]{$\left( {-\dfrac{2023}{2},{2022}}\right)$}
    \begin{jieda}
        $f\left( x\right)  = \dfrac{x}{x + 1} + \dfrac{x + 1}{x + 2} + \cdots  + \dfrac{x + {2021}}{x + {2022}}$
        
        $\dfrac{x}{x + 1} + \dfrac{x + {2021}}{x + {2022}} = 2 - \left( {\dfrac{1}{x + 1} + \dfrac{1}{x + {2022}}}\right) = 2 - \dfrac{2021}{{x}^{2} + {2023x} + {2022}}$
        
        $\left\{\begin{array}{l}
             f\left( 0\right)  = \dfrac{0}{1} + \dfrac{1}{2} + \cdots  + \dfrac{2 \times  {21}}{2022}  \\
             f\left( {-{2023}}\right)  = \dfrac{2023}{2022} + \dfrac{2023}{2021} + \cdots  + \dfrac{2}{1} 
        \end{array}\right.$，有$f\left( 0\right)  + f\left( {-{2023}}\right)  = \dfrac{2}{1} + \dfrac{4}{2} + \cdots  + \dfrac{4044}{2022} = {4044}$
        故 $f\left( x\right)$ 应关于 $\left( {-\dfrac{2023}{2},{2022}}\right)$ 对称
        
        
        检验：$\left\{\begin{array}{l}
             f\left( x\right)  = \dfrac{x}{x + 1} + \dfrac{x + 1}{x + 2} + \cdots  + \dfrac{x + {2021}}{x + {2022}}  \\
             f\left( {-{2023} - x}\right)  = \dfrac{-{2023} - x}{-{2022} - x} + \dfrac{-{2022} - x}{-{2021} - x} + \cdots  + \dfrac{-2 - x}{-1 - x} \\ \quad \quad \quad \quad  \quad \quad  = \dfrac{x + {2023}}{x + {2022}} + \dfrac{x + {2022}}{x + {2021}} + \cdots  + \dfrac{x + 2}{x + 1}
        \end{array}\right.$
        
        故$f\left( x\right)  + f\left( {-{2023} - x}\right)  = \dfrac{2\left( {x + 1}\right) }{x + 1} + \dfrac{2\left( {x + 2}\right) }{x + 2} + \cdots  + \dfrac{2\left( {x + {2022}}\right) }{\left( x + {2022}\right) }$
        $= 2 \times  {2022}$
        $= {4044}$
    \end{jieda} 
    \item 已知函数 $f\left( x\right)  = \dfrac{x}{x + 1} + \dfrac{x + 1}{x + 2} + \dfrac{x + 2}{x + 3}$ ,给出下列判断:
    \begin{dingautolist}{172}
        \item 函数 $f\left( x\right)$ 的值域为 $\mathbf{R}$
        \item $f\left( x\right)$ 在定义域内有三个零点
        \item $f\left( x\right)$ 的图象是中心对称图象
    \end{dingautolist} 其中,正确判断的有\blkda[4]{\ding{172}\ding{173}\ding{174}}
    \begin{jieda}
        由题意可知,函数 $f\left( x\right)  = \dfrac{x}{x + 1} + \dfrac{x + 1}{x + 2} + \dfrac{x + 2}{x + 3} = \left( {1 - \dfrac{1}{x + 1}}\right)  + \left( {1 - \dfrac{1}{x + 2}}\right)  +$  $\left( {1 - \dfrac{1}{x + 3}}\right)  = 3 - \left( {\dfrac{1}{x + 1} + \dfrac{1}{x + 2} + \dfrac{1}{x + 3}}\right)$ ,其定义域为 $\{ x \mid  x \neq   - 1,x \neq   - 2,x \neq   - 3\}$ .
        \begin{dingautolist}{172}
            \item 显然函数在 $x \in  \left( {-2, - 1}\right)$ 内连续且严格单调递增,当 $x \Longrightarrow   - {2}^{ + }$ 时, $f\left( x\right)  \Longrightarrow$  $- \infty$ ,当 $x \Longrightarrow   - {1}^{ - }$ 时, $f\left( x\right)  \Longrightarrow   + \infty$ ,可知函数的值域是 $\mathbf{R}$
            \item 因为 $f\left( x\right)  = \dfrac{x}{x + 1} + \dfrac{x + 1}{x + 2} + \dfrac{x + 2}{x + 3} = 3 - \left( {\dfrac{1}{x + 1} + \dfrac{1}{x + 2} + \dfrac{1}{x + 3}}\right)$ ,显然函数在 $\left( {-\infty , - 3}\right) ,\left( {-3, - 2}\right) ,\left( {-2, - 1}\right) ,\left( {-1, + \infty }\right)$ 四个区间内连续且严格单调递增，分析各区间端点的极限值,根据零点存在性定理,可知 $f\left( x\right)$ 在定义域内有三个零点
            \item 因为 $f\left( x\right)  = 3 - \left( {\dfrac{1}{x + 1} + \dfrac{1}{x + 2} + \dfrac{1}{x + 3}}\right)$ ,所以 $f\left( {-4 - x}\right)  = 3 + \left( {\dfrac{1}{x + 1} + \dfrac{1}{x + 2} + }\right.$
            
            $\left. \dfrac{1}{x + 3}\right)  = 6 - \left\lbrack  {3 - \left( {\dfrac{1}{x + 1} + \dfrac{1}{x + 2} + \dfrac{1}{x + 3}}\right) }\right\rbrack   = 6 - f\left( x\right)$ ,得 $f\left( x\right)  + f\left( {-4 - x}\right)  = 6$ ,即函数的图象关于点$(-2,3)$中心对称

            \ps 对于对称中心可以通过带特殊值或者根据三个权重相同的部分各自的对称中心先猜出对称中心后证明.
        \end{dingautolist}
    \end{jieda}
\end{enumerate}
\newpage
\quad
\newpage
\section*{函数中的恒成立与存在性问题、方程的根与函数的零点}
\subsection*{C Level}
\begin{enumerate}
    \item 对于实数$a, b$定义运算“*”：$a*b = \left\{  
         \begin{array}{ll}  
            a^2 - ab, & a \leq b, \\  
            b^2 - ab, & a > b.
         \end{array}
    \right.$ 设$f(x) = (2x-1) * (x-1)$，且关于$x$的方程$f(x) = m (m \in \mathbf{R})$恰有三个互不相等的实数根$x_1, x_2, x_3$，则$x_1x_2x_3$的取值范围是\blkda[3]{$\left(-\dfrac{\sqrt{3} - 1}{16}, 0\right)$}
    \begin{jieda}
        根据运算定义规则，首先求$2x-1$和$x-1$的交点，为$(0, -1)$，\\
        故$f(x) = \left\{ 
                 \begin{array}{ll}  
                    (2x-1)^2 - (2x-1)(x-1), & x \leq 0, \\  
                    (x-1)^2 - (2x-1)(x-1), & x > 0.
                 \end{array}
                \right. $ \\
        即$f(x) = \left\{ 
                 \begin{array}{ll}  
                    2x^2-x, & x \leq 0, \\  
                    x-x^2, & x > 0.
                 \end{array}
                \right. $ \\
        观察函数图像可知，在满足题意的情况下，$y = m \left(0 < m < \dfrac{1}{4}\right)$与$f(x)$在$y$轴左侧有一个交点，在$y$轴右侧有两个交点
        
        由韦达定理知$x_2x_3 = m$，$x_1 = \dfrac{1-\sqrt{1+8m}}{4}$，故$x_1x_2x_3 = -\dfrac{1}{4}(m \cdot (\sqrt{1+8m} - 1))$
        
        因为$-\dfrac{1}{4}(m \cdot (\sqrt{1+8m} - 1))$在$m \in \left(0, \dfrac{1}{4}\right)$单调递减，故其取值范围是$\left(-\dfrac{\sqrt{3} - 1}{16}, 0\right)$
    \end{jieda}
    \item 用$C(A)$表示非空集合$A$中的元素个数，定义$|A-B| = \left \{ \begin{array}{ll}
        C(A) - C(B), & C(A) \geq C(B) \\
        C(B) - C(A), & C(A) < C(B)
    \end{array}\right.$，若$A = \{1,2\}, B = \{x | |x^2+2x-3| = a\}$，且$|A-B| = 1$，由$a$的所有可能值构成的集合为$S$，那么$C(S) = $\blkda{1}
    \begin{jieda}
        $|x^2+2x-3| = a$即$y = a$与$y = |x^2+2x-3|$的交点横坐标，而$|A-B| = 1$即两者交点数为3或1. 易知两者交点个数为0,2,3,4，故考虑其有三个交点的情况，即$a = 4$，故$S$中只有一个元素.
    \end{jieda}
    \item 用$C(A)$表示非空集合A中元素个数，定义$|A-B| = \left \{ \begin{array}{ll}
         C(A)-C(B),& C(A) \geq C(B)  \\
         C(B)-C(A),& C(A) < C(B)
    \end{array}\right.$，
    已知集合$A = \{x | x^2 + x = 0\}, B = \{x | (3x^2+ax)(x^2+ax+2) = 0\}$，且$|A-B|=1$，设实数$a$的所有可能取值构成集合$S$，则$C(S) = $\blkda{5}
    \begin{jieda}
        $A = \{0, -1\}, C(A) = 2$,因为$|A-B| = 1$，所以$C(B) = 1 \mbox{或} 3$。 \\
        方程$(3x^2+ax)(x^2+ax+2) = 0$的根包括$0, -\dfrac{a}{3}, \dfrac{-a+\sqrt{a^2-8}}{2}, \dfrac{-a-\sqrt{a^2-8}}{2}$ 
        \begin{enumerate}
            \item 当$a \in (-2\sqrt{2}, 2\sqrt{2})$时，后两者并非实根，因此$C(B) \neq 3$，当$a = 0$的时候$B = \{0\}, C(B) = 1$
            \item 当$a \in \{-2\sqrt{2}, 2\sqrt{2}\}$时，$B = \{0, -\dfrac{a}{3}, -\dfrac{a}{2}\}, C(B) = 3$
            \item 当$a \in (-\infty, -2\sqrt{2}) \cup (2\sqrt{2}, +\infty)$时，若有$-\dfrac{a}{3} = \dfrac{-a \pm \sqrt{a^2-8}}{2}$可使得$C(B) = 3$，解得$a = \pm 3$
        \end{enumerate}
        综上$S = \{0, 2\sqrt{2}, -2\sqrt{2}, 3, -3\}$，因此$C(S) = 5$
    \end{jieda} 
    
    \item 已知方程$|x| = ax+1$有一负解，且无正解，则实数$a$的取值范围是\blkda{$[1, +\infty)$ }
    \begin{jieda}
        将方程的根转化为$y = |x|$与$y = ax+1$的交点的横坐标.$y = ax+1$过$(0, 1)$，斜率为$a$，故当$a \in [1, +\infty)$时，两函数在$x > 0$时无交点，而在$x < 0$时有一个交点，满足要求.
    \end{jieda}
    
    \item 设$f(x)$是定义在$\bR$上的奇函数, 且对于任意的$x\zr$, $f(1+x)-f(1-x)=0$恒成立; 当$x\in[0,1]$时, $f(x)=2x$. 
    若方程$f(x)=ax$恰好有5个不同的解, 则实数$a$的取值范围为\blkda[2.5]{$\left(-\dfrac{2}{3},-\dfrac{2}{7}\right)\cup\left\{\dfrac{2}{5}\right\}$}. 
    \begin{jieda}
    $f(x)$图像关于原点对称, 关于直线$x=1$对称, 可知周期$T=4$. \\
    $y = ax$的图像是过原点的直线，随着$a$的变化，其斜率变化. 绘制草图后可以判断要分为$a>0$与$a<0$两种情形, 
    不难求得$a=\dfrac{2}{5}$或$a\in\left(-\dfrac{2}{3},-\dfrac{2}{7}\right)$. 
    \end{jieda}
    
    \item 已知关于$x$的方程$(x^2-1)^2-|x^2-1|+k = 0$. 下列语句中说法正确的有\blkda{\ding{172} \ding{173} \ding{174} \ding{175}}
    \begin{dingautolist}{172}
        \item 存在实数$k$使得方程恰有2个不同的实数解.
        \item 存在实数$k$使得方程恰有4个不同的实数解.
        \item 存在实数$k$使得方程恰有5个不同的实数解.
        \item 存在实数$k$使得方程恰有8个不同的实数解.
    \end{dingautolist}
    \begin{jieda}
        换元$t = |x^2-1|$，参变分离有$k = t - t^2$
        
        对于内层方程$t = |x^2-1|$：
        \begin{dingautolist}{172}
            \item 若$t < 0$，则方程无解
            \item 若$t = 0$，则方程有2个解
            \item 若$t \in (0, 1)$，则方程有4个解
            \item 若$t = 1$，则方程有3个解
            \item 若$t > 1$，则方程有2个解
        \end{dingautolist}
        对于外层方程$k = t - t^2$：
        \begin{dingautolist}{172}
            \item 若$k > \dfrac{1}{4}$，则外层方程无解，复合方程无解
            \item 若$k = \dfrac{1}{4}$，则存在唯一的$t = \dfrac{1}{2}$使外层方程成立，对复合方程本身有4个解
            \item 若$0 < k < \dfrac{1}{4}$，则外层方程两解$t_1, t_2 \in (0, 1)$，复合方程有8个解
            \item 若$k = 0$，则外层方程两解$t_1 = 0, t_2 = 1$，复合方程有5个解
            \item 若$k < 0$，则外层方程两解$t_1 < 0, t_2 > 1$，复合方程有2个解
        \end{dingautolist}
    \end{jieda}
    
    \item 已知函数$f(x) = x - 2$，$g(x) = mx^2-2mx+1(m \in \mathbf{R})$. 
    \begin{enumerate}
        \item 对任意的$x \in \mathbf{R}$，不等式$g(x) > f(x)$恒成立，求$m$取值范围.
        \item 对任意的$x_1 \in [1, 2], x_2 \in [-1, 0]$，不等式$g(x_1) > f(x_2)$恒成立，求$m$取值范围.
        \item 对任意的$x_1 \in [1, 2]$，存在$x_2 \in [3, 4]$，使得$g(x_1) = f(x_2)$，求$m$取值范围.
        \item 对任意的$x_1 \in [1, 2]$，存在$x_2 \in [0, 2]$，使得$g(x_1) > f(x_2)$，求$m$取值范围.
    \end{enumerate}
    
    \begin{jieda}
        \begin{enumerate}
            \item 问题转化为$g(x) - f(x) > 0$在$\mathbf{R}$上恒成立，即$mx^2-(2m+1)x+3 > 0$恒成立. $m = 0$不满足要求，$m < 0$不满足要求，$m > 0$时应有$\Delta < 0$，即$m \in \left(\dfrac{2-\sqrt{3}}{2}, \dfrac{2+\sqrt{3}}{2}\right)$
            \item 问题转化为$g(x)$在$[1, 2]$上的最小值大于$f(x)$在$[3, 4]$上的最大值. $f(x)$在$[-1, 0]$上的最大值为$-2$，故问题转化为$g(x) + 2 > 0$在$[1, 2]$上恒成立，设$h(x) = g(x) + 2 = mx^2-2mx + 3 = m(x-1)^2-(m-3)$，下面分析$h(x)$的最值
            \begin{dingautolist}{172}
                \item $m > 0$，$h(x)$在$[1, 2]$内单调递增，因此最小值为$h(1) = 3-m$
                \item $m = 0$，$h(x) = 3$，最小值为$3$
                \item $m < 0$，$h(x)$在$[1, 2]$内单调递减，因此最小值为$h(2) = 3$
            \end{dingautolist}
            故当$m \in (-\infty, 3)$时，满足要求.
            \item 问题转化为$g(x)$在$[1, 2]$上的值域是$f(x)$在$[3, 4]$上值域的子集. $f(x)$在$[3, 4]$上值域为$[1, 2]$，$g(x) = m(x-1)^2 + (1-m)$，下面分析$g(x)$的值域
            \begin{dingautolist}{172}
                \item $m > 0$，$g(x)$在$[1, 2]$内单调递增，值域为$[1-m, 1]$
                \item $m = 0$，$g(x) = 1$，值域为$\{1\}$
                \item $m < 0$，$g(x)$在$[1, 2]$内单调递增，值域为$[1, 1-m]$
            \end{dingautolist}
            故当$m \in [-1, 0]$时满足要求.
            \item 问题转化为$g(x)$在$[1, 2]$上的最小值大于$f(x)$在$[0, 2]$上的最小值. $f(x)$在$[0, 2]$上最小值为-2，下面分析$g(x)$的最值
            \begin{dingautolist}{172}
                \item $m > 0$，$g(x)$在$[1, 2]$内单调递增，因此最小值为$g(1) = 1-m$
                \item $m = 0$，$g(x) = -1$，最小值为$-1$
                \item $m < 0$，$g(x)$在$[1, 2]$内单调递减，因此最小值为$g(2) = 1$
            \end{dingautolist}
            故当$m \in (-\infty, 3)$时，满足要求.
        \end{enumerate}
    \end{jieda}
    \item 已知函数 $f\left( x\right)  = x,g\left( x\right)  = a{x}^{2} - x$ ,其中 $a > 0$ ,若对任意的 ${x}_{1} \in  \left\lbrack  {1,3}\right\rbrack$ ,总存在 ${x}_{2} \in$  $\left\lbrack  {1,4}\right\rbrack$ ,使得 $f\left( {x}_{1}\right) f\left( {x}_{2}\right)  = g\left( {x}_{1}\right) g\left( {x}_{2}\right)$ 成立,则实数 $a$ 的取值范围是\blkda[4]{$\left\lbrack  {\dfrac{5}{4},\dfrac{4}{3}}\right\rbrack$}.
    \begin{jieda}
        $x \in [1, 4], f(x) \neq 0$欲满足目标式必有$x \in [1, 3], g(x) \neq 0$
    
        设$h(x) = \dfrac{g(x)}{f(x)} = ax-1$，则$f\left( {x}_{1}\right) f\left( {x}_{2}\right)  = g\left( {x}_{1}\right) g\left( {x}_{2}\right) \Longleftrightarrow \dfrac{1}{h(x_1)} = h(x_2)$

        设$m(x) = \dfrac{1}{h(x)}$，则问题转化为“对任意的 ${x}_{1} \in  \left\lbrack  {1,3}\right\rbrack$ ,总存在 ${x}_{2} \in$  $\left\lbrack  {1,4}\right\rbrack$ ,使得$m(x_1) = h(x_2)$”，即$m(x)$在$[1,3]$的值域是$h(x)$在$[1,4]$的值域的子集.

        \begin{dingautolist}{172}
            \item $a > 1$ \\
            $\left[\dfrac{1}{3a-1}, \dfrac{1}{a-1}\right] \subseteq \left[a-1, 4a-1\right]$，故有$\left\{\begin{array}{l}
             \dfrac{1}{3a-1} \geq a - 1 \\
             \dfrac{1}{a-1} \leq 4a-1
            \end{array}\right.$，解得$\dfrac{5}{4} \leq  a \leq  \dfrac{4}{3}$
            \item $a \in \left(\dfrac{1}{3}, 1\right)$\\
            $\left(-\infty, \dfrac{1}{a-1}\right] \cup \left[\dfrac{1}{3a-1}, +\infty\right) \subseteq \left[a-1, 4a-1\right]$，无解
            \item $a\in \left(0, \dfrac{1}{3}\right)$\\
            $\left[\dfrac{1}{3a-1}, \dfrac{1}{a-1}\right] \subseteq \left[a-1, 4a-1\right]$，故有$\left\{\begin{array}{l}
             \dfrac{1}{3a-1} \geq a - 1 \\
             \dfrac{1}{a-1} \leq 4a-1
            \end{array}\right.$，无解
        \end{dingautolist}
        综上，$a \in \left\lbrack  {\dfrac{5}{4},\dfrac{4}{3}}\right\rbrack$
        
    \end{jieda}
\end{enumerate}
\clearpage
\subsection*{D Level}
\begin{enumerate}
    \item 已知$f(x) = x^2 + \dfrac{1}{x^2} + k \left(x + \dfrac{1}{x}\right)+11, x \in (0, +\infty)$，若$f(x) = 0$有两个解，则实数$k$的取值范围是\blkda[4]{$\left(-\infty, -\dfrac{13}{2}\right) \cup \{-6\}$}
    \begin{jieda}
        首先考虑换元$t = x + \dfrac{1}{x}$，$f(x) = g(t) = t^2+k\cdot t + 9$，当$t = 2$时，只对应一个$x$，当$t > 2$时，对应着两个$x$. 故若$f(x)$在$(0, +\infty)$，有两个零点可以转化为$g(t)$在$(2, +\infty)$有一个零点，进一步转化为$y = -k$与$y = t + \dfrac{9}{t}$在$(2, +\infty)$有一个交点. 因此$k \in \left(-\infty, -\dfrac{13}{2}\right) \cup \{-6\}$ \\
        \ps \textbf{换元后及时判断新元的取值范围，并且分析清楚换元前后的对应关系！}
    \end{jieda}
    \item 若关于$x$的方程$\dfrac{|x|}{x+6} = kx^2$有四个不同的实数解，求实数$k$的取值范围.
    \begin{jieda}
        易知必有一根为0 \\
        参变分离有$\dfrac{1}{k} = \dfrac{(x+6)x^2}{|x|} = \left \{ \begin{array}{ll}
            \dfrac{(x+6)x^2}{x}, & x > 0 \\
            -\dfrac{(x+6)x^2}{x},  & x < 0
        \end{array}\right.$
        由于原方程必有一根为0，故可以转化为方程$\dfrac{1}{k} = \left \{ \begin{array}{ll}
            (x+6)x, & x > 0 \\
            -(x+6)x,  & x < 0
        \end{array}\right.$有三个不等实根.因此$\dfrac{1}{k} \in (0, 9)$，即$k \in \left(\dfrac{1}{9}, +\infty\right)$
    \end{jieda}
    \item 已知函数$f(x) = \left\{  
         \begin{array}{ll}  
            x^3, & x \geq 0, \\  
            -x, & x < 0.
         \end{array}
    \right.$ 若函数$g(x) = f(x) - |kx^2-2x| (k \in \mathbf{R})$恰有4零点，$k$的取值范围是\blkda[3.5]{$(-\infty, 0) \cup (2\sqrt{2}, +\infty)$}.
    \begin{jieda}
        容易判断，$g(x)$有一零点为0. 设$F(x) = \left\{  
         \begin{array}{ll}  
            x^2, & x \geq 0, \\  
            1, & x < 0.
         \end{array}
        \right.$，$H(x) = |kx-2|$，故其余三个零点是函数$G(x) = F(x) - H(x)$的零点.易知$k = 0$时$G(x)$没有三个零点. 
        \begin{dingautolist}{172}
            
            \item 当$k < 0$时，$G(x)$在$\left(-\infty, \dfrac{2}{k}\right)$单调增，在$\left(\dfrac{2}{k}, 0\right)$单调减，在$\left(0, -\dfrac{k}{2}\right)$单调减，在$\left(-\dfrac{k}{2}, +\infty\right)$单调增.$G\left(\dfrac{3}{k}\right) = 0, G\left(\dfrac{1}{k}\right) = 0, G(0) = -2, G(2-k) = 2-2k$，故$G(x)$有三个零点，分别是$\dfrac{3}{k}, \dfrac{1}{k}$，以及一个位于$(0, 2-k)$的零点.
            \item 当$k > 0$时，当$x < 0$，$G(x) = kx-1$，无零点. \\
            当$x \in \left(0, \dfrac{2}{k}\right)$，$G(x)$单调增,$G\left(\dfrac{1}{k}\right) = 0$. \\
            当$x \geq \dfrac{2}{k}$，$G(x) = x^2 - kx + 2$，要使得$G(x)$在$\left[\dfrac{2}{k}, +\infty\right)$有两个零点，必有$\Delta > 0$，即$k > 2\sqrt{2}$. 抛物线的对称轴为$x = \dfrac{k}{2}$，因为当$k > 2\sqrt{2}$时$\dfrac{k}{2} > \dfrac{2}{k}$，又因为$G\left(\dfrac{2}{k}\right) > 0$，所以当$k > 2\sqrt{2}$时，$G(x)$在$\left[\dfrac{2}{k}, +\infty\right)$有两个零点.
        \end{dingautolist}
        综上，$k \in (-\infty, 0) \cup (2\sqrt{2}, +\infty)$ \\
        \ps 从数形结合角变形$H(x)$来看会更直观.
\end{jieda}
    \item 已知二次函数 $f(x)=-\dfrac{1}{2} x^2+x$，是否存在实数 $m, n(m<n)$, 使得 $f(x)$ 的定义域为 $[m, n]$, 值域为 $[k m, k n](k>1)$.
    \begin{jieda}
    $f(x)=-\dfrac{1}{2} x^2+x=-\dfrac{1}{2} \times(x-1)^2+\dfrac{1}{2} \leq \dfrac{1}{2}$
    $\therefore k n \leq \dfrac{1}{2}, \therefore m<n<\dfrac{1}{2}$
    又 $\because$ 函数的对称轴为 $x=1$
    $\therefore f(x)$ 在 $[m, n]$ 单调递增则有 $f(m)=k m, f(n)=k n$
    故$m, n$是$-\dfrac{1}{2}x^2 + x = kx$的两根，故$m = 2(1-k), n = 0$ \\
    \ps 根据值域判断函数在定义域内的单调性是简化，解决问题的突破口。\\
    常规解决方法是分析单调性，根据最值列方程组.
    \end{jieda}
    \item 若函数$f(x)$在区间$[a, b]$上的值域恰是$\left[\dfrac{1}{b}, \dfrac{1}{a}\right]$，则称区间$[a, b]$是$f(x)$的一个倒域区间. 求函数$g(x) = \left \{ \begin{array}{ll}
        -x^2+2x & 0 \leq x \leq 2 \\
        x^2 + 2x & -2 \leq x < 0
    \end{array}\right.$的所有倒域区间.
    \begin{jieda}
        $g(x)$的定义域为$[-2, 2]$，$g(x)$值域为$[-1, 1]$，故有$-2 \leq a < b \leq 2, -1 \leq \dfrac{1}{b} < \dfrac{1}{a} \leq 1$，解得$1 \leq a < b \leq 2$或$-2 \leq a < b \leq -1$.
        \begin{dingautolist}{172}
            \item  $1 \leq a < b \leq 2$，此时$g(x)$单调递减，有$\left \{ \begin{array}{l}
                 -a^2 + 2a = \dfrac{1}{a}  \\
                 -b^2 + 2b = \dfrac{1}{b} 
            \end{array} \right.$，即$a, b$是$-x^2 + 2x = \dfrac{1}{x}$的两根，即$a, b$是$-x^3 + 2x^2 -1 = 0$的两根，可以看出有一根为$1$，即$a, b$是$(x-1)(-x^2+x+1) = 0$的两根，故$a = 1, b = \dfrac{1+\sqrt{5}}{2}$
            \item $-2 \leq a < b \leq -1$，此种情况和前一种对称，可得$a = -\dfrac{1+\sqrt{5}}{2}, b = -1$
        \end{dingautolist}
        因此两个倒域区间为$\left[1, \dfrac{1+\sqrt{5}}{2}\right]$和$\left[-\dfrac{1+\sqrt{5}}{2}, -1\right]$
    \end{jieda}
    \item 若函数$f(x)$在区间$[m, n]$上的值域恰是$[m, n]$，则称区间$[m, n]$是$f(x)$的一个等域区间. 若区间$[m,n]$是函数$y = \dfrac{(a+1)x - \dfrac{1}{a}}{ax}(a \neq 0)$的一个等域区间.求$n-m$的最大值.
    \begin{jieda}
        $\dfrac{(a+1)x - \dfrac{1}{a}}{ax} = \dfrac{a+1}{a} - \dfrac{1}{a^2x}$，故函数为$y = \dfrac{a+1}{a} - \dfrac{1}{a^2x} (a \neq 0)$，易知函数的定义域为$(-\infty, 0) \cup (0, +\infty)$，故必有$m < n < 0$或$0 < m < n$. 故易知函数在$[m, n]$内严格单调递增，故$\left \{ \begin{array}{l}
                \dfrac{a+1}{a} - \dfrac{1}{a^2m} = m \\
                \dfrac{a+1}{a} - \dfrac{1}{a^2n} = n
            \end{array} \right.$，即$m, n$是$a^2x^2 - (a+1)ax + 1 = 0$的两根. 由韦达定理有$\left \{ \begin{array}{l}
                m+n = \dfrac{a+1}{a}\\
                m \cdot n = \dfrac{1}{a^2}
            \end{array} \right.$，由$\Delta > 0$得$((a+1)^2 - 4)a^2 > 0$，故$a \in (-\infty, -3]\cup [1, +\infty)$ \\
            $n - m = \sqrt{(m+n)^2 - 4mn} = \sqrt{(\dfrac{a+1}{a})^2 - \dfrac{4}{a^2}} = \sqrt{\dfrac{a^2+2a-3}{a^2}} \leq \dfrac{2\sqrt{3}}{3}$，当$a = 3$时取等号.
    \end{jieda}
\end{enumerate}

\section*{指数运算、对数运算}
\subsection*{A Level}
\begin{enumerate}
    \item 把下列指数式化为对数式,把对数式化为指数式 $\left( {a > 0\text{ 且 }a \neq  1}\right)$：
    \begin{dingautolist}{172}
        \item ${a}^{1} = a$
        \item ${a}^{0} = 1$
        \item ${\log }_{{a}}N = b$
        \item ${\log }_{a}\sqrt[3]{{a}^{2}} = \dfrac{2}{3}$
    \end{dingautolist}
    \begin{jieda}
        \begin{dingautolist}{172}
            \item ${\log }_{a}a = 1$
            \item ${\log }_{a}1 = 0$
            \item ${a}^{b} = N$
            \item ${a}^{\frac{2}{3}} = \sqrt[3]{{a}^{2}}$
        \end{dingautolist}
    \end{jieda}
    \item 求下列各式的值:
    \begin{enumerate*}[i.]
        \item ${2}^{{\log }_{2}8}$
        \item ${3}^{{\log }_{3}9}$
        \item ${10}^{\lg 5}$
        \item ${e}^{\ln 7}$
        \item ${\log }_{6}({6}^{2})$
        \item $\lg {10}^{-5}$
        \item $\lg {10}^{6}$
        \item $\ln {e}^{3}$
    \end{enumerate*}
    \begin{jieda}
        \begin{enumerate*}[i.]
            \item 8
            \item 9
            \item 5
            \item 7
            \item 2
            \item $-5$
            \item 6
            \item 3
        \end{enumerate*}
    \end{jieda}
    \item 求下列各式的值:
    \begin{enumerate*}[i.]
        \item ${\log }_{2}\dfrac{1}{16}$
        \item ${\log }_{\frac{1}{2}}{32}$
        \item ${\log }_{0.4}1$
        \item ${\log }_{2.5}{6.25}$
        \item ${4}^{{\log }_{4}3}$
        \item ${9}^{{\log }_{3}7 + 1}$
    \end{enumerate*}
    \begin{jieda}
        \begin{enumerate*}[i.]
            \item $-4$
            \item $-5$
            \item 0
            \item 2
            \item 3
            \item 441
        \end{enumerate*}
    \end{jieda}
    \item 已知 $\lg a < 1$ ,化简： $\sqrt{{\left( \lg a\right) }^{2} - \lg \dfrac{{a}^{2}}{10}} = $\blkda[4]{$1 - \lg a$} .
    \begin{jieda}
        由 $\lg \dfrac{{a}^{2}}{10} = \lg {a}^{2} - \lg {10} = 2\lg a - 1$ ,且 $\lg a < 1$ ,
        
        于是 $\sqrt{{\left( \lg a\right) }^{2} - \lg \dfrac{{a}^{2}}{10}} = \sqrt{{\left( \lg a\right) }^{2} - 2\lg a + 1} = \sqrt{{\left( \lg a - 1\right) }^{2}} = \left| {\lg a - 1}\right|  = 1 - \lg a$
    \end{jieda}
    \item 求值: ${7}^{1 + {\log }_{7}5} = $\blkda[4]{35}
    \begin{jieda}
        ${7}^{1 + {\log }_{7}5} = 7 \cdot  {7}^{{\log }_{7}5} = 7 \times  5 = {35}$ .
    \end{jieda}
    
    \item 计算下列各式的值.
    \begin{enumerate}
        \item ${\log }_{2}{10} \times  \lg 4 =$ \blkda[4]{2}
        \item ${\log }_{2}3 \times  {\log }_{\sqrt{3}}8 =$ \blkda[4]{6}
        \item ${\log }_{3}4 \times  {\log }_{4}5 \times  {\log }_{5}6 \times  {\log }_{6}7 \times  {\log }_{7}8 \times  {\log }_{8}9 =$ \blkda[4]{2}
    \end{enumerate}
    \begin{jieda}
        \begin{enumerate}
            \item ${\log }_{2}{10} \times  \lg 4 = \dfrac{1}{\lg 2} \times  2\lg 2 = 2$
            \item ${\log }_{2}3 \times  {\log }_{\sqrt{3}}8 = \dfrac{\lg 3}{\lg 2} \times  \dfrac{\lg {2}^{3}}{\lg {3}^{\frac{1}{2}}} = \dfrac{\lg 3}{\lg 2} \times  \dfrac{3\lg 2}{\dfrac{1}{2} \times  \lg 3} = 6$
            \item ${\log }_{3}4 \times  {\log }_{4}5 \times  {\log }_{5}6 \times  {\log }_{6}7 \times  {\log }_{7}8 \times  {\log }_{8}9 = \dfrac{\lg 4}{\lg 3} \times  \dfrac{\lg 5}{\lg 4} \times  \dfrac{\lg 6}{\lg 5} \times  \dfrac{\lg 7}{\lg 6} \times  \dfrac{\lg 8}{\lg 7} \times  \dfrac{\lg 9}{\lg 8} = \dfrac{\lg 9}{\lg 3} = \dfrac{\lg {3}^{2}}{\lg 3} = \dfrac{\lg {3}^{2}}{\lg 3} = 2$
        \end{enumerate}
    \end{jieda}
    \item $\left( {{\log }_{3}2 + {\log }_{9}2}\right) \left( {{\log }_{4}3 + {\log }_{8}3}\right)  =$ \blkda[4]{$\dfrac{5}{4}$}.
    \begin{jieda}
        $\left( {{\log }_{3}2 + {\log }_{9}2}\right) \left( {{\log }_{4}3 + {\log }_{8}3}\right)  = \left( {\dfrac{\lg 2}{\lg 3} + \dfrac{\lg 2}{\lg 9}}\right) \left( {\dfrac{\lg 3}{\lg 4} + \dfrac{\lg 3}{\lg 8}}\right)  \\
        = \left( {\dfrac{\lg 2}{\lg 3} + \dfrac{\lg 2}{2\lg 3}}\right) \left( {\dfrac{\lg 3}{2\lg 2} + \dfrac{\lg 3}{3\lg 2}}\right) = \dfrac{3\lg 2}{2\lg 3} \cdot  \dfrac{5\lg 3}{6\lg 2} = \dfrac{5}{4}$
    \end{jieda}
    
    \item 计算： $\left( {{\log }_{2}5 + {\log }_{4}{125}}\right) \left( {{\log }_{5}4 + {\log }_{25}{64}}\right) = $\blkda[4]{$\dfrac{25}{2}$}
    \begin{jieda}
        原式 $= \left( {{\log }_{2}5 + {\log }_{{2}^{2}}{5}^{3}}\right) \left( {{\log }_{5}{2}^{2} + {\log }_{5^2}{2}^{6}}\right)$

        $= \left( {{\log }_{2}5 + \dfrac{3}{2}{\log }_{2}5}\right) \left( {2{\log }_{5}2 + 3{\log }_{5}2}\right)$
        
        $= \dfrac{5}{2}{\log }_{2}5 \times  5{\log }_{5}2 = \dfrac{25}{2}{\log }_{2}5 \times  {\log }_{5}2 = \dfrac{25}{2}$ .
    \end{jieda}
    \item 化简: ${\log }_{2}\sqrt{\dfrac{7}{48}} + {\log }_{2}{12} - \dfrac{1}{2}{\log }_{2}{42} - 1$ .
    \begin{jieda}
        原式 $= {\log }_{2}\dfrac{\sqrt{7}}{\sqrt{48}} + {\log }_{2}{12} - {\log }_{2}\sqrt{42} - {\log }_{2}2$
        
        $= {\log }_{2}\dfrac{\sqrt{7} \times  {12}}{\sqrt{48} \times  \sqrt{42} \times  2}$
        
        $= {\log }_{2}\dfrac{\sqrt{7} \times  {12}}{4\sqrt{3} \times  \sqrt{7} \times  \sqrt{6} \times  2}$
        
        $= {\log }_{2}\dfrac{1}{2\sqrt{2}} = {\log }_{2}{2}^{-\frac{3}{2}} =  - \dfrac{3}{2}$ .
    \end{jieda}
    
    \item 计算: ${9}^{\frac{1}{2}{\log }_{3}4} + {5}^{1 + {\log }_{5}2} + {\log }_{3}\sqrt{3} + {{e}}^{\ln 3}$ .
    \begin{jieda}
        ${9}^{\frac{1}{2}{\log }_{3}4} + {5}^{1 + {\log }_{5}2} + {\log }_{3}\sqrt{3} + {{e}}^{\ln 3} = {3}^{{\log }_{3}4} + {5}^{{\log }_{5}\left( {2 \times  5}\right) } + \dfrac{1}{2}{\log }_{3}3 + 3 = 4 + {10} + \dfrac{1}{2} + 3 = {17.5}$ .
    \end{jieda}
    
    \item 设 ${2}^{x} = {3}^{y} = {72}$ ,则 $\dfrac{3}{x} + \dfrac{2}{y} =$ \blkda[4]{1}.
    \begin{jieda}
        由 ${2}^{x} = {3}^{y} = {72}$ ,可得 $x = {\log }_{2}{72},y = {\log }_{3}{72}$ ,
        
        所以 $\dfrac{3}{x} + \dfrac{2}{y} = 3{\log }_{72}2 + 2{\log }_{72}3 = {\log }_{72}9 + {\log }_{72}8 = {\log }_{72}{72} = 1$ .
    \end{jieda}
    \item 已知 $\lg a = {2.31},\lg b = {1.31}$ ,则 $\dfrac{a}{b} =$ \blkda[4]{10}.
    \begin{jieda}
        因为 $\lg a = {2.31},\lg b = {1.31}$，所以 $\lg a - \lg b = \lg \dfrac{a}{b} = 1$，所以 $\dfrac{a}{b} = {10}$ .
    \end{jieda}
\end{enumerate}
\subsection*{B Level}
\begin{enumerate}
    
    \item 已知 ${\log }_{8}9 = a,{\log }_{3}5 = b$ ,试用 $a,b$ 表示 $\lg 5$ .
    \begin{jieda}
        ${\log }_{8}9 = \dfrac{{\log }_{3}9}{{\log }_{3}8} = \dfrac{2}{3{\log }_{3}2} = a,\therefore {\log }_{3}2 = \dfrac{2}{3a}$
        
        $\lg 5 = \dfrac{{\log }_{3}5}{{\log }_{3}{10}} = \dfrac{{\log }_{3}5}{{\log }_{3}5 + {\log }_{3}2} = \dfrac{b}{b + \dfrac{2}{3a}} = \dfrac{3ab}{2 + {3ab}}.$
    \end{jieda}
    \item 若$\log_427=m$, $\log_325=n$, 则$\lg2$可用$m,n$表示为\blkda{$\dfrac{3}{3+mn}$}. 
        \begin{jieda}
            $m = \dfrac{3log_23}{2}$，$n = \dfrac{2log_25}{log_23}$，$lg2 = \dfrac{1}{1+log_25}$，故$log_25 = \dfrac{m\cdot n}{2}$，$lg2 = \dfrac{3}{3+mn}$ \\
            \ps 对数运算表示类的题目，首先要统一底数，选取题干中的一个质因子作为底数，然后将所有的对数式都拆彻底（底数和真数互质），然后再凑.
        \end{jieda}
    \item 若 ${\log }_{5}x + {\log }_{5}y = 2{\log }_{5}\left( {x - {2y}}\right)$ ,则 $\dfrac{x}{y} =$ \blkda[4]{4}.
    \begin{jieda}
        因为 ${\log }_{5}x + {\log }_{5}y = 2{\log }_{5}\left( {x - {2y}}\right)$ ,故 ${\log }_{5}{xy} = {\log }_{5}{\left( x - 2y\right) }^{2}$ ,
        
        所以 $\left\{  \begin{matrix} x > 0,y > 0 \\  x - {2y} > 0 \\  {xy} = {\left( x - 2y\right) }^{2} \end{matrix}\right.$ ,由 ${xy} = {\left( x - 2y\right) }^{2}$ 得 ${x}^{2} - {5xy} + 4{y}^{2} = 0 \Longrightarrow  x = y$ 或 $x = {4y}$ ,
        
        又 $x > 0,y > 0,x - {2y} > 0$ ,所以舍去 $x = y$ ,故 $x = {4y}$ ,则 $\dfrac{x}{y} = 4$ .
    \end{jieda}
    \item 若 $a > 0$ 且 $a \neq  1,b > 0,c > 0$ ,给出以下式子:
    \begin{dingautolist}{172}
        \item ${\log }_{a}\dfrac{b}{c} = \dfrac{{\log }_{a}b}{{\log }_{a}c}$
        \item ${\log }_{a}\left( {b \cdot  c}\right)  = {\log }_{a}\left( {b + c}\right)$
        \item ${\log }_{a}\left( {b \cdot  c}\right)  = {\log }_{a}b + {\log }_{a}c$
        \item ${\log }_{a}\left( {b - c}\right)  = \dfrac{{\log }_{a}b}{{\log }_{a}c}$
        \item ${\log }_{a}\left( {b + c}\right)  = {\log }_{a}b \cdot  {\log }_{a}c$
        \item ${\log }_{a}\dfrac{b}{c} = {\log }_{a}b - {\log }_{a}c$
    \end{dingautolist}
    则正确的式子为\blkda[4]{\ding{174}\ding{177}}(填序号). 
    \item 通常我们以分贝(dB)为单位来表示声音大小的等级,如果强度为 $v$ 的声音对应的分贝数为 $f\left( v\right) {{dB}}$ ,那么满足: $f\left( v\right)  = {10} \times  \lg \dfrac{v}{1 \times  {10}^{-{12}}}$ ,若在地铁中多人外放电子设备加上行车噪音,车厢内的声音的分贝能达到 ${90}{{dB}}$ ,则 ${90}{{dB}}$ 的声音与 ${50}{{dB}}$ 的声音强度之比为\blkda[4]{${10000}$}.
    \begin{jieda}
        当 $f\left( v\right)  = {90}$ 时,则有 ${10} \times  \lg \dfrac{v}{1 \times  {10}^{-{12}}} = {90}$ ,解得 ${v}_{1} = {10}^{-3}$ ,
        
        当 $f\left( v\right)  = {50}$ 时,则有 ${10} \times  \lg \dfrac{v}{1 \times  {10}^{-{12}}} = {50}$ ,解得 ${v}_{2} = {10}^{-7}$ ,
        
        所以当 $\dfrac{{v}_{1}}{{v}_{2}} = \dfrac{{10}^{-3}}{{10}^{-7}} = {10}^{4} = {10000}$ .
    \end{jieda}
    \item \begin{enumerate}
        \item 设 ${3}^{x} = {4}^{y} = {36}$ ,求 $\dfrac{2}{x} + \dfrac{1}{y}$ 的值
        \item 若 $x{\log }_{2}3 = 1$ ,求 ${3}^{x} + {9}^{-x}$ 的值
    \end{enumerate}
    \begin{jieda}
        \begin{enumerate}
            \item 由已知得 $x = {\log }_{3}{36},y = {\log }_{4}{36}$ ,

            所以 $\dfrac{2}{x} + \dfrac{1}{y} = \dfrac{2}{{\log }_{3}{36}} + \dfrac{1}{{\log }_{4}{36}} = 2{\log }_{36}3 + {\log }_{36}4 = {\log }_{36}9 + {\log }_{36}4 = {\log }_{36}{36} = 1$
            \item 因为 $x{\log }_{2}3 = 1$ ,所以 ${\log }_{2}{3}^{x} = 1$ ,所以 ${3}^{x} = 2$ ,

            所以 ${3}^{x} + {9}^{-x} = {3}^{x} + {\left( {3}^{x}\right) }^{-2} = 2 + \dfrac{1}{4} = \dfrac{9}{4}$ .
        \end{enumerate}
    \end{jieda} 
    
    \item 方程 $\dfrac{50}{{10}^{x} + 1} - 5 = {10}^{x + 1}$ 的实数解为\blkda[4]{$\lg \dfrac{3}{2}$}.
    \begin{jieda}
        令 ${10}^{x} = t$ ,则 $\dfrac{50}{t + 1} - 5 = {10t}$ ,整理得 $2{t}^{2} + {3t} - 9 = \left( {t + 3}\right) \left( {{2t} - 3}\right)  = 0$ ,解得 $t =  - 3$ (舍去) 或 $t = \dfrac{3}{2}$ ,所以 ${10}^{x} = \dfrac{3}{2}$ ,故 $x = \lg \dfrac{3}{2}$ .
    \end{jieda}
    
    \item 方程 ${5}^{{5x} + 1} = {3}^{x - 1}$ 的解集是\blkda[4]{$\left\{\dfrac{lg3+lg5}{lg3-5lg5}\right\}$}
    \begin{jieda}
        两侧同时以$5$取对数有$5x+1 = (x-1)log_53$，整理可得$(5-log_53)x = -(1+log_53)$，故$x = \dfrac{1+log_53}{log_53-5} = \dfrac{lg3+lg5}{lg3-5lg5}$
    \end{jieda}
    \item 方程 ${2}^{2 + x} - {2}^{2 - x} = {15}$ 的解集是\blkda[4]{$\left\{2\right\}$}
    \begin{jieda}
        设$t = 2^x, t > 0$故$4t-\dfrac{4}{t} = 15$，解得$t = 4$或$-\dfrac{1}{4}$（舍去），故$x = 2$
    \end{jieda}
    \item 方程 ${2}^{\left| x + 1\right| } = 3$ 的解集是\blkda[4]{$\left\{  {{\log }_{2}\dfrac{3}{2}, - {\log }_26}\right\}$}
    \begin{jieda}
        ${2}^{\left| x + 1\right| } = 3 \Longleftrightarrow |x+1| = log_23 \Longleftrightarrow x+1 = \pm log_23$，故解集为$\left\{  {{\log }_{2}\dfrac{3}{2}, - {\log }_26}\right\}$
    \end{jieda}
    
    \item 若 ${\log }_{3}\left( {{\log }_{4}x}\right)  = 1$ ,则 $x =$ \blkda[4]{64}.
    \begin{jieda}
        ${\log }_{3}\left( {{\log }_{4}x}\right)  = 1$ ,则 ${\log }_{4}x = 3$ ,故 $x = {64}$ .
    \end{jieda} 

    \item 方程 ${\log }_{7}\left\lbrack  {{\log }_{2}\left( {{\log }_{x}{16}}\right) }\right\rbrack   = 0$ 的解集是\blkda[4]{$\{4\}$}
    \begin{jieda}
        ${\log }_{7}\left\lbrack  {{\log }_{2}\left( {{\log }_{x}{16}}\right) }\right\rbrack   = 0 \Longleftrightarrow {{\log }_{2}\left( {{\log }_{x}{16}}\right) } = 1 \Longleftrightarrow log_x16 = 2 \Longleftrightarrow x = 4$
    \end{jieda}
    \item 已知 ${\log }_{{x}^{2} - 1}\left( {2{x}^{2} - {3x} + 1}\right)  = 1$，则$x = $\blkda[4]{2}
    \begin{jieda}
    
        ${\log }_{{x}^{2} - 1}\left( {2{x}^{2} - {3x} + 1}\right)  = 1 \Longleftrightarrow  \left\{\begin{array}{l}
             {x}^{2} - 1 = 2{x}^{2} - {3x} + 1  \\
             x^2-1 \in (0, 1) \cup (1, +\infty) \\
             2{x}^{2} - {3x} + 1 > 0
        \end{array}\right . \Longleftrightarrow \left\{\begin{array}{l}
             {x}^{2} - {3x} + 2 = 0  \\
             x^2 \in (1, 2) \cup (2, +\infty) \\
             2{x}^{2} - {3x} + 1 > 0
        \end{array}\right .$ ,

        解得: $x = 1$ (舍去) 或 $x = 2$
    \end{jieda}
    
    \item 方程 ${\log }_{4}\left( {{3x} - 1}\right)  = {\log }_{4}\left( {x - 1}\right)  + {\log }_{4}\left( {3 + x}\right)$ 的解集是\blkda[4]{$\{2\}$}
    \begin{jieda}
        ${\log }_{4}\left( {{3x} - 1}\right)  = {\log }_{4}\left( {x - 1}\right)  + {\log }_{4}\left( {3 + x}\right) \Longleftrightarrow \left\{\begin{array}{l}
             x > 1  \\
             x > -3 \\
             3x-1 = (x-1)(3+x)
        \end{array}\right.$

        解得: $x = -1$ (舍去) 或 $x = 2$
    \end{jieda}
    \item 方程 ${\log }_{2}\left( {x - 1}\right)  = 2 - {\log }_{2}\left( {x + 1}\right)$ 的解集是\blkda[4]{$\{\sqrt{5}\}$}
    \begin{jieda}
        ${\log }_{2}\left( {x - 1}\right)  = 2 - {\log }_{2}\left( {x + 1}\right) \Longleftrightarrow \left\{\begin{array}{l}
             x > -1  \\
             x > 1 \\
             x-1 = \dfrac{4}{x+1}
        \end{array}\right.$，解得$x = \sqrt{5}$（负舍）
    \end{jieda}
    \item ${\log }_{2}\left( {x + 4}\right)  + {\log }_{2}\left( {x - 1}\right)  = 1 + {\log }_{2}\left( {x + 8}\right)$ 的解是\blkda[4]{$x = 4$}.
    \begin{jieda}
        由 $x + 4 > 0,x - 1 > 0,x + 8 > 0$ 得 $x$ 取值范围是 $x > 1$
        
        $\because {\log }_{2}\left( {x + 4}\right)  + {\log }_{2}\left( {x - 1}\right)  = 1 + {\log }_{2}\left( {x + 8}\right) ,\therefore {\log }_{2}\left( {x + 4}\right)  + {\log }_{2}\left( {x - 1}\right)  = {\log }_{2}2 + {\log }_{2}\left( {x + 8}\right)$
        
        $\therefore {\log }_{2}\left( {x + 4}\right) \left( {x - 1}\right)  = {\log }_{2}2\left( {x + 8}\right)$ 即 $\left( {x + 4}\right) \left( {x - 1}\right)  = 2\left( {x + 8}\right)$
        
        化简得 ${x}^{2} + x - {20} = 0$ ,解得 $x =  - 5$ (舍) 或 $x = 4$
        
        即 ${\log }_{2}\left( {x + 4}\right)  + {\log }_{2}\left( {x - 1}\right)  = 1 + {\log }_{2}\left( {x + 8}\right)$ 的解是 $x = 4$
    \end{jieda}
    
\end{enumerate}
\subsection*{C Level}
\begin{enumerate}
    \item 解方程 ${4}^{x} + {4}^{-x} - 5\left( {{2}^{x} + {2}^{-x}}\right)  + 6 = 0$
    \begin{jieda}
        设$t = {2}^{x} + {2}^{-x}, t \geq 2$，则${4}^{x} + {4}^{-x} = t^2 - 2$，故有$t^2-5t+4=0$，解得$t = 4$或$1$（舍去）

        因此${2}^{x} + {2}^{-x} = 4$，设$p = 2^x$，故$p + \dfrac{1}{p} = 4$，解得$p = 2 \pm \sqrt{3}$，故原方程解集为$\left\{log_2(2+\sqrt{3}), log_2(2-\sqrt{3})\right\}$
    \end{jieda}
    \item 求方程 ${9}^{x} + {4}^{x} = \dfrac{5}{2} \cdot  {6}^{x}$ 的解集
    \begin{jieda}
        ${9}^{x} + {4}^{x} = \dfrac{5}{2} \cdot  {6}^{x} \Longleftrightarrow \left(\dfrac{3}{2}\right)^x+\left(\dfrac{2}{3}\right)^x = \dfrac{5}{2}$，故$\left(\dfrac{3}{2}\right)^x = 2$或$\dfrac{1}{2}$，因此解集为$\left\{\log_{\frac{3}{2}}2, log_{\frac{3}{2}}\dfrac{1}{2}\right\}$
    \end{jieda}
    
    \item 解方程：${\log }_{2}\left( {{4}^{x} + 4}\right)  = x + {\log }_{2}\left( {{2}^{x + 1} - 3}\right)$ .
    \begin{jieda}
        ${\log }_{2}\left( {{4}^{x} + 4}\right)  = x + {\log }_{2}\left( {{2}^{x + 1} - 3}\right) \Longleftrightarrow {\log }_{2}\left( {{4}^{x} + 4}\right)  = {\log }_{2}[\left( {{2}^{x + 1} - 3}\right)\cdot \left(2^x\right)] \\
        \Longleftrightarrow 4^x+4 = \left( {{2}^{x + 1} - 3}\right)\cdot \left(2^x\right)$， 令$t = 2^x, t > 0$，则$t^2+4 = (2t-3)t$，解得$t = 4$或$-1$（舍去）
    
        故原方程解集为：$\{2\}$
    \end{jieda}
    
    \item 解方程: ${\log }_{3}\left( {{3}^{x} - 1}\right)  \cdot  {\log }_{3}\left( {{3}^{x - 1} - \dfrac{1}{3}}\right)  = 2$ .
    \begin{jieda}
        令$t = {\log }_{3}\left( {{3}^{x - 1} - \dfrac{1}{3}}\right)$，则有$t(t+1) = 2$，解得$t \in \{1, -2\}$
        
        故原方程解集为：$\{log_310, log_34 - 1\}$
    \end{jieda}
\end{enumerate}
\section*{幂函数及其性质}
\subsection*{A Level}
\begin{enumerate}
    \item 下列函数中, 是奇函数又是严格增函数的为 \dotfill{(\kh{B})}
    \xx{$y = x + 1$}{$y = {x}^{\frac{1}{3}}$}{$y = \dfrac{1}{x}$}{$y =  - {x}^{3}$}
    \item 下列函数中,在 $\left( {0, + \infty }\right)$ 上单调递增的是 \dotfill{(\kh{C})}
    \xx{$y =  - {x}^{2} + 1$}{$y = \left| {x - 1}\right|$}{$y = {x}^{3}$}{$y = \dfrac{1}{x}$}
    \item 图中 ${C}_{1},{C}_{2},{C}_{3}$ 分别为幂函数 $y = {x}^{{\alpha }_{1}},y = {x}^{{\alpha }_{2}},y = {x}^{{\alpha }_{3}}$ 在第一象限内的图象,则 ${\alpha }_{1}$ , ${\alpha }_{2},{\alpha }_{3}$ 依次可以是\dotfill{(\kh{D})}
    
    \begin{center}
    \includegraphics[width=0.2\textwidth]{bodyxb/src/bo_d2em7r3ef24c73b07a3g_1_242_562_316_304_0.jpg}
    \end{center}
    \xx{$\dfrac{1}{2},3, - 1$}{$- 1,3,\dfrac{1}{2}$}{$\dfrac{1}{2}, - 1,3$}{$- 1,\dfrac{1}{2},3$}
    \begin{jieda}
        由题图知: ${\alpha }_{1} < 0,0 < {\alpha }_{2} < 1,{\alpha }_{3} > 1$ ,

        所以 ${\alpha }_{1},{\alpha }_{2},{\alpha }_{3}$ 依次可以是 $- 1,\dfrac{1}{2},3$ .
    \end{jieda}
    \item 现有下列函数: 
    \begin{enumerate*}
        \item $y = {x}^{3}$
        \item $y = {\left( \dfrac{1}{2}\right) }^{x}$
        \item $y = 4{x}^{2}$
        \item $y = {x}^{5} + 1$
        \item $y = {\left( x - 1\right) }^{2}$
        \item $y = x$
        \item $y = {a}^{x}\left( {a > 1}\right)$
    \end{enumerate*}
    其中幂函数的序号为\blkda[4]{(1)(6)}
    \item 若幂函数 $y = {x}^{a}$ 的图象经过点 $\left( {\sqrt[4]{3},3}\right)$ ,则实数 $a =$ \blkda[4]{4}.
    \begin{jieda}
        因为幂函数 $y = {x}^{a}$ 的图象经过点 $\left( {\sqrt[4]{3},3}\right)$ ,所以 $3 = {\left( \sqrt[4]{3}\right) }^{a}$，所以 $3 = {3}^{\frac{a}{4}}$ ,所以 $a = 4$
    \end{jieda} 
    \item 函数 $y = {\left( x + 2\right) }^{\frac{1}{2}}$ 的定义域为\blkda[4]{$\lbrack  - 2, + \infty )$}.
    \begin{jieda}
        要使 $y = {\left( x + 2\right) }^{\frac{1}{2}} = \sqrt{x + 2}$ 有意义,则 $x + 2 \geq  0$ ,解得 $x \geq   - 2$ .
    \end{jieda}
    \item 若幂函数 $y = {x}^{-{m}^{2} + {2m}}$ ( $m$ 为整数)的定义域为 $\mathbf{R}$ ,则 $m$ 的值为\blkda[4]{1}.
    \begin{jieda}
        若幂函数 $y = {x}^{-{m}^{2} + {2m}}$ ( $m$ 为整数) 的定义域为 $\mathbf{R}$ ,则 $- {m}^{2} + {2m} > 0$ ,解得$0 < m < 2$，而 $m$ 是整数,则只能 $m = 1$ ,经检验符合题意.
    \end{jieda}
    \item 若幂函数 $y = \left( {{a}^{2} - a - 5}\right) {x}^{a}$ 的图像关于 $y$ 轴对称,则实数 $a =$ \blkda[4]{$-2$}.
    \begin{jieda}
        由幂函数可得 ${a}^{2} - a - 5 = 1$ ,解得 $a = 3$ 或 $a =  - 2$ ,
        
        又因为函数图像关于 $y$ 轴对称,则 $a$ 为偶数,所以 $a =  - 2$ .
    \end{jieda}
    \item 幂函数 $y = f\left( x\right)$ 的图像经过点 $\left( {\dfrac{1}{2},4}\right)$ ,则 $f\left( \dfrac{\sqrt{2}}{2}\right)$ 的值为\blkda[4]{2}.
    \begin{jieda}
        设幂函数 $f\left( x\right)  = {x}^{\alpha }$ ,将 $\left( {\dfrac{1}{2},4}\right)$ 代入 $f\left( x\right)$ ,可得: $\alpha  =  - 2$，所以 $f\left( x\right)  = {x}^{-2}$ ,所以 $f\left( \dfrac{\sqrt{2}}{2}\right)  = {\left( \dfrac{\sqrt{2}}{2}\right) }^{-2} = 2$
    \end{jieda}
    \item 若幂函数 $y = f\left( x\right)$ 的图像经过点 $\left( {-2, - \dfrac{1}{8}}\right)$ ,则满足 $f\left( x\right)  = {27}$ 的 $x$ 的值是\blkda{$\dfrac{1}{3}$}
    \begin{jieda}
        $y = x^{-3}$
    \end{jieda}
    \item 当 $\alpha  \in  \mathbf{R}$ 时,函数 $y = {x}^{\alpha } - 2$ 的图象恒过定点 $A$ ,则点 $A$ 的坐标为\blkda[4]{$(1, -1)$}.
    \begin{jieda}
        由于对任意的 $\alpha  \in  \mathbf{R},y = {x}^{\alpha }$ 恒经过点(1,1),所以函数 $y = {x}^{\alpha } - 2$ 的图象恒过定点 $A\left( {1, - 1}\right)$ 
    \end{jieda}
    
    \item 已知 $f\left( x\right)  = {\left( 2x - 1\right) }^{n} + 1$ ,则函数 $y = f\left( x\right)$ 的图象恒过的定点 $P$ 的坐标为\blkda[4]{$(1,2)$}.
    \begin{jieda}
        令 ${2x} - 1 = 1$ ,得 $x = 1,y = 2$，故函数 $f\left( x\right)$ 图象过定点 $P\left( {1,2}\right)$
    \end{jieda}
    
    \item 若幂函数 $y = \left( {{m}^{2} - m - 1}\right) {x}^{m}$ 为偶函数,则 $m =$ \blkda[4]{2}.
    \begin{jieda}
        $\because$ 函数 $y = \left( {{m}^{2} - m - 1}\right) {x}^{m}$ 为幂函数,
        
        $\therefore {m}^{2} - m - 1 = 1$ ,解得 $m = 2$ 或 $m =  - 1$ ,
        
        又 $\because y = {x}^{m}$ 为偶函数，$\therefore m = 2$
    \end{jieda}
    
    \item 已知集合 $A = \left\{  {y \mid  y = {x}^{2},x \in  R}\right\}$ ,集合 $B = \left\{  {y \mid  y =  - {x}^{2} + {3x} - 1,x \in  R}\right\}$ ,集合 $C$ 为函数$f\left( x\right)  = \sqrt{-{x}^{2} + {4x} + m - 7}$ 的定义域.
    \begin{enumerate}
        \item 求 $A \cap  B$
        \item 若 $A \cup  C \subseteq  A$ ,求实数 $m$ 的取值范围
    \end{enumerate}
    \begin{jieda}
        集合 $A = \left\{  {y \mid  y = {x}^{2},x \in  R}\right\}   = \{ y \mid  y \geq  0\}$ ,集合 $B = \left\{  {y \mid  y =  - {x}^{2} + {3x} - 1}\right.,x \in  R\}  = \left\{  {y\left| {y =  - {\left( x - \dfrac{3}{2}\right) }^{2} + \dfrac{5}{4} \leq  \dfrac{5}{4}}\right. }\right\}  ,$
        \begin{enumerate}
            \item $A \cap  B = \{ y \mid  y \geq  0\}  \cap  \left\{  {y\left| {y =  - {\left( x - \dfrac{3}{2}\right) }^{2} + \dfrac{5}{4} \leq  \dfrac{5}{4}}\right. }\right\}   = \left\{ y \left |  0 \leq  y \leq  \dfrac{5}{4}\right\}\right.$
            \item 集合 $C$ 为函数 $f\left( x\right)  = \sqrt{-{x}^{2} + {4x} + m - 7}$ 的定义域,
            
            $\therefore  - {x}^{2} + {4x} + m - 7 \geq  0$ . 解得: $2 - \sqrt{m - 3} \leq  x \leq  2 + \sqrt{m - 3}$ 且 $m \geq  3$ .
            
            所以 $C = \left\{  {x\left| {2 - \sqrt{m - 3} \leq  x \leq  2 + \sqrt{m - 3}\text{ 且 }m \geq  3}\right. }\right\}$ ;
            
            若 $A \cup  C \subseteq  A$ ,所以 $C \subseteq  A$ ,则 $\left\{  {\begin{array}{l} 2 - \sqrt{m - 3} \geq  0, \\  m \geq  3, \end{array}} \right.$，解得: $3 \leq m \leq  7$
        \end{enumerate}
    \end{jieda}
\end{enumerate}
\subsection*{B Level}
\begin{enumerate}
    \item 下列幂函数中,其图象关于 $y$ 轴对称且过点 $\left( {0,0}\right) \text{ 、 }\left( {1,1}\right)$ 的是 \dotfill{(\kh{B})}  
    \xx{$y = {x}^{\frac{1}{2}}$}{$y = {x}^{4}$}{$y = {x}^{-2}$}{$y = {x}^{\frac{1}{3}}$}
    \item 已知幂函数 $y = f\left( x\right)$ 经过点 $\left( {3,\sqrt{3}}\right)$ ,则 $f\left( x\right)$ \dotfill{(\kh{D})}  
    \xx{是偶函数,且在 $\left( {0, + \infty }\right)$ 上是增函数}{是偶函数,且在 $\left( {0, + \infty }\right)$ 上是减函数}{是奇函数,且在 $\left( {0, + \infty }\right)$ 上是减函数}{是非奇非偶函数,且在 $\left( {0, + \infty }\right)$ 上是增函数}
    \begin{jieda}
        设幂函数的解析式为 $y = {x}^{\alpha }$ ,
        
        将点 $\left( {3,\sqrt{3}}\right)$ 的坐标代入解析式得 ${3}^{\alpha } = \sqrt{3}$ ,解得 $\alpha  = \frac{1}{2}$ ,
        
        $\therefore y = {x}^{\frac{1}{2}}$ ,函数的定义域为 $\lbrack 0, + \infty )$ ,是非奇非偶函数,且在 $\left( {0, + \infty }\right)$ 上是增函数
    \end{jieda}
    \item 如图,下列 3 个幂函数的图象,则其图象对应的函数可能是\dotfill{(\kh{A})} 

    \threetu{\includegraphics[width=0.6\textwidth]{bodyxb/src/bo_d2f9lh3ef24c73b0d7eg_1_240_194_185_170_0.jpg}}{\includegraphics[width=0.6\textwidth]{bodyxb/src/bo_d2f9lh3ef24c73b0d7eg_1_425_194_175_169_0.jpg}}{\includegraphics[width=0.6\textwidth]{bodyxb/src/bo_d2f9lh3ef24c73b0d7eg_1_603_193_175_166_0.jpg}}
    
    \xx{\ding{172}$y = {x}^{-1}$\ding{173}$y = {x}^{\frac{1}{2}}$\ding{174}$y = {x}^{\frac{1}{3}}$}
    {\ding{172}$y = {x}^{-1}$\ding{173}$y = {x}^{\frac{1}{3}}$\ding{174}$y = {x}^{\frac{1}{2}}$}
    {\ding{172}$y = {x}^{\frac{1}{3}}$\ding{173}$y = {x}^{\frac{1}{2}}$\ding{174}$y = {x}^{-1}$}
    {\ding{172}$y = {x}^{\frac{1}{3}}$\ding{173}$y = {x}^{-1}$\ding{174}$y = {x}^{\frac{1}{2}}$}
    \begin{jieda}
        由函数 $y = {x}^{-1} = \frac{1}{x}$ 是反比例函数,其对应图象为\ding{172};
        
        函数 $y = {x}^{\frac{1}{2}} = \sqrt{x}$ 的定义域为 $\left( {0, + \infty }\right)$ ,应为图\ding{173}:
        
        因为 $y = {x}^{\frac{1}{3}}$ 的定义域为 $\mathbf{R}$ 且为奇函数,故应为图\ding{174}.
    \end{jieda}
    \item 设集合 $A = (1, +\infty)$, $B = \left\{  {y \mid  y = {x}^{k},x \in  A}\right\}$ (其中 $k$ 是常数),则 “ $k < 0$ ” 是 “ $A \cap  B = \varnothing$ ” 的\dotfill{(\kh{A})}
    \xx{充分非必要条件}{必要非充分条件}{充分必要条件}{既非充分又非必要条件}
    \begin{jieda}
        $k < 0$时$B = (0, 1)$可知充分性，$k = 0$时$B = \{1\}$，可知非必要性.
    \end{jieda}
    \item 下列命题中正确的是\dotfill{(\kh{D})}  
    \xx{当 $\alpha  = 0$ 时,函数 $y = {x}^{\alpha }$ 的图像是一条直线}{幂函数的图像都经过(0,0)和(1,1)点}{幂函数 $y = {x}^{-\frac{3}{2}}$ 的定义域为 $\lbrack 0, + \infty )$}{幂函数的图像不可能出现在第四象限}
    \begin{jieda}
        对于 ${A},\alpha  = 0$ 时,函数 $y = {x}^{0}$ 的图像是一条直线除去(0,1)点,故 ${A}$ 错误;
        
        对于 ${B}$ ,幂函数的图像都经过(1,1)点,当指数大于 0 时,都经过(0,0)点,当指数小于 0 时, 不经过(0,0)点,故 ${B}$ 错误;
        
        对于 ${C}$ ,函数 $y = {x}^{-\frac{3}{2}} = \dfrac{1}{{x}^{\frac{3}{2}}} = \dfrac{1}{\sqrt{{x}^{3}}}$ ,故定义域为 $\left( {0, + \infty }\right)$ ,故错误;
        
        对于 ${D}$ ,由幂函数的性质,幂函数的图像一定过第一象限,不可能出现在第四象限,故正确.
    \end{jieda}
    \item 已知幂函数 $y = f\left( x\right)$ 的图像过点 $\left( {3,\dfrac{\sqrt{3}}{3}}\right)$ ,则该幂函数的值域是\blkda[4]{$\left( {0, + \infty }\right)$}.
    \begin{jieda}
        设幂函数 $f\left( x\right)  = {x}^{\alpha },\alpha  \in  \mathbf{R}$ ,
        
        代入点 $\left( {3,\frac{\sqrt{3}}{3}}\right)$ 可得 ${3}^{\alpha } = \frac{\sqrt{3}}{3} = {3}^{-\frac{1}{2}}$ ,即 $\alpha  =  - \frac{1}{2}$ ,
        
        可得 $f\left( x\right)  = {x}^{-\frac{1}{2}} = \frac{1}{\sqrt{x}}$ ,
        
        因为 $\sqrt{x} > 0$ ,可得 $f\left( x\right)  = \frac{1}{\sqrt{x}} > 0$ ,所以该幂函数的值域是 $\left( {0, + \infty }\right)$ .
    \end{jieda}
    \item 已知幂函数 $f\left( x\right)  = \left( {{m}^{2} - {5m} + 7}\right) {x}^{m}$ 的图象关于 $y$ 轴对称,则实数 $m$ 的值是\blkda[4]{2}.
    \begin{jieda}
        由 $f\left( x\right)$ 为幂函数,则 ${m}^{2} - {5m} + 7 = 1$ ,解得 $m = 2$ ,或 $m = 3$ ,
        
        当 $m = 2$ 时, $f\left( x\right)  = {x}^{2}$ ,其图象关于 $y$ 轴对称,
        
        当 $m = 3$ 时, $f\left( x\right)  = {x}^{3}$ ,其图象关于(0,0)对称,
        
        因此 $m = 2$
    \end{jieda} 
    \item 已知函数 $f\left( x\right)  = {x}^{\frac{1}{3}}$ ,则关于 $t$ 的表达式 $f\left( {{t}^{2} - {2t}}\right)  + f\left( {2{t}^{2} - 1}\right)  < 0$ 的解集为\blkda[4]{$\left( {-\dfrac{1}{3},1}\right)$}.
    \begin{jieda}
        由题意可知, $f\left( x\right)$ 的定义域为 $\left( {-\infty , + \infty }\right)$ ,
        
        所以 $f\left( {-x}\right)  = {\left( -x\right) }^{\frac{1}{3}} =  - {x}^{\frac{1}{3}} =  - f\left( x\right)$ ,
        
        所以函数 $f\left( x\right)$ 是奇函数,
        
        由幂函数的性质知,函数 $f\left( x\right)  = {x}^{\frac{1}{3}}$ 在函数 $\left( {-\infty , + \infty }\right)$ 上单调递增,
        
        由 $f\left( {{t}^{2} - {2t}}\right)  + f\left( {2{t}^{2} - 1}\right)  < 0$ ,得 $f\left( {{t}^{2} - {2t}}\right)  <  - f\left( {2{t}^{2} - 1}\right)$ ,即 $f\left( {{t}^{2} - {2t}}\right)  < f\left( {1 - 2{t}^{2}}\right)$ ,
        
        所以 ${t}^{2} - {2t} < 1 - 2{t}^{2}$ ,即 $3{t}^{2} - {2t} - 1 < 0$ ,解得 $- \frac{1}{3} < t < 1$ ,
        
        所以关于 $t$ 的表达式 $f\left( {{t}^{2} - {2t}}\right)  + f\left( {2{t}^{2} - 1}\right)  < 0$ 的解集为 $\left( {-\dfrac{1}{3},1}\right)$ .
    \end{jieda}
    \item 不等式 ${\left( 2x + 1\right) }^{\frac{2}{3}} < {\left( x - 3\right) }^{\frac{2}{3}}$ 的解为\blkda[4]{$\left( {-4,\dfrac{2}{3}}\right)$}.
    \begin{jieda}
        幂函数 $f\left( x\right)  = {x}^{\frac{2}{3}} = \sqrt[3]{{x}^{2}}$ 的定义域为 $\mathbf{R}$ ,且函数在 $\lbrack 0, + \infty )$ 上单调递增,
        
        又 $f\left( {-x}\right)  = \sqrt[3]{{\left( -x\right) }^{2}} = \sqrt[3]{{x}^{2}} = f\left( x\right)$ ,则 $f\left( x\right)$ 为偶函数,所以 $f\left( x\right)$ 在 $\left( {-\infty ,0}\right)$ 上单调递减,
        
        则由不等式 ${\left( 2x + 1\right) }^{\frac{2}{3}} < {\left( x - 3\right) }^{\frac{2}{3}}$ 可得 $\left| {{2x} + 1}\right|  < \left| {x - 3}\right|$ ,平方后整理得 $3{x}^{2} + {10x} - 8 < 0$ ,
        
        即 $\left( {{3x} - 2}\right) \left( {x + 4}\right)  < 0$ ,解得 $- 4 < x < \dfrac{2}{3}$ ,则不等式的解集为 $\left( {-4,\dfrac{2}{3}}\right)$ .
    \end{jieda}
    \item 已知幂函数 $f\left( x\right)  = {x}^{a}$ 的图象过点 $\left( {2,\dfrac{1}{4}}\right)$ ,若 $f\left( {{2m} + 1}\right)  < f\left( 3\right)$ ,则 $m$ 的取值范围是\blkda[4]{$\left( {-\infty , - 2}\right)  \cup  \left( {1, + \infty }\right)$}.
    \begin{jieda}
        由幂函数 $f\left( x\right)  = {x}^{a}$ 的图象过点 $\left( {2,\dfrac{1}{4}}\right)$ 得 ${2}^{a} = \dfrac{1}{4}$ ,解得 $a =  - 2$ ,
        
        则 $f\left( x\right)  = {x}^{-2} = \dfrac{1}{{x}^{2}}$ ,定义域为 $\left( {-\infty ,0}\right)  \cup  \left( {0, + \infty }\right)$ .
        
        由 $f\left( {-x}\right)  = \dfrac{1}{{\left( -x\right) }^{2}} = f\left( x\right)$ 可得 $f\left( x\right)$ 为偶函数,
        
        又幂函数的单调性可知,函数 $f\left( x\right)  = \dfrac{1}{{x}^{2}}$ 在 $\left( {0, + \infty }\right)$ 上单调递减.
        
        于是 $f\left( {{2m} + 1}\right)  < f\left( 3\right)$ 等价于 $\left| {{2m} + 1}\right|  > 3$ ,解得 $m <  - 2$ 或 $m > 1$ .
        
        所以 $m$ 的取值范围是 $\left( {-\infty , - 2}\right)  \cup  \left( {1, + \infty }\right)$ .
    \end{jieda}
    \item 已知幂函数 $f\left( x\right)  = {x}^{-{m}^{2} + {4m}}\left( {m \in  Z}\right)$ 的图象关于 $y$ 轴对称,且在区间 $\left( {0, + \infty }\right)$ 上是严格增函数.
    \begin{enumerate}
        \item 求 $m$ 的值
        \item 求满足不等式 $f\left( {{2a} - 1}\right)  < f\left( {a + 1}\right)$ 的实数 $a$ 的取值范围.
    \end{enumerate}
    \begin{jieda}
        \begin{enumerate}
            \item 因为幂函数 $f\left( x\right)  = {x}^{-{m}^{2} + {4m}}$ 在区间 $\left( {0, + \infty }\right)$ 上是严格增函数,
            
            所以 $- {m}^{2} + {4m} > 0$ ,解得 $0 < m < 4$ ,
            
            又因为 $m \in  \mathbf{Z}$ ,所以 $m = 1$ 或 $m = 2$ 或 $m = 3$ ,
            
            当 $m = 1$ 或 $m = 3$ 时, $f\left( x\right)  = {x}^{3}$ 为奇函数,图象关于原点对称 (舍);
            
            当 $m = 2$ 时, $f\left( x\right)  = {x}^{4}$ 为偶函数,图象关于 $y$ 轴对称,符合题意;
            
            综上所述, $m = 2$ .
            \item 由(1)得 $f\left( x\right)  = {x}^{4}$ 为偶函数,且在区间 $\left( {0, + \infty }\right)$ 上是严格增函数,
            
            则由 $f\left( {{2a} - 1}\right)  < f\left( {a + 1}\right)$ 得 $\left| {{2a} - 1}\right|  < \left| {a + 1}\right|$ ,
            
            即 ${\left( 2a - 1\right) }^{2} < {\left( a + 1\right) }^{2}$ ,即 ${a}^{2} - {2a} < 0$ ,解得 $0 < a < 2$ ,
            
            所以满足 $f\left( {{2a} - 1}\right)  < f\left( {a + 1}\right)$ 的实数 $a$ 的取值范围为 $0 < a < 2$ .
        \end{enumerate}
    \end{jieda}
    \item 已知幂函数 $f\left( x\right)  = \left( {{m}^{2} + {4m} + 4}\right) {x}^{m + 2}$ 在 $\left( {0, + \infty }\right)$ 上单调递减.
    \begin{enumerate}
        \item 求 $m$ 的值
        \item 若 ${\left( 2a - 1\right) }^{-m} < {\left( a + 3\right) }^{-m}$ ,求 $a$ 的取值范围.
    \end{enumerate}
    \begin{jieda}
        \begin{enumerate}
            \item 由幂函数的定义可得 ${m}^{2} + {4m} + 4 = 1$ ,即 ${m}^{2} + {4m} + 3 = 0$ ,解得 $m =  - 1$ 或 $m =  - 3$ .
            
            因为 $f\left( x\right)$ 在 $\left( {0, + \infty }\right)$ 上单调递减,所以 $m + 2 < 0$ ,即 $m <  - 2$，则 $m =  - 3$ .

            \item 设 $g\left( x\right)  = {x}^{3},g\left( x\right)$ 是 $\mathbf{R}$ 上的增函数.
            
            由( 1 )可知 ${\left( 2a - 1\right) }^{-m} < {\left( a + 3\right) }^{-m}$ ,即 ${\left( 2a - 1\right) }^{3} < {\left( a + 3\right) }^{3}$ ,
            
            则 ${2a} - 1 < a + 3$ ,解得 $a < 4$ ,
            
            即 $a$ 的取值范围为 $\left( {-\infty ,4}\right)$ .
        \end{enumerate}
    \end{jieda}
\end{enumerate}
\subsection*{C Level}
\begin{enumerate}
    \item 若关于 $x$ 的不等式 ${x}^{\frac{4}{3}} + {2x} + a + \left| {{x}^{\frac{4}{3}} - {2x} - a}\right|  \geq  2$ 的解集为 $\mathbf{R}$ ,则实数 $a$ 能取到的最小值为\blkda[4]{3}.
    \begin{jieda}
        ${x}^{\frac{4}{3}} + {2x} + a + \left| {{x}^{\frac{4}{3}} - {2x} - a}\right|  \geq  2 \Longleftrightarrow \left| {{x}^{\frac{4}{3}} - {2x} - a}\right| \geq -{x}^{\frac{4}{3}} - {2x} + (2-a)$

        即${x}^{\frac{4}{3}} - {2x} - a \geq -{x}^{\frac{4}{3}} - {2x} + (2-a)$ 或 ${x}^{\frac{4}{3}} - {2x} - a \leq{x}^{\frac{4}{3}} + {2x} + (a-2)$

        即${x}^{\frac{4}{3}} \geq 1$或$2x+a-1 \geq 0$

        因此应有两个不等式的解集的并集为$\mathbf{R}$，即$(-\infty, -1] \cup [1, +\infty) \cup \left[\dfrac{1-a}{2}, +\infty\right) = \mathbf{R}$，故$a \geq  3$
    \end{jieda}

    \item 设函数 $f\left( x\right)  = a{x}^{2} + {bx} + c\left( {a,b,c \in  \mathbf{R},a > 0}\right)$ ,则 “ $f\left( {f\left( {-\dfrac{b}{2a}}\right) }\right)  < 0$ ” 是 “ $f\left( x\right)$ 与 $f\left( {f\left( x\right) }\right)$ 都恰有两个零点”的\dotfill{(\kh{C})}
    \xx{充分不必要条件}{必要不充分条件}{充要条件}{既不充分也不必要条件}
    \begin{jieda}
        设 $\left( {{x}_{0},{y}_{0}}\right)$ 为 $y = f\left( x\right)$ 图象的顶点, ${x}_{1}\text{ 、 }{x}_{2}$ 且 ${x}_{1} < {x}_{2}$ 为 $y = f\left( x\right)$ 的两个零点 (若有). 因为 $f\left( {f\left( {-\dfrac{b}{2a}}\right) }\right)  < 0$ ,即 $f\left( {y}_{0}\right)  < 0$ ,故 ${x}_{1},{x}_{2}$ 必存在,且 ${x}_{1} < {y}_{0} < {x}_{2}$ . 故此时显然 $f\left( x\right)$ 恰有两个零点; 再令 $f\left( {f\left( x\right) }\right)  = 0$ ,解得 $f\left( x\right)  = {x}_{1}$ 或 $f\left( x\right)  = {x}_{2}$
        
        根据 $y = f\left( x\right)$ 的图象可知,当 $k > {y}_{0}$ 时, $f\left( x\right)  = k$ 有两解; 当 $k = {y}_{0}$ 时, $f\left( x\right)  = k$ 有一解; 当 $k < {y}_{0}$ 时, $f\left( x\right)  = k$ 无解; 因此 $f\left( x\right)  = {x}_{1}$ 无解, $f\left( x\right)  = {x}_{2}$ 有两解,综上, $f\left( {f\left( x\right) }\right)$ 恰有两个零点,满足充分性. 
        
        若 $f\left( x\right)$ 有两个零点,则 ${x}_{1}\text{ 、 }{x}_{2}$ 必存在. 由 $f\left( {f\left( x\right) }\right)  = 0$ 得 $f\left( x\right)  = {x}_{1}$ 或 $f\left( x\right)  = {x}_{2}$ ,因为 $f\left( {f\left( x\right) }\right)  = 0$ 也只有两个零点,因此 $f\left( x\right)  = {x}_{1}$ 无解, $f\left( x\right)  = {x}_{2}$ 有两个不等实根,即 ${x}_{1} <$  ${y}_{0} < {x}_{2}$ ,所以 $f\left( {f\left( {-\dfrac{b}{2a}}\right) }\right)  < 0$ ,必要性得证.
    \end{jieda}
\end{enumerate}

\newpage
\quad
\newpage
\section*{指数函数及其性质}
\subsection*{A Level}
\begin{enumerate}
    \item 设 $a > 0,s \in  \mathbf{R}$ . 下列各项中,能推出 ${a}^{s} > a$ 的一项是\dotfill{(\kh{D})}
    \xx{$a > 1$ ,且 $s > 0$}{$a > 1$ ,且 $s < 0$}{$0 < a < 1$ ,且 $s > 0$}{$0 < a < 1$ ,且 $s < 0$}
    \begin{jieda}
        $\because a > 0,{a}^{s} > a,\therefore {a}^{s - 1} > 1 = {a}^{0}$ ,
        
        当 $a \in  \left( {0,1}\right)$ 时, $y = {a}^{x}$ 定义域上严格单调递减,
        
        此时若 $s - 1 < 0$ ,则一定有 ${a}^{s - 1} > 1 = {a}^{0}$ 成立,故 ${D}$ 正确, ${C}$ 错误;
        
        当 $a \in  \left( {1, + \infty }\right)$ 时, $y = {a}^{x}$ 定义域上严格单调递增,要满足 ${a}^{s - 1} > 1 = {a}^{0}$ ,需 $s > 1$ ,即 ${A}\text{ 、 }{B}$ 错误.
    \end{jieda}
    \item  函数 $y = {\left( \dfrac{1}{2}\right) }^{{x}^{2} - {2x}}$ 的值域为\dotfill{(\kh{A})}
    \xx{$(0,2\rbrack$}{$\left( {0, + \infty }\right)$}{$\lbrack 2, + \infty )$}{$\lbrack 1, + \infty )$}
    \begin{jieda}
        依题意，令 $t = {x}^{2} - {2x}$ ,则 $t = {x}^{2} - {2x} = {\left( x - 1\right) }^{2} - 1 \geq   - 1$，因为 $y = {\left( \dfrac{1}{2}\right) }^{t}$ 单调递减,且 $y = {\left( \dfrac{1}{2}\right) }^{t} > 0$
        
        所以 $y = {\left( \dfrac{1}{2}\right) }^{t} \leq  {\left( \dfrac{1}{2}\right) }^{-1} = 2$ ,
        
        所以 $y \in  (0,2\rbrack$ .
    \end{jieda}
    \item 下列函数是指数函数的序号为\blkda[4]{\ding{173}\ding{174}\ding{176}}. (请填入全部正确的序号)
    \begin{dingautolist}{172}
        \item $y = {\left( -4\right) }^{x}$
        \item $y = {\left( \dfrac{1}{4}\right) }^{x}$
        \item $y = {4}^{x}$
        \item $y = {x}^{-4}$
        \item $y = {\left( \sqrt{4}\right) }^{x}$
    \end{dingautolist}

    \item 若函数 $y = \left( {{m}^{2} - {2m} - 2}\right)  \cdot  {m}^{x}$ 是指数函数,则 $m$ 等于\dotfill{(\kh{C})}
    \xx{$-1$ 或 $3$}{$-1$}{3}{$\dfrac{1}{3}$}
    \begin{jieda}
        因为函数 $y = \left( {{m}^{2} - {2m} - 2}\right)  \cdot  {m}^{x}$ 是指数函数，所以 $\left\{  {\begin{array}{l} {m}^{2} - {2m} - 2 = 1 \\  m > 0 \\  m \neq  1 \end{array} \Longrightarrow  m = 3}\right.$ .
    \end{jieda}
    
    \item 若函数 $y = \left( {{a}^{2} - {3a} + 3}\right) {a}^{x}$ 为指数函数,则 $a =$ \blkda[4]{2}.
    \begin{jieda}
        因为函数 $y = \left( {{a}^{2} - {3a} + 3}\right) {a}^{x}$ 为指数函数,
        
        所以 ${a}^{2} - {3a} + 3 = 1$ 且 $a > 0$ 且 $a \neq  1$ ,解得 $a = 2$ .
    \end{jieda}
    \item 函数 $y = {a}^{x} - 1\left( {a > 0\text{ 且 }a \neq  1}\right)$ 的图像一定过点\blkda[4]{$(0, 0)$}.
    \begin{jieda}
        函数 $y = {a}^{x} - 1\left( {a > 0\text{ 且 }a \neq  1}\right)$ ,令 $x = 0$ 可得 $y = {a}^{0} - 1 = 0$
    \end{jieda} 
    \item 函数 $y = {a}^{x + 2} - 3\left( {a > 0\text{ ,且 }a \neq  1}\right)$ 的图象过定点 $A$ ,则点 $A$ 的坐标是\blkda[4]{$(-2, -2)$}.
    \begin{jieda}
        因为 $y = {a}^{x + 2} - 3\left( {a > 0\text{ ,且 }a \neq  1}\right)$ 的图象过定点 $A$，令 $x + 2 = 0$ ,则 $x =  - 2,y = {a}^{0} - 3 =  - 2$，所以点 $A$ 的坐标为$(-2, - 2)$.
    \end{jieda}
    \item 已知 $f\left( x\right)  = \left\{  \begin{matrix} {2}^{x},x > 0 \\  1,x \leq  0 \end{matrix}\right.$ ,则 $f\left( x\right)$ 的值域是\blkda[4]{$\lbrack 1, + \infty )$}
    \begin{jieda}
        当 $x > 0$ 时,根据指数函数的图象与性质知 $f\left( x\right)  = {2}^{x} > 1$ ,
        
        当 $x \leq  0$ 时, $f\left( x\right)  = 1$ .
        
        综上: $y = f\left( x\right)$ 的值域为 $\lbrack 1, + \infty )$ .
    \end{jieda}
    \item 函数 $y = {3}^{x}$ 在 $\left\lbrack  {\dfrac{1}{2},2}\right\rbrack$ 上的最小值是\blkda[4]{$\sqrt{3}$}.
    \begin{jieda}
        因函数 $y = {3}^{x}$ 是 $\mathbf{R}$ 上的增函数,且 $x \in  \left\lbrack  {\dfrac{1}{2},2}\right\rbrack$，则 $x = \dfrac{1}{2}$ 时,函数 $y = {3}^{x}$ 取得最小值为 ${3}^{\frac{1}{2}} = \sqrt{3}$ .
    \end{jieda}
    \item 函数 $y = \sqrt{{3}^{x} - 1}$ 的定义域是\blkda[4]{$\lbrack 0, + \infty )$}.
    \begin{jieda}
        由题知, ${3}^{x} - 1 \geq  0$ ,解得 $x \geq  0$，所以函数 $y = \sqrt{{3}^{x} - 1}$ 的定义域是 $\lbrack 0, + \infty )$
    \end{jieda}
    \item 函数 $y = \sqrt{{\left( \dfrac{1}{2}\right) }^{x} - 1}$ 的定义域为\blkda[4]{$( - \infty ,0\rbrack$}.
    \begin{jieda}
        由题设 ${\left( \dfrac{1}{2}\right) }^{x} - 1 \geq  0$ ,故 $x \leq  {\log }_{\frac{1}{2}}1 = 0$ ,故定义域为 $( - \infty ,0\rbrack$ .
    \end{jieda}
    \item 若关于 $x$ 的方程 ${\left( \dfrac{3}{4}\right) }^{x} = a$ 有负根,则实数 $a$ 的取值范围\blkda[4]{$\left( {1, + \infty }\right)$}.
    \begin{jieda}
        关于 $x$ 的方程 ${\left( \dfrac{3}{4}\right) }^{x} = a$ 有负根等价于指数函数 $y = {\left( \dfrac{3}{4}\right) }^{x}$ 与 $y = a$ 在第二象限有交点
        \begin{center}
        \includegraphics[width=0.3\textwidth]{bodyxb/src/bo_d2ff4jbef24c73b0g1j0_4_240_195_458_309_0.jpg}
        \end{center}
        当 $a > 1$ 时, $y = {\left( \dfrac{3}{4}\right) }^{x}$ 与 $y = a$ 在第二象限有交点，所以实数 $a$ 的取值范围 $\left( {1, + \infty }\right)$ .
    \end{jieda}
\end{enumerate}
\subsection*{B Level}
\begin{enumerate}
    \item 已知函数 $f\left( x\right)$ 是定义域为 $\mathbf{R}$ 的奇函数,当 $x < 0$ 时, $f\left( x\right)  = {2}^{x}$ ,则 $f\left( {{\log }_{2}7}\right)  =$ \blkda{$-\dfrac{1}{7}$}
    \begin{jieda}
        $log_27 > 0$，故$f\left( {{\log }_{2}7}\right)  = - f(-{{\log }_{2}7}) = - f\left({{\log }_{2}\dfrac{1}{7}}\right) = -\dfrac{1}{7}$
    \end{jieda}
    \item 已知常数 $a > 0$ ,函数 $f\left( x\right)  = \dfrac{{2}^{x}}{{2}^{x} + {ax}}$ 的图像经过点 $P\left( {p,\dfrac{6}{5}}\right) \text{ 、 }Q\left( {q, - \dfrac{1}{5}}\right)$ ,若 ${2}^{p + q} =$  ${25pq}$ ,则实数 $a =$ \blkda{5}
    \begin{jieda}
        由条件可知 $\dfrac{{2}^{p}}{{2}^{p} + {ap}} = \dfrac{6}{5}$ ,则 $\dfrac{1}{1 + \dfrac{ap}{{2}^{p}}} = \dfrac{6}{5}$ ,得 $\dfrac{ap}{{2}^{p}} =  - \dfrac{1}{6}$
        
        由 $\dfrac{{2}^{q}}{{2}^{q} + {aq}} =  - \dfrac{1}{5}$ ,得 $\dfrac{1}{1 + \dfrac{aq}{{2}^{q}}} = - \dfrac{1}{5}$，则 $\dfrac{aq}{{2}^{q}} =  - 6$
        
        故得 $\dfrac{{a}^{2}{pq}}{{2}^{p + q}} = 1$ ,又因为 ${2}^{p + q} = {25pq}$ ,所以 $\dfrac{{a}^{2}}{25} = 1$ ,又 $a > 0$ ,得 $a = 5$
    \end{jieda}  
    \item 不等式 ${2}^{{x}^{2} - {2x} - 3} < {\left( \dfrac{1}{2}\right) }^{{3x} - 3}$ 与不等式 ${x}^{2} + {ax} + b < 0$ 解集相同,则 $a + b =$ \blkda[4]{$-5$}.
    \begin{jieda}
        $\because {2}^{{x}^{2} - {2x} - 3} < {\left( \dfrac{1}{2}\right) }^{{3x} - 3} = {2}^{3 - {3x}}$ ,
        
        $\because y = {2}^{x}$ 在 $\mathbf{R}$ 上单调递增,
        
        $\therefore {x}^{2} - {2x} - 3 < 3 - {3x}$ ,即 ${x}^{2} + x - 6 < 0$ ,
        
        $\therefore a = 1,b =  - 6$ ,
        
        $\therefore a + b =  - 5$ .
    \end{jieda}
    \item 高斯函数属于初等函数,以大数学家约翰·卡尔·弗里德里希·高斯的名字命名,高斯函数应用范围很广,在自然科学、社会科学、数学以及工程学等领域都能看到它的身影,设 $x \in  \mathbf{R}$ , 用 $\left\lbrack  x\right\rbrack$ 表示不超过 $x$ 的最大整数,则 $y = \left\lbrack  x\right\rbrack$ 称为高斯函数,例如: $\left\lbrack  {-{3.1}}\right\rbrack   =  - 4,\left\lbrack  {4.8}\right\rbrack   = 4$，则函数 $f\left( x\right)  = \left\lbrack  {\dfrac{{2}^{x + 1}}{1 + {2}^{x}} - \dfrac{1}{3}}\right\rbrack$ 的值域为\blkda{$\{-1, 0, 1\}$}
    \begin{jieda}
        先求$g(x) = \dfrac{{2}^{x + 1}}{1 + {2}^{x}} - \dfrac{1}{3}$的值域：

        $g(x) = \dfrac{5}{3} - \dfrac{2}{1+2^x}$，$x \in \mathbf{R}$，$2^x = (0, +\infty)$，$1+2^x \in (1, +\infty)$，$\dfrac{1}{1+2^x} \in (0, 1)$，$\dfrac{2}{1+2^x} \in (0, 2)$，$g(x) \in \left(-\dfrac{1}{3}, \dfrac{5}{3}\right)$

        故$f(x) \in \{-1, 0, 1\}$
    \end{jieda}
    \item 已知函数 $f\left( x\right)  = {4}^{x} - \dfrac{5}{2} \cdot  {2}^{x + 1} - 6$ ,其中 $x \in  \left\lbrack  {0,3}\right\rbrack$ ,则 $f\left( x\right)$ 的最大值为\blkda[4]{$18$}
    \begin{jieda}
        设$t = 2^x$，则$t \in [1, 8]$
        
        $f(x) = g(t) = t^2-5t-6 \leq g(8) = 18$
    \end{jieda}

    \item 已知 $- 1 \leq  x \leq  1$ ,求函数 $y = 4 \cdot  {3}^{x} - 2 \cdot  {9}^{x}$ 的最大值为\blkda[4]{2}
    \begin{jieda}
        令$t = 3^x, t \in \left[\dfrac{1}{3}, 3\right]$，$y = 4t-2t^2 \leq 2$
    \end{jieda}
  
    \item 函数$y=3\cdot4^x-20\cdot2^x+3\cdot4^{-x}-20\cdot2^{-x}$的最小值为\blkda[4]{$-\dfrac{118}{3}$} 
    \begin{jieda}
    令$t=2^x+2^{-x}\in[2,+\infty)$, 则
    $y = 3t^2-20t-6 = 3(t-\dfrac{10}{3})^2-\dfrac{118}{3}$ \\
    故$y_{\min}=-\dfrac{118}{3}$, 当$t=\dfrac{10}{3}$即$2^x = \dfrac{1}{3}\mbox{或}3, x=log_23 \mbox{或}-log_23$时取到. 
    \end{jieda}
    \item 若关于 $x$ 的方程 $\left| {{3}^{x} - 1}\right|  = k$ 有两解,则实数 $k$ 的取值范围为\blkda[4]{$(0, 1)$}
    \begin{jieda}
        数形结合
    \end{jieda}
    
    \item 若实数 $x$ 、 $y$ 满足 $x + {2y} = 1$ ,则 ${2}^{x} + {4}^{y}$ 的最小值为\blkda[4]{$2\sqrt{2}$}.
    \begin{jieda}
        ${2}^{x} + {4}^{y} \geq  2\sqrt{{2}^{x} \times  {4}^{y}} = 2\sqrt{{2}^{x} \times  {2}^{2y}} = 2\sqrt{{2}^{x + {2y}}} = 2\sqrt{2}$ ,当且仅当 $x = {2y}$，即 $x = \dfrac{1}{2},y = \dfrac{1}{4}$ 时取到等号.
    \end{jieda}
    \item 设 $t \in  {R}$ ,若在区间(1,2)上,关于 $x$ 的不等式 ${2}^{x} > \dfrac{1}{x + t}$ 有意义且能恒成立,则 $t$ 的取值范围为\blkda[4]{$\left( -\infty , - 2 \right \rbrack \cup \left\lbrack  {-\dfrac{1}{2}, + \infty  }\right)$}.
    \begin{jieda}
        由题意得 ${2}^{x} > \dfrac{1}{x + t}$ 在$(1,2)$上有意义,故 $x + t \neq  0$ 在(1,2)上恒成立，故 $t =  - x \in  \left( {-\infty , - 2\rbrack \cup \lbrack  - 1, + \infty }\right)$ 

        \begin{dingautolist}{172}
            \item 当 $t \in  ( - \infty , - 2\rbrack$ 时, $x + t < 0$ ,而 ${2}^{x} > 0$ ,满足 ${2}^{x} > \dfrac{1}{x + t}$ ,符合题意
            \item 当 $t \in  \lbrack  - 1, + \infty )$ 时, $x + t > 0,{2}^{x} > \dfrac{1}{x + t} \Longrightarrow  t > \dfrac{1}{{2}^{x}} - x$ 在$(1,2)$上恒成立

            令 $g\left( x\right)  = \dfrac{1}{{2}^{x}} - x,x \in  \left( {1,2}\right)$，$g\left( x\right)$ 在 $x \in  \left( {1,2}\right)$ 上单调递减,
            
            故 $g\left( x\right)  = \dfrac{1}{{2}^{x}} - x \in  \left( {-\dfrac{7}{4}, - \dfrac{1}{2}}\right)$，故 $t \geq   - \dfrac{1}{2}$
        \end{dingautolist}
        综上, $t$ 的取值范围是 $\left( -\infty , - 2 \right \rbrack \cup \left\lbrack  {-\dfrac{1}{2}, + \infty  }\right)$ 
    \end{jieda}
    \item 若函数 $f\left( x\right)  = \left\{  \begin{array}{l} {a}^{x},x \geq  1 \\  \left( {4 - \dfrac{a}{2}}\right) x + 2,x < 1 \end{array}\right.$ 在区间 $\left( {-\infty , + \infty }\right)$ 严格单调递增函数,则实数 $a$ 的取值范围是\blkda[4]{$[4,8)$}. 
    \begin{jieda}
        由题意,函数 $f\left( x\right)  = \left\{  \begin{array}{l} {a}^{x},x \geq  1 \\  \left( {4 - \dfrac{a}{2}}\right) x + 2,x < 1 \end{array}\right.$ 是 $\left( {-\infty , + \infty }\right)$ 上的单调递增函数,
        
        则满足 $\left\{  \begin{array}{l} a > 1 \\  4 - \dfrac{a}{2} > 0 \\  \left( {4 - \dfrac{a}{2}}\right)  \times  1 + 2 \leq  {a}^{1} \end{array}\right.$ ,解得 $4 \leq  a < 8$
    \end{jieda}
    
    \item 若 $m \in  \mathbf{R}$ , $f\left( x\right)  = \left\{  \begin{array}{l} {2}^{x},x \geq  0 \\  \dfrac{1}{{2}^{x}},x < 0 \end{array}\right.$ ,则满足 $f\left( {m - 2}\right)  \geq  f\left( {m + 3}\right)$ 的 $m$ 的最大值为\blkda[4]{$- \dfrac{1}{2}$}.
    \begin{jieda}
        $f\left( x\right)$ 为偶函数，由指数函数的单调性可知, $f\left( x\right)$ 在 $\left( {0, + \infty }\right)$ 上单调递增,
        
        因此,不等式 $f\left( {m - 2}\right)  \geq  f\left( {m + 3}\right)$ 等价于 $\left| {m - 2}\right|  \geq  \left| {m + 3}\right|$ ,
        
        即 ${\left( m - 2\right) }^{2} \geq  {\left( m + 3\right) }^{2}$ ,解得 $m \leq   - \dfrac{1}{2}$ .
        
        故 $m$ 的最大值为 $- \dfrac{1}{2}$
    \end{jieda} 
    
    \item 设函数 $f\left( x\right)  = \left\{  \begin{array}{l} {2}^{-\left| {x - a}\right| },x \leq  1 \\   - \dfrac{1}{2}x + 1,x > 1 \end{array}\right.$ ,若 $f\left( 1\right)$ 是函数 $f\left( x\right)$ 的最大值,则实数 $a$ 的取值范围为\blkda[4]{$\left\lbrack  {1,2}\right\rbrack$}.
    \begin{jieda}
        当 $x > 1$ 时 $f\left( x\right)  =  - \dfrac{1}{2}x + 1$ 函数单调递减且 $f\left( x\right)  < \dfrac{1}{2}$
        
        当 $x \leq  1$ 时 $f\left( x\right)  = {2}^{-\left| {x - a}\right| } = {\left( \dfrac{1}{2}\right) }^{\left| x - a\right| }$，结合图像可得在 $x < a$ 单调递增，在 $x > a$ 时函数单调递减，故必有$a \geq 1$，且 ${\left( \dfrac{1}{2}\right) }^{a - 1} \geq  \dfrac{1}{2}$ ,解得 $1 \leq  a \leq  2$       
    \end{jieda}
    \item 已知 $0 < a < 2$ ,函数 $y = \left\{  \begin{matrix} \left( {a - 2}\right) x + {4a} + 1, & x \leq  2 \\  2{a}^{x - 1}, & x > 2 \end{matrix}\right.$ ,若该函数存在最小值,则实数 $a$ 的取值范围是\blkda[4]{$\left\{ a \left |  0 < a \leq  \dfrac{1}{2} \text{或} a = 1\right\}\right.$}.
    \begin{jieda}
        由题意,令 $g\left( x\right)  = \left( {a - 2}\right) x + {4a} + 1,x \in  ( - \infty ,2\rbrack ,h\left( x\right)  = 2{a}^{x - 1},x \in  \left( {2, + \infty }\right)$

        \begin{dingautolist}{172}
            \item 当 $0 < a < 1$ 时, $g\left( x\right)$ 在 $( - \infty ,2\rbrack$ 上单调递减, $h\left( x\right)$ 在 $\left( {2, + \infty }\right)$ 上单调递减,则 $h\left( x\right)$ 在 $\left( {2, + \infty }\right)$ 上的值域为$(0,2a)$,
            
            因为 $f\left( x\right)$ 存在最小值,故需 $g\left( 2\right)  = \left( {a - 2}\right)  \times  2 + {4a} + 1 \leq  0$ ,解得 $a \leq  \dfrac{1}{2}$ ,
            
            结合 $0 < a < 1$ ,此时 $0 < a \leq  \dfrac{1}{2}$
            \item 当 $1 < a < 2$ 时, $g\left( x\right)$ 在 $( - \infty ,2\rbrack$ 上单调递减, $h\left( x\right)$ 在 $\left( {2, + \infty }\right)$ 上单调递增,则 $h\left( x\right)$ 在 $\left( {2, + \infty }\right)$ 上的值域为 $\left( {{2a}, + \infty }\right)$ ,
            
            因为 $f\left( x\right)$ 存在最小值,故需 $g\left( 2\right)  \leq  {2a}$ ,即 $\left( {a - 2}\right)  \times  2 + {4a} + 1 \leq  {2a}$ ,解得 $a \leq  \dfrac{3}{4}$ ,
            
            这与 $1 < a < 2$ 矛盾

            \item 当 $a = 1$ 时, $g\left( x\right)  =  - x + 5$ 在 $( - \infty ,2\rbrack$ 上单调递减,且在 $( - \infty ,2\rbrack$ 上的值域为 $\lbrack 3, + \infty ),h\left( x\right)  = 2$ , 此时存在最小值 2
        \end{dingautolist}
        故实数 $a$ 的取值范围为 $\left\{  {a\left| {0 < a \leq  \dfrac{1}{2}}\text{或} a = 1\right. }\right\}$ .
    \end{jieda}
\end{enumerate}
\subsection*{C Level}
\begin{enumerate}
    \item 已知函数 $f\left( x\right)  = 3 \cdot  {2}^{x} + 2$ ,对于任意的 ${x}_{2} \in  \left\lbrack  {0,1}\right\rbrack$ ,都存在 ${x}_{1} \in  \left\lbrack  {0,1}\right\rbrack$ ,使得 $f\left( {x}_{1}\right)  + {2f}\left( {{x}_{2} + m}\right)  = {13}$ 成立,则实数 $m$ 的取值范围为\blkda[4]{$\left\lbrack  {{\log }_{2}\dfrac{1}{6},{\log }_{2}\dfrac{1}{3}}\right\rbrack$}.
    \begin{jieda}
        对于任意的 ${x}_{2} \in  \left\lbrack  {0,1}\right\rbrack$ ,都存在 ${x}_{1} \in  \left\lbrack  {0,1}\right\rbrack$ ,使得 ${f}\left( {{x}_{2} + m}\right)  = \dfrac{{13} - f\left( {x}_{1}\right) }{2}$ 成立，故等式左侧的取值范围是等式右侧取值范围的子集.
        
        $f\left( {x}_{1}\right)  \in  \left\lbrack  {5,8}\right\rbrack  $，故$\dfrac{{13} - f\left( {x}_{1}\right) }{2} \in  \left\lbrack  {\dfrac{5}{2},4}\right\rbrack$
        
        $f\left( {{x}_{2} + m}\right)  = 3 \cdot  {2}^{{x}_{2} + m} + 2 \in  \left\lbrack  {3 \cdot  {2}^{m} + 2,3 \cdot  {2}^{1 + m} + 2}\right\rbrack$ ,
        
        由题意得 $\left\{  {\begin{array}{l} 3 \cdot  {2}^{m} + 2 \geq  \dfrac{5}{2} \\  3 \cdot  {2}^{m + 1} + 2 \leq  4 \end{array} \Longrightarrow  \left\{  {\begin{array}{l} {2}^{m} \geq  \dfrac{1}{6} \\  {2}^{m + 1} \leq  \dfrac{2}{3} \end{array} \Longrightarrow  {\log }_{2}\dfrac{1}{6} \leq  m \leq  {\log }_{2}\dfrac{1}{3}}\right. }\right.$
    \end{jieda}
    \item 已知函数 $f\left( x\right)  = \left\{  \begin{array}{ll} \dfrac{{x}^{2} + 4}{x} & x \geq  2 \\  {2}^{\left| x - a\right| } & x < 2 \end{array}\right.$ ,若对任意的 ${x}_{1} \in  \lbrack 2, + \infty )$ ,都存在唯一的 ${x}_{2} \in  \left( {-\infty ,2}\right)$ ,满足 $f\left( {x}_{2}\right)  = f\left( {x}_{1}\right)$ ,则实数 $a$ 的取值范围是\blkda[4]{$0 \leq  a < 4$}.
    \begin{jieda}
        设函数 $g\left( x\right)  = \dfrac{{x}^{2} + 4}{x}$，其在$x \in [2, +\infty)$的函数值取值范围为 ${A}$ ,函数 $h\left( x\right)  = {2}^{\left| x - a\right| }$，其在$x \in (-\infty, 2)$只对应一个自变量的函数值的取值范围为 $B$ , 因为对任意的 ${x}_{1} \in  \lbrack 2, + \infty )$ ,都存在唯一的 ${x}_{2} \in  \left( {-\infty ,2}\right)$ ,满足 $f\left( {x}_{2}\right)  = f\left( {x}_{1}\right)$，则 $A \subseteq  B$.
        
        $g\left( x\right)  = \dfrac{{x}^{2} + 4}{x} = x + \dfrac{4}{x}$，所以 $A = \lbrack 4, + \infty )$
        
         $h\left( x\right)  = {2}^{\left| x - a\right| }$在$(-\infty, a)$严格减，在$(a, +\infty)$严格增，结合图像可知
         \begin{dingautolist}{172}
             \item 当 $a \geq  2$ 时, $B = \left( {{2}^{a - 2}, + \infty }\right)$

             $\therefore {2}^{a - 2} < 4$ ,解得 $2 \leq  a < 4$

             \item 当 $a < 2$ 时, $B = \left( {{2}^{2-a}, + \infty }\right)$
             
             $\therefore {2}^{2 - a} \leq  4$ ,解得 $0 \leq  a < 2$
         \end{dingautolist}
        综合得 $0 \leq  a < 4$ 
    \end{jieda}
    \item 已知函数 $f\left( x\right)  = \dfrac{a \cdot  {2}^{x} + 1}{{2}^{x} - 1}$ 
    \begin{enumerate}
        \item 当 $a = 1$ 时,解方程 $\lg f\left( {2x}\right)  - \lg f\left( x\right)  = 1 - \lg {18}$
        \item 当 $x \in  (0,1\rbrack$ 时, $\left| {f\left( {2x}\right)  - f\left( x\right) }\right|  \geq  1$ 恒成立,求实数 $a$ 的取值范围
    \end{enumerate}
    \begin{jieda}
        \begin{enumerate}
            \item 当 $a = 1$ 时, $f\left( x\right)  = \dfrac{a \cdot  {2}^{x} + 1}{{2}^{x} - 1} = \dfrac{{2}^{x} + 1}{{2}^{x} - 1},f\left( {2x}\right)  = \dfrac{{2}^{2x} + 1}{{2}^{2x} - 1} = \dfrac{{\left( {2}^{x}\right) }^{2} + 1}{{\left( {2}^{x}\right) }^{2} - 1}$ 
            
            所以方程 $\lg f\left( {2x}\right)  - \lg f\left( x\right)  = 1 - \lg {18}$ 化为 $\lg \dfrac{f\left( {2x}\right) }{f\left( x\right) } = \lg \dfrac{10}{18}$ 且 $f\left( {2x}\right)  > 0,f\left( x\right)  > 0$
            
            即 $\dfrac{f\left( {2x}\right) }{f\left( x\right) } = \dfrac{5}{9}$ 且 $\dfrac{{\left( {2}^{x}\right) }^{2} + 1}{{\left( {2}^{x}\right) }^{2} - 1} > 0$ ,
            $\dfrac{{2}^{x} + 1}{{2}^{x} - 1} > 0$
            
            所以 $\dfrac{\dfrac{{\left( {2}^{x}\right) }^{2} + 1}{{\left( {2}^{x}\right) }^{2} - 1}}{\dfrac{{2}^{x} + 1}{{2}^{x} - 1}} = \dfrac{5}{9}$ ,即 $\dfrac{{\left( {2}^{x}\right) }^{2} + 1}{{\left( {2}^{x} + 1\right) }^{2}} = \dfrac{5}{9}$ 
            
            令 ${2}^{x} = t$ ,则原方程化为 $\dfrac{{t}^{2} + 1}{{\left( t + 1\right) }^{2}} = \dfrac{5}{9}$ ,整理得 $2{t}^{2} - {5t} +$
            $2 = 0$ ,解得 $t = 2$ 或 $t = \dfrac{1}{2}$ ,即 ${2}^{x} = 2$ 或 ${2}^{x} = \dfrac{1}{2}$ ,解得 $x = 1$ 或 $x =  - 1$
            
            当 $x =  - 1$ 时, $\dfrac{{\left( {2}^{x}\right) }^{2} + 1}{{\left( {2}^{x}\right) }^{2} - 1} < 0,\dfrac{{2}^{x} + 1}{{2}^{x} - 1} <0$,故舍去,故原方程的解为 $x = 1$
            
            \item 由 $\left| {f\left( {2x}\right)  - f\left( x\right) }\right|  \geq  1$ ,得 $\left| {\dfrac{a \cdot  {\left( {2}^{x}\right) }^{2} + 1}{{\left( {2}^{x}\right) }^{2} - 1} - \dfrac{a \cdot  {2}^{x} + 1}{{2}^{x} - 1}}\right|  \geq  1$ ,即
            $\left| \dfrac{-a \cdot  {2}^{x} - {2}^{x}}{{\left( {2}^{x}\right) }^{2} - 1}\right|  \geq  1$
            
            令 ${2}^{x} = t$ ,当 $x \in  (0,1\rbrack$ 时, $t \in  (1,2\rbrack$
            
            所以当 $x \in  (0,1\rbrack$ 时,
            $\left| {f\left( {2x}\right)  - f\left( x\right) }\right|  \geq  1$ 恒成立，等价于当 $t \in  (1,2\rbrack$ 时, $\left| \dfrac{\left( {a + 1}\right)  \cdot  t}{{t}^{2} - 1}\right|  \geq  1$ 恒成立，即 $\left| {a + 1}\right|  \geq  \dfrac{{t}^{2} - 1}{t}$ 在
            $t \in  (1,2\rbrack$ 时恒成立
            
            令 $g\left( t\right)  = \dfrac{{t}^{2} - 1}{t} = t - \dfrac{1}{t}$，$g\left( t\right)$在 $(1,2\rbrack$ 上严格增
            
            $g\left( 2\right)  = 2 - \dfrac{1}{2} = \dfrac{3}{2},g\left( 1\right)  = 1 - 1 = 0$，则 $0 < g\left( t\right)  \leq  \dfrac{3}{2}$
            
            所以 $\left| {a + 1}\right|  \geq  \dfrac{3}{2}$ ,解得 $a \leq   - \dfrac{5}{2}$ 或 $a \geq  \dfrac{1}{2}$
        \end{enumerate}
    \end{jieda}
\end{enumerate}

\section*{对数函数及其性质}
\subsection*{A Level}
\begin{enumerate}
    \item 已知函数 $f\left( x\right)  = {4}^{x} + {\log }_{2}x$ ,则 $f\left( \dfrac{1}{2}\right)  =$ \blkda[4]{1}.
    \begin{jieda}
        函数 $f\left( x\right)  = {4}^{x} + {\log }_{2}x$ ,所以 $f\left( \frac{1}{2}\right)  = {4}^{\frac{1}{2}} + {\log }_{2}\frac{1}{2} = 2 - 1 = 1$ .
    \end{jieda}
    
    \item 已知函数 $f\left( x\right)  = \left\{  \begin{array}{l} {\log }_{3}\left( {1 - x}\right) ,x < 1, \\  1 + {3}^{x - 1},x \geq  1, \end{array}\right.$ 则 $f\left( {-8}\right)  + f\left( {{\log }_{3}{15}}\right)  =$ \blkda[4]{8}.
    \begin{jieda}
        由题意，$f\left( {-8}\right)  + f\left( {{\log }_{3}{15}}\right)  = {\log }_{3}9 + 1 + {3}^{{\log }_{3}{15} - 1} = {\log }_{3}9 + 1 + {3}^{{\log }_{3}5} = 2 + 1 + 5 = 8$
    \end{jieda}
    
    \item 函数 $y = \lg \left( {{x}^{2} - {5x} + 4}\right)$ 的定义域为\blkda[4]{$\left( {4, + \infty }\right)  \cup  \left( {-\infty ,1}\right)$}.
    \begin{jieda}
        由对数函数定义可得 ${x}^{2} - {5x} + 4 > 0$ ,解得 $x > 4$ 或 $x < 1$ ,
        
        所以函数定义域为 $\left( {4, + \infty }\right)  \cup  \left( {-\infty ,1}\right)$ .
    \end{jieda}
    \item 在对数式中 ${\log }_{\left( x - 1\right) }\left( {3 - x}\right)$ ,实数 $x$ 的取值范围是\blkda[4]{$1 < x < 3$ 且 $x \neq  2$}.
    \begin{jieda}
        由对数式 ${\log }_{\left( x - 1\right) }\left( {3 - x}\right)$ 可知:
        
        $\left\{  \begin{array}{l} x - 1 > 0 \\  x - 1 \neq  1 \\  3 - x > 0 \end{array}\right.$ ,解之得: $1 < x < 3$ 且 $x \neq  2$
    \end{jieda}
    
    \item 函数 $y = {\log }_{3}\left( {{x}^{2} - {2x} + 4}\right)$ 的值域是\blkda[4]{$\lbrack 1, + \infty )$}.
    \begin{jieda}
        易知 ${x}^{2} - {2x} + 4 = {\left( x - 1\right) }^{2} + 3 \geq  3$， 又 $y = {\log }_{3}x$ 定义域上单调递增，所以 $y = {\log }_{3}\left( {{x}^{2} - {2x} + 4}\right)  \geq  {\log }_{3}3 = 1$ .
    \end{jieda}
    
    \item 已知函数 $f\left( x\right)  = 1 + {\log }_{a}\left( {{2x} - 3}\right) \left( {a > 0,a \neq  1}\right)$ 恒过定点$(m, n)$,则 $m + n =$\blkda[4]{3}
    \begin{jieda}
        令 ${2x} - 3 = 1$ ,解得 $x = 2$ ,此时 $f\left( 2\right)  = 1 + {\log }_{a}1 = 1$，所以 $f\left( x\right)$ 恒过定点(2,1),则 $m = 2,n = 1$，所以 $m + n = 3$ .
    \end{jieda}
        \item 已知函数 $f\left( x\right)  = {\log }_{a}\left( {x - b}\right)$ ( $a > 0$ 且 $a \neq  1,a,b$ 为常数) 的图象如图,则下列结论正确的是\dotfill{(\kh{D})}
    
    \begin{center}
    \includegraphics[width=0.3\textwidth]{bodyxb/src/bo_d2gik7bef24c73b0ne3g_0_239_1183_471_383_0.jpg}
    \end{center}

    \xx{$a > 0,b <  - 1$}{$a > 0, - 1 < b < 0$}{$0 < a < 1,b <  - 1$}{$0 < a < 1, - 1 < b < 0$}
    \begin{jieda}
        因为函数 $f\left( x\right)  = {\log }_{a}\left( {x - b}\right)$ 为减函数,所以 $0 < a < 1$
        
        又因为函数图象与 $x$ 轴的交点在正半轴,所以 $x = 1 + b > 0$ ,即 $b >  - 1$
        
        又因为函数图象与 $y$ 轴有交点,所以 $b < 0$ ,所以 $- 1 < b < 0$
    \end{jieda}
    
    \item 根据下列不等式,比较正数 $m$ 及 $n$ 的大小:
    \begin{dingautolist}{172}
        \item ${\log }_{3}m < {\log }_{3}n$
        \item ${\log }_{a}m < {\log }_{a}n$ ( $a > 0$ 且 $a \neq  1$ )
        \item ${\log }_{m}N < {\log }_{n}N\left( {0 < m < 1,0 < n < 1,0 < N < 1}\right)$ 
    \end{dingautolist}
    \begin{jieda}
        \begin{dingautolist}{172}
            \item $n > m$
            \item 当 $a > 1$ 时, $n > m$ ；当 $0 < a < 1$ 时, $m > n$
            \item $n > m$
        \end{dingautolist}
    \end{jieda}
    
    \item 不等式 $\lg \left( {x + 1}\right)  > 1$ 的解集为\blkda[4]{$\left( {9, + \infty }\right)$}.
    \begin{jieda}
        由不等式 $\lg \left( {x + 1}\right)  > 1$ ,得 $x + 1 > {10}$ ,解得 $x > 9$ ,
        
        所以不等式 $\lg \left( {x + 1}\right)  > 1$ 的解集为 $\left( {9, + \infty }\right)$ .
    \end{jieda} 
    \item 已知 $a \in  R,{\log }_{2}\left( {{4a} + 3}\right)  < 2$ ,则 $a$ 的取值范围为\blkda[4]{$\left( {-\dfrac{3}{4},\dfrac{1}{4}}\right)$}
    \begin{jieda}
        由 ${\log }_{2}\left( {{4a} + 3}\right)  < 2$ ,得 $0 < {4a} + 3 < 4$ ,解得 $- \frac{3}{4} < a < \frac{1}{4}$ ,
        
        所以 $a$ 的取值范围为 $\left( {-\dfrac{3}{4},\dfrac{1}{4}}\right)$ .
    \end{jieda}
    
    \item 方程 $\lg \left( {2 - x}\right)  + \lg \left( {3 - x}\right)  = \lg {12}$ 的解是\blkda[4]{$x =  - 1$}.
    \begin{jieda}
        由方程 $\lg \left( {2 - x}\right)  + \lg \left( {3 - x}\right)  = \lg {12}$ ,可得 $\lg \left\lbrack  {\left( {2 - x}\right) \left( {3 - x}\right) }\right\rbrack   = \lg {12}$
        
        $\therefore \left\{  \begin{array}{l} 2 - x > 0 \\  3 - x > 0 \\  \left( {2 - x}\right) \left( {3 - x}\right)  = {12} \end{array}\right.$ ,解得 $x =  - 1$ .
    \end{jieda}
    \item 不等式 ${\log }_{2}\left( {x + {14}}\right)  + {\log }_{2}\left( {x + 2}\right)  > 3 + {\log }_{2}\left( {x + 6}\right)$ 的解集为\blkda[4]{$\left( {2, + \infty }\right)$}.
    \begin{jieda}
        由 ${\log }_{2}\left( {x + {14}}\right)  + {\log }_{2}\left( {x + 2}\right)  > 3 + {\log }_{2}\left( {x + 6}\right)$，可得 ${\log }_{2}\left( {x + {14}}\right) \left( {x + 2}\right)  > {\log }_{2}8\left( {x + 6}\right)$，又 $y = {\log }_{2}x$ 在 $\left( {0, + \infty }\right)$ 上单调递增，所以 $\left\{  \begin{array}{l} \left( {x + {14}}\right) \left( {x + 2}\right)  > 8\left( {x + 6}\right) \\  x + {14} > 0 \\  x + 2 > 0 \\  x + 6 > 0 \end{array}\right.$ ,解不等式组可得 $x > 2$，所以不等式 ${\log }_{2}\left( {x + {14}}\right)  + {\log }_{2}\left( {x + 2}\right)  > 3 + {\log }_{2}\left( {x + 6}\right)$ 的解集为 $\left( {2, + \infty }\right)$ .
    \end{jieda}
\end{enumerate}
\subsection*{B Level}
\begin{enumerate}
    \item 下列函数中,在区间 $\left( {0, + \infty }\right)$ 上单调递增的是 \dotfill{(\kh{C})}
    \xx{$f\left( x\right)  =  - \ln x$}{$f\left( x\right)  = \dfrac{1}{{2}^{x}}$}{$f\left( x\right)  =  - \dfrac{1}{x}$}{$f\left( x\right)  = {3}^{\left| x - 1\right| }$}
    \item 函数 $y = \dfrac{\ln \left| x\right| }{{x}^{2} + 2}$ 的图像大致为\dotfill{(\kh{B})}
    \xx{\includegraphics[width=0.37\textwidth]{bodyxb/src/bo_d2gj7jref24c73b0nip0_0_330_545_457_391_0.jpg}}{\includegraphics[width=0.37\textwidth]{bodyxb/src/bo_d2gj7jref24c73b0nip0_0_865_548_461_389_0.jpg}}{\includegraphics[width=0.37\textwidth]{bodyxb/src/bo_d2gj7jref24c73b0nip0_0_332_937_453_388_0.jpg}}{\includegraphics[width=0.37\textwidth]{bodyxb/src/bo_d2gj7jref24c73b0nip0_0_867_937_444_385_0.jpg}}
    \begin{jieda}
        设 $y = f\left( x\right)  = \dfrac{\ln \left| x\right| }{{x}^{2} + 2}$ ,则函数 $f\left( x\right)$ 的定义域为 $\{ x \mid  x \neq  0\}$ ,关于原点对称,
        
        又 $f\left( {-x}\right)  = \dfrac{\ln \left| {-x}\right| }{{\left( -x\right) }^{2} + 2} = f\left( x\right)$ ,所以函数 $f\left( x\right)$ 为偶函数,排除 ${{AC}}$ ;
        
        当 $x \in  \left( {0,1}\right)$ 时, $\ln \left| x\right| < 0, {x}^{2} + 2 >  0$ ,所以 $f\left( x\right)  < 0$ ,排除 ${D}$ .
    \end{jieda}
    \item 在天文学中, 天体的明暗程度可以用星等或亮度来描述. 两颗星的星等与亮度满足 ${m}_{2} - {m}_{1} = \dfrac{5}{2}\lg \dfrac{{E}_{1}}{{E}_{2}}$ ,其中星等为 ${m_k}$ 的星的亮度为 ${E_k}\left( {k = 1,2}\right)$ . 已知太阳的星等是-26.7,天狼星的星等是 -1.45 ,则太阳与天狼星的亮度的比值为\blkda[4]{${10}^{10.1}$}
    \begin{jieda}
        两颗星的星等与亮度满足 ${m}_{2} - {m}_{1} = \dfrac{5}{2}\lg \dfrac{{E}_{1}}{{E}_{2}}$ ,令 ${m}_{2} =  - {1.45},{m}_{1} =  - {26.7}$，$\lg \dfrac{{E}_{1}}{{E}_{2}} = \dfrac{2}{5} \cdot  \left( {{m}_{2} - {m}_{1}}\right)  = \dfrac{2}{5}\left( {-{1.45} + {26.7}}\right)  = {10.1},\dfrac{{E}_{1}}{{E}_{2}} = {10}^{10.1}$ .
    \end{jieda}
    
    \item 若 $f\left( x\right)$ 是奇函数,且当 $x > 0$ 时, $f\left( x\right)  = 2 + \lg x$ ,则 $f\left( {-{10}}\right)  =$ \blkda[4]{$-3$}.
    \begin{jieda}
        由题意得 $f\left( {-{10}}\right)  =  - f\left( {10}\right)  =  - \left( {2 + \lg {10}}\right)  =  - 3$ .
    \end{jieda}
    
    \item 已知函数 $f\left( x\right)  = \left\{  \begin{array}{l}  - {x}^{2} - {2ax} - a,x < 0 \\  {{e}}^{x} + \ln \left( {x + 1}\right) ,x \geq  0 \end{array}\right.$ 在 $\mathbf{R}$ 上单调递增,则 $a$ 的取值范围是 \blkda[4]{$\left\lbrack  {-1,0}\right\rbrack$}
    \begin{jieda}
        因为 $f\left( x\right)$ 在 ${R}$ 上单调递增,且 $x \geq  0$ 时, $f\left( x\right)  = {{e}}^{x} + \ln \left( {x + 1}\right)$ 单调递增,
        
        则需满足 $\left\{  \begin{array}{l}  - \dfrac{-{2a}}{2 \times  \left( {-1}\right) } \geq  0 \\   - a \leq  {{e}}^{0} + \ln 1 \end{array}\right.$ ,解得 $- 1 \leq  a \leq  0$ ,
        
        即 $a$ 的范围是 $\left\lbrack  {-1,0}\right\rbrack$ .
    \end{jieda}
    
    \item 已知 $f\left( x\right)  = {\log }_{a}x\left( {a > 0\text{ 且 }a \neq  1}\right)$ 满足 $f\left( \dfrac{1}{2}\right)  + f\left( 4\right)  = 1$ 
    \begin{enumerate}
        \item 求 $f\left( 1\right)  + f\left( 2\right)  + f\left( 4\right)  + \cdots  + f\left( {2}^{10}\right)$ 的值
        \item 若 $g\left( x\right)  = f\left( {2 - x}\right)  + f\left( {1 + x}\right)$ ,求 $g\left( x\right)$ 的单调区间
    \end{enumerate}
    \begin{jieda}
        \begin{enumerate}
            \item 由已知 $f\left( x\right)  = {\log }_{a}x\left( {a > 0\text{且}a \neq  1}\right)$ 满足 $f\left( \dfrac{1}{2}\right)  + f\left( 4\right)  = 1$ ,则 ${\log }_{a}\dfrac{1}{2} + {\log }_{a}4 = 1$ ,即 $- {\log }_{a}2 + 2{\log }_{a}2 = 1$ ,得 ${\log }_{a}2 = 1$ ,所以 $a = 2$ ,得 $f\left( x\right)  = {\log }_{2}x$ ,则 $f\left( 1\right)  + f\left( 2\right)  + f\left( 4\right)  + \cdots  + f\left( {2}^{10}\right)  = {\log }_{2}1 + {\log }_{2}2 + {\log }_{2}4 + {\log }_{2}8 + \cdots  + {\log }_{2}{2}^{10} = 1 + 2 + 3 + \cdots  + {10} = {55}$
            
            \item 由 $\left\{  \begin{array}{l} 2 - x > 0, \\  1 + x > 0, \end{array}\right.$ 得 $- 1 < x < 2$，故$g(x)$的定义域为 $\left( {-1,2}\right)$
            
            $g\left( x\right)  = f\left( {2 - x}\right)  + f\left( {1 + x}\right)  = {\log }_{2}\left( {2 - x}\right)  + {\log }_{2}\left( {1 + x}\right)  = {\log }_{2}\left( {-{x}^{2} + x + 2 }\right)$，根据同增异减可知 $g\left( x\right)$ 的单调增区间为 $\left( {-1,\dfrac{1}{2}}\right\rbrack$ ,单调减区间为 $\left( {\dfrac{1}{2},2}\right)$
        \end{enumerate}
    \end{jieda}
    
    \item 设正实数 $a\text{ 、 }b\text{ 、 }c$ 分别满足 $a \cdot  {2}^{a} = 1,b \cdot  {\log }_{2}b = 1,c \cdot  {\log }_{3}c = 1$ ,则 $a\text{ 、 }b\text{ 、 }c$ 的大小关系为\dotfill(\kh{C})
    \xx{$a > b > c$}{$b > a > c$}{$c > b > a$}{$a > c > b$}
    \begin{jieda}
        $2^a = \dfrac{1}{a}$，$log_2b = \dfrac{1}{b}$，$log_3c = \dfrac{1}{c}$，根据函数图像即可判断.
    \end{jieda}
    
    \item 已知函数 $f\left( x\right)  = \left| {\lg x}\right|$ ,若 $f\left( {\lg m}\right)  > f\left( 2\right)$ ,则实数 $m$ 的取值范围是\blkda[4]{$\left( {1,\sqrt{10}}\right)  \cup  \left( {{100}, + \infty }\right)$}.
    \begin{jieda}
        如图,画出 $f\left( x\right)  = \left| {\lg x}\right|$ 的大致图象,易知 $f\left( \dfrac{1}{2}\right)  = f\left( 2\right)$ ,所以 $0 < \lg m < \dfrac{1}{2}$ 或 $\lg m > 2$ , 解得 $1 < m < \sqrt{10}$ 或 $m > {100}$ ,即实数 $m$ 的取值范围是 $\left( {1,\sqrt{10}}\right)  \cup  \left( {{100}, + \infty }\right)$
        
        \begin{center}
        \includegraphics[width=0.3\textwidth]{bodyxb/src/bo_d2gj7jref24c73b0nip0_3_245_1700_506_377_0.jpg}
        \end{center}
    \end{jieda} .
    \item 已知函数 $f\left( x\right)  = \left| {{\log }_{3}x}\right|$ ,若 $a < b$ ,且 $f\left( a\right)  = f\left( b\right)$ ,则 $a + {2b}$ 的取值范围是\blkda[4]{$\left( {3, + \infty }\right)$}.
    \begin{jieda}
        $f\left( x\right)  = \left| {{\log }_{3}x}\right|$ 的图象如下:
        
        \begin{center}
        \includegraphics[width=0.3\textwidth]{bodyxb/src/bo_d2gj7jref24c73b0nip0_4_241_612_373_252_0.jpg}
        \end{center}
        
        
        因为 $0 < a < b$ 且 $f\left( a\right)  = f\left( b\right)$ ,所以 $\left| {{\log }_{3}a}\right|  = \left| {{\log }_{3}b}\right|$ 且 $0 < a < 1,b > 1$ ,
        
        所以 $- {\log }_{3}a = {\log }_{3}b$ ,所以 ${ab} = 1$ ,故 $a + {2b} = a + \frac{2}{a}$ ,
        
        由对勾函数 $y = x + \dfrac{2}{x}$ 在(0,1)上单调递减,所以 $a + {2b} = a + \dfrac{2}{a} > 1 + 2 = 3$ ,
        
        所以 $a+ 2b$ 的取值范围是 $\left( {3, + \infty }\right)$ .
    \end{jieda}
    \item 已知函数 $f\left( x\right)  = \left| {{\log }_{2}x}\right|$ ,实数 $a,b$ 满足 $0 < a < b$ ,且 $f\left( a\right)  = f\left( b\right)$ ,若 $f\left( x\right)$ 在 $\left\lbrack  {{a}^{2},b}\right\rbrack$ 上的最大值为 2,则 $\dfrac{1}{a} + b =$ \blkda[4]{4}.
    \begin{jieda}
        如图所示: 根据函数 $f\left( x\right)  = \left| {{\log }_{2}x}\right|$ 的图象

        \begin{center}
        \includegraphics[width=0.3\textwidth]{bodyxb/src/bo_d2gj7jref24c73b0nip0_4_242_1769_257_197_0.jpg}
        \end{center}
        
        
        得 $0 < a < 1 < b$ ,所以 $0 < {a}^{2} < a < 1$ . 结合函数图象,
        
        易知当 $x = {a}^{2}$ 时 $f\left( x\right)$ 在 $\left\lbrack  {{a}^{2},b}\right\rbrack$ 上取得最大值,所以 $f\left( {a}^{2}\right)  = \left| {{\log }_{2}{a}^{2}}\right|  = 2$
        
        又 $0 < a < 1$ ,所以 $a = \dfrac{1}{2}$ ,
        
        再结合 $f\left( a\right)  = f\left( b\right)$ ,可得 $b = 2$ ,所以 $\dfrac{1}{a} + b = 2 + 2 = 4$ .
    \end{jieda}
    
    \item 已知函数 $y = {{e}}^{\left| \ln x\right| } - \left| {x - 2}\right|$ 与 $y = {ax}$ 的图像有 3 个不同公共点 (其中 ${e}$ 为自然对数的底数),则实数 $a$ 的取值范围是\blkda{$(0, 1)$}
    \begin{jieda}
        $f(x) = {{e}}^{\left| \ln x\right| } - \left| {x - 2}\right| = \left\{\begin{array}{ll}
            \dfrac{1}{x} + x - 2 & 0 < x \leq 1 \\
            2x - 2 & 1 < x \leq 2 \\
            2 & x > 2
        \end{array}\right.$
        \begin{center}
        \includegraphics[width=0.25\linewidth]{bodyxb//src/yikeiylian4-10.png}
        \end{center}

        根据函数图像可知$a \in (0, 1)$
    \end{jieda}
    \item 已知函数 $f\left( x\right)  = \left\{  \begin{array}{l}  - {x}^{2} - {2x},x \leq  0 \\  \left| {1 - \ln x}\right| ,x > 0 \end{array}\right.$ ,若 $f\left( x\right)  = m$ 存在四个不相等的实根 ${x}_{1},{x}_{2},{x}_{3},{x}_{4}$ ,且 ${x}_{1} < {x}_{2} < {x}_{3} < {x}_{4}$ ,则 ${x}_{4} - \left( {{x}_{1} + {x}_{2}}\right) {x}_{3}$ 的最小值是\blkda[4]{$2\sqrt{2}{e}$}.
    \begin{jieda}
        作函数 $f\left( x\right)  = \left\{  \begin{array}{l}  - {x}^{2} - {2x},x \leq  0 \\  \left| {1 - \ln x}\right| ,x > 0 \end{array}\right.$ 与 $y = m$ 图像如下:
        
        \begin{center}
        \includegraphics[width=0.4\textwidth]{bodyxb/src/bo_d2gj7jref24c73b0nip0_5_249_1281_576_445_0.jpg}
        \end{center}
        
        
        $y = m\because f\left( x\right)  = m$ 存在四个不相等的实根
        
        ${x}_{1},{x}_{2},{x}_{3},{x}_{4}$ ,且 ${x}_{1} < {x}_{2} < {x}_{3} < {x}_{4}$ ,
        
        则 ${x}_{1} + {x}_{2} =  - 2,{x}_{3}\text{ 、 }{x}_{4} > 0$ ,且 $1 - \ln {x}_{3} =  - \left( {1 - \ln {x}_{4}}\right)$ ,
        
        则 $\ln {x}_{3} + \ln {x}_{4} = 2$ ,即 $\ln( {x}_{3}{x}_{4}) = 2$ ,得 ${x}_{3}{x}_{4} = {{e}}^{2}$ ,
        
        则 ${x}_{4} - \left( {{x}_{1} + {x}_{2}}\right) {x}_{3} = {x}_{4} + 2{x}_{3} \geq  2\sqrt{2{x}_{3}{x}_{4}} = 2\sqrt{2}{e}$ ,
        
        当且仅当 ${x}_{4} = 2{x}_{3}$ 时,即 ${x}_{3} = \dfrac{\sqrt{2}}{2}{e},{x}_{4} = \sqrt{2}{e}$ 时,等号成立
    \end{jieda}
\end{enumerate}
\subsection*{C Level}
\begin{enumerate}
    \item 设函数 $f\left( x\right)  = \left\{  \begin{array}{l} {x}^{2},x \leq  1 \\  \left| {{\log }_{2}\left( {x - 1}\right) }\right| ,x > 1 \end{array}\right.$ ,若 $f\left( {x}_{1}\right)  = f\left( {x}_{2}\right)  = f\left( {x}_{3}\right)  = f\left( {x}_{4}\right)$ (其中$\left. {{x}_{1} < {x}_{2} < {x}_{3} < {x}_{4}}\right)$ ,则 $\dfrac{4}{{x}_{4} + 1} + \left( {{x}_{1} + {x}_{2} + 2}\right) {x}_{3}$ 的取值范围是 \blkda[4]{$\left\lbrack  {4,\dfrac{16}{3}}\right)$}
    \begin{jieda}
        当 $x \leq  1$ 时,函数 $f\left( x\right)  = {x}^{2}$ 在 $\left( {-\infty ,0}\right\rbrack$ 上递减,函数值集合为 $\lbrack 0, + \infty )$ ,在 $\left\lbrack  {0,1}\right\rbrack$ 上递增, 函数值集合为 $\left\lbrack  {0,1}\right\rbrack$ ,

        当 $x > 1$ 时,函数 $f\left( x\right)  = \left| {{\log }_{2}\left( {x - 1}\right) }\right|$ 在 $(1,2\rbrack$ 上递减,函数值集合为 $\lbrack 0, + \infty )$ ,在 $\lbrack 2, + \infty )$ 上递增, 函数值集合为 $\lbrack 0, + \infty )$ ,
        
        作出函数 $f\left( x\right)$ 的图象,如图,设 $f\left( {x}_{1}\right)  = f\left( {x}_{2}\right)  = f\left( {x}_{3}\right)  = f\left( {x}_{4}\right)  = t$ ,
        
        \begin{center}
        \includegraphics[width=0.4\textwidth]{bodyxb/src/bo_d2gj7jref24c73b0nip0_2_241_1692_494_349_0.jpg}
        \end{center}
        
        
        当 $0 < t \leq  1$ 时,直线 $y = t$ 与函数 $y = f\left( x\right)$ 的图象有四个交点,且交点的横坐标分别为
        
        ${x}_{1},{x}_{2},{x}_{3},{x}_{4}$ ,且 ${x}_{1} < {x}_{2} < {x}_{3} < {x}_{4}$ ,
        
        当 $x > 1$ 时,由 $f\left( x\right)  = \left| {{\log }_{2}\left( {x - 1}\right) }\right|  = 1$ ,解得 $x = \dfrac{3}{2}$ 或 $x = 3$ ,于是 $\dfrac{3}{2} \leq  {x}_{3} < 2,2 < {x}_{4} \leq  3$ ,
        
        由 $f\left( {x}_{3}\right)  = f\left( {x}_{4}\right)$ ,得 $- {\log }_{2}\left( {{x}_{3} - 1}\right)  = {\log }_{2}\left( {{x}_{4} - 1}\right)$ ,则 ${x}_{3} - 1 = \dfrac{1}{{x}_{4} - 1}$ ,即 ${x}_{3} = \dfrac{1}{{x}_{4} - 1} + 1$ ,而 ${x}_{1} + {x}_{2} = 0$ ,
        
        因此 $\dfrac{4}{{x}_{4} + 1} + \left( {{x}_{1} + {x}_{2} + 2}\right) {x}_{3} = \dfrac{4}{{x}_{4} + 1} + 2{x}_{3} = \dfrac{4}{{x}_{4} + 1} + \dfrac{2}{{x}_{4} - 1} + 2$ ,
        
        令 $g\left( x\right)  = \dfrac{4}{x + 1} + \dfrac{2}{x - 1} + 2$ ,显然函数 $g\left( x\right)$ 在 $(2,3\rbrack$ 上递减,且 $g\left( 2\right)  = \dfrac{16}{3},g\left( 3\right)  = 4$ ,于是
        
        $g\left( x\right)  \in  \left\lbrack  {4,\dfrac{16}{3}}\right) ,$
        
        所以 $\dfrac{4}{{x}_{4} + 1} + \left( {{x}_{1} + {x}_{2} + 2}\right) {x}_{3}$ 的取值范围是 $\left\lbrack  {4,\dfrac{16}{3}}\right)$ .
    \end{jieda}
    
    \item 已知函数 $f\left( x\right)  = {\log }_{3}\left( {{3x} + {12}}\right) ,g\left( x\right)  = m \cdot  {2}^{x} + 2\left( {m \neq  0}\right)$ ,若 $\forall {x}_{1} \in  \left\lbrack  {0,1}\right\rbrack  ,\exists {x}_{2} \in  \left\lbrack  {-3,5}\right\rbrack$，使得 $g\left( {x}_{1}\right)  = f\left( {x}_{2}\right)$ ,则实数 $m$ 的取值范围是\blkda[4]{$\left\lbrack  {-\dfrac{1}{2},0}\right)  \cup  \left( {0,\dfrac{1}{2}}\right\rbrack$}.
    \begin{jieda}
        当 $x \in  \left\lbrack  {-3,5}\right\rbrack$ 时, $f\left( x\right)  \in  \left\lbrack  {1,3}\right\rbrack$ ；
        
        当 $m > 0$ 时,当 $x \in  \left\lbrack  {0,1}\right\rbrack  ,g\left( x\right)  \in  \left\lbrack  {m + 2,{2m} + 2}\right\rbrack$ ,
        
        又 $\forall {x}_{1} \in  \left\lbrack  {0,1}\right\rbrack  ,\exists {x}_{2} \in  \left\lbrack  {-3,5}\right\rbrack$ ,使得 $g\left( {x}_{1}\right)  = f\left( {x}_{2}\right)$ ,
        
        所以 $\left\lbrack  {m + 2,{2m} + 2}\right\rbrack   \subseteq  \left\lbrack  {1,3}\right\rbrack$ ,
        
        所以 $\left\{  \begin{array}{l} m + 2 \geq  1 \\  {2m} + 2 \leq  3 \end{array}\right.$ ,解得 $0 < m \leq  \dfrac{1}{2}$ ;
        
        当 $m < 0$ 时,当 $x \in  \left\lbrack  {0,1}\right\rbrack  ,g\left( x\right)  \in  \left\lbrack  {{2m} + 2,m + 2}\right\rbrack$ ,
        
        又 $\forall {x}_{1} \in  \left\lbrack  {0,1}\right\rbrack  ,\exists {x}_{2} \in  \left\lbrack  {-3,5}\right\rbrack$ ,使得 $g\left( {x}_{1}\right)  = f\left( {x}_{2}\right)$ ,
        
        所以 $\left\lbrack  {{2m} + 2,m + 2}\right\rbrack   \subseteq  \left\lbrack  {1,3}\right\rbrack$ ,
        
        所以 $\left\{  \begin{array}{l} {2m} + 2 \geq  1 \\  m + 2 \leq  3 \end{array}\right.$ ,解得 $- \dfrac{1}{2} \leq  m < 0$ .
        
        综上,实数 $m$ 的取值范围是 $\left\lbrack  {-\dfrac{1}{2},0}\right)  \cup  \left( {0,\dfrac{1}{2}}\right\rbrack$ .
    \end{jieda}
    \item 已知函数 $f\left( x\right)  = {\log }_{2}\left( {x + a}\right)$
    \begin{enumerate}
        \item 当 $a = 2$ 时,解不等式: $f\left( x\right)  < {\log }_{2}x + 1$
        \item 当 $a = 0$ 时,不等式 ${\left\lbrack  f\left( x\right)  + 2\right\rbrack  }^{2} > m \cdot  f\left( x\right)$ 在 $x \in  \left( {1, + \infty }\right)$ 上恒成立,求实数 $m$ 的取值范围
        \item 是否存在实数 $a$ ,使得函数 $y = \left| {f\left( x\right) }\right|$ 在 $x \in  \left\lbrack  {-1,2}\right\rbrack$ 上的最大值为 ${\log }_{2}3$ ,若存在,请求出 $a$ 的值；若不存在,请说明理由
    \end{enumerate}
    \begin{jieda}
        \begin{enumerate}
            \item 由已知得 ${\log }_{2}\left( {x + 2}\right)  < {\log }_{2}x + 1$，即 ${\log }_{2}\left( {x + 2}\right)  < {\log }_{2}{2x}$ ,因为 $y = {\log }_{2}x$ 是增函数,
            所以 $\left\{  \begin{array}{l} x + 2 > 0 \\  {2x} > 0 \\  x + 2 < {2x} \end{array}\right.$ ,解得 $x > 2$，所以原不等式的解集为 $\{ x \mid  x > 2\}$
            \item 由题意令 $t = {\log }_{2}x$ ,因为 $x > 1$ ,所以 $t > 0$ ,
            所以不等式 ${\left\lbrack  f\left( x\right)  + 2\right\rbrack  }^{2} > m \cdot  f\left( x\right)$ 在 $x \in  \left( {1, + \infty }\right)$ 上恒成立,
            可化为 ${\left( t + 2\right) }^{2} > {mt}$ 在 $\left( {0, + \infty }\right)$ 上恒成立,
            分离参数得 $m < \dfrac{{t}^{2} + {4t} + 4}{t} = t + \dfrac{4}{t} + 4$ ,因为 $t + \dfrac{4}{t} \geq  2\sqrt{t \cdot  \dfrac{4}{t}} = 4$ ,当且仅当 $t = 2$ 时取等号,
            则要使原式恒成立,只需 $m < 4 + 4 = 8$ 即可,即实数 $m$ 的取值范围为 $\left( {-\infty ,8}\right)$
            \item 首先要使函数在 $x \in  \left\lbrack  {-1,2}\right\rbrack$ 上有意义,需 $- 1 + a > 0$ ,所以 $a > 1$ .
            易知函数 $y = \left| {f\left( x\right) }\right|$ 在 $x \in  \left\lbrack  {-1,2}\right\rbrack$ 上的最大值必在端点处产生,
            故只需 $\left\{  \begin{array}{l} \left| {{\log }_{2}\left( {-1 + a}\right) }\right|  = {\log }_{2}3 \\  \left| {{\log }_{2}\left( {2 + a}\right) }\right|  \leq  {\log }_{2}3 \end{array}\right.$ ,或 $\left\{  \begin{array}{l} \left| {{\log }_{2}\left( {-1 + a}\right) }\right|  \leq  {\log }_{2}3 \\  \left| {{\log }_{2}\left( {2 + a}\right) }\right|  = {\log }_{2}3 \end{array}\right.$ ,
            均无解，故不存在 $a$ 使结论成立.
        \end{enumerate}
    \end{jieda} 
\end{enumerate}
\newpage
\quad
\newpage